% =========================================================================
% Figura (clon de «Tema 2. Figuras», diapositiva «Equipotenciales»):
% equipotenciales de los términos monopolar (phi ~ 1/r), dipolar
% (phi ~ cos t / r^2) y cuadrupolar (phi ~ P2(cos t)/r^3), a intervalos
% iguales de potencial. Monopolo: circunferencias analíticas; dipolo y
% cuadrupolo: datos de figuras/datos/{dip,cua}_puntual_eq.tex.
% =========================================================================
\begin{figure}[ht]
\centering
\newcommand{\rotulo}[1]{{\footnotesize #1}\\[3pt]}
\renewcommand{\EQ}[1]{\draw[mmaEQsolo] plot coordinates {#1};}
\colorlet{marcoc}{grispizarra!60}
\begin{minipage}[t]{0.325\textwidth}\centering
\rotulo{Término monopolar ($\ell=0$)}
\begin{tikzpicture}[x=0.198cm, y=0.198cm]
  \panelMMA{10}{white}{0}{%
    \foreach \k in {0,...,12}{\pgfmathsetmacro{\rr}{1/(0.1+0.0194*\k)}\draw[mmaEQsolo] (0,0) circle (\rr);}}
  \marcoMMA{10}{5}{$x$ (u.\,arb.)}{$z$ (u.\,arb.)}
\end{tikzpicture}
\end{minipage}\hfill
\begin{minipage}[t]{0.325\textwidth}\centering
\rotulo{Término dipolar ($\ell=1$)}
\begin{tikzpicture}[x=0.198cm, y=0.198cm]
  \panelMMA{10}{white}{0}{%
    % Dipolo puntual - equipotenciales
% Generado por scripts/generar_datos.py (no editar a mano).
\EQ{(1.224,-0.476) (1.250,-0.522) (1.275,-0.574) (1.296,-0.625) (1.313,-0.675) (1.326,-0.725) (1.337,-0.775) (1.345,-0.825) (1.350,-0.875) (1.353,-0.925) (1.354,-0.975) (1.352,-1.025) (1.347,-1.075) (1.340,-1.125) (1.331,-1.175) (1.319,-1.225) (1.305,-1.275) (1.288,-1.325) (1.269,-1.375) (1.246,-1.425) (1.221,-1.475) (1.194,-1.525) (1.162,-1.575) (1.128,-1.625) (1.089,-1.675) (1.050,-1.721) (1.000,-1.774) (0.950,-1.822) (0.900,-1.865) (0.856,-1.900) (0.800,-1.940) (0.750,-1.972) (0.702,-2.000) (0.656,-2.025) (0.604,-2.050) (0.550,-2.074) (0.500,-2.093) (0.450,-2.111) (0.400,-2.126) (0.350,-2.139) (0.300,-2.151) (0.250,-2.161) (0.200,-2.168) (0.150,-2.174) (0.100,-2.179) (0.050,-2.181) (0.000,-2.182) (-0.050,-2.181) (-0.100,-2.179) (-0.150,-2.174) (-0.200,-2.168) (-0.250,-2.161) (-0.300,-2.151) (-0.350,-2.139) (-0.400,-2.126) (-0.450,-2.111) (-0.500,-2.093) (-0.547,-2.075) (-0.600,-2.052) (-0.650,-2.028) (-0.700,-2.001) (-0.745,-1.975) (-0.800,-1.940) (-0.850,-1.904) (-0.900,-1.865) (-0.947,-1.825) (-0.999,-1.775) (-1.047,-1.725) (-1.089,-1.675) (-1.125,-1.629) (-1.162,-1.575) (-1.194,-1.525) (-1.221,-1.475) (-1.246,-1.425) (-1.269,-1.375) (-1.288,-1.325) (-1.305,-1.275) (-1.319,-1.225) (-1.331,-1.175) (-1.340,-1.125) (-1.347,-1.075) (-1.352,-1.025) (-1.354,-0.975) (-1.353,-0.925) (-1.350,-0.875) (-1.345,-0.825) (-1.337,-0.775) (-1.326,-0.725) (-1.313,-0.675) (-1.296,-0.625) (-1.275,-0.575) (-1.252,-0.525) (-1.225,-0.477)}
\EQ{(1.229,-0.446) (1.263,-0.500) (1.290,-0.550) (1.313,-0.600) (1.333,-0.650) (1.349,-0.700) (1.362,-0.750) (1.373,-0.800) (1.381,-0.850) (1.386,-0.900) (1.389,-0.950) (1.389,-1.000) (1.387,-1.050) (1.383,-1.100) (1.376,-1.150) (1.367,-1.200) (1.355,-1.250) (1.341,-1.300) (1.325,-1.349) (1.306,-1.400) (1.284,-1.450) (1.260,-1.500) (1.233,-1.550) (1.203,-1.600) (1.175,-1.642) (1.132,-1.700) (1.100,-1.740) (1.050,-1.796) (1.000,-1.846) (0.950,-1.891) (0.909,-1.925) (0.850,-1.970) (0.806,-2.000) (0.750,-2.035) (0.700,-2.063) (0.650,-2.089) (0.600,-2.112) (0.550,-2.133) (0.500,-2.152) (0.450,-2.169) (0.400,-2.184) (0.350,-2.197) (0.300,-2.208) (0.250,-2.218) (0.200,-2.225) (0.150,-2.231) (0.100,-2.236) (0.050,-2.238) (0.000,-2.239) (-0.050,-2.238) (-0.100,-2.236) (-0.150,-2.231) (-0.200,-2.225) (-0.250,-2.218) (-0.300,-2.208) (-0.350,-2.197) (-0.400,-2.184) (-0.450,-2.169) (-0.500,-2.152) (-0.550,-2.133) (-0.600,-2.112) (-0.650,-2.089) (-0.700,-2.063) (-0.750,-2.035) (-0.800,-2.004) (-0.843,-1.975) (-0.900,-1.932) (-0.940,-1.900) (-0.996,-1.850) (-1.046,-1.800) (-1.091,-1.750) (-1.125,-1.709) (-1.169,-1.650) (-1.200,-1.604) (-1.233,-1.550) (-1.260,-1.500) (-1.284,-1.450) (-1.306,-1.400) (-1.325,-1.350) (-1.341,-1.300) (-1.355,-1.250) (-1.367,-1.200) (-1.376,-1.150) (-1.383,-1.100) (-1.387,-1.050) (-1.389,-1.000) (-1.389,-0.950) (-1.386,-0.900) (-1.381,-0.850) (-1.373,-0.800) (-1.362,-0.750) (-1.349,-0.700) (-1.333,-0.650) (-1.313,-0.600) (-1.290,-0.550) (-1.263,-0.500) (-1.232,-0.450) (-1.229,-0.446)}
\EQ{(1.244,-0.431) (1.274,-0.475) (1.300,-0.518) (1.325,-0.565) (1.350,-0.619) (1.372,-0.675) (1.388,-0.725) (1.400,-0.775) (1.411,-0.825) (1.419,-0.875) (1.424,-0.925) (1.427,-0.975) (1.427,-1.025) (1.425,-1.075) (1.421,-1.125) (1.414,-1.175) (1.406,-1.225) (1.395,-1.275) (1.381,-1.325) (1.365,-1.375) (1.347,-1.425) (1.326,-1.475) (1.303,-1.525) (1.277,-1.575) (1.250,-1.621) (1.215,-1.675) (1.180,-1.725) (1.141,-1.775) (1.100,-1.822) (1.050,-1.874) (1.000,-1.922) (0.950,-1.965) (0.906,-2.000) (0.850,-2.040) (0.800,-2.073) (0.755,-2.100) (0.709,-2.125) (0.660,-2.150) (0.605,-2.175) (0.550,-2.198) (0.500,-2.216) (0.450,-2.233) (0.400,-2.247) (0.350,-2.260) (0.300,-2.271) (0.250,-2.280) (0.200,-2.287) (0.150,-2.293) (0.100,-2.297) (0.050,-2.299) (0.000,-2.300) (-0.050,-2.299) (-0.100,-2.297) (-0.150,-2.293) (-0.200,-2.287) (-0.250,-2.280) (-0.300,-2.271) (-0.350,-2.260) (-0.400,-2.247) (-0.450,-2.233) (-0.500,-2.216) (-0.550,-2.198) (-0.600,-2.177) (-0.650,-2.155) (-0.700,-2.130) (-0.750,-2.103) (-0.796,-2.075) (-0.850,-2.040) (-0.900,-2.004) (-0.950,-1.965) (-0.997,-1.925) (-1.049,-1.875) (-1.097,-1.825) (-1.141,-1.775) (-1.175,-1.731) (-1.215,-1.675) (-1.248,-1.625) (-1.275,-1.578) (-1.300,-1.531) (-1.325,-1.478) (-1.347,-1.425) (-1.365,-1.375) (-1.381,-1.325) (-1.395,-1.275) (-1.406,-1.225) (-1.414,-1.175) (-1.421,-1.125) (-1.425,-1.075) (-1.427,-1.025) (-1.427,-0.975) (-1.424,-0.925) (-1.419,-0.875) (-1.411,-0.825) (-1.400,-0.775) (-1.388,-0.725) (-1.372,-0.675) (-1.352,-0.625) (-1.330,-0.575) (-1.304,-0.525) (-1.275,-0.477) (-1.244,-0.431)}
\EQ{(1.250,-0.405) (1.285,-0.450) (1.318,-0.500) (1.347,-0.550) (1.373,-0.600) (1.395,-0.650) (1.413,-0.700) (1.429,-0.750) (1.442,-0.800) (1.452,-0.850) (1.460,-0.900) (1.465,-0.950) (1.468,-1.000) (1.468,-1.050) (1.467,-1.100) (1.463,-1.150) (1.457,-1.200) (1.448,-1.250) (1.438,-1.300) (1.425,-1.350) (1.410,-1.400) (1.392,-1.450) (1.372,-1.500) (1.350,-1.550) (1.325,-1.600) (1.300,-1.645) (1.266,-1.700) (1.232,-1.750) (1.200,-1.793) (1.154,-1.850) (1.109,-1.900) (1.059,-1.950) (1.004,-2.000) (0.950,-2.044) (0.909,-2.075) (0.850,-2.116) (0.800,-2.147) (0.751,-2.175) (0.704,-2.200) (0.652,-2.225) (0.600,-2.248) (0.550,-2.268) (0.500,-2.285) (0.450,-2.301) (0.400,-2.315) (0.350,-2.328) (0.300,-2.338) (0.250,-2.347) (0.200,-2.354) (0.150,-2.360) (0.100,-2.364) (0.050,-2.366) (0.000,-2.367) (-0.050,-2.366) (-0.100,-2.364) (-0.150,-2.360) (-0.200,-2.354) (-0.250,-2.347) (-0.300,-2.338) (-0.350,-2.328) (-0.400,-2.315) (-0.450,-2.301) (-0.500,-2.285) (-0.550,-2.268) (-0.600,-2.248) (-0.650,-2.226) (-0.700,-2.202) (-0.750,-2.176) (-0.795,-2.150) (-0.850,-2.116) (-0.900,-2.081) (-0.942,-2.050) (-1.000,-2.003) (-1.050,-1.958) (-1.100,-1.909) (-1.150,-1.854) (-1.195,-1.800) (-1.225,-1.760) (-1.266,-1.700) (-1.297,-1.650) (-1.325,-1.600) (-1.350,-1.550) (-1.372,-1.500) (-1.392,-1.450) (-1.410,-1.400) (-1.425,-1.350) (-1.438,-1.300) (-1.448,-1.250) (-1.457,-1.200) (-1.463,-1.150) (-1.467,-1.100) (-1.468,-1.050) (-1.468,-1.000) (-1.465,-0.950) (-1.460,-0.900) (-1.452,-0.850) (-1.442,-0.800) (-1.429,-0.750) (-1.413,-0.700) (-1.395,-0.650) (-1.373,-0.600) (-1.350,-0.555) (-1.325,-0.512) (-1.285,-0.450) (-1.250,-0.405)}
\EQ{(1.251,-0.374) (1.295,-0.425) (1.325,-0.465) (1.365,-0.525) (1.393,-0.575) (1.418,-0.625) (1.440,-0.675) (1.458,-0.725) (1.473,-0.775) (1.486,-0.825) (1.496,-0.875) (1.504,-0.925) (1.510,-0.975) (1.513,-1.025) (1.514,-1.075) (1.512,-1.125) (1.509,-1.175) (1.503,-1.225) (1.495,-1.275) (1.486,-1.325) (1.473,-1.375) (1.459,-1.425) (1.442,-1.475) (1.423,-1.525) (1.402,-1.575) (1.378,-1.625) (1.352,-1.675) (1.325,-1.721) (1.290,-1.775) (1.255,-1.825) (1.216,-1.875) (1.175,-1.923) (1.127,-1.975) (1.075,-2.025) (1.025,-2.070) (0.975,-2.110) (0.925,-2.147) (0.875,-2.182) (0.825,-2.213) (0.775,-2.242) (0.725,-2.268) (0.675,-2.292) (0.625,-2.314) (0.575,-2.334) (0.525,-2.352) (0.475,-2.369) (0.425,-2.383) (0.375,-2.396) (0.325,-2.407) (0.275,-2.416) (0.225,-2.424) (0.175,-2.430) (0.125,-2.435) (0.075,-2.438) (0.025,-2.440) (-0.050,-2.439) (-0.100,-2.437) (-0.150,-2.433) (-0.200,-2.427) (-0.250,-2.420) (-0.300,-2.412) (-0.350,-2.402) (-0.400,-2.390) (-0.450,-2.376) (-0.500,-2.361) (-0.550,-2.344) (-0.599,-2.325) (-0.650,-2.304) (-0.700,-2.280) (-0.750,-2.255) (-0.800,-2.228) (-0.846,-2.200) (-0.900,-2.165) (-0.950,-2.129) (-1.000,-2.090) (-1.048,-2.050) (-1.100,-2.002) (-1.150,-1.951) (-1.195,-1.900) (-1.236,-1.850) (-1.273,-1.800) (-1.307,-1.750) (-1.337,-1.700) (-1.365,-1.650) (-1.390,-1.600) (-1.413,-1.550) (-1.433,-1.500) (-1.450,-1.453) (-1.466,-1.400) (-1.480,-1.350) (-1.491,-1.300) (-1.500,-1.250) (-1.506,-1.200) (-1.511,-1.150) (-1.513,-1.100) (-1.514,-1.050) (-1.512,-1.000) (-1.507,-0.950) (-1.501,-0.900) (-1.492,-0.850) (-1.480,-0.800) (-1.466,-0.750) (-1.449,-0.700) (-1.429,-0.650) (-1.406,-0.600) (-1.379,-0.550) (-1.350,-0.502) (-1.314,-0.450) (-1.275,-0.401) (-1.251,-0.374)}
\EQ{(1.266,-0.359) (1.305,-0.400) (1.346,-0.450) (1.382,-0.500) (1.414,-0.550) (1.442,-0.600) (1.467,-0.650) (1.488,-0.700) (1.506,-0.750) (1.522,-0.800) (1.534,-0.850) (1.545,-0.900) (1.553,-0.950) (1.559,-1.000) (1.562,-1.050) (1.563,-1.100) (1.563,-1.150) (1.560,-1.200) (1.555,-1.250) (1.547,-1.300) (1.538,-1.350) (1.527,-1.400) (1.513,-1.450) (1.498,-1.500) (1.480,-1.550) (1.460,-1.600) (1.437,-1.650) (1.412,-1.700) (1.385,-1.750) (1.354,-1.800) (1.325,-1.844) (1.285,-1.900) (1.250,-1.944) (1.201,-2.000) (1.154,-2.050) (1.101,-2.100) (1.050,-2.145) (1.000,-2.185) (0.950,-2.222) (0.908,-2.250) (0.850,-2.287) (0.800,-2.315) (0.750,-2.342) (0.700,-2.366) (0.650,-2.388) (0.600,-2.409) (0.550,-2.427) (0.500,-2.443) (0.450,-2.458) (0.400,-2.471) (0.350,-2.483) (0.300,-2.493) (0.250,-2.501) (0.200,-2.508) (0.150,-2.513) (0.100,-2.517) (0.050,-2.519) (0.000,-2.520) (-0.050,-2.519) (-0.100,-2.517) (-0.150,-2.513) (-0.200,-2.508) (-0.250,-2.501) (-0.300,-2.493) (-0.350,-2.483) (-0.400,-2.471) (-0.450,-2.458) (-0.500,-2.443) (-0.550,-2.427) (-0.600,-2.409) (-0.650,-2.388) (-0.700,-2.366) (-0.750,-2.342) (-0.800,-2.315) (-0.850,-2.287) (-0.900,-2.256) (-0.945,-2.225) (-1.000,-2.185) (-1.044,-2.150) (-1.100,-2.101) (-1.150,-2.054) (-1.200,-2.001) (-1.245,-1.950) (-1.285,-1.900) (-1.321,-1.850) (-1.350,-1.807) (-1.385,-1.750) (-1.412,-1.700) (-1.437,-1.650) (-1.460,-1.600) (-1.480,-1.550) (-1.498,-1.500) (-1.513,-1.450) (-1.527,-1.400) (-1.538,-1.350) (-1.547,-1.300) (-1.555,-1.250) (-1.560,-1.200) (-1.563,-1.150) (-1.563,-1.100) (-1.562,-1.050) (-1.559,-1.000) (-1.553,-0.950) (-1.545,-0.900) (-1.534,-0.850) (-1.522,-0.800) (-1.506,-0.750) (-1.488,-0.700) (-1.467,-0.650) (-1.442,-0.600) (-1.414,-0.550) (-1.382,-0.500) (-1.350,-0.456) (-1.305,-0.400) (-1.266,-0.359)}
\EQ{(1.275,-0.337) (1.314,-0.375) (1.350,-0.414) (1.400,-0.475) (1.436,-0.525) (1.467,-0.575) (1.495,-0.625) (1.519,-0.675) (1.540,-0.725) (1.559,-0.775) (1.574,-0.825) (1.587,-0.875) (1.598,-0.925) (1.606,-0.975) (1.612,-1.025) (1.616,-1.075) (1.618,-1.125) (1.618,-1.175) (1.616,-1.225) (1.611,-1.275) (1.605,-1.325) (1.597,-1.375) (1.586,-1.425) (1.574,-1.475) (1.559,-1.525) (1.543,-1.575) (1.524,-1.625) (1.503,-1.675) (1.479,-1.725) (1.454,-1.775) (1.425,-1.825) (1.394,-1.875) (1.360,-1.925) (1.325,-1.973) (1.283,-2.025) (1.250,-2.063) (1.216,-2.100) (1.175,-2.141) (1.125,-2.187) (1.081,-2.225) (1.025,-2.269) (0.983,-2.300) (0.925,-2.339) (0.875,-2.370) (0.825,-2.398) (0.775,-2.425) (0.725,-2.449) (0.675,-2.471) (0.625,-2.492) (0.575,-2.510) (0.525,-2.527) (0.475,-2.542) (0.425,-2.555) (0.375,-2.567) (0.325,-2.578) (0.275,-2.586) (0.225,-2.594) (0.175,-2.599) (0.125,-2.604) (0.075,-2.607) (0.025,-2.608) (-0.050,-2.608) (-0.100,-2.605) (-0.150,-2.602) (-0.200,-2.597) (-0.250,-2.590) (-0.300,-2.582) (-0.350,-2.573) (-0.400,-2.562) (-0.450,-2.549) (-0.500,-2.535) (-0.550,-2.519) (-0.600,-2.501) (-0.650,-2.482) (-0.700,-2.460) (-0.750,-2.437) (-0.800,-2.412) (-0.850,-2.384) (-0.900,-2.355) (-0.946,-2.325) (-1.000,-2.288) (-1.050,-2.250) (-1.100,-2.209) (-1.139,-2.175) (-1.191,-2.125) (-1.239,-2.075) (-1.275,-2.034) (-1.323,-1.975) (-1.360,-1.925) (-1.394,-1.875) (-1.425,-1.825) (-1.450,-1.781) (-1.475,-1.734) (-1.500,-1.681) (-1.524,-1.625) (-1.543,-1.575) (-1.559,-1.525) (-1.574,-1.475) (-1.586,-1.425) (-1.597,-1.375) (-1.605,-1.325) (-1.611,-1.275) (-1.616,-1.225) (-1.618,-1.175) (-1.618,-1.125) (-1.616,-1.075) (-1.612,-1.025) (-1.606,-0.975) (-1.598,-0.925) (-1.587,-0.875) (-1.574,-0.825) (-1.559,-0.775) (-1.540,-0.725) (-1.519,-0.675) (-1.495,-0.625) (-1.467,-0.575) (-1.436,-0.525) (-1.400,-0.475) (-1.360,-0.425) (-1.325,-0.387) (-1.289,-0.350) (-1.275,-0.337)}
\EQ{(1.275,-0.308) (1.323,-0.350) (1.374,-0.400) (1.419,-0.450) (1.450,-0.489) (1.493,-0.550) (1.524,-0.600) (1.550,-0.647) (1.575,-0.698) (1.597,-0.750) (1.616,-0.800) (1.632,-0.850) (1.645,-0.900) (1.656,-0.950) (1.665,-1.000) (1.672,-1.050) (1.676,-1.100) (1.679,-1.150) (1.679,-1.200) (1.678,-1.250) (1.674,-1.300) (1.669,-1.350) (1.662,-1.400) (1.652,-1.450) (1.641,-1.500) (1.628,-1.550) (1.613,-1.600) (1.595,-1.650) (1.576,-1.700) (1.554,-1.750) (1.530,-1.800) (1.504,-1.850) (1.475,-1.900) (1.444,-1.950) (1.410,-2.000) (1.375,-2.047) (1.333,-2.100) (1.289,-2.150) (1.250,-2.191) (1.200,-2.240) (1.162,-2.275) (1.102,-2.325) (1.050,-2.365) (1.001,-2.400) (0.950,-2.434) (0.900,-2.464) (0.850,-2.492) (0.800,-2.518) (0.750,-2.543) (0.700,-2.565) (0.650,-2.585) (0.600,-2.604) (0.550,-2.621) (0.500,-2.636) (0.450,-2.650) (0.400,-2.662) (0.350,-2.672) (0.300,-2.682) (0.250,-2.689) (0.200,-2.696) (0.150,-2.700) (0.100,-2.704) (0.050,-2.706) (0.000,-2.707) (-0.050,-2.706) (-0.100,-2.704) (-0.150,-2.700) (-0.200,-2.696) (-0.250,-2.689) (-0.300,-2.682) (-0.350,-2.672) (-0.400,-2.662) (-0.450,-2.650) (-0.500,-2.636) (-0.550,-2.621) (-0.600,-2.604) (-0.650,-2.585) (-0.700,-2.565) (-0.750,-2.543) (-0.800,-2.518) (-0.850,-2.492) (-0.900,-2.464) (-0.950,-2.434) (-1.000,-2.401) (-1.050,-2.365) (-1.100,-2.327) (-1.150,-2.285) (-1.189,-2.250) (-1.241,-2.200) (-1.289,-2.150) (-1.325,-2.109) (-1.373,-2.050) (-1.410,-2.000) (-1.444,-1.950) (-1.475,-1.900) (-1.504,-1.850) (-1.530,-1.800) (-1.554,-1.750) (-1.575,-1.702) (-1.595,-1.650) (-1.613,-1.600) (-1.628,-1.550) (-1.641,-1.500) (-1.652,-1.450) (-1.662,-1.400) (-1.669,-1.350) (-1.674,-1.300) (-1.678,-1.250) (-1.679,-1.200) (-1.679,-1.150) (-1.676,-1.100) (-1.672,-1.050) (-1.665,-1.000) (-1.656,-0.950) (-1.645,-0.900) (-1.632,-0.850) (-1.616,-0.800) (-1.597,-0.750) (-1.576,-0.700) (-1.552,-0.650) (-1.525,-0.601) (-1.493,-0.550) (-1.458,-0.500) (-1.425,-0.458) (-1.375,-0.401) (-1.325,-0.352) (-1.275,-0.308)}
\EQ{(1.275,-0.280) (1.325,-0.319) (1.375,-0.362) (1.414,-0.400) (1.450,-0.438) (1.500,-0.497) (1.539,-0.550) (1.572,-0.600) (1.600,-0.648) (1.625,-0.696) (1.650,-0.750) (1.670,-0.800) (1.688,-0.850) (1.703,-0.900) (1.716,-0.950) (1.726,-1.000) (1.734,-1.050) (1.741,-1.100) (1.745,-1.150) (1.747,-1.200) (1.748,-1.250) (1.746,-1.300) (1.743,-1.350) (1.738,-1.400) (1.731,-1.450) (1.722,-1.500) (1.711,-1.550) (1.698,-1.600) (1.683,-1.650) (1.667,-1.700) (1.648,-1.750) (1.627,-1.800) (1.604,-1.850) (1.579,-1.900) (1.552,-1.950) (1.525,-1.995) (1.489,-2.050) (1.454,-2.100) (1.416,-2.150) (1.375,-2.200) (1.330,-2.250) (1.282,-2.300) (1.229,-2.350) (1.175,-2.397) (1.125,-2.437) (1.075,-2.475) (1.025,-2.509) (0.975,-2.542) (0.925,-2.571) (0.875,-2.599) (0.825,-2.625) (0.775,-2.649) (0.725,-2.671) (0.675,-2.691) (0.625,-2.710) (0.575,-2.727) (0.525,-2.742) (0.475,-2.756) (0.425,-2.768) (0.375,-2.779) (0.325,-2.789) (0.275,-2.797) (0.225,-2.804) (0.175,-2.809) (0.125,-2.813) (0.075,-2.816) (0.025,-2.817) (-0.050,-2.817) (-0.100,-2.815) (-0.150,-2.811) (-0.200,-2.807) (-0.250,-2.800) (-0.300,-2.793) (-0.350,-2.784) (-0.400,-2.774) (-0.450,-2.762) (-0.500,-2.749) (-0.550,-2.735) (-0.600,-2.719) (-0.650,-2.701) (-0.700,-2.681) (-0.750,-2.660) (-0.800,-2.637) (-0.850,-2.612) (-0.900,-2.586) (-0.950,-2.557) (-1.000,-2.526) (-1.050,-2.492) (-1.100,-2.456) (-1.141,-2.425) (-1.200,-2.376) (-1.250,-2.331) (-1.300,-2.282) (-1.350,-2.228) (-1.396,-2.175) (-1.436,-2.125) (-1.472,-2.075) (-1.506,-2.025) (-1.537,-1.975) (-1.566,-1.925) (-1.592,-1.875) (-1.616,-1.825) (-1.638,-1.775) (-1.658,-1.725) (-1.675,-1.675) (-1.691,-1.625) (-1.705,-1.575) (-1.716,-1.525) (-1.726,-1.475) (-1.734,-1.425) (-1.741,-1.375) (-1.745,-1.325) (-1.747,-1.275) (-1.748,-1.225) (-1.746,-1.175) (-1.743,-1.125) (-1.738,-1.075) (-1.730,-1.025) (-1.721,-0.975) (-1.709,-0.925) (-1.695,-0.875) (-1.679,-0.825) (-1.660,-0.775) (-1.639,-0.725) (-1.614,-0.675) (-1.587,-0.625) (-1.556,-0.575) (-1.525,-0.530) (-1.482,-0.475) (-1.438,-0.425) (-1.400,-0.386) (-1.362,-0.350) (-1.301,-0.300) (-1.275,-0.280)}
\EQ{(1.288,-0.262) (1.342,-0.300) (1.400,-0.347) (1.450,-0.391) (1.500,-0.442) (1.550,-0.499) (1.590,-0.550) (1.625,-0.600) (1.656,-0.650) (1.684,-0.700) (1.708,-0.750) (1.730,-0.800) (1.750,-0.850) (1.767,-0.900) (1.781,-0.950) (1.794,-1.000) (1.804,-1.050) (1.812,-1.100) (1.818,-1.150) (1.823,-1.200) (1.825,-1.250) (1.825,-1.300) (1.824,-1.350) (1.821,-1.400) (1.816,-1.450) (1.810,-1.500) (1.801,-1.550) (1.791,-1.600) (1.779,-1.650) (1.765,-1.700) (1.750,-1.748) (1.732,-1.800) (1.712,-1.850) (1.691,-1.900) (1.667,-1.950) (1.641,-2.000) (1.613,-2.050) (1.583,-2.100) (1.550,-2.149) (1.514,-2.200) (1.476,-2.250) (1.434,-2.300) (1.400,-2.339) (1.366,-2.375) (1.325,-2.416) (1.289,-2.450) (1.250,-2.484) (1.201,-2.525) (1.150,-2.565) (1.101,-2.600) (1.050,-2.634) (1.000,-2.666) (0.950,-2.695) (0.900,-2.722) (0.850,-2.748) (0.800,-2.771) (0.750,-2.793) (0.700,-2.813) (0.650,-2.831) (0.600,-2.848) (0.550,-2.864) (0.500,-2.878) (0.450,-2.890) (0.400,-2.901) (0.350,-2.911) (0.300,-2.919) (0.250,-2.926) (0.200,-2.932) (0.150,-2.937) (0.100,-2.940) (0.050,-2.942) (0.000,-2.943) (-0.050,-2.942) (-0.100,-2.940) (-0.150,-2.937) (-0.200,-2.932) (-0.250,-2.926) (-0.300,-2.919) (-0.350,-2.911) (-0.400,-2.901) (-0.450,-2.890) (-0.500,-2.878) (-0.550,-2.864) (-0.600,-2.848) (-0.650,-2.831) (-0.700,-2.813) (-0.750,-2.793) (-0.800,-2.771) (-0.850,-2.748) (-0.900,-2.722) (-0.950,-2.695) (-1.000,-2.666) (-1.050,-2.634) (-1.100,-2.601) (-1.150,-2.565) (-1.200,-2.526) (-1.250,-2.484) (-1.289,-2.450) (-1.325,-2.416) (-1.366,-2.375) (-1.400,-2.339) (-1.434,-2.300) (-1.475,-2.251) (-1.514,-2.200) (-1.550,-2.150) (-1.583,-2.100) (-1.613,-2.050) (-1.641,-2.000) (-1.667,-1.950) (-1.691,-1.900) (-1.712,-1.850) (-1.732,-1.800) (-1.749,-1.750) (-1.765,-1.700) (-1.779,-1.650) (-1.791,-1.600) (-1.801,-1.550) (-1.810,-1.500) (-1.816,-1.450) (-1.821,-1.400) (-1.824,-1.350) (-1.825,-1.300) (-1.825,-1.250) (-1.823,-1.200) (-1.818,-1.150) (-1.812,-1.100) (-1.804,-1.050) (-1.794,-1.000) (-1.781,-0.950) (-1.767,-0.900) (-1.750,-0.851) (-1.730,-0.800) (-1.708,-0.750) (-1.684,-0.700) (-1.656,-0.650) (-1.625,-0.601) (-1.590,-0.550) (-1.551,-0.500) (-1.508,-0.450) (-1.459,-0.400) (-1.404,-0.350) (-1.350,-0.307) (-1.307,-0.275) (-1.288,-0.262)}
\EQ{(1.300,-0.243) (1.350,-0.275) (1.400,-0.310) (1.450,-0.348) (1.500,-0.391) (1.550,-0.439) (1.600,-0.493) (1.647,-0.550) (1.684,-0.600) (1.717,-0.650) (1.747,-0.700) (1.774,-0.750) (1.798,-0.800) (1.819,-0.850) (1.838,-0.900) (1.855,-0.950) (1.869,-1.000) (1.881,-1.050) (1.892,-1.100) (1.900,-1.150) (1.906,-1.200) (1.911,-1.250) (1.914,-1.300) (1.915,-1.350) (1.914,-1.400) (1.911,-1.450) (1.907,-1.500) (1.901,-1.550) (1.894,-1.600) (1.885,-1.650) (1.873,-1.700) (1.861,-1.750) (1.846,-1.800) (1.830,-1.850) (1.812,-1.900) (1.792,-1.950) (1.770,-2.000) (1.746,-2.050) (1.720,-2.100) (1.692,-2.150) (1.661,-2.200) (1.629,-2.250) (1.593,-2.300) (1.555,-2.350) (1.515,-2.400) (1.475,-2.445) (1.425,-2.498) (1.375,-2.548) (1.325,-2.593) (1.275,-2.636) (1.226,-2.675) (1.175,-2.713) (1.125,-2.747) (1.083,-2.775) (1.025,-2.810) (0.975,-2.838) (0.925,-2.865) (0.875,-2.890) (0.825,-2.913) (0.775,-2.934) (0.725,-2.954) (0.675,-2.972) (0.625,-2.989) (0.575,-3.004) (0.525,-3.018) (0.475,-3.030) (0.425,-3.042) (0.375,-3.052) (0.325,-3.060) (0.275,-3.068) (0.225,-3.074) (0.175,-3.079) (0.125,-3.082) (0.075,-3.085) (0.025,-3.086) (-0.050,-3.086) (-0.100,-3.084) (-0.150,-3.081) (-0.200,-3.076) (-0.250,-3.071) (-0.300,-3.064) (-0.350,-3.056) (-0.400,-3.047) (-0.450,-3.036) (-0.500,-3.024) (-0.550,-3.011) (-0.600,-2.996) (-0.650,-2.980) (-0.700,-2.963) (-0.750,-2.944) (-0.800,-2.923) (-0.850,-2.901) (-0.900,-2.877) (-0.950,-2.852) (-0.999,-2.825) (-1.050,-2.795) (-1.100,-2.764) (-1.150,-2.730) (-1.192,-2.700) (-1.250,-2.656) (-1.300,-2.615) (-1.345,-2.575) (-1.398,-2.525) (-1.448,-2.475) (-1.493,-2.425) (-1.535,-2.375) (-1.575,-2.325) (-1.611,-2.275) (-1.645,-2.225) (-1.675,-2.178) (-1.706,-2.125) (-1.733,-2.075) (-1.758,-2.025) (-1.781,-1.975) (-1.802,-1.925) (-1.821,-1.875) (-1.838,-1.825) (-1.854,-1.775) (-1.867,-1.725) (-1.879,-1.675) (-1.889,-1.625) (-1.898,-1.575) (-1.905,-1.525) (-1.910,-1.475) (-1.913,-1.425) (-1.915,-1.375) (-1.914,-1.325) (-1.913,-1.275) (-1.909,-1.225) (-1.903,-1.175) (-1.896,-1.125) (-1.887,-1.075) (-1.875,-1.025) (-1.862,-0.975) (-1.847,-0.925) (-1.829,-0.875) (-1.809,-0.825) (-1.786,-0.775) (-1.761,-0.725) (-1.733,-0.675) (-1.701,-0.625) (-1.666,-0.575) (-1.627,-0.525) (-1.584,-0.475) (-1.536,-0.425) (-1.482,-0.375) (-1.425,-0.328) (-1.375,-0.292) (-1.325,-0.258) (-1.300,-0.243)}
\EQ{(1.300,-0.216) (1.350,-0.244) (1.400,-0.274) (1.450,-0.308) (1.500,-0.345) (1.550,-0.385) (1.595,-0.425) (1.646,-0.475) (1.692,-0.525) (1.725,-0.565) (1.770,-0.625) (1.800,-0.669) (1.834,-0.725) (1.862,-0.775) (1.887,-0.825) (1.909,-0.875) (1.929,-0.925) (1.947,-0.975) (1.962,-1.025) (1.975,-1.075) (1.987,-1.125) (1.997,-1.175) (2.004,-1.225) (2.011,-1.275) (2.015,-1.325) (2.017,-1.375) (2.018,-1.425) (2.017,-1.475) (2.015,-1.525) (2.011,-1.575) (2.005,-1.625) (1.998,-1.675) (1.989,-1.725) (1.979,-1.775) (1.967,-1.825) (1.953,-1.875) (1.937,-1.925) (1.920,-1.975) (1.901,-2.025) (1.880,-2.075) (1.858,-2.125) (1.833,-2.175) (1.807,-2.225) (1.778,-2.275) (1.750,-2.321) (1.714,-2.375) (1.679,-2.425) (1.641,-2.475) (1.601,-2.525) (1.557,-2.575) (1.510,-2.625) (1.475,-2.661) (1.434,-2.700) (1.378,-2.750) (1.325,-2.794) (1.285,-2.825) (1.225,-2.869) (1.180,-2.900) (1.125,-2.935) (1.075,-2.965) (1.025,-2.993) (0.975,-3.020) (0.925,-3.045) (0.875,-3.068) (0.825,-3.089) (0.775,-3.109) (0.725,-3.128) (0.675,-3.145) (0.625,-3.161) (0.576,-3.175) (0.525,-3.188) (0.476,-3.200) (0.425,-3.211) (0.375,-3.220) (0.325,-3.229) (0.275,-3.236) (0.225,-3.241) (0.175,-3.246) (0.125,-3.249) (0.075,-3.252) (0.025,-3.253) (-0.050,-3.252) (-0.100,-3.251) (-0.150,-3.248) (-0.200,-3.244) (-0.250,-3.239) (-0.300,-3.232) (-0.350,-3.225) (-0.400,-3.216) (-0.450,-3.206) (-0.500,-3.195) (-0.550,-3.182) (-0.600,-3.168) (-0.650,-3.153) (-0.700,-3.137) (-0.750,-3.119) (-0.798,-3.100) (-0.850,-3.079) (-0.900,-3.056) (-0.950,-3.032) (-1.000,-3.007) (-1.050,-2.980) (-1.100,-2.950) (-1.150,-2.919) (-1.200,-2.886) (-1.250,-2.851) (-1.300,-2.814) (-1.348,-2.775) (-1.400,-2.731) (-1.450,-2.685) (-1.486,-2.650) (-1.525,-2.610) (-1.575,-2.555) (-1.621,-2.500) (-1.661,-2.450) (-1.697,-2.400) (-1.731,-2.350) (-1.763,-2.300) (-1.793,-2.250) (-1.820,-2.200) (-1.846,-2.150) (-1.869,-2.100) (-1.891,-2.050) (-1.911,-2.000) (-1.929,-1.950) (-1.945,-1.900) (-1.960,-1.850) (-1.973,-1.800) (-1.984,-1.750) (-1.994,-1.700) (-2.002,-1.650) (-2.008,-1.600) (-2.013,-1.550) (-2.017,-1.500) (-2.018,-1.450) (-2.018,-1.400) (-2.016,-1.350) (-2.013,-1.300) (-2.008,-1.250) (-2.001,-1.200) (-1.992,-1.150) (-1.982,-1.100) (-1.969,-1.050) (-1.955,-1.000) (-1.938,-0.950) (-1.919,-0.900) (-1.898,-0.850) (-1.875,-0.801) (-1.850,-0.753) (-1.825,-0.710) (-1.787,-0.650) (-1.752,-0.600) (-1.713,-0.550) (-1.675,-0.506) (-1.625,-0.454) (-1.575,-0.407) (-1.525,-0.364) (-1.475,-0.326) (-1.425,-0.291) (-1.375,-0.259) (-1.325,-0.230) (-1.300,-0.216)}
\EQ{(1.300,-0.191) (1.350,-0.215) (1.400,-0.241) (1.450,-0.270) (1.499,-0.300) (1.550,-0.335) (1.600,-0.372) (1.650,-0.414) (1.691,-0.450) (1.742,-0.500) (1.788,-0.550) (1.825,-0.594) (1.868,-0.650) (1.900,-0.696) (1.934,-0.750) (1.963,-0.800) (1.989,-0.850) (2.013,-0.900) (2.034,-0.950) (2.053,-1.000) (2.070,-1.050) (2.085,-1.100) (2.098,-1.150) (2.109,-1.200) (2.118,-1.250) (2.126,-1.300) (2.132,-1.350) (2.137,-1.400) (2.139,-1.450) (2.141,-1.500) (2.140,-1.550) (2.138,-1.600) (2.135,-1.650) (2.130,-1.700) (2.123,-1.750) (2.115,-1.800) (2.106,-1.850) (2.095,-1.900) (2.082,-1.950) (2.068,-2.000) (2.052,-2.050) (2.035,-2.100) (2.016,-2.150) (1.995,-2.200) (1.972,-2.250) (1.950,-2.296) (1.925,-2.344) (1.900,-2.389) (1.864,-2.450) (1.831,-2.500) (1.800,-2.546) (1.760,-2.600) (1.725,-2.645) (1.679,-2.700) (1.634,-2.750) (1.600,-2.786) (1.562,-2.825) (1.525,-2.860) (1.481,-2.900) (1.425,-2.947) (1.375,-2.987) (1.325,-3.025) (1.275,-3.060) (1.225,-3.093) (1.175,-3.124) (1.132,-3.150) (1.087,-3.175) (1.040,-3.200) (0.990,-3.225) (0.936,-3.250) (0.879,-3.275) (0.825,-3.297) (0.775,-3.315) (0.725,-3.333) (0.675,-3.349) (0.625,-3.364) (0.575,-3.377) (0.525,-3.390) (0.475,-3.401) (0.425,-3.411) (0.375,-3.420) (0.325,-3.427) (0.275,-3.434) (0.225,-3.439) (0.175,-3.444) (0.125,-3.447) (0.075,-3.449) (0.025,-3.450) (-0.037,-3.450) (-0.100,-3.448) (-0.150,-3.445) (-0.200,-3.442) (-0.250,-3.437) (-0.300,-3.431) (-0.350,-3.424) (-0.400,-3.415) (-0.450,-3.406) (-0.500,-3.395) (-0.550,-3.384) (-0.600,-3.371) (-0.650,-3.356) (-0.700,-3.341) (-0.748,-3.325) (-0.800,-3.306) (-0.850,-3.287) (-0.900,-3.266) (-0.950,-3.244) (-1.000,-3.220) (-1.050,-3.195) (-1.100,-3.168) (-1.150,-3.139) (-1.200,-3.109) (-1.250,-3.077) (-1.300,-3.043) (-1.350,-3.006) (-1.391,-2.975) (-1.450,-2.927) (-1.500,-2.883) (-1.550,-2.836) (-1.587,-2.800) (-1.625,-2.760) (-1.675,-2.705) (-1.721,-2.650) (-1.760,-2.600) (-1.797,-2.550) (-1.831,-2.500) (-1.864,-2.450) (-1.894,-2.400) (-1.922,-2.350) (-1.948,-2.300) (-1.972,-2.250) (-1.995,-2.200) (-2.016,-2.150) (-2.035,-2.100) (-2.052,-2.050) (-2.068,-2.000) (-2.082,-1.950) (-2.095,-1.900) (-2.106,-1.850) (-2.115,-1.800) (-2.123,-1.750) (-2.130,-1.700) (-2.135,-1.650) (-2.138,-1.600) (-2.140,-1.550) (-2.141,-1.500) (-2.139,-1.450) (-2.137,-1.400) (-2.132,-1.350) (-2.126,-1.300) (-2.118,-1.250) (-2.109,-1.200) (-2.098,-1.150) (-2.085,-1.100) (-2.070,-1.050) (-2.053,-1.000) (-2.034,-0.950) (-2.013,-0.900) (-1.989,-0.850) (-1.963,-0.800) (-1.934,-0.750) (-1.903,-0.700) (-1.868,-0.650) (-1.830,-0.600) (-1.788,-0.550) (-1.750,-0.508) (-1.700,-0.459) (-1.650,-0.414) (-1.603,-0.375) (-1.550,-0.335) (-1.500,-0.301) (-1.450,-0.270) (-1.400,-0.241) (-1.350,-0.215) (-1.300,-0.191)}
\EQ{(1.300,-0.165) (1.350,-0.186) (1.400,-0.208) (1.450,-0.233) (1.500,-0.259) (1.550,-0.288) (1.600,-0.319) (1.646,-0.350) (1.700,-0.390) (1.744,-0.425) (1.800,-0.474) (1.850,-0.522) (1.900,-0.575) (1.943,-0.625) (1.975,-0.665) (2.018,-0.725) (2.050,-0.773) (2.082,-0.825) (2.109,-0.875) (2.134,-0.925) (2.157,-0.975) (2.178,-1.025) (2.197,-1.075) (2.213,-1.125) (2.228,-1.175) (2.242,-1.225) (2.253,-1.275) (2.263,-1.325) (2.271,-1.375) (2.277,-1.425) (2.282,-1.475) (2.286,-1.525) (2.288,-1.575) (2.288,-1.625) (2.288,-1.675) (2.285,-1.725) (2.281,-1.775) (2.276,-1.825) (2.269,-1.875) (2.261,-1.925) (2.252,-1.975) (2.241,-2.025) (2.228,-2.075) (2.215,-2.125) (2.199,-2.175) (2.182,-2.225) (2.164,-2.275) (2.144,-2.325) (2.122,-2.375) (2.100,-2.423) (2.075,-2.473) (2.050,-2.520) (2.025,-2.565) (1.989,-2.625) (1.956,-2.675) (1.925,-2.720) (1.885,-2.775) (1.850,-2.820) (1.805,-2.875) (1.761,-2.925) (1.725,-2.964) (1.690,-3.000) (1.650,-3.039) (1.611,-3.075) (1.554,-3.125) (1.500,-3.170) (1.450,-3.209) (1.400,-3.245) (1.358,-3.275) (1.300,-3.313) (1.250,-3.344) (1.200,-3.373) (1.152,-3.400) (1.104,-3.425) (1.053,-3.450) (1.000,-3.475) (0.950,-3.497) (0.900,-3.517) (0.850,-3.536) (0.800,-3.554) (0.750,-3.571) (0.700,-3.587) (0.650,-3.601) (0.600,-3.614) (0.550,-3.626) (0.500,-3.637) (0.450,-3.647) (0.400,-3.656) (0.350,-3.664) (0.300,-3.670) (0.250,-3.676) (0.200,-3.680) (0.150,-3.684) (0.100,-3.687) (0.050,-3.688) (0.000,-3.689) (-0.050,-3.688) (-0.100,-3.687) (-0.150,-3.684) (-0.200,-3.680) (-0.250,-3.676) (-0.300,-3.670) (-0.350,-3.664) (-0.400,-3.656) (-0.450,-3.647) (-0.500,-3.637) (-0.550,-3.626) (-0.600,-3.614) (-0.650,-3.601) (-0.700,-3.587) (-0.750,-3.571) (-0.800,-3.554) (-0.850,-3.536) (-0.900,-3.517) (-0.950,-3.497) (-1.000,-3.475) (-1.050,-3.452) (-1.100,-3.427) (-1.150,-3.401) (-1.197,-3.375) (-1.250,-3.344) (-1.300,-3.313) (-1.350,-3.280) (-1.394,-3.250) (-1.450,-3.209) (-1.494,-3.175) (-1.550,-3.129) (-1.600,-3.085) (-1.638,-3.050) (-1.675,-3.015) (-1.714,-2.975) (-1.750,-2.937) (-1.800,-2.881) (-1.846,-2.825) (-1.885,-2.775) (-1.922,-2.725) (-1.956,-2.675) (-1.989,-2.625) (-2.019,-2.575) (-2.047,-2.525) (-2.074,-2.475) (-2.099,-2.425) (-2.122,-2.375) (-2.144,-2.325) (-2.164,-2.275) (-2.182,-2.225) (-2.199,-2.175) (-2.215,-2.125) (-2.228,-2.075) (-2.241,-2.025) (-2.252,-1.975) (-2.261,-1.925) (-2.269,-1.875) (-2.276,-1.825) (-2.281,-1.775) (-2.285,-1.725) (-2.288,-1.675) (-2.288,-1.625) (-2.288,-1.575) (-2.286,-1.525) (-2.282,-1.475) (-2.277,-1.425) (-2.271,-1.375) (-2.263,-1.325) (-2.253,-1.275) (-2.242,-1.225) (-2.228,-1.175) (-2.213,-1.125) (-2.197,-1.075) (-2.178,-1.025) (-2.157,-0.975) (-2.134,-0.925) (-2.109,-0.875) (-2.082,-0.825) (-2.051,-0.775) (-2.018,-0.725) (-1.982,-0.675) (-1.950,-0.634) (-1.900,-0.575) (-1.853,-0.525) (-1.801,-0.475) (-1.750,-0.430) (-1.700,-0.390) (-1.650,-0.353) (-1.600,-0.319) (-1.550,-0.288) (-1.500,-0.259) (-1.450,-0.233) (-1.400,-0.208) (-1.350,-0.186) (-1.300,-0.165)}
\EQ{(1.300,-0.141) (1.350,-0.158) (1.400,-0.177) (1.450,-0.197) (1.500,-0.220) (1.550,-0.243) (1.600,-0.269) (1.650,-0.297) (1.697,-0.325) (1.750,-0.359) (1.800,-0.394) (1.841,-0.425) (1.900,-0.473) (1.950,-0.517) (2.000,-0.566) (2.050,-0.619) (2.099,-0.675) (2.138,-0.725) (2.174,-0.775) (2.208,-0.825) (2.239,-0.875) (2.267,-0.925) (2.293,-0.975) (2.317,-1.025) (2.339,-1.075) (2.359,-1.125) (2.377,-1.175) (2.394,-1.225) (2.409,-1.275) (2.422,-1.325) (2.433,-1.375) (2.443,-1.425) (2.451,-1.475) (2.458,-1.525) (2.464,-1.575) (2.468,-1.625) (2.470,-1.675) (2.472,-1.725) (2.472,-1.775) (2.470,-1.825) (2.467,-1.875) (2.463,-1.925) (2.458,-1.975) (2.451,-2.025) (2.443,-2.075) (2.434,-2.125) (2.423,-2.175) (2.411,-2.225) (2.398,-2.275) (2.383,-2.325) (2.367,-2.375) (2.350,-2.423) (2.330,-2.475) (2.310,-2.525) (2.288,-2.575) (2.265,-2.625) (2.240,-2.675) (2.213,-2.725) (2.184,-2.775) (2.154,-2.825) (2.125,-2.871) (2.089,-2.925) (2.053,-2.975) (2.015,-3.025) (1.975,-3.075) (1.933,-3.125) (1.888,-3.175) (1.850,-3.215) (1.800,-3.265) (1.763,-3.300) (1.725,-3.335) (1.679,-3.375) (1.625,-3.420) (1.575,-3.459) (1.525,-3.496) (1.475,-3.531) (1.425,-3.565) (1.375,-3.597) (1.328,-3.625) (1.285,-3.650) (1.239,-3.675) (1.191,-3.700) (1.141,-3.725) (1.087,-3.750) (1.031,-3.775) (0.975,-3.798) (0.925,-3.817) (0.875,-3.835) (0.825,-3.852) (0.775,-3.868) (0.725,-3.883) (0.675,-3.897) (0.625,-3.909) (0.575,-3.921) (0.525,-3.932) (0.475,-3.941) (0.425,-3.950) (0.375,-3.958) (0.325,-3.964) (0.275,-3.970) (0.225,-3.975) (0.175,-3.978) (0.125,-3.981) (0.075,-3.983) (0.025,-3.984) (-0.050,-3.984) (-0.100,-3.982) (-0.150,-3.980) (-0.200,-3.977) (-0.250,-3.972) (-0.300,-3.967) (-0.350,-3.961) (-0.400,-3.954) (-0.450,-3.946) (-0.500,-3.937) (-0.550,-3.926) (-0.600,-3.915) (-0.650,-3.903) (-0.700,-3.890) (-0.750,-3.876) (-0.800,-3.860) (-0.850,-3.844) (-0.900,-3.826) (-0.950,-3.808) (-1.000,-3.788) (-1.050,-3.767) (-1.100,-3.744) (-1.150,-3.721) (-1.200,-3.696) (-1.250,-3.669) (-1.300,-3.641) (-1.350,-3.612) (-1.400,-3.581) (-1.448,-3.550) (-1.500,-3.514) (-1.550,-3.478) (-1.600,-3.440) (-1.650,-3.400) (-1.700,-3.357) (-1.750,-3.312) (-1.790,-3.275) (-1.840,-3.225) (-1.875,-3.189) (-1.910,-3.150) (-1.950,-3.105) (-1.995,-3.050) (-2.034,-3.000) (-2.071,-2.950) (-2.100,-2.908) (-2.139,-2.850) (-2.170,-2.800) (-2.199,-2.750) (-2.225,-2.703) (-2.250,-2.654) (-2.275,-2.603) (-2.299,-2.550) (-2.320,-2.500) (-2.340,-2.450) (-2.358,-2.400) (-2.375,-2.351) (-2.391,-2.300) (-2.405,-2.250) (-2.417,-2.200) (-2.429,-2.150) (-2.439,-2.100) (-2.447,-2.050) (-2.455,-2.000) (-2.461,-1.950) (-2.466,-1.900) (-2.469,-1.850) (-2.471,-1.800) (-2.472,-1.750) (-2.471,-1.700) (-2.469,-1.650) (-2.466,-1.600) (-2.461,-1.550) (-2.455,-1.500) (-2.447,-1.450) (-2.438,-1.400) (-2.427,-1.350) (-2.415,-1.300) (-2.401,-1.250) (-2.386,-1.200) (-2.369,-1.150) (-2.350,-1.101) (-2.329,-1.050) (-2.306,-1.000) (-2.281,-0.950) (-2.253,-0.900) (-2.225,-0.852) (-2.192,-0.800) (-2.157,-0.750) (-2.125,-0.708) (-2.078,-0.650) (-2.033,-0.600) (-1.984,-0.550) (-1.931,-0.500) (-1.875,-0.452) (-1.825,-0.413) (-1.775,-0.376) (-1.725,-0.343) (-1.675,-0.312) (-1.625,-0.283) (-1.575,-0.256) (-1.525,-0.231) (-1.475,-0.208) (-1.425,-0.187) (-1.375,-0.167) (-1.327,-0.150) (-1.300,-0.141)}
\EQ{(1.300,-0.117) (1.350,-0.131) (1.400,-0.146) (1.450,-0.163) (1.500,-0.181) (1.549,-0.200) (1.600,-0.221) (1.650,-0.244) (1.700,-0.268) (1.750,-0.293) (1.800,-0.321) (1.849,-0.350) (1.900,-0.382) (1.950,-0.416) (1.997,-0.450) (2.050,-0.492) (2.090,-0.525) (2.145,-0.575) (2.196,-0.625) (2.243,-0.675) (2.287,-0.725) (2.325,-0.772) (2.365,-0.825) (2.400,-0.875) (2.432,-0.925) (2.462,-0.975) (2.490,-1.025) (2.516,-1.075) (2.540,-1.125) (2.562,-1.175) (2.582,-1.225) (2.600,-1.275) (2.617,-1.325) (2.632,-1.375) (2.646,-1.425) (2.658,-1.475) (2.669,-1.525) (2.679,-1.575) (2.687,-1.625) (2.693,-1.675) (2.699,-1.725) (2.703,-1.775) (2.706,-1.825) (2.707,-1.875) (2.708,-1.925) (2.707,-1.975) (2.705,-2.025) (2.701,-2.075) (2.697,-2.125) (2.691,-2.175) (2.684,-2.225) (2.676,-2.275) (2.667,-2.325) (2.656,-2.375) (2.644,-2.425) (2.631,-2.475) (2.617,-2.525) (2.602,-2.575) (2.585,-2.625) (2.567,-2.675) (2.547,-2.725) (2.527,-2.775) (2.504,-2.825) (2.481,-2.875) (2.456,-2.925) (2.430,-2.975) (2.402,-3.025) (2.375,-3.070) (2.341,-3.125) (2.308,-3.175) (2.275,-3.223) (2.237,-3.275) (2.200,-3.323) (2.158,-3.375) (2.125,-3.414) (2.075,-3.470) (2.025,-3.523) (1.975,-3.574) (1.925,-3.621) (1.875,-3.667) (1.836,-3.700) (1.776,-3.750) (1.725,-3.790) (1.677,-3.825) (1.625,-3.862) (1.575,-3.896) (1.530,-3.925) (1.475,-3.959) (1.425,-3.988) (1.375,-4.016) (1.325,-4.043) (1.275,-4.068) (1.225,-4.092) (1.175,-4.115) (1.125,-4.137) (1.075,-4.158) (1.025,-4.177) (0.975,-4.196) (0.925,-4.213) (0.875,-4.229) (0.825,-4.245) (0.775,-4.259) (0.725,-4.272) (0.675,-4.285) (0.625,-4.296) (0.575,-4.307) (0.525,-4.317) (0.475,-4.325) (0.425,-4.333) (0.375,-4.340) (0.325,-4.346) (0.275,-4.351) (0.225,-4.356) (0.175,-4.359) (0.125,-4.362) (0.075,-4.363) (0.025,-4.364) (-0.050,-4.364) (-0.100,-4.363) (-0.150,-4.361) (-0.200,-4.358) (-0.250,-4.354) (-0.300,-4.349) (-0.350,-4.343) (-0.400,-4.337) (-0.450,-4.329) (-0.500,-4.321) (-0.550,-4.312) (-0.600,-4.302) (-0.650,-4.291) (-0.700,-4.279) (-0.750,-4.266) (-0.800,-4.252) (-0.850,-4.237) (-0.900,-4.221) (-0.950,-4.204) (-1.000,-4.187) (-1.050,-4.168) (-1.100,-4.147) (-1.150,-4.126) (-1.200,-4.104) (-1.250,-4.080) (-1.300,-4.056) (-1.350,-4.030) (-1.400,-4.002) (-1.448,-3.975) (-1.500,-3.944) (-1.550,-3.912) (-1.600,-3.879) (-1.642,-3.850) (-1.700,-3.808) (-1.744,-3.775) (-1.800,-3.730) (-1.850,-3.688) (-1.894,-3.650) (-1.948,-3.600) (-1.999,-3.550) (-2.047,-3.500) (-2.093,-3.450) (-2.137,-3.400) (-2.175,-3.354) (-2.218,-3.300) (-2.250,-3.257) (-2.291,-3.200) (-2.325,-3.150) (-2.357,-3.100) (-2.387,-3.050) (-2.416,-3.000) (-2.443,-2.950) (-2.469,-2.900) (-2.493,-2.850) (-2.516,-2.800) (-2.537,-2.750) (-2.557,-2.700) (-2.575,-2.653) (-2.593,-2.600) (-2.610,-2.550) (-2.624,-2.500) (-2.638,-2.450) (-2.650,-2.400) (-2.662,-2.350) (-2.672,-2.300) (-2.680,-2.250) (-2.688,-2.200) (-2.694,-2.150) (-2.699,-2.100) (-2.703,-2.050) (-2.706,-2.000) (-2.707,-1.950) (-2.708,-1.900) (-2.707,-1.850) (-2.704,-1.800) (-2.701,-1.750) (-2.696,-1.700) (-2.690,-1.650) (-2.683,-1.600) (-2.674,-1.550) (-2.664,-1.500) (-2.652,-1.450) (-2.639,-1.400) (-2.625,-1.351) (-2.609,-1.300) (-2.591,-1.250) (-2.572,-1.200) (-2.551,-1.150) (-2.528,-1.100) (-2.503,-1.050) (-2.476,-1.000) (-2.450,-0.954) (-2.416,-0.900) (-2.383,-0.850) (-2.350,-0.805) (-2.308,-0.750) (-2.275,-0.711) (-2.225,-0.655) (-2.175,-0.604) (-2.125,-0.557) (-2.075,-0.513) (-2.029,-0.475) (-1.975,-0.434) (-1.927,-0.400) (-1.875,-0.366) (-1.825,-0.335) (-1.775,-0.307) (-1.725,-0.280) (-1.675,-0.255) (-1.625,-0.232) (-1.575,-0.211) (-1.525,-0.191) (-1.475,-0.172) (-1.425,-0.155) (-1.375,-0.139) (-1.325,-0.124) (-1.300,-0.117)}
\EQ{(1.300,-0.093) (1.350,-0.104) (1.400,-0.116) (1.450,-0.130) (1.500,-0.144) (1.550,-0.159) (1.600,-0.175) (1.650,-0.193) (1.700,-0.211) (1.750,-0.231) (1.800,-0.252) (1.850,-0.275) (1.900,-0.299) (1.950,-0.324) (1.997,-0.350) (2.041,-0.375) (2.100,-0.412) (2.150,-0.444) (2.194,-0.475) (2.250,-0.517) (2.292,-0.550) (2.350,-0.599) (2.400,-0.645) (2.450,-0.693) (2.500,-0.746) (2.548,-0.800) (2.589,-0.850) (2.625,-0.896) (2.664,-0.950) (2.697,-1.000) (2.725,-1.044) (2.759,-1.100) (2.786,-1.150) (2.812,-1.200) (2.836,-1.250) (2.858,-1.300) (2.879,-1.350) (2.898,-1.400) (2.916,-1.450) (2.932,-1.500) (2.947,-1.550) (2.960,-1.600) (2.972,-1.650) (2.983,-1.700) (2.993,-1.750) (3.001,-1.800) (3.008,-1.850) (3.014,-1.900) (3.019,-1.950) (3.023,-2.000) (3.025,-2.050) (3.027,-2.100) (3.027,-2.150) (3.026,-2.200) (3.025,-2.250) (3.022,-2.300) (3.018,-2.350) (3.013,-2.400) (3.006,-2.450) (2.999,-2.500) (2.991,-2.550) (2.981,-2.600) (2.971,-2.650) (2.959,-2.700) (2.947,-2.750) (2.933,-2.800) (2.918,-2.850) (2.902,-2.900) (2.885,-2.950) (2.866,-3.000) (2.847,-3.050) (2.826,-3.100) (2.804,-3.150) (2.781,-3.200) (2.756,-3.250) (2.730,-3.300) (2.703,-3.350) (2.675,-3.399) (2.645,-3.450) (2.613,-3.500) (2.580,-3.550) (2.550,-3.594) (2.509,-3.650) (2.475,-3.696) (2.432,-3.750) (2.400,-3.789) (2.350,-3.846) (2.301,-3.900) (2.253,-3.950) (2.203,-4.000) (2.151,-4.050) (2.100,-4.096) (2.050,-4.139) (2.007,-4.175) (1.950,-4.220) (1.900,-4.258) (1.850,-4.295) (1.807,-4.325) (1.750,-4.363) (1.700,-4.395) (1.651,-4.425) (1.600,-4.455) (1.550,-4.483) (1.500,-4.510) (1.450,-4.536) (1.400,-4.561) (1.350,-4.585) (1.300,-4.607) (1.250,-4.629) (1.200,-4.649) (1.150,-4.669) (1.100,-4.687) (1.050,-4.705) (1.000,-4.722) (0.950,-4.737) (0.900,-4.752) (0.850,-4.766) (0.800,-4.779) (0.750,-4.792) (0.700,-4.803) (0.650,-4.814) (0.600,-4.824) (0.550,-4.833) (0.500,-4.841) (0.450,-4.848) (0.400,-4.855) (0.350,-4.861) (0.300,-4.866) (0.250,-4.870) (0.200,-4.873) (0.150,-4.876) (0.100,-4.878) (0.050,-4.879) (0.000,-4.880) (-0.050,-4.879) (-0.100,-4.878) (-0.150,-4.876) (-0.200,-4.873) (-0.250,-4.870) (-0.300,-4.866) (-0.350,-4.861) (-0.400,-4.855) (-0.450,-4.848) (-0.500,-4.841) (-0.550,-4.833) (-0.600,-4.824) (-0.650,-4.814) (-0.700,-4.803) (-0.750,-4.792) (-0.800,-4.779) (-0.850,-4.766) (-0.900,-4.752) (-0.950,-4.737) (-1.000,-4.722) (-1.050,-4.705) (-1.100,-4.687) (-1.150,-4.669) (-1.198,-4.650) (-1.250,-4.629) (-1.300,-4.607) (-1.350,-4.585) (-1.400,-4.561) (-1.450,-4.536) (-1.500,-4.510) (-1.550,-4.483) (-1.600,-4.455) (-1.650,-4.426) (-1.700,-4.395) (-1.750,-4.363) (-1.800,-4.330) (-1.843,-4.300) (-1.900,-4.258) (-1.944,-4.225) (-2.000,-4.181) (-2.050,-4.139) (-2.095,-4.100) (-2.150,-4.051) (-2.200,-4.003) (-2.250,-3.954) (-2.300,-3.901) (-2.347,-3.850) (-2.390,-3.800) (-2.425,-3.758) (-2.472,-3.700) (-2.509,-3.650) (-2.546,-3.600) (-2.575,-3.558) (-2.613,-3.500) (-2.645,-3.450) (-2.675,-3.400) (-2.700,-3.356) (-2.725,-3.310) (-2.750,-3.262) (-2.775,-3.212) (-2.800,-3.159) (-2.825,-3.102) (-2.847,-3.050) (-2.866,-3.000) (-2.885,-2.950) (-2.902,-2.900) (-2.918,-2.850) (-2.933,-2.800) (-2.947,-2.750) (-2.959,-2.700) (-2.971,-2.650) (-2.981,-2.600) (-2.991,-2.550) (-2.999,-2.500) (-3.006,-2.450) (-3.013,-2.400) (-3.018,-2.350) (-3.022,-2.300) (-3.025,-2.250) (-3.026,-2.200) (-3.027,-2.150) (-3.027,-2.100) (-3.025,-2.050) (-3.023,-2.000) (-3.019,-1.950) (-3.014,-1.900) (-3.008,-1.850) (-3.001,-1.800) (-2.993,-1.750) (-2.983,-1.700) (-2.972,-1.650) (-2.960,-1.600) (-2.947,-1.550) (-2.932,-1.500) (-2.916,-1.450) (-2.898,-1.400) (-2.879,-1.350) (-2.858,-1.300) (-2.836,-1.250) (-2.812,-1.200) (-2.786,-1.150) (-2.759,-1.100) (-2.729,-1.050) (-2.700,-1.004) (-2.664,-0.950) (-2.628,-0.900) (-2.589,-0.850) (-2.550,-0.802) (-2.504,-0.750) (-2.457,-0.700) (-2.406,-0.650) (-2.351,-0.600) (-2.300,-0.557) (-2.261,-0.525) (-2.200,-0.479) (-2.158,-0.450) (-2.100,-0.412) (-2.050,-0.381) (-2.000,-0.352) (-1.951,-0.325) (-1.902,-0.300) (-1.850,-0.275) (-1.800,-0.252) (-1.750,-0.231) (-1.700,-0.211) (-1.650,-0.193) (-1.600,-0.175) (-1.550,-0.159) (-1.500,-0.144) (-1.450,-0.130) (-1.400,-0.116) (-1.350,-0.104) (-1.300,-0.093)}
\EQ{(1.300,-0.070) (1.350,-0.078) (1.400,-0.087) (1.450,-0.097) (1.500,-0.107) (1.550,-0.118) (1.600,-0.130) (1.650,-0.143) (1.700,-0.157) (1.750,-0.171) (1.800,-0.187) (1.850,-0.203) (1.900,-0.220) (1.950,-0.239) (2.000,-0.258) (2.050,-0.279) (2.098,-0.300) (2.150,-0.324) (2.200,-0.348) (2.250,-0.374) (2.298,-0.400) (2.342,-0.425) (2.400,-0.460) (2.450,-0.492) (2.500,-0.525) (2.550,-0.561) (2.600,-0.598) (2.650,-0.638) (2.694,-0.675) (2.750,-0.724) (2.800,-0.772) (2.850,-0.822) (2.899,-0.875) (2.943,-0.925) (2.985,-0.975) (3.024,-1.025) (3.060,-1.075) (3.095,-1.125) (3.125,-1.171) (3.159,-1.225) (3.188,-1.275) (3.216,-1.325) (3.241,-1.375) (3.266,-1.425) (3.289,-1.475) (3.310,-1.525) (3.330,-1.575) (3.349,-1.625) (3.367,-1.675) (3.383,-1.725) (3.398,-1.775) (3.412,-1.825) (3.425,-1.875) (3.437,-1.925) (3.447,-1.975) (3.456,-2.025) (3.465,-2.075) (3.472,-2.125) (3.479,-2.175) (3.484,-2.225) (3.488,-2.275) (3.491,-2.325) (3.494,-2.375) (3.495,-2.425) (3.496,-2.475) (3.495,-2.525) (3.494,-2.575) (3.491,-2.625) (3.488,-2.675) (3.483,-2.725) (3.478,-2.775) (3.472,-2.825) (3.465,-2.875) (3.457,-2.925) (3.448,-2.975) (3.438,-3.025) (3.427,-3.075) (3.415,-3.125) (3.403,-3.175) (3.389,-3.225) (3.374,-3.275) (3.359,-3.325) (3.342,-3.375) (3.325,-3.423) (3.306,-3.475) (3.286,-3.525) (3.265,-3.575) (3.243,-3.625) (3.220,-3.675) (3.196,-3.725) (3.171,-3.775) (3.145,-3.825) (3.118,-3.875) (3.089,-3.925) (3.059,-3.975) (3.028,-4.025) (3.000,-4.068) (2.962,-4.125) (2.927,-4.175) (2.890,-4.225) (2.852,-4.275) (2.812,-4.325) (2.775,-4.370) (2.728,-4.425) (2.683,-4.475) (2.637,-4.525) (2.600,-4.563) (2.563,-4.600) (2.525,-4.637) (2.485,-4.675) (2.430,-4.725) (2.375,-4.773) (2.325,-4.815) (2.281,-4.850) (2.225,-4.894) (2.183,-4.925) (2.125,-4.967) (2.078,-5.000) (2.025,-5.035) (1.975,-5.067) (1.925,-5.098) (1.881,-5.125) (1.825,-5.157) (1.775,-5.185) (1.725,-5.212) (1.675,-5.238) (1.625,-5.262) (1.575,-5.286) (1.525,-5.309) (1.475,-5.331) (1.425,-5.352) (1.375,-5.373) (1.325,-5.392) (1.275,-5.411) (1.225,-5.428) (1.175,-5.445) (1.125,-5.462) (1.075,-5.477) (1.025,-5.492) (0.975,-5.505) (0.925,-5.519) (0.875,-5.531) (0.825,-5.543) (0.775,-5.553) (0.725,-5.564) (0.675,-5.573) (0.625,-5.582) (0.575,-5.590) (0.525,-5.597) (0.475,-5.604) (0.425,-5.610) (0.375,-5.616) (0.325,-5.620) (0.275,-5.624) (0.225,-5.628) (0.175,-5.630) (0.125,-5.632) (0.075,-5.634) (0.025,-5.634) (-0.050,-5.634) (-0.100,-5.633) (-0.150,-5.631) (-0.200,-5.629) (-0.250,-5.626) (-0.300,-5.622) (-0.350,-5.618) (-0.400,-5.613) (-0.450,-5.607) (-0.500,-5.601) (-0.550,-5.594) (-0.600,-5.586) (-0.650,-5.578) (-0.700,-5.569) (-0.750,-5.559) (-0.800,-5.548) (-0.850,-5.537) (-0.899,-5.525) (-0.950,-5.512) (-1.000,-5.499) (-1.050,-5.484) (-1.100,-5.469) (-1.150,-5.453) (-1.200,-5.437) (-1.250,-5.420) (-1.300,-5.401) (-1.350,-5.382) (-1.400,-5.363) (-1.450,-5.342) (-1.500,-5.320) (-1.550,-5.298) (-1.599,-5.275) (-1.650,-5.250) (-1.700,-5.225) (-1.747,-5.200) (-1.793,-5.175) (-1.838,-5.150) (-1.900,-5.114) (-1.950,-5.083) (-2.000,-5.051) (-2.050,-5.019) (-2.100,-4.985) (-2.149,-4.950) (-2.200,-4.913) (-2.249,-4.875) (-2.300,-4.835) (-2.343,-4.800) (-2.400,-4.751) (-2.450,-4.707) (-2.500,-4.661) (-2.538,-4.625) (-2.575,-4.588) (-2.613,-4.550) (-2.650,-4.511) (-2.700,-4.457) (-2.750,-4.400) (-2.792,-4.350) (-2.825,-4.309) (-2.871,-4.250) (-2.908,-4.200) (-2.944,-4.150) (-2.975,-4.106) (-3.012,-4.050) (-3.044,-4.000) (-3.074,-3.950) (-3.100,-3.906) (-3.125,-3.862) (-3.150,-3.816) (-3.175,-3.768) (-3.200,-3.718) (-3.225,-3.665) (-3.250,-3.610) (-3.275,-3.551) (-3.296,-3.500) (-3.315,-3.450) (-3.333,-3.400) (-3.350,-3.351) (-3.366,-3.300) (-3.382,-3.250) (-3.396,-3.200) (-3.409,-3.150) (-3.421,-3.100) (-3.433,-3.050) (-3.443,-3.000) (-3.453,-2.950) (-3.461,-2.900) (-3.469,-2.850) (-3.475,-2.800) (-3.481,-2.750) (-3.486,-2.700) (-3.490,-2.650) (-3.492,-2.600) (-3.494,-2.550) (-3.495,-2.500) (-3.496,-2.450) (-3.495,-2.400) (-3.493,-2.350) (-3.490,-2.300) (-3.486,-2.250) (-3.481,-2.200) (-3.475,-2.150) (-3.469,-2.100) (-3.461,-2.050) (-3.452,-2.000) (-3.442,-1.950) (-3.431,-1.900) (-3.419,-1.850) (-3.405,-1.800) (-3.391,-1.750) (-3.375,-1.700) (-3.358,-1.650) (-3.340,-1.600) (-3.321,-1.550) (-3.300,-1.501) (-3.277,-1.450) (-3.254,-1.400) (-3.229,-1.350) (-3.202,-1.300) (-3.175,-1.252) (-3.144,-1.200) (-3.112,-1.150) (-3.078,-1.100) (-3.042,-1.050) (-3.004,-1.000) (-2.964,-0.950) (-2.925,-0.904) (-2.876,-0.850) (-2.828,-0.800) (-2.777,-0.750) (-2.725,-0.702) (-2.675,-0.659) (-2.634,-0.625) (-2.575,-0.579) (-2.525,-0.543) (-2.475,-0.508) (-2.425,-0.475) (-2.375,-0.444) (-2.325,-0.415) (-2.275,-0.387) (-2.225,-0.361) (-2.175,-0.336) (-2.125,-0.312) (-2.075,-0.290) (-2.025,-0.269) (-1.975,-0.248) (-1.925,-0.230) (-1.875,-0.212) (-1.825,-0.195) (-1.775,-0.179) (-1.725,-0.164) (-1.676,-0.150) (-1.625,-0.137) (-1.575,-0.124) (-1.525,-0.113) (-1.475,-0.102) (-1.425,-0.092) (-1.375,-0.082) (-1.325,-0.074) (-1.300,-0.070)}
\EQ{(1.300,-0.046) (1.350,-0.052) (1.400,-0.058) (1.450,-0.064) (1.500,-0.071) (1.550,-0.079) (1.600,-0.086) (1.650,-0.095) (1.700,-0.104) (1.750,-0.113) (1.800,-0.123) (1.850,-0.134) (1.900,-0.145) (1.950,-0.157) (2.000,-0.170) (2.050,-0.183) (2.100,-0.197) (2.150,-0.212) (2.200,-0.227) (2.250,-0.243) (2.300,-0.260) (2.350,-0.278) (2.400,-0.297) (2.450,-0.317) (2.500,-0.337) (2.550,-0.359) (2.600,-0.381) (2.650,-0.405) (2.700,-0.429) (2.750,-0.455) (2.800,-0.482) (2.850,-0.510) (2.900,-0.539) (2.950,-0.570) (2.998,-0.600) (3.050,-0.635) (3.100,-0.670) (3.141,-0.700) (3.200,-0.745) (3.250,-0.785) (3.298,-0.825) (3.350,-0.871) (3.400,-0.917) (3.450,-0.966) (3.500,-1.017) (3.550,-1.071) (3.598,-1.125) (3.640,-1.175) (3.675,-1.219) (3.718,-1.275) (3.750,-1.319) (3.789,-1.375) (3.822,-1.425) (3.850,-1.470) (3.883,-1.525) (3.912,-1.575) (3.939,-1.625) (3.965,-1.675) (3.990,-1.725) (4.014,-1.775) (4.036,-1.825) (4.057,-1.875) (4.077,-1.925) (4.096,-1.975) (4.114,-2.025) (4.131,-2.075) (4.147,-2.125) (4.162,-2.175) (4.176,-2.225) (4.189,-2.275) (4.201,-2.325) (4.213,-2.375) (4.223,-2.425) (4.232,-2.475) (4.241,-2.525) (4.249,-2.575) (4.256,-2.625) (4.262,-2.675) (4.267,-2.725) (4.271,-2.775) (4.275,-2.825) (4.278,-2.875) (4.280,-2.925) (4.281,-2.975) (4.281,-3.025) (4.281,-3.075) (4.280,-3.125) (4.278,-3.175) (4.275,-3.225) (4.272,-3.275) (4.268,-3.325) (4.263,-3.375) (4.257,-3.425) (4.250,-3.475) (4.243,-3.525) (4.235,-3.575) (4.226,-3.625) (4.217,-3.675) (4.206,-3.725) (4.195,-3.775) (4.184,-3.825) (4.171,-3.875) (4.157,-3.925) (4.143,-3.975) (4.128,-4.025) (4.112,-4.075) (4.096,-4.125) (4.078,-4.175) (4.060,-4.225) (4.041,-4.275) (4.021,-4.325) (4.000,-4.375) (3.979,-4.425) (3.956,-4.475) (3.933,-4.525) (3.908,-4.575) (3.883,-4.625) (3.857,-4.675) (3.830,-4.725) (3.802,-4.775) (3.775,-4.821) (3.743,-4.875) (3.711,-4.925) (3.679,-4.975) (3.650,-5.019) (3.611,-5.075) (3.576,-5.125) (3.539,-5.175) (3.501,-5.225) (3.462,-5.275) (3.425,-5.321) (3.380,-5.375) (3.336,-5.425) (3.300,-5.466) (3.250,-5.520) (3.200,-5.573) (3.150,-5.623) (3.100,-5.672) (3.050,-5.720) (3.000,-5.766) (2.961,-5.800) (2.904,-5.850) (2.850,-5.895) (2.800,-5.935) (2.750,-5.974) (2.700,-6.013) (2.650,-6.049) (2.600,-6.085) (2.550,-6.120) (2.505,-6.150) (2.450,-6.186) (2.400,-6.218) (2.350,-6.249) (2.306,-6.275) (2.250,-6.308) (2.200,-6.336) (2.150,-6.363) (2.100,-6.390) (2.050,-6.415) (2.000,-6.440) (1.950,-6.464) (1.900,-6.488) (1.850,-6.510) (1.800,-6.532) (1.750,-6.553) (1.700,-6.574) (1.650,-6.593) (1.600,-6.612) (1.550,-6.631) (1.500,-6.648) (1.450,-6.665) (1.400,-6.682) (1.350,-6.698) (1.300,-6.713) (1.250,-6.727) (1.200,-6.741) (1.150,-6.754) (1.100,-6.767) (1.050,-6.779) (1.000,-6.791) (0.950,-6.801) (0.900,-6.812) (0.850,-6.821) (0.800,-6.831) (0.750,-6.839) (0.700,-6.847) (0.650,-6.854) (0.600,-6.861) (0.550,-6.868) (0.500,-6.873) (0.450,-6.879) (0.400,-6.883) (0.350,-6.887) (0.300,-6.891) (0.250,-6.894) (0.200,-6.896) (0.150,-6.898) (0.100,-6.900) (0.050,-6.900) (0.000,-6.901) (-0.050,-6.900) (-0.100,-6.900) (-0.150,-6.898) (-0.200,-6.896) (-0.250,-6.894) (-0.300,-6.891) (-0.350,-6.887) (-0.400,-6.883) (-0.450,-6.879) (-0.500,-6.873) (-0.550,-6.868) (-0.600,-6.861) (-0.650,-6.854) (-0.700,-6.847) (-0.750,-6.839) (-0.800,-6.831) (-0.850,-6.821) (-0.900,-6.812) (-0.950,-6.801) (-1.000,-6.791) (-1.050,-6.779) (-1.100,-6.767) (-1.150,-6.754) (-1.200,-6.741) (-1.250,-6.727) (-1.300,-6.713) (-1.350,-6.698) (-1.400,-6.682) (-1.450,-6.665) (-1.500,-6.648) (-1.550,-6.631) (-1.600,-6.612) (-1.650,-6.593) (-1.697,-6.575) (-1.750,-6.553) (-1.800,-6.532) (-1.850,-6.510) (-1.900,-6.488) (-1.950,-6.464) (-2.000,-6.440) (-2.050,-6.415) (-2.100,-6.390) (-2.150,-6.363) (-2.200,-6.336) (-2.250,-6.308) (-2.300,-6.279) (-2.348,-6.250) (-2.400,-6.218) (-2.450,-6.186) (-2.500,-6.154) (-2.543,-6.125) (-2.600,-6.085) (-2.649,-6.050) (-2.700,-6.013) (-2.749,-5.975) (-2.800,-5.935) (-2.844,-5.900) (-2.900,-5.853) (-2.950,-5.810) (-2.990,-5.775) (-3.044,-5.725) (-3.097,-5.675) (-3.148,-5.625) (-3.198,-5.575) (-3.246,-5.525) (-3.292,-5.475) (-3.336,-5.425) (-3.375,-5.380) (-3.421,-5.325) (-3.462,-5.275) (-3.500,-5.227) (-3.539,-5.175) (-3.575,-5.126) (-3.611,-5.075) (-3.646,-5.025) (-3.675,-4.981) (-3.711,-4.925) (-3.743,-4.875) (-3.773,-4.825) (-3.800,-4.778) (-3.825,-4.734) (-3.850,-4.688) (-3.875,-4.641) (-3.900,-4.592) (-3.925,-4.541) (-3.950,-4.488) (-3.975,-4.433) (-4.000,-4.376) (-4.021,-4.325) (-4.041,-4.275) (-4.060,-4.225) (-4.078,-4.175) (-4.096,-4.125) (-4.112,-4.075) (-4.128,-4.025) (-4.143,-3.975) (-4.157,-3.925) (-4.171,-3.875) (-4.184,-3.825) (-4.195,-3.775) (-4.206,-3.725) (-4.217,-3.675) (-4.226,-3.625) (-4.235,-3.575) (-4.243,-3.525) (-4.250,-3.475) (-4.257,-3.425) (-4.263,-3.375) (-4.268,-3.325) (-4.272,-3.275) (-4.275,-3.225) (-4.278,-3.175) (-4.280,-3.125) (-4.281,-3.075) (-4.281,-3.025) (-4.281,-2.975) (-4.280,-2.925) (-4.278,-2.875) (-4.275,-2.825) (-4.271,-2.775) (-4.267,-2.725) (-4.262,-2.675) (-4.256,-2.625) (-4.249,-2.575) (-4.241,-2.525) (-4.232,-2.475) (-4.223,-2.425) (-4.213,-2.375) (-4.201,-2.325) (-4.189,-2.275) (-4.176,-2.225) (-4.162,-2.175) (-4.147,-2.125) (-4.131,-2.075) (-4.114,-2.025) (-4.096,-1.975) (-4.077,-1.925) (-4.057,-1.875) (-4.036,-1.825) (-4.014,-1.775) (-3.990,-1.725) (-3.965,-1.675) (-3.939,-1.625) (-3.912,-1.575) (-3.883,-1.525) (-3.853,-1.475) (-3.825,-1.430) (-3.789,-1.375) (-3.754,-1.325) (-3.718,-1.275) (-3.680,-1.225) (-3.640,-1.175) (-3.600,-1.127) (-3.554,-1.075) (-3.508,-1.025) (-3.459,-0.975) (-3.408,-0.925) (-3.355,-0.875) (-3.300,-0.827) (-3.250,-0.785) (-3.207,-0.750) (-3.150,-0.707) (-3.107,-0.675) (-3.050,-0.635) (-3.000,-0.602) (-2.950,-0.570) (-2.900,-0.539) (-2.850,-0.510) (-2.800,-0.482) (-2.750,-0.455) (-2.700,-0.429) (-2.650,-0.405) (-2.600,-0.381) (-2.550,-0.359) (-2.500,-0.337) (-2.450,-0.317) (-2.400,-0.297) (-2.350,-0.278) (-2.300,-0.260) (-2.250,-0.243) (-2.200,-0.227) (-2.150,-0.212) (-2.100,-0.197) (-2.050,-0.183) (-2.000,-0.170) (-1.950,-0.157) (-1.900,-0.145) (-1.850,-0.134) (-1.800,-0.123) (-1.750,-0.113) (-1.700,-0.104) (-1.650,-0.095) (-1.600,-0.086) (-1.550,-0.079) (-1.500,-0.071) (-1.450,-0.064) (-1.400,-0.058) (-1.350,-0.052) (-1.300,-0.046)}
\EQ{(1.300,-0.023) (1.350,-0.026) (1.400,-0.029) (1.450,-0.032) (1.500,-0.035) (1.550,-0.039) (1.600,-0.043) (1.650,-0.047) (1.700,-0.052) (1.750,-0.056) (1.800,-0.061) (1.850,-0.067) (1.900,-0.072) (1.950,-0.078) (2.000,-0.084) (2.050,-0.091) (2.100,-0.098) (2.150,-0.105) (2.200,-0.112) (2.250,-0.120) (2.300,-0.128) (2.350,-0.137) (2.400,-0.146) (2.450,-0.155) (2.500,-0.165) (2.550,-0.175) (2.600,-0.186) (2.650,-0.197) (2.700,-0.209) (2.750,-0.220) (2.800,-0.233) (2.850,-0.246) (2.900,-0.259) (2.950,-0.273) (3.000,-0.287) (3.050,-0.302) (3.100,-0.318) (3.150,-0.334) (3.199,-0.350) (3.250,-0.367) (3.300,-0.385) (3.350,-0.403) (3.400,-0.422) (3.450,-0.442) (3.500,-0.462) (3.550,-0.483) (3.600,-0.504) (3.646,-0.525) (3.700,-0.550) (3.750,-0.573) (3.800,-0.598) (3.850,-0.623) (3.900,-0.649) (3.949,-0.675) (3.994,-0.700) (4.038,-0.725) (4.100,-0.761) (4.150,-0.792) (4.200,-0.823) (4.250,-0.856) (4.300,-0.889) (4.350,-0.923) (4.400,-0.959) (4.450,-0.996) (4.500,-1.034) (4.550,-1.073) (4.600,-1.113) (4.644,-1.150) (4.700,-1.198) (4.750,-1.243) (4.800,-1.289) (4.838,-1.325) (4.875,-1.362) (4.913,-1.400) (4.950,-1.438) (4.985,-1.475) (5.025,-1.519) (5.074,-1.575) (5.117,-1.625) (5.150,-1.665) (5.198,-1.725) (5.237,-1.775) (5.274,-1.825) (5.310,-1.875) (5.345,-1.925) (5.375,-1.970) (5.411,-2.025) (5.443,-2.075) (5.473,-2.125) (5.500,-2.171) (5.525,-2.214) (5.550,-2.259) (5.575,-2.306) (5.598,-2.350) (5.623,-2.400) (5.648,-2.450) (5.671,-2.500) (5.694,-2.550) (5.716,-2.600) (5.737,-2.650) (5.757,-2.700) (5.777,-2.750) (5.796,-2.800) (5.814,-2.850) (5.831,-2.900) (5.848,-2.950) (5.863,-3.000) (5.879,-3.050) (5.893,-3.100) (5.907,-3.150) (5.920,-3.200) (5.933,-3.250) (5.944,-3.300) (5.956,-3.350) (5.966,-3.400) (5.976,-3.450) (5.985,-3.500) (5.994,-3.550) (6.002,-3.600) (6.010,-3.650) (6.017,-3.700) (6.023,-3.750) (6.029,-3.800) (6.034,-3.850) (6.038,-3.900) (6.042,-3.950) (6.046,-4.000) (6.049,-4.050) (6.051,-4.100) (6.053,-4.150) (6.054,-4.200) (6.054,-4.250) (6.054,-4.300) (6.054,-4.350) (6.053,-4.400) (6.051,-4.450) (6.049,-4.500) (6.047,-4.550) (6.043,-4.600) (6.040,-4.650) (6.035,-4.700) (6.031,-4.750) (6.025,-4.800) (6.019,-4.850) (6.013,-4.900) (6.006,-4.950) (5.998,-5.000) (5.990,-5.050) (5.982,-5.100) (5.973,-5.150) (5.963,-5.200) (5.953,-5.250) (5.942,-5.300) (5.931,-5.350) (5.919,-5.400) (5.906,-5.450) (5.893,-5.500) (5.880,-5.550) (5.866,-5.600) (5.851,-5.650) (5.836,-5.700) (5.820,-5.750) (5.804,-5.800) (5.787,-5.850) (5.769,-5.900) (5.751,-5.950) (5.732,-6.000) (5.713,-6.050) (5.693,-6.100) (5.673,-6.150) (5.652,-6.200) (5.630,-6.250) (5.608,-6.300) (5.585,-6.350) (5.561,-6.400) (5.537,-6.450) (5.512,-6.500) (5.487,-6.550) (5.461,-6.600) (5.434,-6.650) (5.406,-6.700) (5.378,-6.750) (5.350,-6.799) (5.320,-6.850) (5.289,-6.900) (5.258,-6.950) (5.226,-7.000) (5.194,-7.050) (5.160,-7.100) (5.126,-7.150) (5.091,-7.200) (5.056,-7.250) (5.025,-7.292) (4.981,-7.350) (4.950,-7.391) (4.904,-7.450) (4.864,-7.500) (4.825,-7.547) (4.781,-7.600) (4.738,-7.650) (4.700,-7.693) (4.650,-7.749) (4.603,-7.800) (4.555,-7.850) (4.507,-7.900) (4.457,-7.950) (4.407,-8.000) (4.355,-8.050) (4.302,-8.100) (4.250,-8.147) (4.200,-8.192) (4.162,-8.225) (4.104,-8.275) (4.050,-8.320) (4.000,-8.361) (3.952,-8.400) (3.900,-8.441) (3.855,-8.475) (3.800,-8.517) (3.755,-8.550) (3.700,-8.589) (3.650,-8.625) (3.600,-8.659) (3.550,-8.693) (3.501,-8.725) (3.450,-8.758) (3.400,-8.790) (3.350,-8.821) (3.302,-8.850) (3.250,-8.881) (3.200,-8.910) (3.150,-8.939) (3.100,-8.967) (3.050,-8.994) (3.000,-9.021) (2.950,-9.047) (2.900,-9.072) (2.850,-9.097) (2.800,-9.122) (2.750,-9.146) (2.700,-9.169) (2.650,-9.192) (2.600,-9.214) (2.550,-9.236) (2.500,-9.257) (2.450,-9.278) (2.400,-9.298) (2.350,-9.318) (2.300,-9.337) (2.250,-9.356) (2.200,-9.374) (2.150,-9.392) (2.100,-9.410) (2.050,-9.427) (2.000,-9.443) (1.950,-9.459) (1.900,-9.475) (1.850,-9.490) (1.800,-9.504) (1.750,-9.519) (1.700,-9.533) (1.650,-9.546) (1.600,-9.559) (1.550,-9.571) (1.500,-9.583) (1.450,-9.595) (1.400,-9.606) (1.350,-9.617) (1.300,-9.628) (1.250,-9.638) (1.200,-9.647) (1.150,-9.656) (1.100,-9.665) (1.050,-9.674) (1.000,-9.682) (0.950,-9.689) (0.900,-9.696) (0.850,-9.703) (0.800,-9.710) (0.750,-9.716) (0.700,-9.721) (0.650,-9.726) (0.600,-9.731) (0.550,-9.736) (0.500,-9.740) (0.450,-9.743) (0.400,-9.747) (0.350,-9.750) (0.300,-9.752) (0.250,-9.754) (0.200,-9.756) (0.150,-9.757) (0.100,-9.758) (0.050,-9.759) (0.000,-9.759) (-0.050,-9.759) (-0.100,-9.758) (-0.150,-9.757) (-0.200,-9.756) (-0.250,-9.754) (-0.300,-9.752) (-0.350,-9.750) (-0.400,-9.747) (-0.450,-9.743) (-0.500,-9.740) (-0.550,-9.736) (-0.600,-9.731) (-0.650,-9.726) (-0.700,-9.721) (-0.750,-9.716) (-0.800,-9.710) (-0.850,-9.703) (-0.900,-9.696) (-0.950,-9.689) (-1.000,-9.682) (-1.050,-9.674) (-1.100,-9.665) (-1.150,-9.656) (-1.200,-9.647) (-1.250,-9.638) (-1.300,-9.628) (-1.350,-9.617) (-1.400,-9.606) (-1.450,-9.595) (-1.500,-9.583) (-1.550,-9.571) (-1.600,-9.559) (-1.650,-9.546) (-1.700,-9.533) (-1.750,-9.519) (-1.800,-9.504) (-1.850,-9.490) (-1.899,-9.475) (-1.950,-9.459) (-2.000,-9.443) (-2.050,-9.427) (-2.100,-9.410) (-2.150,-9.392) (-2.198,-9.375) (-2.250,-9.356) (-2.300,-9.337) (-2.350,-9.318) (-2.400,-9.298) (-2.450,-9.278) (-2.500,-9.257) (-2.550,-9.236) (-2.600,-9.214) (-2.650,-9.192) (-2.700,-9.169) (-2.750,-9.146) (-2.800,-9.122) (-2.850,-9.097) (-2.900,-9.072) (-2.950,-9.047) (-3.000,-9.021) (-3.050,-8.994) (-3.100,-8.967) (-3.150,-8.939) (-3.200,-8.910) (-3.250,-8.881) (-3.300,-8.851) (-3.344,-8.825) (-3.400,-8.790) (-3.450,-8.758) (-3.500,-8.726) (-3.550,-8.693) (-3.600,-8.659) (-3.650,-8.625) (-3.700,-8.589) (-3.750,-8.553) (-3.800,-8.517) (-3.850,-8.479) (-3.900,-8.441) (-3.950,-8.401) (-4.000,-8.361) (-4.045,-8.325) (-4.100,-8.279) (-4.150,-8.236) (-4.191,-8.200) (-4.247,-8.150) (-4.300,-8.101) (-4.350,-8.055) (-4.400,-8.007) (-4.450,-7.957) (-4.500,-7.907) (-4.550,-7.856) (-4.600,-7.803) (-4.649,-7.750) (-4.694,-7.700) (-4.738,-7.650) (-4.775,-7.607) (-4.823,-7.550) (-4.864,-7.500) (-4.900,-7.455) (-4.943,-7.400) (-4.975,-7.359) (-5.019,-7.300) (-5.050,-7.258) (-5.091,-7.200) (-5.125,-7.152) (-5.160,-7.100) (-5.194,-7.050) (-5.225,-7.002) (-5.258,-6.950) (-5.289,-6.900) (-5.320,-6.850) (-5.349,-6.800) (-5.375,-6.755) (-5.400,-6.711) (-5.425,-6.666) (-5.450,-6.620) (-5.474,-6.575) (-5.500,-6.525) (-5.525,-6.475) (-5.549,-6.425) (-5.573,-6.375) (-5.596,-6.325) (-5.619,-6.275) (-5.641,-6.225) (-5.662,-6.175) (-5.683,-6.125) (-5.703,-6.075) (-5.723,-6.025) (-5.742,-5.975) (-5.760,-5.925) (-5.778,-5.875) (-5.795,-5.825) (-5.812,-5.775) (-5.828,-5.725) (-5.843,-5.675) (-5.858,-5.625) (-5.873,-5.575) (-5.887,-5.525) (-5.900,-5.475) (-5.913,-5.425) (-5.925,-5.375) (-5.936,-5.325) (-5.947,-5.275) (-5.958,-5.225) (-5.968,-5.175) (-5.977,-5.125) (-5.986,-5.075) (-5.994,-5.025) (-6.002,-4.975) (-6.009,-4.925) (-6.016,-4.875) (-6.022,-4.825) (-6.028,-4.775) (-6.033,-4.725) (-6.038,-4.675) (-6.042,-4.625) (-6.045,-4.575) (-6.048,-4.525) (-6.050,-4.475) (-6.052,-4.425) (-6.054,-4.375) (-6.054,-4.325) (-6.055,-4.275) (-6.054,-4.225) (-6.053,-4.175) (-6.052,-4.125) (-6.050,-4.075) (-6.047,-4.025) (-6.044,-3.975) (-6.040,-3.925) (-6.036,-3.875) (-6.031,-3.825) (-6.026,-3.775) (-6.020,-3.725) (-6.013,-3.675) (-6.006,-3.625) (-5.998,-3.575) (-5.990,-3.525) (-5.981,-3.475) (-5.971,-3.425) (-5.961,-3.375) (-5.950,-3.325) (-5.939,-3.275) (-5.926,-3.225) (-5.914,-3.175) (-5.900,-3.125) (-5.886,-3.075) (-5.871,-3.025) (-5.856,-2.975) (-5.839,-2.925) (-5.822,-2.875) (-5.805,-2.825) (-5.786,-2.775) (-5.767,-2.725) (-5.747,-2.675) (-5.727,-2.625) (-5.705,-2.575) (-5.683,-2.525) (-5.660,-2.475) (-5.636,-2.425) (-5.611,-2.375) (-5.585,-2.325) (-5.559,-2.275) (-5.531,-2.225) (-5.503,-2.175) (-5.475,-2.128) (-5.443,-2.075) (-5.411,-2.025) (-5.378,-1.975) (-5.350,-1.933) (-5.310,-1.875) (-5.275,-1.827) (-5.237,-1.775) (-5.200,-1.727) (-5.158,-1.675) (-5.125,-1.634) (-5.075,-1.576) (-5.030,-1.525) (-4.985,-1.475) (-4.950,-1.438) (-4.913,-1.400) (-4.875,-1.362) (-4.838,-1.325) (-4.800,-1.289) (-4.758,-1.250) (-4.702,-1.200) (-4.650,-1.155) (-4.600,-1.113) (-4.553,-1.075) (-4.500,-1.034) (-4.456,-1.000) (-4.400,-0.959) (-4.352,-0.925) (-4.300,-0.889) (-4.250,-0.856) (-4.203,-0.825) (-4.150,-0.792) (-4.100,-0.761) (-4.050,-0.732) (-4.000,-0.703) (-3.950,-0.676) (-3.902,-0.650) (-3.854,-0.625) (-3.805,-0.600) (-3.754,-0.575) (-3.701,-0.550) (-3.650,-0.527) (-3.600,-0.504) (-3.550,-0.483) (-3.500,-0.462) (-3.450,-0.442) (-3.400,-0.422) (-3.350,-0.403) (-3.300,-0.385) (-3.250,-0.367) (-3.200,-0.350) (-3.150,-0.334) (-3.100,-0.318) (-3.050,-0.302) (-3.000,-0.287) (-2.950,-0.273) (-2.900,-0.259) (-2.850,-0.246) (-2.800,-0.233) (-2.750,-0.220) (-2.700,-0.209) (-2.650,-0.197) (-2.600,-0.186) (-2.550,-0.175) (-2.500,-0.165) (-2.450,-0.155) (-2.400,-0.146) (-2.350,-0.137) (-2.300,-0.128) (-2.250,-0.120) (-2.200,-0.112) (-2.150,-0.105) (-2.100,-0.098) (-2.050,-0.091) (-2.000,-0.084) (-1.950,-0.078) (-1.900,-0.072) (-1.850,-0.067) (-1.800,-0.061) (-1.750,-0.056) (-1.700,-0.052) (-1.650,-0.047) (-1.600,-0.043) (-1.550,-0.039) (-1.500,-0.035) (-1.450,-0.032) (-1.400,-0.029) (-1.350,-0.026) (-1.302,-0.023)}
\EQ{(1.300,0.023) (1.350,0.026) (1.400,0.029) (1.450,0.032) (1.500,0.035) (1.550,0.039) (1.600,0.043) (1.650,0.047) (1.700,0.052) (1.750,0.056) (1.800,0.061) (1.850,0.067) (1.900,0.072) (1.950,0.078) (2.000,0.084) (2.050,0.091) (2.100,0.098) (2.150,0.105) (2.200,0.112) (2.250,0.120) (2.300,0.128) (2.350,0.137) (2.400,0.146) (2.450,0.155) (2.500,0.165) (2.550,0.175) (2.600,0.186) (2.650,0.197) (2.700,0.209) (2.750,0.220) (2.800,0.233) (2.850,0.246) (2.900,0.259) (2.950,0.273) (3.000,0.287) (3.050,0.302) (3.100,0.318) (3.150,0.334) (3.199,0.350) (3.250,0.367) (3.300,0.385) (3.350,0.403) (3.400,0.422) (3.450,0.442) (3.500,0.462) (3.550,0.483) (3.600,0.504) (3.646,0.525) (3.700,0.550) (3.750,0.573) (3.800,0.598) (3.850,0.623) (3.900,0.649) (3.949,0.675) (3.994,0.700) (4.038,0.725) (4.100,0.761) (4.150,0.792) (4.200,0.823) (4.250,0.856) (4.300,0.889) (4.350,0.923) (4.400,0.959) (4.450,0.996) (4.500,1.034) (4.550,1.073) (4.600,1.113) (4.644,1.150) (4.700,1.198) (4.750,1.243) (4.800,1.289) (4.838,1.325) (4.875,1.362) (4.913,1.400) (4.950,1.438) (4.985,1.475) (5.025,1.519) (5.074,1.575) (5.117,1.625) (5.150,1.665) (5.198,1.725) (5.237,1.775) (5.274,1.825) (5.310,1.875) (5.345,1.925) (5.375,1.970) (5.411,2.025) (5.443,2.075) (5.473,2.125) (5.500,2.171) (5.525,2.214) (5.550,2.259) (5.575,2.306) (5.598,2.350) (5.623,2.400) (5.648,2.450) (5.671,2.500) (5.694,2.550) (5.716,2.600) (5.737,2.650) (5.757,2.700) (5.777,2.750) (5.796,2.800) (5.814,2.850) (5.831,2.900) (5.848,2.950) (5.863,3.000) (5.879,3.050) (5.893,3.100) (5.907,3.150) (5.920,3.200) (5.933,3.250) (5.944,3.300) (5.956,3.350) (5.966,3.400) (5.976,3.450) (5.985,3.500) (5.994,3.550) (6.002,3.600) (6.010,3.650) (6.017,3.700) (6.023,3.750) (6.029,3.800) (6.034,3.850) (6.038,3.900) (6.042,3.950) (6.046,4.000) (6.049,4.050) (6.051,4.100) (6.053,4.150) (6.054,4.200) (6.054,4.250) (6.054,4.300) (6.054,4.350) (6.053,4.400) (6.051,4.450) (6.049,4.500) (6.047,4.550) (6.043,4.600) (6.040,4.650) (6.035,4.700) (6.031,4.750) (6.025,4.800) (6.019,4.850) (6.013,4.900) (6.006,4.950) (5.998,5.000) (5.990,5.050) (5.982,5.100) (5.973,5.150) (5.963,5.200) (5.953,5.250) (5.942,5.300) (5.931,5.350) (5.919,5.400) (5.906,5.450) (5.893,5.500) (5.880,5.550) (5.866,5.600) (5.851,5.650) (5.836,5.700) (5.820,5.750) (5.804,5.800) (5.787,5.850) (5.769,5.900) (5.751,5.950) (5.732,6.000) (5.713,6.050) (5.693,6.100) (5.673,6.150) (5.652,6.200) (5.630,6.250) (5.608,6.300) (5.585,6.350) (5.561,6.400) (5.537,6.450) (5.512,6.500) (5.487,6.550) (5.461,6.600) (5.434,6.650) (5.406,6.700) (5.378,6.750) (5.350,6.799) (5.320,6.850) (5.289,6.900) (5.258,6.950) (5.226,7.000) (5.194,7.050) (5.160,7.100) (5.126,7.150) (5.091,7.200) (5.056,7.250) (5.025,7.292) (4.981,7.350) (4.950,7.391) (4.904,7.450) (4.864,7.500) (4.825,7.547) (4.781,7.600) (4.738,7.650) (4.700,7.693) (4.650,7.749) (4.603,7.800) (4.555,7.850) (4.507,7.900) (4.457,7.950) (4.407,8.000) (4.355,8.050) (4.302,8.100) (4.250,8.147) (4.200,8.192) (4.162,8.225) (4.104,8.275) (4.050,8.320) (4.000,8.361) (3.952,8.400) (3.900,8.441) (3.855,8.475) (3.800,8.517) (3.755,8.550) (3.700,8.589) (3.650,8.625) (3.600,8.659) (3.550,8.693) (3.501,8.725) (3.450,8.758) (3.400,8.790) (3.350,8.821) (3.302,8.850) (3.250,8.881) (3.200,8.910) (3.150,8.939) (3.100,8.967) (3.050,8.994) (3.000,9.021) (2.950,9.047) (2.900,9.072) (2.850,9.097) (2.800,9.122) (2.750,9.146) (2.700,9.169) (2.650,9.192) (2.600,9.214) (2.550,9.236) (2.500,9.257) (2.450,9.278) (2.400,9.298) (2.350,9.318) (2.300,9.337) (2.250,9.356) (2.200,9.374) (2.150,9.392) (2.100,9.410) (2.050,9.427) (2.000,9.443) (1.950,9.459) (1.900,9.475) (1.850,9.490) (1.800,9.504) (1.750,9.519) (1.700,9.533) (1.650,9.546) (1.600,9.559) (1.550,9.571) (1.500,9.583) (1.450,9.595) (1.400,9.606) (1.350,9.617) (1.300,9.628) (1.250,9.638) (1.200,9.647) (1.150,9.656) (1.100,9.665) (1.050,9.674) (1.000,9.682) (0.950,9.689) (0.900,9.696) (0.850,9.703) (0.800,9.710) (0.750,9.716) (0.700,9.721) (0.650,9.726) (0.600,9.731) (0.550,9.736) (0.500,9.740) (0.450,9.743) (0.400,9.747) (0.350,9.750) (0.300,9.752) (0.250,9.754) (0.200,9.756) (0.150,9.757) (0.100,9.758) (0.050,9.759) (0.000,9.759) (-0.050,9.759) (-0.100,9.758) (-0.150,9.757) (-0.200,9.756) (-0.250,9.754) (-0.300,9.752) (-0.350,9.750) (-0.400,9.747) (-0.450,9.743) (-0.500,9.740) (-0.550,9.736) (-0.600,9.731) (-0.650,9.726) (-0.700,9.721) (-0.750,9.716) (-0.800,9.710) (-0.850,9.703) (-0.900,9.696) (-0.950,9.689) (-1.000,9.682) (-1.050,9.674) (-1.100,9.665) (-1.150,9.656) (-1.200,9.647) (-1.250,9.638) (-1.300,9.628) (-1.350,9.617) (-1.400,9.606) (-1.450,9.595) (-1.500,9.583) (-1.550,9.571) (-1.600,9.559) (-1.650,9.546) (-1.700,9.533) (-1.750,9.519) (-1.800,9.504) (-1.850,9.490) (-1.899,9.475) (-1.950,9.459) (-2.000,9.443) (-2.050,9.427) (-2.100,9.410) (-2.150,9.392) (-2.198,9.375) (-2.250,9.356) (-2.300,9.337) (-2.350,9.318) (-2.400,9.298) (-2.450,9.278) (-2.500,9.257) (-2.550,9.236) (-2.600,9.214) (-2.650,9.192) (-2.700,9.169) (-2.750,9.146) (-2.800,9.122) (-2.850,9.097) (-2.900,9.072) (-2.950,9.047) (-3.000,9.021) (-3.050,8.994) (-3.100,8.967) (-3.150,8.939) (-3.200,8.910) (-3.250,8.881) (-3.300,8.851) (-3.344,8.825) (-3.400,8.790) (-3.450,8.758) (-3.500,8.726) (-3.550,8.693) (-3.600,8.659) (-3.650,8.625) (-3.700,8.589) (-3.750,8.553) (-3.800,8.517) (-3.850,8.479) (-3.900,8.441) (-3.950,8.401) (-4.000,8.361) (-4.045,8.325) (-4.100,8.279) (-4.150,8.236) (-4.191,8.200) (-4.247,8.150) (-4.300,8.101) (-4.350,8.055) (-4.400,8.007) (-4.450,7.957) (-4.500,7.907) (-4.550,7.856) (-4.600,7.803) (-4.649,7.750) (-4.694,7.700) (-4.738,7.650) (-4.775,7.607) (-4.823,7.550) (-4.864,7.500) (-4.900,7.455) (-4.943,7.400) (-4.975,7.359) (-5.019,7.300) (-5.050,7.258) (-5.091,7.200) (-5.125,7.152) (-5.160,7.100) (-5.194,7.050) (-5.225,7.002) (-5.258,6.950) (-5.289,6.900) (-5.320,6.850) (-5.349,6.800) (-5.375,6.755) (-5.400,6.711) (-5.425,6.666) (-5.450,6.620) (-5.474,6.575) (-5.500,6.525) (-5.525,6.475) (-5.549,6.425) (-5.573,6.375) (-5.596,6.325) (-5.619,6.275) (-5.641,6.225) (-5.662,6.175) (-5.683,6.125) (-5.703,6.075) (-5.723,6.025) (-5.742,5.975) (-5.760,5.925) (-5.778,5.875) (-5.795,5.825) (-5.812,5.775) (-5.828,5.725) (-5.843,5.675) (-5.858,5.625) (-5.873,5.575) (-5.887,5.525) (-5.900,5.475) (-5.913,5.425) (-5.925,5.375) (-5.936,5.325) (-5.947,5.275) (-5.958,5.225) (-5.968,5.175) (-5.977,5.125) (-5.986,5.075) (-5.994,5.025) (-6.002,4.975) (-6.009,4.925) (-6.016,4.875) (-6.022,4.825) (-6.028,4.775) (-6.033,4.725) (-6.038,4.675) (-6.042,4.625) (-6.045,4.575) (-6.048,4.525) (-6.050,4.475) (-6.052,4.425) (-6.054,4.375) (-6.054,4.325) (-6.055,4.275) (-6.054,4.225) (-6.053,4.175) (-6.052,4.125) (-6.050,4.075) (-6.047,4.025) (-6.044,3.975) (-6.040,3.925) (-6.036,3.875) (-6.031,3.825) (-6.026,3.775) (-6.020,3.725) (-6.013,3.675) (-6.006,3.625) (-5.998,3.575) (-5.990,3.525) (-5.981,3.475) (-5.971,3.425) (-5.961,3.375) (-5.950,3.325) (-5.939,3.275) (-5.926,3.225) (-5.914,3.175) (-5.900,3.125) (-5.886,3.075) (-5.871,3.025) (-5.856,2.975) (-5.839,2.925) (-5.822,2.875) (-5.805,2.825) (-5.786,2.775) (-5.767,2.725) (-5.747,2.675) (-5.727,2.625) (-5.705,2.575) (-5.683,2.525) (-5.660,2.475) (-5.636,2.425) (-5.611,2.375) (-5.585,2.325) (-5.559,2.275) (-5.531,2.225) (-5.503,2.175) (-5.475,2.128) (-5.443,2.075) (-5.411,2.025) (-5.378,1.975) (-5.350,1.933) (-5.310,1.875) (-5.275,1.827) (-5.237,1.775) (-5.200,1.727) (-5.158,1.675) (-5.125,1.634) (-5.075,1.576) (-5.030,1.525) (-4.985,1.475) (-4.950,1.438) (-4.913,1.400) (-4.875,1.362) (-4.838,1.325) (-4.800,1.289) (-4.758,1.250) (-4.702,1.200) (-4.650,1.155) (-4.600,1.113) (-4.553,1.075) (-4.500,1.034) (-4.456,1.000) (-4.400,0.959) (-4.352,0.925) (-4.300,0.889) (-4.250,0.856) (-4.203,0.825) (-4.150,0.792) (-4.100,0.761) (-4.050,0.732) (-4.000,0.703) (-3.950,0.676) (-3.902,0.650) (-3.854,0.625) (-3.805,0.600) (-3.754,0.575) (-3.701,0.550) (-3.650,0.527) (-3.600,0.504) (-3.550,0.483) (-3.500,0.462) (-3.450,0.442) (-3.400,0.422) (-3.350,0.403) (-3.300,0.385) (-3.250,0.367) (-3.200,0.350) (-3.150,0.334) (-3.100,0.318) (-3.050,0.302) (-3.000,0.287) (-2.950,0.273) (-2.900,0.259) (-2.850,0.246) (-2.800,0.233) (-2.750,0.220) (-2.700,0.209) (-2.650,0.197) (-2.600,0.186) (-2.550,0.175) (-2.500,0.165) (-2.450,0.155) (-2.400,0.146) (-2.350,0.137) (-2.300,0.128) (-2.250,0.120) (-2.200,0.112) (-2.150,0.105) (-2.100,0.098) (-2.050,0.091) (-2.000,0.084) (-1.950,0.078) (-1.900,0.072) (-1.850,0.067) (-1.800,0.061) (-1.750,0.056) (-1.700,0.052) (-1.650,0.047) (-1.600,0.043) (-1.550,0.039) (-1.500,0.035) (-1.450,0.032) (-1.400,0.029) (-1.350,0.026) (-1.302,0.023)}
\EQ{(1.300,0.046) (1.350,0.052) (1.400,0.058) (1.450,0.064) (1.500,0.071) (1.550,0.079) (1.600,0.086) (1.650,0.095) (1.700,0.104) (1.750,0.113) (1.800,0.123) (1.850,0.134) (1.900,0.145) (1.950,0.157) (2.000,0.170) (2.050,0.183) (2.100,0.197) (2.150,0.212) (2.200,0.227) (2.250,0.243) (2.300,0.260) (2.350,0.278) (2.400,0.297) (2.450,0.317) (2.500,0.337) (2.550,0.359) (2.600,0.381) (2.650,0.405) (2.700,0.429) (2.750,0.455) (2.800,0.482) (2.850,0.510) (2.900,0.539) (2.950,0.570) (2.998,0.600) (3.050,0.635) (3.100,0.670) (3.141,0.700) (3.200,0.745) (3.250,0.785) (3.298,0.825) (3.350,0.871) (3.400,0.917) (3.450,0.966) (3.500,1.017) (3.550,1.071) (3.598,1.125) (3.640,1.175) (3.675,1.219) (3.718,1.275) (3.750,1.319) (3.789,1.375) (3.822,1.425) (3.850,1.470) (3.883,1.525) (3.912,1.575) (3.939,1.625) (3.965,1.675) (3.990,1.725) (4.014,1.775) (4.036,1.825) (4.057,1.875) (4.077,1.925) (4.096,1.975) (4.114,2.025) (4.131,2.075) (4.147,2.125) (4.162,2.175) (4.176,2.225) (4.189,2.275) (4.201,2.325) (4.213,2.375) (4.223,2.425) (4.232,2.475) (4.241,2.525) (4.249,2.575) (4.256,2.625) (4.262,2.675) (4.267,2.725) (4.271,2.775) (4.275,2.825) (4.278,2.875) (4.280,2.925) (4.281,2.975) (4.281,3.025) (4.281,3.075) (4.280,3.125) (4.278,3.175) (4.275,3.225) (4.272,3.275) (4.268,3.325) (4.263,3.375) (4.257,3.425) (4.250,3.475) (4.243,3.525) (4.235,3.575) (4.226,3.625) (4.217,3.675) (4.206,3.725) (4.195,3.775) (4.184,3.825) (4.171,3.875) (4.157,3.925) (4.143,3.975) (4.128,4.025) (4.112,4.075) (4.096,4.125) (4.078,4.175) (4.060,4.225) (4.041,4.275) (4.021,4.325) (4.000,4.375) (3.979,4.425) (3.956,4.475) (3.933,4.525) (3.908,4.575) (3.883,4.625) (3.857,4.675) (3.830,4.725) (3.802,4.775) (3.775,4.821) (3.743,4.875) (3.711,4.925) (3.679,4.975) (3.650,5.019) (3.611,5.075) (3.576,5.125) (3.539,5.175) (3.501,5.225) (3.462,5.275) (3.425,5.321) (3.380,5.375) (3.336,5.425) (3.300,5.466) (3.250,5.520) (3.200,5.573) (3.150,5.623) (3.100,5.672) (3.050,5.720) (3.000,5.766) (2.961,5.800) (2.904,5.850) (2.850,5.895) (2.800,5.935) (2.750,5.974) (2.700,6.013) (2.650,6.049) (2.600,6.085) (2.550,6.120) (2.505,6.150) (2.450,6.186) (2.400,6.218) (2.350,6.249) (2.306,6.275) (2.250,6.308) (2.200,6.336) (2.150,6.363) (2.100,6.390) (2.050,6.415) (2.000,6.440) (1.950,6.464) (1.900,6.488) (1.850,6.510) (1.800,6.532) (1.750,6.553) (1.700,6.574) (1.650,6.593) (1.600,6.612) (1.550,6.631) (1.500,6.648) (1.450,6.665) (1.400,6.682) (1.350,6.698) (1.300,6.713) (1.250,6.727) (1.200,6.741) (1.150,6.754) (1.100,6.767) (1.050,6.779) (1.000,6.791) (0.950,6.801) (0.900,6.812) (0.850,6.821) (0.800,6.831) (0.750,6.839) (0.700,6.847) (0.650,6.854) (0.600,6.861) (0.550,6.868) (0.500,6.873) (0.450,6.879) (0.400,6.883) (0.350,6.887) (0.300,6.891) (0.250,6.894) (0.200,6.896) (0.150,6.898) (0.100,6.900) (0.050,6.900) (0.000,6.901) (-0.050,6.900) (-0.100,6.900) (-0.150,6.898) (-0.200,6.896) (-0.250,6.894) (-0.300,6.891) (-0.350,6.887) (-0.400,6.883) (-0.450,6.879) (-0.500,6.873) (-0.550,6.868) (-0.600,6.861) (-0.650,6.854) (-0.700,6.847) (-0.750,6.839) (-0.800,6.831) (-0.850,6.821) (-0.900,6.812) (-0.950,6.801) (-1.000,6.791) (-1.050,6.779) (-1.100,6.767) (-1.150,6.754) (-1.200,6.741) (-1.250,6.727) (-1.300,6.713) (-1.350,6.698) (-1.400,6.682) (-1.450,6.665) (-1.500,6.648) (-1.550,6.631) (-1.600,6.612) (-1.650,6.593) (-1.697,6.575) (-1.750,6.553) (-1.800,6.532) (-1.850,6.510) (-1.900,6.488) (-1.950,6.464) (-2.000,6.440) (-2.050,6.415) (-2.100,6.390) (-2.150,6.363) (-2.200,6.336) (-2.250,6.308) (-2.300,6.279) (-2.348,6.250) (-2.400,6.218) (-2.450,6.186) (-2.500,6.154) (-2.543,6.125) (-2.600,6.085) (-2.649,6.050) (-2.700,6.013) (-2.749,5.975) (-2.800,5.935) (-2.844,5.900) (-2.900,5.853) (-2.950,5.810) (-2.990,5.775) (-3.044,5.725) (-3.097,5.675) (-3.148,5.625) (-3.198,5.575) (-3.246,5.525) (-3.292,5.475) (-3.336,5.425) (-3.375,5.380) (-3.421,5.325) (-3.462,5.275) (-3.500,5.227) (-3.539,5.175) (-3.575,5.126) (-3.611,5.075) (-3.646,5.025) (-3.675,4.981) (-3.711,4.925) (-3.743,4.875) (-3.773,4.825) (-3.800,4.778) (-3.825,4.734) (-3.850,4.688) (-3.875,4.641) (-3.900,4.592) (-3.925,4.541) (-3.950,4.488) (-3.975,4.433) (-4.000,4.376) (-4.021,4.325) (-4.041,4.275) (-4.060,4.225) (-4.078,4.175) (-4.096,4.125) (-4.112,4.075) (-4.128,4.025) (-4.143,3.975) (-4.157,3.925) (-4.171,3.875) (-4.184,3.825) (-4.195,3.775) (-4.206,3.725) (-4.217,3.675) (-4.226,3.625) (-4.235,3.575) (-4.243,3.525) (-4.250,3.475) (-4.257,3.425) (-4.263,3.375) (-4.268,3.325) (-4.272,3.275) (-4.275,3.225) (-4.278,3.175) (-4.280,3.125) (-4.281,3.075) (-4.281,3.025) (-4.281,2.975) (-4.280,2.925) (-4.278,2.875) (-4.275,2.825) (-4.271,2.775) (-4.267,2.725) (-4.262,2.675) (-4.256,2.625) (-4.249,2.575) (-4.241,2.525) (-4.232,2.475) (-4.223,2.425) (-4.213,2.375) (-4.201,2.325) (-4.189,2.275) (-4.176,2.225) (-4.162,2.175) (-4.147,2.125) (-4.131,2.075) (-4.114,2.025) (-4.096,1.975) (-4.077,1.925) (-4.057,1.875) (-4.036,1.825) (-4.014,1.775) (-3.990,1.725) (-3.965,1.675) (-3.939,1.625) (-3.912,1.575) (-3.883,1.525) (-3.853,1.475) (-3.825,1.430) (-3.789,1.375) (-3.754,1.325) (-3.718,1.275) (-3.680,1.225) (-3.640,1.175) (-3.600,1.127) (-3.554,1.075) (-3.508,1.025) (-3.459,0.975) (-3.408,0.925) (-3.355,0.875) (-3.300,0.827) (-3.250,0.785) (-3.207,0.750) (-3.150,0.707) (-3.107,0.675) (-3.050,0.635) (-3.000,0.602) (-2.950,0.570) (-2.900,0.539) (-2.850,0.510) (-2.800,0.482) (-2.750,0.455) (-2.700,0.429) (-2.650,0.405) (-2.600,0.381) (-2.550,0.359) (-2.500,0.337) (-2.450,0.317) (-2.400,0.297) (-2.350,0.278) (-2.300,0.260) (-2.250,0.243) (-2.200,0.227) (-2.150,0.212) (-2.100,0.197) (-2.050,0.183) (-2.000,0.170) (-1.950,0.157) (-1.900,0.145) (-1.850,0.134) (-1.800,0.123) (-1.750,0.113) (-1.700,0.104) (-1.650,0.095) (-1.600,0.086) (-1.550,0.079) (-1.500,0.071) (-1.450,0.064) (-1.400,0.058) (-1.350,0.052) (-1.300,0.046)}
\EQ{(1.300,0.070) (1.350,0.078) (1.400,0.087) (1.450,0.097) (1.500,0.107) (1.550,0.118) (1.600,0.130) (1.650,0.143) (1.700,0.157) (1.750,0.171) (1.800,0.187) (1.850,0.203) (1.900,0.220) (1.950,0.239) (2.000,0.258) (2.050,0.279) (2.098,0.300) (2.150,0.324) (2.200,0.348) (2.250,0.374) (2.298,0.400) (2.342,0.425) (2.400,0.460) (2.450,0.492) (2.500,0.525) (2.550,0.561) (2.600,0.598) (2.650,0.638) (2.694,0.675) (2.750,0.724) (2.800,0.772) (2.850,0.822) (2.899,0.875) (2.943,0.925) (2.985,0.975) (3.024,1.025) (3.060,1.075) (3.095,1.125) (3.125,1.171) (3.159,1.225) (3.188,1.275) (3.216,1.325) (3.241,1.375) (3.266,1.425) (3.289,1.475) (3.310,1.525) (3.330,1.575) (3.349,1.625) (3.367,1.675) (3.383,1.725) (3.398,1.775) (3.412,1.825) (3.425,1.875) (3.437,1.925) (3.447,1.975) (3.456,2.025) (3.465,2.075) (3.472,2.125) (3.479,2.175) (3.484,2.225) (3.488,2.275) (3.491,2.325) (3.494,2.375) (3.495,2.425) (3.496,2.475) (3.495,2.525) (3.494,2.575) (3.491,2.625) (3.488,2.675) (3.483,2.725) (3.478,2.775) (3.472,2.825) (3.465,2.875) (3.457,2.925) (3.448,2.975) (3.438,3.025) (3.427,3.075) (3.415,3.125) (3.403,3.175) (3.389,3.225) (3.374,3.275) (3.359,3.325) (3.342,3.375) (3.325,3.423) (3.306,3.475) (3.286,3.525) (3.265,3.575) (3.243,3.625) (3.220,3.675) (3.196,3.725) (3.171,3.775) (3.145,3.825) (3.118,3.875) (3.089,3.925) (3.059,3.975) (3.028,4.025) (3.000,4.068) (2.962,4.125) (2.927,4.175) (2.890,4.225) (2.852,4.275) (2.812,4.325) (2.775,4.370) (2.728,4.425) (2.683,4.475) (2.637,4.525) (2.600,4.563) (2.563,4.600) (2.525,4.637) (2.485,4.675) (2.430,4.725) (2.375,4.773) (2.325,4.815) (2.281,4.850) (2.225,4.894) (2.183,4.925) (2.125,4.967) (2.078,5.000) (2.025,5.035) (1.975,5.067) (1.925,5.098) (1.881,5.125) (1.825,5.157) (1.775,5.185) (1.725,5.212) (1.675,5.238) (1.625,5.262) (1.575,5.286) (1.525,5.309) (1.475,5.331) (1.425,5.352) (1.375,5.373) (1.325,5.392) (1.275,5.411) (1.225,5.428) (1.175,5.445) (1.125,5.462) (1.075,5.477) (1.025,5.492) (0.975,5.505) (0.925,5.519) (0.875,5.531) (0.825,5.543) (0.775,5.553) (0.725,5.564) (0.675,5.573) (0.625,5.582) (0.575,5.590) (0.525,5.597) (0.475,5.604) (0.425,5.610) (0.375,5.616) (0.325,5.620) (0.275,5.624) (0.225,5.628) (0.175,5.630) (0.125,5.632) (0.075,5.634) (0.025,5.634) (-0.050,5.634) (-0.100,5.633) (-0.150,5.631) (-0.200,5.629) (-0.250,5.626) (-0.300,5.622) (-0.350,5.618) (-0.400,5.613) (-0.450,5.607) (-0.500,5.601) (-0.550,5.594) (-0.600,5.586) (-0.650,5.578) (-0.700,5.569) (-0.750,5.559) (-0.800,5.548) (-0.850,5.537) (-0.899,5.525) (-0.950,5.512) (-1.000,5.499) (-1.050,5.484) (-1.100,5.469) (-1.150,5.453) (-1.200,5.437) (-1.250,5.420) (-1.300,5.401) (-1.350,5.382) (-1.400,5.363) (-1.450,5.342) (-1.500,5.320) (-1.550,5.298) (-1.599,5.275) (-1.650,5.250) (-1.700,5.225) (-1.747,5.200) (-1.793,5.175) (-1.838,5.150) (-1.900,5.114) (-1.950,5.083) (-2.000,5.051) (-2.050,5.019) (-2.100,4.985) (-2.149,4.950) (-2.200,4.913) (-2.249,4.875) (-2.300,4.835) (-2.343,4.800) (-2.400,4.751) (-2.450,4.707) (-2.500,4.661) (-2.538,4.625) (-2.575,4.588) (-2.613,4.550) (-2.650,4.511) (-2.700,4.457) (-2.750,4.400) (-2.792,4.350) (-2.825,4.309) (-2.871,4.250) (-2.908,4.200) (-2.944,4.150) (-2.975,4.106) (-3.012,4.050) (-3.044,4.000) (-3.074,3.950) (-3.100,3.906) (-3.125,3.862) (-3.150,3.816) (-3.175,3.768) (-3.200,3.718) (-3.225,3.665) (-3.250,3.610) (-3.275,3.551) (-3.296,3.500) (-3.315,3.450) (-3.333,3.400) (-3.350,3.351) (-3.366,3.300) (-3.382,3.250) (-3.396,3.200) (-3.409,3.150) (-3.421,3.100) (-3.433,3.050) (-3.443,3.000) (-3.453,2.950) (-3.461,2.900) (-3.469,2.850) (-3.475,2.800) (-3.481,2.750) (-3.486,2.700) (-3.490,2.650) (-3.492,2.600) (-3.494,2.550) (-3.495,2.500) (-3.496,2.450) (-3.495,2.400) (-3.493,2.350) (-3.490,2.300) (-3.486,2.250) (-3.481,2.200) (-3.475,2.150) (-3.469,2.100) (-3.461,2.050) (-3.452,2.000) (-3.442,1.950) (-3.431,1.900) (-3.419,1.850) (-3.405,1.800) (-3.391,1.750) (-3.375,1.700) (-3.358,1.650) (-3.340,1.600) (-3.321,1.550) (-3.300,1.501) (-3.277,1.450) (-3.254,1.400) (-3.229,1.350) (-3.202,1.300) (-3.175,1.252) (-3.144,1.200) (-3.112,1.150) (-3.078,1.100) (-3.042,1.050) (-3.004,1.000) (-2.964,0.950) (-2.925,0.904) (-2.876,0.850) (-2.828,0.800) (-2.777,0.750) (-2.725,0.702) (-2.675,0.659) (-2.634,0.625) (-2.575,0.579) (-2.525,0.543) (-2.475,0.508) (-2.425,0.475) (-2.375,0.444) (-2.325,0.415) (-2.275,0.387) (-2.225,0.361) (-2.175,0.336) (-2.125,0.312) (-2.075,0.290) (-2.025,0.269) (-1.975,0.248) (-1.925,0.230) (-1.875,0.212) (-1.825,0.195) (-1.775,0.179) (-1.725,0.164) (-1.676,0.150) (-1.625,0.137) (-1.575,0.124) (-1.525,0.113) (-1.475,0.102) (-1.425,0.092) (-1.375,0.082) (-1.325,0.074) (-1.300,0.070)}
\EQ{(1.300,0.093) (1.350,0.104) (1.400,0.116) (1.450,0.130) (1.500,0.144) (1.550,0.159) (1.600,0.175) (1.650,0.193) (1.700,0.211) (1.750,0.231) (1.800,0.252) (1.850,0.275) (1.900,0.299) (1.950,0.324) (1.997,0.350) (2.041,0.375) (2.100,0.412) (2.150,0.444) (2.194,0.475) (2.250,0.517) (2.292,0.550) (2.350,0.599) (2.400,0.645) (2.450,0.693) (2.500,0.746) (2.548,0.800) (2.589,0.850) (2.625,0.896) (2.664,0.950) (2.697,1.000) (2.725,1.044) (2.759,1.100) (2.786,1.150) (2.812,1.200) (2.836,1.250) (2.858,1.300) (2.879,1.350) (2.898,1.400) (2.916,1.450) (2.932,1.500) (2.947,1.550) (2.960,1.600) (2.972,1.650) (2.983,1.700) (2.993,1.750) (3.001,1.800) (3.008,1.850) (3.014,1.900) (3.019,1.950) (3.023,2.000) (3.025,2.050) (3.027,2.100) (3.027,2.150) (3.026,2.200) (3.025,2.250) (3.022,2.300) (3.018,2.350) (3.013,2.400) (3.006,2.450) (2.999,2.500) (2.991,2.550) (2.981,2.600) (2.971,2.650) (2.959,2.700) (2.947,2.750) (2.933,2.800) (2.918,2.850) (2.902,2.900) (2.885,2.950) (2.866,3.000) (2.847,3.050) (2.826,3.100) (2.804,3.150) (2.781,3.200) (2.756,3.250) (2.730,3.300) (2.703,3.350) (2.675,3.399) (2.645,3.450) (2.613,3.500) (2.580,3.550) (2.550,3.594) (2.509,3.650) (2.475,3.696) (2.432,3.750) (2.400,3.789) (2.350,3.846) (2.301,3.900) (2.253,3.950) (2.203,4.000) (2.151,4.050) (2.100,4.096) (2.050,4.139) (2.007,4.175) (1.950,4.220) (1.900,4.258) (1.850,4.295) (1.807,4.325) (1.750,4.363) (1.700,4.395) (1.651,4.425) (1.600,4.455) (1.550,4.483) (1.500,4.510) (1.450,4.536) (1.400,4.561) (1.350,4.585) (1.300,4.607) (1.250,4.629) (1.200,4.649) (1.150,4.669) (1.100,4.687) (1.050,4.705) (1.000,4.722) (0.950,4.737) (0.900,4.752) (0.850,4.766) (0.800,4.779) (0.750,4.792) (0.700,4.803) (0.650,4.814) (0.600,4.824) (0.550,4.833) (0.500,4.841) (0.450,4.848) (0.400,4.855) (0.350,4.861) (0.300,4.866) (0.250,4.870) (0.200,4.873) (0.150,4.876) (0.100,4.878) (0.050,4.879) (0.000,4.880) (-0.050,4.879) (-0.100,4.878) (-0.150,4.876) (-0.200,4.873) (-0.250,4.870) (-0.300,4.866) (-0.350,4.861) (-0.400,4.855) (-0.450,4.848) (-0.500,4.841) (-0.550,4.833) (-0.600,4.824) (-0.650,4.814) (-0.700,4.803) (-0.750,4.792) (-0.800,4.779) (-0.850,4.766) (-0.900,4.752) (-0.950,4.737) (-1.000,4.722) (-1.050,4.705) (-1.100,4.687) (-1.150,4.669) (-1.198,4.650) (-1.250,4.629) (-1.300,4.607) (-1.350,4.585) (-1.400,4.561) (-1.450,4.536) (-1.500,4.510) (-1.550,4.483) (-1.600,4.455) (-1.650,4.426) (-1.700,4.395) (-1.750,4.363) (-1.800,4.330) (-1.843,4.300) (-1.900,4.258) (-1.944,4.225) (-2.000,4.181) (-2.050,4.139) (-2.095,4.100) (-2.150,4.051) (-2.200,4.003) (-2.250,3.954) (-2.300,3.901) (-2.347,3.850) (-2.390,3.800) (-2.425,3.758) (-2.472,3.700) (-2.509,3.650) (-2.546,3.600) (-2.575,3.558) (-2.613,3.500) (-2.645,3.450) (-2.675,3.400) (-2.700,3.356) (-2.725,3.310) (-2.750,3.262) (-2.775,3.212) (-2.800,3.159) (-2.825,3.102) (-2.847,3.050) (-2.866,3.000) (-2.885,2.950) (-2.902,2.900) (-2.918,2.850) (-2.933,2.800) (-2.947,2.750) (-2.959,2.700) (-2.971,2.650) (-2.981,2.600) (-2.991,2.550) (-2.999,2.500) (-3.006,2.450) (-3.013,2.400) (-3.018,2.350) (-3.022,2.300) (-3.025,2.250) (-3.026,2.200) (-3.027,2.150) (-3.027,2.100) (-3.025,2.050) (-3.023,2.000) (-3.019,1.950) (-3.014,1.900) (-3.008,1.850) (-3.001,1.800) (-2.993,1.750) (-2.983,1.700) (-2.972,1.650) (-2.960,1.600) (-2.947,1.550) (-2.932,1.500) (-2.916,1.450) (-2.898,1.400) (-2.879,1.350) (-2.858,1.300) (-2.836,1.250) (-2.812,1.200) (-2.786,1.150) (-2.759,1.100) (-2.729,1.050) (-2.700,1.004) (-2.664,0.950) (-2.628,0.900) (-2.589,0.850) (-2.550,0.802) (-2.504,0.750) (-2.457,0.700) (-2.406,0.650) (-2.351,0.600) (-2.300,0.557) (-2.261,0.525) (-2.200,0.479) (-2.158,0.450) (-2.100,0.412) (-2.050,0.381) (-2.000,0.352) (-1.951,0.325) (-1.902,0.300) (-1.850,0.275) (-1.800,0.252) (-1.750,0.231) (-1.700,0.211) (-1.650,0.193) (-1.600,0.175) (-1.550,0.159) (-1.500,0.144) (-1.450,0.130) (-1.400,0.116) (-1.350,0.104) (-1.300,0.093)}
\EQ{(1.300,0.117) (1.350,0.131) (1.400,0.146) (1.450,0.163) (1.500,0.181) (1.549,0.200) (1.600,0.221) (1.650,0.244) (1.700,0.268) (1.750,0.293) (1.800,0.321) (1.849,0.350) (1.900,0.382) (1.950,0.416) (1.997,0.450) (2.050,0.492) (2.090,0.525) (2.145,0.575) (2.196,0.625) (2.243,0.675) (2.287,0.725) (2.325,0.772) (2.365,0.825) (2.400,0.875) (2.432,0.925) (2.462,0.975) (2.490,1.025) (2.516,1.075) (2.540,1.125) (2.562,1.175) (2.582,1.225) (2.600,1.275) (2.617,1.325) (2.632,1.375) (2.646,1.425) (2.658,1.475) (2.669,1.525) (2.679,1.575) (2.687,1.625) (2.693,1.675) (2.699,1.725) (2.703,1.775) (2.706,1.825) (2.707,1.875) (2.708,1.925) (2.707,1.975) (2.705,2.025) (2.701,2.075) (2.697,2.125) (2.691,2.175) (2.684,2.225) (2.676,2.275) (2.667,2.325) (2.656,2.375) (2.644,2.425) (2.631,2.475) (2.617,2.525) (2.602,2.575) (2.585,2.625) (2.567,2.675) (2.547,2.725) (2.527,2.775) (2.504,2.825) (2.481,2.875) (2.456,2.925) (2.430,2.975) (2.402,3.025) (2.375,3.070) (2.341,3.125) (2.308,3.175) (2.275,3.223) (2.237,3.275) (2.200,3.323) (2.158,3.375) (2.125,3.414) (2.075,3.470) (2.025,3.523) (1.975,3.574) (1.925,3.621) (1.875,3.667) (1.836,3.700) (1.776,3.750) (1.725,3.790) (1.677,3.825) (1.625,3.862) (1.575,3.896) (1.530,3.925) (1.475,3.959) (1.425,3.988) (1.375,4.016) (1.325,4.043) (1.275,4.068) (1.225,4.092) (1.175,4.115) (1.125,4.137) (1.075,4.158) (1.025,4.177) (0.975,4.196) (0.925,4.213) (0.875,4.229) (0.825,4.245) (0.775,4.259) (0.725,4.272) (0.675,4.285) (0.625,4.296) (0.575,4.307) (0.525,4.317) (0.475,4.325) (0.425,4.333) (0.375,4.340) (0.325,4.346) (0.275,4.351) (0.225,4.356) (0.175,4.359) (0.125,4.362) (0.075,4.363) (0.025,4.364) (-0.050,4.364) (-0.100,4.363) (-0.150,4.361) (-0.200,4.358) (-0.250,4.354) (-0.300,4.349) (-0.350,4.343) (-0.400,4.337) (-0.450,4.329) (-0.500,4.321) (-0.550,4.312) (-0.600,4.302) (-0.650,4.291) (-0.700,4.279) (-0.750,4.266) (-0.800,4.252) (-0.850,4.237) (-0.900,4.221) (-0.950,4.204) (-1.000,4.187) (-1.050,4.168) (-1.100,4.147) (-1.150,4.126) (-1.200,4.104) (-1.250,4.080) (-1.300,4.056) (-1.350,4.030) (-1.400,4.002) (-1.448,3.975) (-1.500,3.944) (-1.550,3.912) (-1.600,3.879) (-1.642,3.850) (-1.700,3.808) (-1.744,3.775) (-1.800,3.730) (-1.850,3.688) (-1.894,3.650) (-1.948,3.600) (-1.999,3.550) (-2.047,3.500) (-2.093,3.450) (-2.137,3.400) (-2.175,3.354) (-2.218,3.300) (-2.250,3.257) (-2.291,3.200) (-2.325,3.150) (-2.357,3.100) (-2.387,3.050) (-2.416,3.000) (-2.443,2.950) (-2.469,2.900) (-2.493,2.850) (-2.516,2.800) (-2.537,2.750) (-2.557,2.700) (-2.575,2.653) (-2.593,2.600) (-2.610,2.550) (-2.624,2.500) (-2.638,2.450) (-2.650,2.400) (-2.662,2.350) (-2.672,2.300) (-2.680,2.250) (-2.688,2.200) (-2.694,2.150) (-2.699,2.100) (-2.703,2.050) (-2.706,2.000) (-2.707,1.950) (-2.708,1.900) (-2.707,1.850) (-2.704,1.800) (-2.701,1.750) (-2.696,1.700) (-2.690,1.650) (-2.683,1.600) (-2.674,1.550) (-2.664,1.500) (-2.652,1.450) (-2.639,1.400) (-2.625,1.351) (-2.609,1.300) (-2.591,1.250) (-2.572,1.200) (-2.551,1.150) (-2.528,1.100) (-2.503,1.050) (-2.476,1.000) (-2.450,0.954) (-2.416,0.900) (-2.383,0.850) (-2.350,0.805) (-2.308,0.750) (-2.275,0.711) (-2.225,0.655) (-2.175,0.604) (-2.125,0.557) (-2.075,0.513) (-2.029,0.475) (-1.975,0.434) (-1.927,0.400) (-1.875,0.366) (-1.825,0.335) (-1.775,0.307) (-1.725,0.280) (-1.675,0.255) (-1.625,0.232) (-1.575,0.211) (-1.525,0.191) (-1.475,0.172) (-1.425,0.155) (-1.375,0.139) (-1.325,0.124) (-1.300,0.117)}
\EQ{(1.300,0.141) (1.350,0.158) (1.400,0.177) (1.450,0.197) (1.500,0.220) (1.550,0.243) (1.600,0.269) (1.650,0.297) (1.697,0.325) (1.750,0.359) (1.800,0.394) (1.841,0.425) (1.900,0.473) (1.950,0.517) (2.000,0.566) (2.050,0.619) (2.099,0.675) (2.138,0.725) (2.174,0.775) (2.208,0.825) (2.239,0.875) (2.267,0.925) (2.293,0.975) (2.317,1.025) (2.339,1.075) (2.359,1.125) (2.377,1.175) (2.394,1.225) (2.409,1.275) (2.422,1.325) (2.433,1.375) (2.443,1.425) (2.451,1.475) (2.458,1.525) (2.464,1.575) (2.468,1.625) (2.470,1.675) (2.472,1.725) (2.472,1.775) (2.470,1.825) (2.467,1.875) (2.463,1.925) (2.458,1.975) (2.451,2.025) (2.443,2.075) (2.434,2.125) (2.423,2.175) (2.411,2.225) (2.398,2.275) (2.383,2.325) (2.367,2.375) (2.350,2.423) (2.330,2.475) (2.310,2.525) (2.288,2.575) (2.265,2.625) (2.240,2.675) (2.213,2.725) (2.184,2.775) (2.154,2.825) (2.125,2.871) (2.089,2.925) (2.053,2.975) (2.015,3.025) (1.975,3.075) (1.933,3.125) (1.888,3.175) (1.850,3.215) (1.800,3.265) (1.763,3.300) (1.725,3.335) (1.679,3.375) (1.625,3.420) (1.575,3.459) (1.525,3.496) (1.475,3.531) (1.425,3.565) (1.375,3.597) (1.328,3.625) (1.285,3.650) (1.239,3.675) (1.191,3.700) (1.141,3.725) (1.087,3.750) (1.031,3.775) (0.975,3.798) (0.925,3.817) (0.875,3.835) (0.825,3.852) (0.775,3.868) (0.725,3.883) (0.675,3.897) (0.625,3.909) (0.575,3.921) (0.525,3.932) (0.475,3.941) (0.425,3.950) (0.375,3.958) (0.325,3.964) (0.275,3.970) (0.225,3.975) (0.175,3.978) (0.125,3.981) (0.075,3.983) (0.025,3.984) (-0.050,3.984) (-0.100,3.982) (-0.150,3.980) (-0.200,3.977) (-0.250,3.972) (-0.300,3.967) (-0.350,3.961) (-0.400,3.954) (-0.450,3.946) (-0.500,3.937) (-0.550,3.926) (-0.600,3.915) (-0.650,3.903) (-0.700,3.890) (-0.750,3.876) (-0.800,3.860) (-0.850,3.844) (-0.900,3.826) (-0.950,3.808) (-1.000,3.788) (-1.050,3.767) (-1.100,3.744) (-1.150,3.721) (-1.200,3.696) (-1.250,3.669) (-1.300,3.641) (-1.350,3.612) (-1.400,3.581) (-1.448,3.550) (-1.500,3.514) (-1.550,3.478) (-1.600,3.440) (-1.650,3.400) (-1.700,3.357) (-1.750,3.312) (-1.790,3.275) (-1.840,3.225) (-1.875,3.189) (-1.910,3.150) (-1.950,3.105) (-1.995,3.050) (-2.034,3.000) (-2.071,2.950) (-2.100,2.908) (-2.139,2.850) (-2.170,2.800) (-2.199,2.750) (-2.225,2.703) (-2.250,2.654) (-2.275,2.603) (-2.299,2.550) (-2.320,2.500) (-2.340,2.450) (-2.358,2.400) (-2.375,2.351) (-2.391,2.300) (-2.405,2.250) (-2.417,2.200) (-2.429,2.150) (-2.439,2.100) (-2.447,2.050) (-2.455,2.000) (-2.461,1.950) (-2.466,1.900) (-2.469,1.850) (-2.471,1.800) (-2.472,1.750) (-2.471,1.700) (-2.469,1.650) (-2.466,1.600) (-2.461,1.550) (-2.455,1.500) (-2.447,1.450) (-2.438,1.400) (-2.427,1.350) (-2.415,1.300) (-2.401,1.250) (-2.386,1.200) (-2.369,1.150) (-2.350,1.101) (-2.329,1.050) (-2.306,1.000) (-2.281,0.950) (-2.253,0.900) (-2.225,0.852) (-2.192,0.800) (-2.157,0.750) (-2.125,0.708) (-2.078,0.650) (-2.033,0.600) (-1.984,0.550) (-1.931,0.500) (-1.875,0.452) (-1.825,0.413) (-1.775,0.376) (-1.725,0.343) (-1.675,0.312) (-1.625,0.283) (-1.575,0.256) (-1.525,0.231) (-1.475,0.208) (-1.425,0.187) (-1.375,0.167) (-1.327,0.150) (-1.300,0.141)}
\EQ{(1.300,0.165) (1.350,0.186) (1.400,0.208) (1.450,0.233) (1.500,0.259) (1.550,0.288) (1.600,0.319) (1.646,0.350) (1.700,0.390) (1.744,0.425) (1.800,0.474) (1.850,0.522) (1.900,0.575) (1.943,0.625) (1.975,0.665) (2.018,0.725) (2.050,0.773) (2.082,0.825) (2.109,0.875) (2.134,0.925) (2.157,0.975) (2.178,1.025) (2.197,1.075) (2.213,1.125) (2.228,1.175) (2.242,1.225) (2.253,1.275) (2.263,1.325) (2.271,1.375) (2.277,1.425) (2.282,1.475) (2.286,1.525) (2.288,1.575) (2.288,1.625) (2.288,1.675) (2.285,1.725) (2.281,1.775) (2.276,1.825) (2.269,1.875) (2.261,1.925) (2.252,1.975) (2.241,2.025) (2.228,2.075) (2.215,2.125) (2.199,2.175) (2.182,2.225) (2.164,2.275) (2.144,2.325) (2.122,2.375) (2.100,2.423) (2.075,2.473) (2.050,2.520) (2.025,2.565) (1.989,2.625) (1.956,2.675) (1.925,2.720) (1.885,2.775) (1.850,2.820) (1.805,2.875) (1.761,2.925) (1.725,2.964) (1.690,3.000) (1.650,3.039) (1.611,3.075) (1.554,3.125) (1.500,3.170) (1.450,3.209) (1.400,3.245) (1.358,3.275) (1.300,3.313) (1.250,3.344) (1.200,3.373) (1.152,3.400) (1.104,3.425) (1.053,3.450) (1.000,3.475) (0.950,3.497) (0.900,3.517) (0.850,3.536) (0.800,3.554) (0.750,3.571) (0.700,3.587) (0.650,3.601) (0.600,3.614) (0.550,3.626) (0.500,3.637) (0.450,3.647) (0.400,3.656) (0.350,3.664) (0.300,3.670) (0.250,3.676) (0.200,3.680) (0.150,3.684) (0.100,3.687) (0.050,3.688) (0.000,3.689) (-0.050,3.688) (-0.100,3.687) (-0.150,3.684) (-0.200,3.680) (-0.250,3.676) (-0.300,3.670) (-0.350,3.664) (-0.400,3.656) (-0.450,3.647) (-0.500,3.637) (-0.550,3.626) (-0.600,3.614) (-0.650,3.601) (-0.700,3.587) (-0.750,3.571) (-0.800,3.554) (-0.850,3.536) (-0.900,3.517) (-0.950,3.497) (-1.000,3.475) (-1.050,3.452) (-1.100,3.427) (-1.150,3.401) (-1.197,3.375) (-1.250,3.344) (-1.300,3.313) (-1.350,3.280) (-1.394,3.250) (-1.450,3.209) (-1.494,3.175) (-1.550,3.129) (-1.600,3.085) (-1.638,3.050) (-1.675,3.015) (-1.714,2.975) (-1.750,2.937) (-1.800,2.881) (-1.846,2.825) (-1.885,2.775) (-1.922,2.725) (-1.956,2.675) (-1.989,2.625) (-2.019,2.575) (-2.047,2.525) (-2.074,2.475) (-2.099,2.425) (-2.122,2.375) (-2.144,2.325) (-2.164,2.275) (-2.182,2.225) (-2.199,2.175) (-2.215,2.125) (-2.228,2.075) (-2.241,2.025) (-2.252,1.975) (-2.261,1.925) (-2.269,1.875) (-2.276,1.825) (-2.281,1.775) (-2.285,1.725) (-2.288,1.675) (-2.288,1.625) (-2.288,1.575) (-2.286,1.525) (-2.282,1.475) (-2.277,1.425) (-2.271,1.375) (-2.263,1.325) (-2.253,1.275) (-2.242,1.225) (-2.228,1.175) (-2.213,1.125) (-2.197,1.075) (-2.178,1.025) (-2.157,0.975) (-2.134,0.925) (-2.109,0.875) (-2.082,0.825) (-2.051,0.775) (-2.018,0.725) (-1.982,0.675) (-1.950,0.634) (-1.900,0.575) (-1.853,0.525) (-1.801,0.475) (-1.750,0.430) (-1.700,0.390) (-1.650,0.353) (-1.600,0.319) (-1.550,0.288) (-1.500,0.259) (-1.450,0.233) (-1.400,0.208) (-1.350,0.186) (-1.300,0.165)}
\EQ{(1.300,0.191) (1.350,0.215) (1.400,0.241) (1.450,0.270) (1.499,0.300) (1.550,0.335) (1.600,0.372) (1.650,0.414) (1.691,0.450) (1.742,0.500) (1.788,0.550) (1.825,0.594) (1.868,0.650) (1.900,0.696) (1.934,0.750) (1.963,0.800) (1.989,0.850) (2.013,0.900) (2.034,0.950) (2.053,1.000) (2.070,1.050) (2.085,1.100) (2.098,1.150) (2.109,1.200) (2.118,1.250) (2.126,1.300) (2.132,1.350) (2.137,1.400) (2.139,1.450) (2.141,1.500) (2.140,1.550) (2.138,1.600) (2.135,1.650) (2.130,1.700) (2.123,1.750) (2.115,1.800) (2.106,1.850) (2.095,1.900) (2.082,1.950) (2.068,2.000) (2.052,2.050) (2.035,2.100) (2.016,2.150) (1.995,2.200) (1.972,2.250) (1.950,2.296) (1.925,2.344) (1.900,2.389) (1.864,2.450) (1.831,2.500) (1.800,2.546) (1.760,2.600) (1.725,2.645) (1.679,2.700) (1.634,2.750) (1.600,2.786) (1.562,2.825) (1.525,2.860) (1.481,2.900) (1.425,2.947) (1.375,2.987) (1.325,3.025) (1.275,3.060) (1.225,3.093) (1.175,3.124) (1.132,3.150) (1.087,3.175) (1.040,3.200) (0.990,3.225) (0.936,3.250) (0.879,3.275) (0.825,3.297) (0.775,3.315) (0.725,3.333) (0.675,3.349) (0.625,3.364) (0.575,3.377) (0.525,3.390) (0.475,3.401) (0.425,3.411) (0.375,3.420) (0.325,3.427) (0.275,3.434) (0.225,3.439) (0.175,3.444) (0.125,3.447) (0.075,3.449) (0.025,3.450) (-0.037,3.450) (-0.100,3.448) (-0.150,3.445) (-0.200,3.442) (-0.250,3.437) (-0.300,3.431) (-0.350,3.424) (-0.400,3.415) (-0.450,3.406) (-0.500,3.395) (-0.550,3.384) (-0.600,3.371) (-0.650,3.356) (-0.700,3.341) (-0.748,3.325) (-0.800,3.306) (-0.850,3.287) (-0.900,3.266) (-0.950,3.244) (-1.000,3.220) (-1.050,3.195) (-1.100,3.168) (-1.150,3.139) (-1.200,3.109) (-1.250,3.077) (-1.300,3.043) (-1.350,3.006) (-1.391,2.975) (-1.450,2.927) (-1.500,2.883) (-1.550,2.836) (-1.587,2.800) (-1.625,2.760) (-1.675,2.705) (-1.721,2.650) (-1.760,2.600) (-1.797,2.550) (-1.831,2.500) (-1.864,2.450) (-1.894,2.400) (-1.922,2.350) (-1.948,2.300) (-1.972,2.250) (-1.995,2.200) (-2.016,2.150) (-2.035,2.100) (-2.052,2.050) (-2.068,2.000) (-2.082,1.950) (-2.095,1.900) (-2.106,1.850) (-2.115,1.800) (-2.123,1.750) (-2.130,1.700) (-2.135,1.650) (-2.138,1.600) (-2.140,1.550) (-2.141,1.500) (-2.139,1.450) (-2.137,1.400) (-2.132,1.350) (-2.126,1.300) (-2.118,1.250) (-2.109,1.200) (-2.098,1.150) (-2.085,1.100) (-2.070,1.050) (-2.053,1.000) (-2.034,0.950) (-2.013,0.900) (-1.989,0.850) (-1.963,0.800) (-1.934,0.750) (-1.903,0.700) (-1.868,0.650) (-1.830,0.600) (-1.788,0.550) (-1.750,0.508) (-1.700,0.459) (-1.650,0.414) (-1.603,0.375) (-1.550,0.335) (-1.500,0.301) (-1.450,0.270) (-1.400,0.241) (-1.350,0.215) (-1.300,0.191)}
\EQ{(1.300,0.216) (1.350,0.244) (1.400,0.274) (1.450,0.308) (1.500,0.345) (1.550,0.385) (1.595,0.425) (1.646,0.475) (1.692,0.525) (1.725,0.565) (1.770,0.625) (1.800,0.669) (1.834,0.725) (1.862,0.775) (1.887,0.825) (1.909,0.875) (1.929,0.925) (1.947,0.975) (1.962,1.025) (1.975,1.075) (1.987,1.125) (1.997,1.175) (2.004,1.225) (2.011,1.275) (2.015,1.325) (2.017,1.375) (2.018,1.425) (2.017,1.475) (2.015,1.525) (2.011,1.575) (2.005,1.625) (1.998,1.675) (1.989,1.725) (1.979,1.775) (1.967,1.825) (1.953,1.875) (1.937,1.925) (1.920,1.975) (1.901,2.025) (1.880,2.075) (1.858,2.125) (1.833,2.175) (1.807,2.225) (1.778,2.275) (1.750,2.321) (1.714,2.375) (1.679,2.425) (1.641,2.475) (1.601,2.525) (1.557,2.575) (1.510,2.625) (1.475,2.661) (1.434,2.700) (1.378,2.750) (1.325,2.794) (1.285,2.825) (1.225,2.869) (1.180,2.900) (1.125,2.935) (1.075,2.965) (1.025,2.993) (0.975,3.020) (0.925,3.045) (0.875,3.068) (0.825,3.089) (0.775,3.109) (0.725,3.128) (0.675,3.145) (0.625,3.161) (0.576,3.175) (0.525,3.188) (0.476,3.200) (0.425,3.211) (0.375,3.220) (0.325,3.229) (0.275,3.236) (0.225,3.241) (0.175,3.246) (0.125,3.249) (0.075,3.252) (0.025,3.253) (-0.050,3.252) (-0.100,3.251) (-0.150,3.248) (-0.200,3.244) (-0.250,3.239) (-0.300,3.232) (-0.350,3.225) (-0.400,3.216) (-0.450,3.206) (-0.500,3.195) (-0.550,3.182) (-0.600,3.168) (-0.650,3.153) (-0.700,3.137) (-0.750,3.119) (-0.798,3.100) (-0.850,3.079) (-0.900,3.056) (-0.950,3.032) (-1.000,3.007) (-1.050,2.980) (-1.100,2.950) (-1.150,2.919) (-1.200,2.886) (-1.250,2.851) (-1.300,2.814) (-1.348,2.775) (-1.400,2.731) (-1.450,2.685) (-1.486,2.650) (-1.525,2.610) (-1.575,2.555) (-1.621,2.500) (-1.661,2.450) (-1.697,2.400) (-1.731,2.350) (-1.763,2.300) (-1.793,2.250) (-1.820,2.200) (-1.846,2.150) (-1.869,2.100) (-1.891,2.050) (-1.911,2.000) (-1.929,1.950) (-1.945,1.900) (-1.960,1.850) (-1.973,1.800) (-1.984,1.750) (-1.994,1.700) (-2.002,1.650) (-2.008,1.600) (-2.013,1.550) (-2.017,1.500) (-2.018,1.450) (-2.018,1.400) (-2.016,1.350) (-2.013,1.300) (-2.008,1.250) (-2.001,1.200) (-1.992,1.150) (-1.982,1.100) (-1.969,1.050) (-1.955,1.000) (-1.938,0.950) (-1.919,0.900) (-1.898,0.850) (-1.875,0.801) (-1.850,0.753) (-1.825,0.710) (-1.787,0.650) (-1.752,0.600) (-1.713,0.550) (-1.675,0.506) (-1.625,0.454) (-1.575,0.407) (-1.525,0.364) (-1.475,0.326) (-1.425,0.291) (-1.375,0.259) (-1.325,0.230) (-1.300,0.216)}
\EQ{(1.300,0.243) (1.350,0.275) (1.400,0.310) (1.450,0.348) (1.500,0.391) (1.550,0.439) (1.600,0.493) (1.647,0.550) (1.684,0.600) (1.717,0.650) (1.747,0.700) (1.774,0.750) (1.798,0.800) (1.819,0.850) (1.838,0.900) (1.855,0.950) (1.869,1.000) (1.881,1.050) (1.892,1.100) (1.900,1.150) (1.906,1.200) (1.911,1.250) (1.914,1.300) (1.915,1.350) (1.914,1.400) (1.911,1.450) (1.907,1.500) (1.901,1.550) (1.894,1.600) (1.885,1.650) (1.873,1.700) (1.861,1.750) (1.846,1.800) (1.830,1.850) (1.812,1.900) (1.792,1.950) (1.770,2.000) (1.746,2.050) (1.720,2.100) (1.692,2.150) (1.661,2.200) (1.629,2.250) (1.593,2.300) (1.555,2.350) (1.515,2.400) (1.475,2.445) (1.425,2.498) (1.375,2.548) (1.325,2.593) (1.275,2.636) (1.226,2.675) (1.175,2.713) (1.125,2.747) (1.083,2.775) (1.025,2.810) (0.975,2.838) (0.925,2.865) (0.875,2.890) (0.825,2.913) (0.775,2.934) (0.725,2.954) (0.675,2.972) (0.625,2.989) (0.575,3.004) (0.525,3.018) (0.475,3.030) (0.425,3.042) (0.375,3.052) (0.325,3.060) (0.275,3.068) (0.225,3.074) (0.175,3.079) (0.125,3.082) (0.075,3.085) (0.025,3.086) (-0.050,3.086) (-0.100,3.084) (-0.150,3.081) (-0.200,3.076) (-0.250,3.071) (-0.300,3.064) (-0.350,3.056) (-0.400,3.047) (-0.450,3.036) (-0.500,3.024) (-0.550,3.011) (-0.600,2.996) (-0.650,2.980) (-0.700,2.963) (-0.750,2.944) (-0.800,2.923) (-0.850,2.901) (-0.900,2.877) (-0.950,2.852) (-0.999,2.825) (-1.050,2.795) (-1.100,2.764) (-1.150,2.730) (-1.192,2.700) (-1.250,2.656) (-1.300,2.615) (-1.345,2.575) (-1.398,2.525) (-1.448,2.475) (-1.493,2.425) (-1.535,2.375) (-1.575,2.325) (-1.611,2.275) (-1.645,2.225) (-1.675,2.178) (-1.706,2.125) (-1.733,2.075) (-1.758,2.025) (-1.781,1.975) (-1.802,1.925) (-1.821,1.875) (-1.838,1.825) (-1.854,1.775) (-1.867,1.725) (-1.879,1.675) (-1.889,1.625) (-1.898,1.575) (-1.905,1.525) (-1.910,1.475) (-1.913,1.425) (-1.915,1.375) (-1.914,1.325) (-1.913,1.275) (-1.909,1.225) (-1.903,1.175) (-1.896,1.125) (-1.887,1.075) (-1.875,1.025) (-1.862,0.975) (-1.847,0.925) (-1.829,0.875) (-1.809,0.825) (-1.786,0.775) (-1.761,0.725) (-1.733,0.675) (-1.701,0.625) (-1.666,0.575) (-1.627,0.525) (-1.584,0.475) (-1.536,0.425) (-1.482,0.375) (-1.425,0.328) (-1.375,0.292) (-1.325,0.258) (-1.300,0.243)}
\EQ{(1.288,0.262) (1.342,0.300) (1.400,0.347) (1.450,0.391) (1.500,0.442) (1.550,0.499) (1.590,0.550) (1.625,0.600) (1.656,0.650) (1.684,0.700) (1.708,0.750) (1.730,0.800) (1.750,0.850) (1.767,0.900) (1.781,0.950) (1.794,1.000) (1.804,1.050) (1.812,1.100) (1.818,1.150) (1.823,1.200) (1.825,1.250) (1.825,1.300) (1.824,1.350) (1.821,1.400) (1.816,1.450) (1.810,1.500) (1.801,1.550) (1.791,1.600) (1.779,1.650) (1.765,1.700) (1.750,1.748) (1.732,1.800) (1.712,1.850) (1.691,1.900) (1.667,1.950) (1.641,2.000) (1.613,2.050) (1.583,2.100) (1.550,2.149) (1.514,2.200) (1.476,2.250) (1.434,2.300) (1.400,2.339) (1.366,2.375) (1.325,2.416) (1.289,2.450) (1.250,2.484) (1.201,2.525) (1.150,2.565) (1.101,2.600) (1.050,2.634) (1.000,2.666) (0.950,2.695) (0.900,2.722) (0.850,2.748) (0.800,2.771) (0.750,2.793) (0.700,2.813) (0.650,2.831) (0.600,2.848) (0.550,2.864) (0.500,2.878) (0.450,2.890) (0.400,2.901) (0.350,2.911) (0.300,2.919) (0.250,2.926) (0.200,2.932) (0.150,2.937) (0.100,2.940) (0.050,2.942) (0.000,2.943) (-0.050,2.942) (-0.100,2.940) (-0.150,2.937) (-0.200,2.932) (-0.250,2.926) (-0.300,2.919) (-0.350,2.911) (-0.400,2.901) (-0.450,2.890) (-0.500,2.878) (-0.550,2.864) (-0.600,2.848) (-0.650,2.831) (-0.700,2.813) (-0.750,2.793) (-0.800,2.771) (-0.850,2.748) (-0.900,2.722) (-0.950,2.695) (-1.000,2.666) (-1.050,2.634) (-1.100,2.601) (-1.150,2.565) (-1.200,2.526) (-1.250,2.484) (-1.289,2.450) (-1.325,2.416) (-1.366,2.375) (-1.400,2.339) (-1.434,2.300) (-1.475,2.251) (-1.514,2.200) (-1.550,2.150) (-1.583,2.100) (-1.613,2.050) (-1.641,2.000) (-1.667,1.950) (-1.691,1.900) (-1.712,1.850) (-1.732,1.800) (-1.749,1.750) (-1.765,1.700) (-1.779,1.650) (-1.791,1.600) (-1.801,1.550) (-1.810,1.500) (-1.816,1.450) (-1.821,1.400) (-1.824,1.350) (-1.825,1.300) (-1.825,1.250) (-1.823,1.200) (-1.818,1.150) (-1.812,1.100) (-1.804,1.050) (-1.794,1.000) (-1.781,0.950) (-1.767,0.900) (-1.750,0.851) (-1.730,0.800) (-1.708,0.750) (-1.684,0.700) (-1.656,0.650) (-1.625,0.601) (-1.590,0.550) (-1.551,0.500) (-1.508,0.450) (-1.459,0.400) (-1.404,0.350) (-1.350,0.307) (-1.307,0.275) (-1.288,0.262)}
\EQ{(1.275,0.280) (1.325,0.319) (1.375,0.362) (1.414,0.400) (1.450,0.438) (1.500,0.497) (1.539,0.550) (1.572,0.600) (1.600,0.648) (1.625,0.696) (1.650,0.750) (1.670,0.800) (1.688,0.850) (1.703,0.900) (1.716,0.950) (1.726,1.000) (1.734,1.050) (1.741,1.100) (1.745,1.150) (1.747,1.200) (1.748,1.250) (1.746,1.300) (1.743,1.350) (1.738,1.400) (1.731,1.450) (1.722,1.500) (1.711,1.550) (1.698,1.600) (1.683,1.650) (1.667,1.700) (1.648,1.750) (1.627,1.800) (1.604,1.850) (1.579,1.900) (1.552,1.950) (1.525,1.995) (1.489,2.050) (1.454,2.100) (1.416,2.150) (1.375,2.200) (1.330,2.250) (1.282,2.300) (1.229,2.350) (1.175,2.397) (1.125,2.437) (1.075,2.475) (1.025,2.509) (0.975,2.542) (0.925,2.571) (0.875,2.599) (0.825,2.625) (0.775,2.649) (0.725,2.671) (0.675,2.691) (0.625,2.710) (0.575,2.727) (0.525,2.742) (0.475,2.756) (0.425,2.768) (0.375,2.779) (0.325,2.789) (0.275,2.797) (0.225,2.804) (0.175,2.809) (0.125,2.813) (0.075,2.816) (0.025,2.817) (-0.050,2.817) (-0.100,2.815) (-0.150,2.811) (-0.200,2.807) (-0.250,2.800) (-0.300,2.793) (-0.350,2.784) (-0.400,2.774) (-0.450,2.762) (-0.500,2.749) (-0.550,2.735) (-0.600,2.719) (-0.650,2.701) (-0.700,2.681) (-0.750,2.660) (-0.800,2.637) (-0.850,2.612) (-0.900,2.586) (-0.950,2.557) (-1.000,2.526) (-1.050,2.492) (-1.100,2.456) (-1.141,2.425) (-1.200,2.376) (-1.250,2.331) (-1.300,2.282) (-1.350,2.228) (-1.396,2.175) (-1.436,2.125) (-1.472,2.075) (-1.506,2.025) (-1.537,1.975) (-1.566,1.925) (-1.592,1.875) (-1.616,1.825) (-1.638,1.775) (-1.658,1.725) (-1.675,1.675) (-1.691,1.625) (-1.705,1.575) (-1.716,1.525) (-1.726,1.475) (-1.734,1.425) (-1.741,1.375) (-1.745,1.325) (-1.747,1.275) (-1.748,1.225) (-1.746,1.175) (-1.743,1.125) (-1.738,1.075) (-1.730,1.025) (-1.721,0.975) (-1.709,0.925) (-1.695,0.875) (-1.679,0.825) (-1.660,0.775) (-1.639,0.725) (-1.614,0.675) (-1.587,0.625) (-1.556,0.575) (-1.525,0.530) (-1.482,0.475) (-1.438,0.425) (-1.400,0.386) (-1.362,0.350) (-1.301,0.300) (-1.275,0.280)}
\EQ{(1.275,0.308) (1.323,0.350) (1.374,0.400) (1.419,0.450) (1.450,0.489) (1.493,0.550) (1.524,0.600) (1.550,0.647) (1.575,0.698) (1.597,0.750) (1.616,0.800) (1.632,0.850) (1.645,0.900) (1.656,0.950) (1.665,1.000) (1.672,1.050) (1.676,1.100) (1.679,1.150) (1.679,1.200) (1.678,1.250) (1.674,1.300) (1.669,1.350) (1.662,1.400) (1.652,1.450) (1.641,1.500) (1.628,1.550) (1.613,1.600) (1.595,1.650) (1.576,1.700) (1.554,1.750) (1.530,1.800) (1.504,1.850) (1.475,1.900) (1.444,1.950) (1.410,2.000) (1.375,2.047) (1.333,2.100) (1.289,2.150) (1.250,2.191) (1.200,2.240) (1.162,2.275) (1.102,2.325) (1.050,2.365) (1.001,2.400) (0.950,2.434) (0.900,2.464) (0.850,2.492) (0.800,2.518) (0.750,2.543) (0.700,2.565) (0.650,2.585) (0.600,2.604) (0.550,2.621) (0.500,2.636) (0.450,2.650) (0.400,2.662) (0.350,2.672) (0.300,2.682) (0.250,2.689) (0.200,2.696) (0.150,2.700) (0.100,2.704) (0.050,2.706) (0.000,2.707) (-0.050,2.706) (-0.100,2.704) (-0.150,2.700) (-0.200,2.696) (-0.250,2.689) (-0.300,2.682) (-0.350,2.672) (-0.400,2.662) (-0.450,2.650) (-0.500,2.636) (-0.550,2.621) (-0.600,2.604) (-0.650,2.585) (-0.700,2.565) (-0.750,2.543) (-0.800,2.518) (-0.850,2.492) (-0.900,2.464) (-0.950,2.434) (-1.000,2.401) (-1.050,2.365) (-1.100,2.327) (-1.150,2.285) (-1.189,2.250) (-1.241,2.200) (-1.289,2.150) (-1.325,2.109) (-1.373,2.050) (-1.410,2.000) (-1.444,1.950) (-1.475,1.900) (-1.504,1.850) (-1.530,1.800) (-1.554,1.750) (-1.575,1.702) (-1.595,1.650) (-1.613,1.600) (-1.628,1.550) (-1.641,1.500) (-1.652,1.450) (-1.662,1.400) (-1.669,1.350) (-1.674,1.300) (-1.678,1.250) (-1.679,1.200) (-1.679,1.150) (-1.676,1.100) (-1.672,1.050) (-1.665,1.000) (-1.656,0.950) (-1.645,0.900) (-1.632,0.850) (-1.616,0.800) (-1.597,0.750) (-1.576,0.700) (-1.552,0.650) (-1.525,0.601) (-1.493,0.550) (-1.458,0.500) (-1.425,0.458) (-1.375,0.401) (-1.325,0.352) (-1.275,0.308)}
\EQ{(1.275,0.337) (1.314,0.375) (1.350,0.414) (1.400,0.475) (1.436,0.525) (1.467,0.575) (1.495,0.625) (1.519,0.675) (1.540,0.725) (1.559,0.775) (1.574,0.825) (1.587,0.875) (1.598,0.925) (1.606,0.975) (1.612,1.025) (1.616,1.075) (1.618,1.125) (1.618,1.175) (1.616,1.225) (1.611,1.275) (1.605,1.325) (1.597,1.375) (1.586,1.425) (1.574,1.475) (1.559,1.525) (1.543,1.575) (1.524,1.625) (1.503,1.675) (1.479,1.725) (1.454,1.775) (1.425,1.825) (1.394,1.875) (1.360,1.925) (1.325,1.973) (1.283,2.025) (1.250,2.063) (1.216,2.100) (1.175,2.141) (1.125,2.187) (1.081,2.225) (1.025,2.269) (0.983,2.300) (0.925,2.339) (0.875,2.370) (0.825,2.398) (0.775,2.425) (0.725,2.449) (0.675,2.471) (0.625,2.492) (0.575,2.510) (0.525,2.527) (0.475,2.542) (0.425,2.555) (0.375,2.567) (0.325,2.578) (0.275,2.586) (0.225,2.594) (0.175,2.599) (0.125,2.604) (0.075,2.607) (0.025,2.608) (-0.050,2.608) (-0.100,2.605) (-0.150,2.602) (-0.200,2.597) (-0.250,2.590) (-0.300,2.582) (-0.350,2.573) (-0.400,2.562) (-0.450,2.549) (-0.500,2.535) (-0.550,2.519) (-0.600,2.501) (-0.650,2.482) (-0.700,2.460) (-0.750,2.437) (-0.800,2.412) (-0.850,2.384) (-0.900,2.355) (-0.946,2.325) (-1.000,2.288) (-1.050,2.250) (-1.100,2.209) (-1.139,2.175) (-1.191,2.125) (-1.239,2.075) (-1.275,2.034) (-1.323,1.975) (-1.360,1.925) (-1.394,1.875) (-1.425,1.825) (-1.450,1.781) (-1.475,1.734) (-1.500,1.681) (-1.524,1.625) (-1.543,1.575) (-1.559,1.525) (-1.574,1.475) (-1.586,1.425) (-1.597,1.375) (-1.605,1.325) (-1.611,1.275) (-1.616,1.225) (-1.618,1.175) (-1.618,1.125) (-1.616,1.075) (-1.612,1.025) (-1.606,0.975) (-1.598,0.925) (-1.587,0.875) (-1.574,0.825) (-1.559,0.775) (-1.540,0.725) (-1.519,0.675) (-1.495,0.625) (-1.467,0.575) (-1.436,0.525) (-1.400,0.475) (-1.360,0.425) (-1.325,0.387) (-1.289,0.350) (-1.275,0.337)}
\EQ{(1.266,0.359) (1.305,0.400) (1.346,0.450) (1.382,0.500) (1.414,0.550) (1.442,0.600) (1.467,0.650) (1.488,0.700) (1.506,0.750) (1.522,0.800) (1.534,0.850) (1.545,0.900) (1.553,0.950) (1.559,1.000) (1.562,1.050) (1.563,1.100) (1.563,1.150) (1.560,1.200) (1.555,1.250) (1.547,1.300) (1.538,1.350) (1.527,1.400) (1.513,1.450) (1.498,1.500) (1.480,1.550) (1.460,1.600) (1.437,1.650) (1.412,1.700) (1.385,1.750) (1.354,1.800) (1.325,1.844) (1.285,1.900) (1.250,1.944) (1.201,2.000) (1.154,2.050) (1.101,2.100) (1.050,2.145) (1.000,2.185) (0.950,2.222) (0.908,2.250) (0.850,2.287) (0.800,2.315) (0.750,2.342) (0.700,2.366) (0.650,2.388) (0.600,2.409) (0.550,2.427) (0.500,2.443) (0.450,2.458) (0.400,2.471) (0.350,2.483) (0.300,2.493) (0.250,2.501) (0.200,2.508) (0.150,2.513) (0.100,2.517) (0.050,2.519) (0.000,2.520) (-0.050,2.519) (-0.100,2.517) (-0.150,2.513) (-0.200,2.508) (-0.250,2.501) (-0.300,2.493) (-0.350,2.483) (-0.400,2.471) (-0.450,2.458) (-0.500,2.443) (-0.550,2.427) (-0.600,2.409) (-0.650,2.388) (-0.700,2.366) (-0.750,2.342) (-0.800,2.315) (-0.850,2.287) (-0.900,2.256) (-0.945,2.225) (-1.000,2.185) (-1.044,2.150) (-1.100,2.101) (-1.150,2.054) (-1.200,2.001) (-1.245,1.950) (-1.285,1.900) (-1.321,1.850) (-1.350,1.807) (-1.385,1.750) (-1.412,1.700) (-1.437,1.650) (-1.460,1.600) (-1.480,1.550) (-1.498,1.500) (-1.513,1.450) (-1.527,1.400) (-1.538,1.350) (-1.547,1.300) (-1.555,1.250) (-1.560,1.200) (-1.563,1.150) (-1.563,1.100) (-1.562,1.050) (-1.559,1.000) (-1.553,0.950) (-1.545,0.900) (-1.534,0.850) (-1.522,0.800) (-1.506,0.750) (-1.488,0.700) (-1.467,0.650) (-1.442,0.600) (-1.414,0.550) (-1.382,0.500) (-1.350,0.456) (-1.305,0.400) (-1.266,0.359)}
\EQ{(1.251,0.374) (1.295,0.425) (1.325,0.465) (1.365,0.525) (1.393,0.575) (1.418,0.625) (1.440,0.675) (1.458,0.725) (1.473,0.775) (1.486,0.825) (1.496,0.875) (1.504,0.925) (1.510,0.975) (1.513,1.025) (1.514,1.075) (1.512,1.125) (1.509,1.175) (1.503,1.225) (1.495,1.275) (1.486,1.325) (1.473,1.375) (1.459,1.425) (1.442,1.475) (1.423,1.525) (1.402,1.575) (1.378,1.625) (1.352,1.675) (1.325,1.721) (1.290,1.775) (1.255,1.825) (1.216,1.875) (1.175,1.923) (1.127,1.975) (1.075,2.025) (1.025,2.070) (0.975,2.110) (0.925,2.147) (0.875,2.182) (0.825,2.213) (0.775,2.242) (0.725,2.268) (0.675,2.292) (0.625,2.314) (0.575,2.334) (0.525,2.352) (0.475,2.369) (0.425,2.383) (0.375,2.396) (0.325,2.407) (0.275,2.416) (0.225,2.424) (0.175,2.430) (0.125,2.435) (0.075,2.438) (0.025,2.440) (-0.050,2.439) (-0.100,2.437) (-0.150,2.433) (-0.200,2.427) (-0.250,2.420) (-0.300,2.412) (-0.350,2.402) (-0.400,2.390) (-0.450,2.376) (-0.500,2.361) (-0.550,2.344) (-0.599,2.325) (-0.650,2.304) (-0.700,2.280) (-0.750,2.255) (-0.800,2.228) (-0.846,2.200) (-0.900,2.165) (-0.950,2.129) (-1.000,2.090) (-1.048,2.050) (-1.100,2.002) (-1.150,1.951) (-1.195,1.900) (-1.236,1.850) (-1.273,1.800) (-1.307,1.750) (-1.337,1.700) (-1.365,1.650) (-1.390,1.600) (-1.413,1.550) (-1.433,1.500) (-1.450,1.453) (-1.466,1.400) (-1.480,1.350) (-1.491,1.300) (-1.500,1.250) (-1.506,1.200) (-1.511,1.150) (-1.513,1.100) (-1.514,1.050) (-1.512,1.000) (-1.507,0.950) (-1.501,0.900) (-1.492,0.850) (-1.480,0.800) (-1.466,0.750) (-1.449,0.700) (-1.429,0.650) (-1.406,0.600) (-1.379,0.550) (-1.350,0.502) (-1.314,0.450) (-1.275,0.401) (-1.251,0.374)}
\EQ{(1.250,0.405) (1.285,0.450) (1.318,0.500) (1.347,0.550) (1.373,0.600) (1.395,0.650) (1.413,0.700) (1.429,0.750) (1.442,0.800) (1.452,0.850) (1.460,0.900) (1.465,0.950) (1.468,1.000) (1.468,1.050) (1.467,1.100) (1.463,1.150) (1.457,1.200) (1.448,1.250) (1.438,1.300) (1.425,1.350) (1.410,1.400) (1.392,1.450) (1.372,1.500) (1.350,1.550) (1.325,1.600) (1.300,1.645) (1.266,1.700) (1.232,1.750) (1.200,1.793) (1.154,1.850) (1.109,1.900) (1.059,1.950) (1.004,2.000) (0.950,2.044) (0.909,2.075) (0.850,2.116) (0.800,2.147) (0.751,2.175) (0.704,2.200) (0.652,2.225) (0.600,2.248) (0.550,2.268) (0.500,2.285) (0.450,2.301) (0.400,2.315) (0.350,2.328) (0.300,2.338) (0.250,2.347) (0.200,2.354) (0.150,2.360) (0.100,2.364) (0.050,2.366) (0.000,2.367) (-0.050,2.366) (-0.100,2.364) (-0.150,2.360) (-0.200,2.354) (-0.250,2.347) (-0.300,2.338) (-0.350,2.328) (-0.400,2.315) (-0.450,2.301) (-0.500,2.285) (-0.550,2.268) (-0.600,2.248) (-0.650,2.226) (-0.700,2.202) (-0.750,2.176) (-0.795,2.150) (-0.850,2.116) (-0.900,2.081) (-0.942,2.050) (-1.000,2.003) (-1.050,1.958) (-1.100,1.909) (-1.150,1.854) (-1.195,1.800) (-1.225,1.760) (-1.266,1.700) (-1.297,1.650) (-1.325,1.600) (-1.350,1.550) (-1.372,1.500) (-1.392,1.450) (-1.410,1.400) (-1.425,1.350) (-1.438,1.300) (-1.448,1.250) (-1.457,1.200) (-1.463,1.150) (-1.467,1.100) (-1.468,1.050) (-1.468,1.000) (-1.465,0.950) (-1.460,0.900) (-1.452,0.850) (-1.442,0.800) (-1.429,0.750) (-1.413,0.700) (-1.395,0.650) (-1.373,0.600) (-1.350,0.555) (-1.325,0.512) (-1.285,0.450) (-1.250,0.405)}
\EQ{(1.244,0.431) (1.274,0.475) (1.300,0.518) (1.325,0.565) (1.350,0.619) (1.372,0.675) (1.388,0.725) (1.400,0.775) (1.411,0.825) (1.419,0.875) (1.424,0.925) (1.427,0.975) (1.427,1.025) (1.425,1.075) (1.421,1.125) (1.414,1.175) (1.406,1.225) (1.395,1.275) (1.381,1.325) (1.365,1.375) (1.347,1.425) (1.326,1.475) (1.303,1.525) (1.277,1.575) (1.250,1.621) (1.215,1.675) (1.180,1.725) (1.141,1.775) (1.100,1.822) (1.050,1.874) (1.000,1.922) (0.950,1.965) (0.906,2.000) (0.850,2.040) (0.800,2.073) (0.755,2.100) (0.709,2.125) (0.660,2.150) (0.605,2.175) (0.550,2.198) (0.500,2.216) (0.450,2.233) (0.400,2.247) (0.350,2.260) (0.300,2.271) (0.250,2.280) (0.200,2.287) (0.150,2.293) (0.100,2.297) (0.050,2.299) (0.000,2.300) (-0.050,2.299) (-0.100,2.297) (-0.150,2.293) (-0.200,2.287) (-0.250,2.280) (-0.300,2.271) (-0.350,2.260) (-0.400,2.247) (-0.450,2.233) (-0.500,2.216) (-0.550,2.198) (-0.600,2.177) (-0.650,2.155) (-0.700,2.130) (-0.750,2.103) (-0.796,2.075) (-0.850,2.040) (-0.900,2.004) (-0.950,1.965) (-0.997,1.925) (-1.049,1.875) (-1.097,1.825) (-1.141,1.775) (-1.175,1.731) (-1.215,1.675) (-1.248,1.625) (-1.275,1.578) (-1.300,1.531) (-1.325,1.478) (-1.347,1.425) (-1.365,1.375) (-1.381,1.325) (-1.395,1.275) (-1.406,1.225) (-1.414,1.175) (-1.421,1.125) (-1.425,1.075) (-1.427,1.025) (-1.427,0.975) (-1.424,0.925) (-1.419,0.875) (-1.411,0.825) (-1.400,0.775) (-1.388,0.725) (-1.372,0.675) (-1.352,0.625) (-1.330,0.575) (-1.304,0.525) (-1.275,0.477) (-1.244,0.431)}
\EQ{(1.229,0.446) (1.263,0.500) (1.290,0.550) (1.313,0.600) (1.333,0.650) (1.349,0.700) (1.362,0.750) (1.373,0.800) (1.381,0.850) (1.386,0.900) (1.389,0.950) (1.389,1.000) (1.387,1.050) (1.383,1.100) (1.376,1.150) (1.367,1.200) (1.355,1.250) (1.341,1.300) (1.325,1.349) (1.306,1.400) (1.284,1.450) (1.260,1.500) (1.233,1.550) (1.203,1.600) (1.175,1.642) (1.132,1.700) (1.100,1.740) (1.050,1.796) (1.000,1.846) (0.950,1.891) (0.909,1.925) (0.850,1.970) (0.806,2.000) (0.750,2.035) (0.700,2.063) (0.650,2.089) (0.600,2.112) (0.550,2.133) (0.500,2.152) (0.450,2.169) (0.400,2.184) (0.350,2.197) (0.300,2.208) (0.250,2.218) (0.200,2.225) (0.150,2.231) (0.100,2.236) (0.050,2.238) (0.000,2.239) (-0.050,2.238) (-0.100,2.236) (-0.150,2.231) (-0.200,2.225) (-0.250,2.218) (-0.300,2.208) (-0.350,2.197) (-0.400,2.184) (-0.450,2.169) (-0.500,2.152) (-0.550,2.133) (-0.600,2.112) (-0.650,2.089) (-0.700,2.063) (-0.750,2.035) (-0.800,2.004) (-0.843,1.975) (-0.900,1.932) (-0.940,1.900) (-0.996,1.850) (-1.046,1.800) (-1.091,1.750) (-1.125,1.709) (-1.169,1.650) (-1.200,1.604) (-1.233,1.550) (-1.260,1.500) (-1.284,1.450) (-1.306,1.400) (-1.325,1.350) (-1.341,1.300) (-1.355,1.250) (-1.367,1.200) (-1.376,1.150) (-1.383,1.100) (-1.387,1.050) (-1.389,1.000) (-1.389,0.950) (-1.386,0.900) (-1.381,0.850) (-1.373,0.800) (-1.362,0.750) (-1.349,0.700) (-1.333,0.650) (-1.313,0.600) (-1.290,0.550) (-1.263,0.500) (-1.232,0.450) (-1.229,0.446)}
\EQ{(1.224,0.476) (1.250,0.522) (1.275,0.574) (1.296,0.625) (1.313,0.675) (1.326,0.725) (1.337,0.775) (1.345,0.825) (1.350,0.875) (1.353,0.925) (1.354,0.975) (1.352,1.025) (1.347,1.075) (1.340,1.125) (1.331,1.175) (1.319,1.225) (1.305,1.275) (1.288,1.325) (1.269,1.375) (1.246,1.425) (1.221,1.475) (1.194,1.525) (1.162,1.575) (1.128,1.625) (1.089,1.675) (1.050,1.721) (1.000,1.774) (0.950,1.822) (0.900,1.865) (0.856,1.900) (0.800,1.940) (0.750,1.972) (0.702,2.000) (0.656,2.025) (0.604,2.050) (0.550,2.074) (0.500,2.093) (0.450,2.111) (0.400,2.126) (0.350,2.139) (0.300,2.151) (0.250,2.161) (0.200,2.168) (0.150,2.174) (0.100,2.179) (0.050,2.181) (0.000,2.182) (-0.050,2.181) (-0.100,2.179) (-0.150,2.174) (-0.200,2.168) (-0.250,2.161) (-0.300,2.151) (-0.350,2.139) (-0.400,2.126) (-0.450,2.111) (-0.500,2.093) (-0.547,2.075) (-0.600,2.052) (-0.650,2.028) (-0.700,2.001) (-0.745,1.975) (-0.800,1.940) (-0.850,1.904) (-0.900,1.865) (-0.947,1.825) (-0.999,1.775) (-1.047,1.725) (-1.089,1.675) (-1.125,1.629) (-1.162,1.575) (-1.194,1.525) (-1.221,1.475) (-1.246,1.425) (-1.269,1.375) (-1.288,1.325) (-1.305,1.275) (-1.319,1.225) (-1.331,1.175) (-1.340,1.125) (-1.347,1.075) (-1.352,1.025) (-1.354,0.975) (-1.353,0.925) (-1.350,0.875) (-1.345,0.825) (-1.337,0.775) (-1.326,0.725) (-1.313,0.675) (-1.296,0.625) (-1.275,0.575) (-1.252,0.525) (-1.225,0.477)}

    \draw[mmaEQsolo] (-10,0) -- (-1.2,0) (1.2,0) -- (10,0);}
  \marcoMMA{10}{5}{$x$ (u.\,arb.)}{$z$ (u.\,arb.)}
\end{tikzpicture}
\end{minipage}\hfill
\begin{minipage}[t]{0.325\textwidth}\centering
\rotulo{Término cuadrupolar ($\ell=2$)}
\begin{tikzpicture}[x=0.198cm, y=0.198cm]
  \panelMMA{10}{white}{0}{% Cuadrupolo puntual - equipotenciales
% Generado por scripts/generar_datos.py (no editar a mano).
\EQ{(-0.220,-0.155) (-0.275,-0.195) (-0.319,-0.225) (-0.375,-0.266) (-0.425,-0.300) (-0.475,-0.336) (-0.525,-0.371) (-0.568,-0.400) (-0.625,-0.441) (-0.675,-0.475) (-0.725,-0.510) (-0.775,-0.544) (-0.821,-0.575) (-0.875,-0.612) (-0.925,-0.646) (-0.969,-0.675) (-1.025,-0.713) (-1.075,-0.745) (-1.121,-0.775) (-1.175,-0.810) (-1.225,-0.841) (-1.275,-0.872) (-1.321,-0.900) (-1.375,-0.933) (-1.425,-0.963) (-1.475,-0.991) (-1.525,-1.020) (-1.575,-1.047) (-1.625,-1.074) (-1.674,-1.100) (-1.723,-1.125) (-1.774,-1.150) (-1.825,-1.175) (-1.875,-1.198) (-1.925,-1.220) (-1.975,-1.241) (-2.025,-1.262) (-2.075,-1.281) (-2.125,-1.300) (-2.175,-1.317) (-2.225,-1.334) (-2.275,-1.349) (-2.325,-1.364) (-2.375,-1.377) (-2.425,-1.389) (-2.475,-1.400) (-2.525,-1.410) (-2.575,-1.419) (-2.625,-1.426) (-2.675,-1.433) (-2.725,-1.438) (-2.775,-1.441) (-2.825,-1.444) (-2.875,-1.445) (-2.925,-1.445) (-2.975,-1.443) (-3.025,-1.440) (-3.075,-1.435) (-3.125,-1.429) (-3.175,-1.421) (-3.225,-1.412) (-3.275,-1.401) (-3.325,-1.388) (-3.375,-1.373) (-3.425,-1.357) (-3.475,-1.338) (-3.525,-1.318) (-3.575,-1.295) (-3.625,-1.270) (-3.675,-1.242) (-3.725,-1.212) (-3.775,-1.179) (-3.825,-1.142) (-3.875,-1.102) (-3.925,-1.058) (-3.975,-1.010) (-4.025,-0.956) (-4.072,-0.900) (-4.110,-0.850) (-4.145,-0.800) (-4.175,-0.753) (-4.206,-0.700) (-4.232,-0.650) (-4.256,-0.600) (-4.277,-0.550) (-4.297,-0.500) (-4.314,-0.450) (-4.329,-0.400) (-4.343,-0.350) (-4.354,-0.300) (-4.364,-0.250) (-4.371,-0.200) (-4.377,-0.150) (-4.382,-0.100) (-4.384,-0.050) (-4.385,0.000) (-4.384,0.050) (-4.382,0.100) (-4.377,0.150) (-4.371,0.200) (-4.364,0.250) (-4.354,0.300) (-4.343,0.350) (-4.329,0.400) (-4.314,0.450) (-4.297,0.500) (-4.277,0.550) (-4.256,0.600) (-4.232,0.650) (-4.206,0.700) (-4.177,0.750) (-4.150,0.792) (-4.110,0.850) (-4.075,0.896) (-4.030,0.950) (-3.985,1.000) (-3.934,1.050) (-3.878,1.100) (-3.825,1.142) (-3.780,1.175) (-3.725,1.212) (-3.675,1.242) (-3.625,1.270) (-3.575,1.295) (-3.525,1.318) (-3.475,1.338) (-3.425,1.357) (-3.375,1.373) (-3.325,1.388) (-3.275,1.401) (-3.225,1.412) (-3.175,1.421) (-3.125,1.429) (-3.075,1.435) (-3.025,1.440) (-2.975,1.443) (-2.925,1.445) (-2.875,1.445) (-2.825,1.444) (-2.775,1.441) (-2.725,1.438) (-2.675,1.433) (-2.625,1.426) (-2.575,1.419) (-2.525,1.410) (-2.475,1.400) (-2.425,1.389) (-2.375,1.377) (-2.325,1.364) (-2.275,1.349) (-2.225,1.334) (-2.175,1.317) (-2.126,1.300) (-2.075,1.281) (-2.025,1.262) (-1.975,1.241) (-1.925,1.220) (-1.875,1.198) (-1.826,1.175) (-1.775,1.151) (-1.725,1.126) (-1.675,1.100) (-1.627,1.075) (-1.580,1.050) (-1.535,1.025) (-1.491,1.000) (-1.447,0.975) (-1.400,0.948) (-1.350,0.918) (-1.300,0.888) (-1.250,0.857) (-1.200,0.825) (-1.150,0.794) (-1.100,0.762) (-1.050,0.729) (-1.007,0.700) (-0.950,0.663) (-0.900,0.629) (-0.858,0.600) (-0.800,0.561) (-0.750,0.527) (-0.700,0.493) (-0.650,0.458) (-0.604,0.425) (-0.550,0.388) (-0.500,0.353) (-0.450,0.318) (-0.400,0.283) (-0.354,0.250) (-0.300,0.213) (-0.250,0.177) (-0.220,0.155)}
\EQ{(0.206,0.144) (0.248,0.175) (0.300,0.213) (0.350,0.248) (0.400,0.283) (0.450,0.318) (0.496,0.350) (0.550,0.388) (0.600,0.423) (0.650,0.458) (0.700,0.493) (0.748,0.525) (0.800,0.561) (0.850,0.595) (0.895,0.625) (0.950,0.663) (1.000,0.696) (1.045,0.725) (1.100,0.762) (1.150,0.794) (1.200,0.825) (1.250,0.857) (1.300,0.888) (1.350,0.918) (1.400,0.948) (1.447,0.975) (1.491,1.000) (1.535,1.025) (1.580,1.050) (1.625,1.074) (1.674,1.100) (1.723,1.125) (1.774,1.150) (1.825,1.175) (1.875,1.198) (1.925,1.220) (1.975,1.241) (2.025,1.262) (2.075,1.281) (2.125,1.300) (2.175,1.317) (2.225,1.334) (2.275,1.349) (2.325,1.364) (2.375,1.377) (2.425,1.389) (2.475,1.400) (2.525,1.410) (2.575,1.419) (2.625,1.426) (2.675,1.433) (2.725,1.438) (2.775,1.441) (2.825,1.444) (2.875,1.445) (2.925,1.445) (2.975,1.443) (3.025,1.440) (3.075,1.435) (3.125,1.429) (3.175,1.421) (3.225,1.412) (3.275,1.401) (3.325,1.388) (3.375,1.373) (3.425,1.357) (3.475,1.338) (3.525,1.318) (3.575,1.295) (3.625,1.270) (3.675,1.242) (3.725,1.212) (3.775,1.179) (3.825,1.142) (3.875,1.102) (3.925,1.058) (3.975,1.010) (4.025,0.956) (4.072,0.900) (4.110,0.850) (4.145,0.800) (4.175,0.753) (4.206,0.700) (4.232,0.650) (4.256,0.600) (4.277,0.550) (4.297,0.500) (4.314,0.450) (4.329,0.400) (4.343,0.350) (4.354,0.300) (4.364,0.250) (4.371,0.200) (4.377,0.150) (4.382,0.100) (4.384,0.050) (4.385,0.000) (4.384,-0.050) (4.382,-0.100) (4.377,-0.150) (4.371,-0.200) (4.364,-0.250) (4.354,-0.300) (4.343,-0.350) (4.329,-0.400) (4.314,-0.450) (4.297,-0.500) (4.277,-0.550) (4.256,-0.600) (4.232,-0.650) (4.206,-0.700) (4.177,-0.750) (4.150,-0.792) (4.110,-0.850) (4.075,-0.896) (4.030,-0.950) (3.985,-1.000) (3.934,-1.050) (3.878,-1.100) (3.825,-1.142) (3.780,-1.175) (3.725,-1.212) (3.675,-1.242) (3.625,-1.270) (3.575,-1.295) (3.525,-1.318) (3.475,-1.338) (3.425,-1.357) (3.375,-1.373) (3.325,-1.388) (3.275,-1.401) (3.225,-1.412) (3.175,-1.421) (3.125,-1.429) (3.075,-1.435) (3.025,-1.440) (2.975,-1.443) (2.925,-1.445) (2.875,-1.445) (2.825,-1.444) (2.775,-1.441) (2.725,-1.438) (2.675,-1.433) (2.625,-1.426) (2.575,-1.419) (2.525,-1.410) (2.475,-1.400) (2.425,-1.389) (2.375,-1.377) (2.325,-1.364) (2.275,-1.349) (2.225,-1.334) (2.175,-1.317) (2.126,-1.300) (2.075,-1.281) (2.025,-1.262) (1.975,-1.241) (1.925,-1.220) (1.875,-1.198) (1.826,-1.175) (1.775,-1.151) (1.725,-1.126) (1.675,-1.100) (1.627,-1.075) (1.580,-1.050) (1.535,-1.025) (1.491,-1.000) (1.447,-0.975) (1.400,-0.948) (1.350,-0.918) (1.300,-0.888) (1.250,-0.857) (1.200,-0.825) (1.150,-0.794) (1.100,-0.762) (1.050,-0.729) (1.007,-0.700) (0.950,-0.663) (0.900,-0.629) (0.858,-0.600) (0.800,-0.561) (0.750,-0.527) (0.700,-0.493) (0.650,-0.458) (0.604,-0.425) (0.550,-0.388) (0.500,-0.353) (0.450,-0.318) (0.400,-0.283) (0.354,-0.250) (0.300,-0.213) (0.250,-0.177) (0.206,-0.144)}
\EQ{(-0.220,-0.155) (-0.275,-0.195) (-0.319,-0.225) (-0.375,-0.266) (-0.425,-0.300) (-0.475,-0.336) (-0.525,-0.371) (-0.568,-0.400) (-0.625,-0.441) (-0.675,-0.475) (-0.725,-0.510) (-0.775,-0.545) (-0.820,-0.575) (-0.875,-0.613) (-0.925,-0.647) (-0.968,-0.675) (-1.025,-0.714) (-1.075,-0.747) (-1.119,-0.775) (-1.175,-0.812) (-1.225,-0.844) (-1.275,-0.875) (-1.325,-0.906) (-1.375,-0.937) (-1.425,-0.967) (-1.475,-0.996) (-1.525,-1.025) (-1.569,-1.050) (-1.614,-1.075) (-1.660,-1.100) (-1.707,-1.125) (-1.755,-1.150) (-1.800,-1.173) (-1.850,-1.197) (-1.900,-1.221) (-1.950,-1.244) (-2.000,-1.266) (-2.050,-1.288) (-2.100,-1.308) (-2.150,-1.327) (-2.200,-1.346) (-2.250,-1.363) (-2.300,-1.380) (-2.350,-1.396) (-2.400,-1.410) (-2.450,-1.424) (-2.500,-1.436) (-2.550,-1.447) (-2.600,-1.458) (-2.650,-1.467) (-2.700,-1.475) (-2.750,-1.481) (-2.800,-1.487) (-2.850,-1.491) (-2.900,-1.494) (-2.950,-1.496) (-3.000,-1.497) (-3.050,-1.496) (-3.100,-1.493) (-3.150,-1.490) (-3.200,-1.485) (-3.250,-1.478) (-3.300,-1.470) (-3.350,-1.460) (-3.400,-1.449) (-3.450,-1.436) (-3.500,-1.421) (-3.550,-1.405) (-3.600,-1.386) (-3.650,-1.365) (-3.700,-1.343) (-3.750,-1.318) (-3.800,-1.290) (-3.850,-1.260) (-3.900,-1.228) (-3.950,-1.192) (-4.000,-1.153) (-4.050,-1.110) (-4.088,-1.075) (-4.125,-1.038) (-4.160,-1.000) (-4.200,-0.954) (-4.243,-0.900) (-4.275,-0.855) (-4.312,-0.800) (-4.342,-0.750) (-4.369,-0.700) (-4.395,-0.650) (-4.417,-0.600) (-4.438,-0.550) (-4.457,-0.500) (-4.473,-0.450) (-4.488,-0.400) (-4.501,-0.350) (-4.512,-0.300) (-4.521,-0.250) (-4.529,-0.200) (-4.534,-0.150) (-4.539,-0.100) (-4.541,-0.050) (-4.542,0.000) (-4.541,0.050) (-4.539,0.100) (-4.534,0.150) (-4.529,0.200) (-4.521,0.250) (-4.512,0.300) (-4.501,0.350) (-4.488,0.400) (-4.473,0.450) (-4.457,0.500) (-4.438,0.550) (-4.417,0.600) (-4.395,0.650) (-4.369,0.700) (-4.342,0.750) (-4.312,0.800) (-4.279,0.850) (-4.243,0.900) (-4.203,0.950) (-4.160,1.000) (-4.125,1.038) (-4.088,1.075) (-4.050,1.110) (-4.003,1.150) (-3.950,1.192) (-3.904,1.225) (-3.850,1.260) (-3.800,1.290) (-3.750,1.318) (-3.700,1.343) (-3.650,1.365) (-3.600,1.386) (-3.550,1.405) (-3.500,1.421) (-3.450,1.436) (-3.400,1.449) (-3.350,1.460) (-3.300,1.470) (-3.250,1.478) (-3.200,1.485) (-3.150,1.490) (-3.100,1.493) (-3.050,1.496) (-3.000,1.497) (-2.950,1.496) (-2.900,1.494) (-2.850,1.491) (-2.800,1.487) (-2.750,1.481) (-2.700,1.475) (-2.650,1.467) (-2.600,1.458) (-2.550,1.447) (-2.500,1.436) (-2.450,1.424) (-2.400,1.410) (-2.350,1.396) (-2.300,1.380) (-2.250,1.363) (-2.200,1.346) (-2.150,1.327) (-2.100,1.308) (-2.050,1.288) (-2.000,1.266) (-1.950,1.244) (-1.900,1.221) (-1.850,1.197) (-1.804,1.175) (-1.755,1.150) (-1.707,1.125) (-1.660,1.100) (-1.614,1.075) (-1.569,1.050) (-1.525,1.025) (-1.482,1.000) (-1.425,0.967) (-1.375,0.937) (-1.325,0.906) (-1.275,0.875) (-1.225,0.844) (-1.175,0.812) (-1.125,0.779) (-1.081,0.750) (-1.025,0.714) (-0.975,0.680) (-0.930,0.650) (-0.875,0.613) (-0.825,0.579) (-0.784,0.550) (-0.725,0.510) (-0.675,0.475) (-0.625,0.441) (-0.575,0.406) (-0.532,0.375) (-0.475,0.336) (-0.425,0.300) (-0.375,0.266) (-0.325,0.230) (-0.284,0.200) (-0.225,0.160) (-0.220,0.155)}
\EQ{(0.206,0.144) (0.248,0.175) (0.300,0.213) (0.350,0.248) (0.400,0.283) (0.450,0.318) (0.496,0.350) (0.550,0.388) (0.600,0.423) (0.650,0.458) (0.700,0.493) (0.747,0.525) (0.800,0.562) (0.850,0.596) (0.894,0.625) (0.950,0.664) (1.000,0.697) (1.043,0.725) (1.100,0.763) (1.150,0.796) (1.196,0.825) (1.250,0.859) (1.300,0.891) (1.350,0.921) (1.397,0.950) (1.450,0.982) (1.500,1.011) (1.550,1.040) (1.600,1.068) (1.650,1.095) (1.700,1.122) (1.750,1.148) (1.800,1.173) (1.850,1.197) (1.900,1.221) (1.950,1.244) (2.000,1.266) (2.050,1.288) (2.100,1.308) (2.150,1.327) (2.200,1.346) (2.250,1.363) (2.300,1.380) (2.350,1.396) (2.400,1.410) (2.450,1.424) (2.500,1.436) (2.550,1.447) (2.600,1.458) (2.650,1.467) (2.700,1.475) (2.750,1.481) (2.800,1.487) (2.850,1.491) (2.900,1.494) (2.950,1.496) (3.000,1.497) (3.050,1.496) (3.100,1.493) (3.150,1.490) (3.200,1.485) (3.250,1.478) (3.300,1.470) (3.350,1.460) (3.400,1.449) (3.450,1.436) (3.500,1.421) (3.550,1.405) (3.600,1.386) (3.650,1.365) (3.700,1.343) (3.750,1.318) (3.800,1.290) (3.850,1.260) (3.900,1.228) (3.950,1.192) (4.000,1.153) (4.050,1.110) (4.088,1.075) (4.125,1.038) (4.160,1.000) (4.200,0.954) (4.243,0.900) (4.275,0.855) (4.312,0.800) (4.342,0.750) (4.369,0.700) (4.395,0.650) (4.417,0.600) (4.438,0.550) (4.457,0.500) (4.473,0.450) (4.488,0.400) (4.501,0.350) (4.512,0.300) (4.521,0.250) (4.529,0.200) (4.534,0.150) (4.539,0.100) (4.541,0.050) (4.542,0.000) (4.541,-0.050) (4.539,-0.100) (4.534,-0.150) (4.529,-0.200) (4.521,-0.250) (4.512,-0.300) (4.501,-0.350) (4.488,-0.400) (4.473,-0.450) (4.457,-0.500) (4.438,-0.550) (4.417,-0.600) (4.395,-0.650) (4.369,-0.700) (4.342,-0.750) (4.312,-0.800) (4.279,-0.850) (4.243,-0.900) (4.203,-0.950) (4.160,-1.000) (4.125,-1.038) (4.088,-1.075) (4.050,-1.110) (4.003,-1.150) (3.950,-1.192) (3.904,-1.225) (3.850,-1.260) (3.800,-1.290) (3.750,-1.318) (3.700,-1.343) (3.650,-1.365) (3.600,-1.386) (3.550,-1.405) (3.500,-1.421) (3.450,-1.436) (3.400,-1.449) (3.350,-1.460) (3.300,-1.470) (3.250,-1.478) (3.200,-1.485) (3.150,-1.490) (3.100,-1.493) (3.050,-1.496) (3.000,-1.497) (2.950,-1.496) (2.900,-1.494) (2.850,-1.491) (2.800,-1.487) (2.750,-1.481) (2.700,-1.475) (2.650,-1.467) (2.600,-1.458) (2.550,-1.447) (2.500,-1.436) (2.450,-1.424) (2.400,-1.410) (2.350,-1.396) (2.300,-1.380) (2.250,-1.363) (2.200,-1.346) (2.150,-1.327) (2.100,-1.308) (2.050,-1.288) (2.000,-1.266) (1.950,-1.244) (1.900,-1.221) (1.850,-1.197) (1.804,-1.175) (1.755,-1.150) (1.707,-1.125) (1.660,-1.100) (1.614,-1.075) (1.569,-1.050) (1.525,-1.025) (1.482,-1.000) (1.425,-0.967) (1.375,-0.937) (1.325,-0.906) (1.275,-0.875) (1.225,-0.844) (1.175,-0.812) (1.125,-0.779) (1.081,-0.750) (1.025,-0.714) (0.975,-0.680) (0.930,-0.650) (0.875,-0.613) (0.825,-0.579) (0.784,-0.550) (0.725,-0.510) (0.675,-0.475) (0.625,-0.441) (0.575,-0.406) (0.532,-0.375) (0.475,-0.336) (0.425,-0.300) (0.375,-0.266) (0.325,-0.230) (0.284,-0.200) (0.225,-0.160) (0.206,-0.144)}
\EQ{(-0.220,-0.155) (-0.275,-0.195) (-0.319,-0.225) (-0.375,-0.266) (-0.425,-0.300) (-0.475,-0.336) (-0.525,-0.371) (-0.568,-0.400) (-0.625,-0.441) (-0.674,-0.475) (-0.725,-0.510) (-0.775,-0.545) (-0.819,-0.575) (-0.875,-0.614) (-0.925,-0.647) (-0.975,-0.681) (-1.025,-0.715) (-1.075,-0.748) (-1.125,-0.781) (-1.175,-0.814) (-1.225,-0.846) (-1.271,-0.875) (-1.325,-0.910) (-1.375,-0.941) (-1.425,-0.971) (-1.473,-1.000) (-1.525,-1.031) (-1.575,-1.060) (-1.625,-1.089) (-1.675,-1.116) (-1.725,-1.144) (-1.775,-1.170) (-1.825,-1.196) (-1.875,-1.221) (-1.925,-1.246) (-1.975,-1.270) (-2.025,-1.293) (-2.075,-1.315) (-2.125,-1.336) (-2.175,-1.357) (-2.222,-1.375) (-2.275,-1.395) (-2.325,-1.413) (-2.375,-1.430) (-2.425,-1.446) (-2.475,-1.461) (-2.525,-1.474) (-2.575,-1.487) (-2.625,-1.499) (-2.675,-1.510) (-2.725,-1.520) (-2.775,-1.529) (-2.825,-1.536) (-2.875,-1.542) (-2.925,-1.548) (-2.975,-1.552) (-3.025,-1.555) (-3.075,-1.556) (-3.125,-1.556) (-3.175,-1.556) (-3.225,-1.553) (-3.275,-1.550) (-3.325,-1.545) (-3.375,-1.538) (-3.425,-1.530) (-3.475,-1.521) (-3.525,-1.510) (-3.575,-1.497) (-3.625,-1.483) (-3.675,-1.467) (-3.725,-1.449) (-3.775,-1.429) (-3.825,-1.407) (-3.875,-1.383) (-3.925,-1.357) (-3.975,-1.329) (-4.021,-1.300) (-4.075,-1.264) (-4.125,-1.227) (-4.175,-1.187) (-4.217,-1.150) (-4.270,-1.100) (-4.318,-1.050) (-4.362,-1.000) (-4.400,-0.952) (-4.439,-0.900) (-4.473,-0.850) (-4.500,-0.807) (-4.533,-0.750) (-4.559,-0.700) (-4.583,-0.650) (-4.605,-0.600) (-4.624,-0.550) (-4.642,-0.500) (-4.658,-0.450) (-4.672,-0.400) (-4.684,-0.350) (-4.695,-0.300) (-4.704,-0.250) (-4.711,-0.200) (-4.717,-0.150) (-4.720,-0.100) (-4.723,-0.050) (-4.724,0.000) (-4.723,0.050) (-4.720,0.100) (-4.717,0.150) (-4.711,0.200) (-4.704,0.250) (-4.695,0.300) (-4.684,0.350) (-4.672,0.400) (-4.658,0.450) (-4.642,0.500) (-4.625,0.548) (-4.605,0.600) (-4.583,0.650) (-4.559,0.700) (-4.533,0.750) (-4.504,0.800) (-4.475,0.847) (-4.439,0.900) (-4.402,0.950) (-4.362,1.000) (-4.325,1.042) (-4.275,1.095) (-4.225,1.143) (-4.175,1.187) (-4.127,1.225) (-4.075,1.264) (-4.025,1.298) (-3.981,1.325) (-3.925,1.357) (-3.875,1.383) (-3.825,1.407) (-3.775,1.429) (-3.725,1.449) (-3.675,1.467) (-3.625,1.483) (-3.575,1.497) (-3.525,1.510) (-3.475,1.521) (-3.425,1.530) (-3.375,1.538) (-3.325,1.545) (-3.275,1.550) (-3.225,1.553) (-3.175,1.556) (-3.125,1.556) (-3.075,1.556) (-3.025,1.555) (-2.975,1.552) (-2.925,1.548) (-2.875,1.542) (-2.825,1.536) (-2.775,1.529) (-2.725,1.520) (-2.675,1.510) (-2.625,1.499) (-2.575,1.487) (-2.525,1.474) (-2.475,1.461) (-2.425,1.446) (-2.375,1.430) (-2.325,1.413) (-2.275,1.395) (-2.225,1.376) (-2.175,1.357) (-2.125,1.336) (-2.075,1.315) (-2.025,1.293) (-1.975,1.270) (-1.925,1.246) (-1.875,1.221) (-1.825,1.196) (-1.775,1.170) (-1.725,1.144) (-1.675,1.116) (-1.625,1.089) (-1.575,1.060) (-1.525,1.031) (-1.475,1.001) (-1.432,0.975) (-1.375,0.941) (-1.325,0.910) (-1.275,0.878) (-1.232,0.850) (-1.175,0.814) (-1.125,0.781) (-1.078,0.750) (-1.025,0.715) (-0.975,0.681) (-0.929,0.650) (-0.875,0.614) (-0.825,0.579) (-0.783,0.550) (-0.725,0.510) (-0.675,0.475) (-0.625,0.441) (-0.575,0.406) (-0.532,0.375) (-0.475,0.336) (-0.425,0.300) (-0.375,0.266) (-0.325,0.230) (-0.284,0.200) (-0.225,0.160) (-0.220,0.155)}
\EQ{(0.206,0.144) (0.248,0.175) (0.300,0.213) (0.350,0.248) (0.400,0.283) (0.450,0.318) (0.496,0.350) (0.550,0.389) (0.600,0.423) (0.650,0.458) (0.700,0.493) (0.747,0.525) (0.800,0.562) (0.850,0.596) (0.893,0.625) (0.950,0.665) (1.000,0.698) (1.050,0.732) (1.100,0.765) (1.150,0.797) (1.193,0.825) (1.250,0.862) (1.300,0.894) (1.350,0.925) (1.400,0.956) (1.450,0.986) (1.500,1.016) (1.550,1.046) (1.600,1.074) (1.646,1.100) (1.691,1.125) (1.737,1.150) (1.785,1.175) (1.833,1.200) (1.883,1.225) (1.934,1.250) (1.987,1.275) (2.042,1.300) (2.099,1.325) (2.150,1.346) (2.200,1.367) (2.250,1.386) (2.300,1.404) (2.350,1.421) (2.400,1.438) (2.450,1.453) (2.500,1.468) (2.550,1.481) (2.600,1.493) (2.650,1.505) (2.700,1.515) (2.750,1.524) (2.800,1.532) (2.850,1.539) (2.900,1.545) (2.950,1.550) (3.000,1.553) (3.050,1.556) (3.100,1.556) (3.150,1.556) (3.200,1.555) (3.250,1.552) (3.300,1.547) (3.350,1.542) (3.400,1.534) (3.450,1.526) (3.500,1.515) (3.550,1.504) (3.600,1.490) (3.650,1.475) (3.700,1.458) (3.750,1.439) (3.800,1.418) (3.850,1.396) (3.900,1.371) (3.950,1.343) (4.000,1.313) (4.050,1.281) (4.094,1.250) (4.150,1.207) (4.189,1.175) (4.244,1.125) (4.294,1.075) (4.340,1.025) (4.375,0.984) (4.421,0.925) (4.450,0.884) (4.489,0.825) (4.519,0.775) (4.546,0.725) (4.571,0.675) (4.594,0.625) (4.615,0.575) (4.633,0.525) (4.650,0.475) (4.665,0.425) (4.678,0.375) (4.690,0.325) (4.699,0.275) (4.708,0.225) (4.714,0.175) (4.719,0.125) (4.722,0.075) (4.723,0.025) (4.723,-0.050) (4.720,-0.100) (4.717,-0.150) (4.711,-0.200) (4.704,-0.250) (4.695,-0.300) (4.684,-0.350) (4.672,-0.400) (4.658,-0.450) (4.642,-0.500) (4.625,-0.548) (4.605,-0.600) (4.583,-0.650) (4.559,-0.700) (4.533,-0.750) (4.504,-0.800) (4.475,-0.847) (4.439,-0.900) (4.402,-0.950) (4.362,-1.000) (4.325,-1.042) (4.275,-1.095) (4.225,-1.143) (4.175,-1.187) (4.127,-1.225) (4.075,-1.264) (4.025,-1.298) (3.981,-1.325) (3.925,-1.357) (3.875,-1.383) (3.825,-1.407) (3.775,-1.429) (3.725,-1.449) (3.675,-1.467) (3.625,-1.483) (3.575,-1.497) (3.525,-1.510) (3.475,-1.521) (3.425,-1.530) (3.375,-1.538) (3.325,-1.545) (3.275,-1.550) (3.225,-1.553) (3.175,-1.556) (3.125,-1.556) (3.075,-1.556) (3.025,-1.555) (2.975,-1.552) (2.925,-1.548) (2.875,-1.542) (2.825,-1.536) (2.775,-1.529) (2.725,-1.520) (2.675,-1.510) (2.625,-1.499) (2.575,-1.487) (2.525,-1.474) (2.475,-1.461) (2.425,-1.446) (2.375,-1.430) (2.325,-1.413) (2.275,-1.395) (2.225,-1.376) (2.175,-1.357) (2.125,-1.336) (2.075,-1.315) (2.025,-1.293) (1.975,-1.270) (1.925,-1.246) (1.875,-1.221) (1.825,-1.196) (1.775,-1.170) (1.725,-1.144) (1.675,-1.116) (1.625,-1.089) (1.575,-1.060) (1.525,-1.031) (1.475,-1.001) (1.432,-0.975) (1.375,-0.941) (1.325,-0.910) (1.275,-0.878) (1.232,-0.850) (1.175,-0.814) (1.125,-0.781) (1.078,-0.750) (1.025,-0.715) (0.975,-0.681) (0.929,-0.650) (0.875,-0.614) (0.825,-0.579) (0.783,-0.550) (0.725,-0.510) (0.675,-0.475) (0.625,-0.441) (0.575,-0.406) (0.532,-0.375) (0.475,-0.336) (0.425,-0.300) (0.375,-0.266) (0.325,-0.230) (0.284,-0.200) (0.225,-0.160) (0.206,-0.144)}
\EQ{(-0.220,-0.155) (-0.275,-0.195) (-0.319,-0.225) (-0.375,-0.266) (-0.425,-0.300) (-0.475,-0.336) (-0.525,-0.371) (-0.568,-0.400) (-0.625,-0.441) (-0.674,-0.475) (-0.725,-0.511) (-0.775,-0.545) (-0.819,-0.575) (-0.875,-0.614) (-0.925,-0.648) (-0.975,-0.682) (-1.025,-0.716) (-1.075,-0.750) (-1.125,-0.783) (-1.175,-0.816) (-1.225,-0.848) (-1.275,-0.881) (-1.325,-0.913) (-1.375,-0.944) (-1.424,-0.975) (-1.475,-1.006) (-1.525,-1.037) (-1.575,-1.066) (-1.625,-1.096) (-1.675,-1.124) (-1.721,-1.150) (-1.766,-1.175) (-1.812,-1.200) (-1.859,-1.225) (-1.907,-1.250) (-1.957,-1.275) (-2.008,-1.300) (-2.060,-1.325) (-2.114,-1.350) (-2.171,-1.375) (-2.225,-1.398) (-2.275,-1.419) (-2.325,-1.438) (-2.375,-1.457) (-2.425,-1.475) (-2.475,-1.492) (-2.525,-1.508) (-2.575,-1.523) (-2.625,-1.538) (-2.675,-1.551) (-2.725,-1.563) (-2.775,-1.574) (-2.825,-1.585) (-2.875,-1.594) (-2.925,-1.602) (-2.975,-1.609) (-3.025,-1.615) (-3.075,-1.620) (-3.125,-1.623) (-3.175,-1.626) (-3.225,-1.627) (-3.275,-1.627) (-3.325,-1.626) (-3.375,-1.624) (-3.425,-1.620) (-3.475,-1.615) (-3.525,-1.609) (-3.575,-1.601) (-3.625,-1.592) (-3.675,-1.581) (-3.725,-1.569) (-3.775,-1.555) (-3.825,-1.539) (-3.875,-1.522) (-3.925,-1.503) (-3.975,-1.482) (-4.025,-1.459) (-4.075,-1.434) (-4.125,-1.407) (-4.175,-1.378) (-4.218,-1.350) (-4.275,-1.311) (-4.323,-1.275) (-4.375,-1.232) (-4.425,-1.188) (-4.464,-1.150) (-4.500,-1.113) (-4.535,-1.075) (-4.575,-1.028) (-4.616,-0.975) (-4.650,-0.928) (-4.685,-0.875) (-4.716,-0.825) (-4.744,-0.775) (-4.770,-0.725) (-4.793,-0.675) (-4.815,-0.625) (-4.835,-0.575) (-4.853,-0.525) (-4.869,-0.475) (-4.883,-0.425) (-4.895,-0.375) (-4.906,-0.325) (-4.916,-0.275) (-4.923,-0.225) (-4.929,-0.175) (-4.934,-0.125) (-4.937,-0.075) (-4.938,-0.025) (-4.938,0.050) (-4.936,0.100) (-4.932,0.150) (-4.926,0.200) (-4.920,0.250) (-4.911,0.300) (-4.901,0.350) (-4.889,0.400) (-4.876,0.450) (-4.861,0.500) (-4.844,0.550) (-4.825,0.600) (-4.804,0.650) (-4.782,0.700) (-4.757,0.750) (-4.730,0.800) (-4.701,0.850) (-4.669,0.900) (-4.635,0.950) (-4.600,0.996) (-4.556,1.050) (-4.512,1.100) (-4.475,1.139) (-4.438,1.175) (-4.400,1.210) (-4.354,1.250) (-4.300,1.292) (-4.255,1.325) (-4.200,1.362) (-4.150,1.393) (-4.100,1.421) (-4.050,1.447) (-4.000,1.471) (-3.950,1.493) (-3.900,1.513) (-3.850,1.531) (-3.800,1.547) (-3.750,1.562) (-3.700,1.575) (-3.650,1.586) (-3.600,1.596) (-3.550,1.605) (-3.500,1.612) (-3.450,1.618) (-3.400,1.622) (-3.350,1.625) (-3.300,1.627) (-3.250,1.627) (-3.200,1.627) (-3.150,1.625) (-3.100,1.622) (-3.050,1.617) (-3.000,1.612) (-2.950,1.605) (-2.900,1.598) (-2.850,1.589) (-2.800,1.580) (-2.750,1.569) (-2.700,1.557) (-2.650,1.544) (-2.600,1.531) (-2.550,1.516) (-2.500,1.500) (-2.450,1.484) (-2.400,1.466) (-2.350,1.448) (-2.300,1.429) (-2.250,1.409) (-2.200,1.388) (-2.150,1.366) (-2.100,1.344) (-2.050,1.320) (-2.000,1.296) (-1.950,1.272) (-1.900,1.246) (-1.850,1.220) (-1.800,1.194) (-1.750,1.167) (-1.700,1.139) (-1.650,1.110) (-1.600,1.081) (-1.550,1.051) (-1.506,1.025) (-1.450,0.991) (-1.400,0.960) (-1.350,0.929) (-1.305,0.900) (-1.250,0.865) (-1.200,0.832) (-1.151,0.800) (-1.100,0.766) (-1.050,0.733) (-1.001,0.700) (-0.950,0.665) (-0.900,0.631) (-0.855,0.600) (-0.800,0.563) (-0.750,0.528) (-0.700,0.493) (-0.650,0.459) (-0.603,0.425) (-0.550,0.389) (-0.500,0.353) (-0.450,0.318) (-0.400,0.283) (-0.354,0.250) (-0.300,0.213) (-0.250,0.177) (-0.220,0.155)}
\EQ{(0.206,0.144) (0.248,0.175) (0.300,0.213) (0.350,0.248) (0.400,0.283) (0.450,0.318) (0.496,0.350) (0.550,0.389) (0.600,0.423) (0.650,0.459) (0.700,0.493) (0.746,0.525) (0.800,0.563) (0.850,0.597) (0.891,0.625) (0.950,0.665) (1.000,0.699) (1.050,0.733) (1.100,0.766) (1.150,0.799) (1.200,0.832) (1.250,0.865) (1.300,0.897) (1.345,0.925) (1.400,0.960) (1.450,0.991) (1.500,1.021) (1.548,1.050) (1.600,1.081) (1.650,1.110) (1.700,1.139) (1.750,1.167) (1.800,1.194) (1.850,1.220) (1.900,1.246) (1.950,1.272) (2.000,1.296) (2.050,1.320) (2.100,1.344) (2.150,1.366) (2.200,1.388) (2.250,1.409) (2.300,1.429) (2.350,1.448) (2.400,1.466) (2.450,1.484) (2.500,1.500) (2.550,1.516) (2.600,1.531) (2.650,1.544) (2.700,1.557) (2.750,1.569) (2.800,1.580) (2.850,1.589) (2.900,1.598) (2.950,1.605) (3.000,1.612) (3.050,1.617) (3.100,1.622) (3.150,1.625) (3.200,1.627) (3.250,1.627) (3.300,1.627) (3.350,1.625) (3.400,1.622) (3.450,1.618) (3.500,1.612) (3.550,1.605) (3.600,1.596) (3.650,1.586) (3.700,1.575) (3.750,1.562) (3.800,1.547) (3.850,1.531) (3.900,1.513) (3.950,1.493) (4.000,1.471) (4.050,1.447) (4.100,1.421) (4.150,1.393) (4.200,1.362) (4.250,1.328) (4.300,1.292) (4.350,1.253) (4.400,1.210) (4.438,1.175) (4.475,1.139) (4.512,1.100) (4.550,1.058) (4.597,1.000) (4.635,0.950) (4.669,0.900) (4.700,0.851) (4.730,0.800) (4.757,0.750) (4.782,0.700) (4.804,0.650) (4.825,0.600) (4.844,0.550) (4.861,0.500) (4.876,0.450) (4.889,0.400) (4.901,0.350) (4.911,0.300) (4.920,0.250) (4.926,0.200) (4.932,0.150) (4.936,0.100) (4.938,0.050) (4.939,0.000) (4.938,-0.050) (4.936,-0.100) (4.932,-0.150) (4.926,-0.200) (4.920,-0.250) (4.911,-0.300) (4.901,-0.350) (4.889,-0.400) (4.876,-0.450) (4.861,-0.500) (4.844,-0.550) (4.825,-0.600) (4.804,-0.650) (4.782,-0.700) (4.757,-0.750) (4.730,-0.800) (4.701,-0.850) (4.669,-0.900) (4.635,-0.950) (4.600,-0.996) (4.556,-1.050) (4.512,-1.100) (4.475,-1.139) (4.438,-1.175) (4.400,-1.210) (4.354,-1.250) (4.300,-1.292) (4.255,-1.325) (4.200,-1.362) (4.150,-1.393) (4.100,-1.421) (4.050,-1.447) (4.000,-1.471) (3.950,-1.493) (3.900,-1.513) (3.850,-1.531) (3.800,-1.547) (3.750,-1.562) (3.700,-1.575) (3.650,-1.586) (3.600,-1.596) (3.550,-1.605) (3.500,-1.612) (3.450,-1.618) (3.400,-1.622) (3.350,-1.625) (3.300,-1.627) (3.250,-1.627) (3.200,-1.627) (3.150,-1.625) (3.100,-1.622) (3.050,-1.617) (3.000,-1.612) (2.950,-1.605) (2.900,-1.598) (2.850,-1.589) (2.800,-1.580) (2.750,-1.569) (2.700,-1.557) (2.650,-1.544) (2.600,-1.531) (2.550,-1.516) (2.500,-1.500) (2.450,-1.484) (2.400,-1.466) (2.350,-1.448) (2.300,-1.429) (2.250,-1.409) (2.200,-1.388) (2.150,-1.366) (2.100,-1.344) (2.050,-1.320) (2.000,-1.296) (1.950,-1.272) (1.900,-1.246) (1.850,-1.220) (1.800,-1.194) (1.750,-1.167) (1.700,-1.139) (1.650,-1.110) (1.600,-1.081) (1.550,-1.051) (1.506,-1.025) (1.450,-0.991) (1.400,-0.960) (1.350,-0.929) (1.305,-0.900) (1.250,-0.865) (1.200,-0.832) (1.151,-0.800) (1.100,-0.766) (1.050,-0.733) (1.001,-0.700) (0.950,-0.665) (0.900,-0.631) (0.855,-0.600) (0.800,-0.563) (0.750,-0.528) (0.700,-0.493) (0.650,-0.459) (0.603,-0.425) (0.550,-0.389) (0.500,-0.353) (0.450,-0.318) (0.400,-0.283) (0.354,-0.250) (0.300,-0.213) (0.250,-0.177) (0.206,-0.144)}
\EQ{(-0.220,-0.155) (-0.275,-0.195) (-0.319,-0.225) (-0.375,-0.266) (-0.425,-0.300) (-0.475,-0.336) (-0.525,-0.371) (-0.567,-0.400) (-0.625,-0.441) (-0.674,-0.475) (-0.725,-0.511) (-0.775,-0.546) (-0.818,-0.575) (-0.875,-0.615) (-0.925,-0.649) (-0.975,-0.683) (-1.025,-0.717) (-1.074,-0.750) (-1.125,-0.785) (-1.175,-0.818) (-1.224,-0.850) (-1.275,-0.884) (-1.325,-0.916) (-1.375,-0.948) (-1.417,-0.975) (-1.475,-1.011) (-1.525,-1.042) (-1.575,-1.073) (-1.620,-1.100) (-1.675,-1.133) (-1.725,-1.162) (-1.775,-1.190) (-1.825,-1.218) (-1.875,-1.246) (-1.925,-1.273) (-1.975,-1.299) (-2.025,-1.325) (-2.075,-1.350) (-2.125,-1.374) (-2.175,-1.398) (-2.225,-1.421) (-2.275,-1.443) (-2.325,-1.465) (-2.375,-1.486) (-2.425,-1.505) (-2.475,-1.525) (-2.525,-1.543) (-2.575,-1.560) (-2.625,-1.577) (-2.675,-1.593) (-2.725,-1.608) (-2.775,-1.622) (-2.825,-1.635) (-2.875,-1.647) (-2.925,-1.658) (-2.975,-1.668) (-3.025,-1.677) (-3.075,-1.686) (-3.125,-1.693) (-3.175,-1.699) (-3.225,-1.704) (-3.275,-1.708) (-3.325,-1.711) (-3.375,-1.713) (-3.425,-1.713) (-3.475,-1.713) (-3.525,-1.711) (-3.575,-1.708) (-3.625,-1.704) (-3.675,-1.698) (-3.725,-1.691) (-3.775,-1.683) (-3.825,-1.674) (-3.875,-1.663) (-3.925,-1.650) (-3.975,-1.636) (-4.025,-1.621) (-4.075,-1.604) (-4.125,-1.585) (-4.175,-1.564) (-4.225,-1.542) (-4.275,-1.518) (-4.325,-1.491) (-4.375,-1.462) (-4.425,-1.431) (-4.472,-1.400) (-4.525,-1.362) (-4.572,-1.325) (-4.625,-1.280) (-4.675,-1.234) (-4.725,-1.185) (-4.775,-1.130) (-4.822,-1.075) (-4.861,-1.025) (-4.897,-0.975) (-4.925,-0.932) (-4.961,-0.875) (-4.989,-0.825) (-5.015,-0.775) (-5.040,-0.725) (-5.062,-0.675) (-5.082,-0.625) (-5.100,-0.577) (-5.117,-0.525) (-5.133,-0.475) (-5.146,-0.425) (-5.158,-0.375) (-5.168,-0.325) (-5.177,-0.275) (-5.184,-0.225) (-5.190,-0.175) (-5.195,-0.125) (-5.197,-0.075) (-5.199,-0.025) (-5.198,0.050) (-5.196,0.100) (-5.193,0.150) (-5.187,0.200) (-5.181,0.250) (-5.173,0.300) (-5.163,0.350) (-5.152,0.400) (-5.140,0.450) (-5.125,0.500) (-5.109,0.550) (-5.092,0.600) (-5.072,0.650) (-5.051,0.700) (-5.028,0.750) (-5.002,0.800) (-4.975,0.850) (-4.946,0.900) (-4.914,0.950) (-4.879,1.000) (-4.842,1.050) (-4.801,1.100) (-4.758,1.150) (-4.710,1.200) (-4.659,1.250) (-4.602,1.300) (-4.550,1.343) (-4.507,1.375) (-4.450,1.415) (-4.400,1.447) (-4.353,1.475) (-4.309,1.500) (-4.260,1.525) (-4.208,1.550) (-4.150,1.575) (-4.100,1.595) (-4.050,1.613) (-4.000,1.629) (-3.950,1.644) (-3.900,1.657) (-3.850,1.668) (-3.800,1.679) (-3.750,1.687) (-3.700,1.695) (-3.650,1.701) (-3.600,1.706) (-3.550,1.710) (-3.500,1.712) (-3.450,1.713) (-3.400,1.713) (-3.350,1.712) (-3.300,1.710) (-3.250,1.706) (-3.200,1.702) (-3.150,1.696) (-3.100,1.689) (-3.050,1.682) (-3.000,1.673) (-2.950,1.663) (-2.900,1.653) (-2.850,1.641) (-2.800,1.628) (-2.750,1.615) (-2.700,1.600) (-2.650,1.585) (-2.600,1.569) (-2.550,1.552) (-2.500,1.534) (-2.450,1.515) (-2.400,1.496) (-2.350,1.475) (-2.300,1.454) (-2.250,1.432) (-2.200,1.410) (-2.150,1.386) (-2.100,1.362) (-2.050,1.338) (-2.000,1.312) (-1.950,1.286) (-1.900,1.259) (-1.850,1.232) (-1.800,1.204) (-1.750,1.176) (-1.705,1.150) (-1.650,1.118) (-1.600,1.088) (-1.550,1.058) (-1.500,1.027) (-1.457,1.000) (-1.400,0.964) (-1.350,0.932) (-1.300,0.900) (-1.250,0.867) (-1.200,0.835) (-1.150,0.801) (-1.100,0.768) (-1.050,0.734) (-1.000,0.700) (-0.950,0.666) (-0.900,0.632) (-0.854,0.600) (-0.800,0.563) (-0.750,0.528) (-0.700,0.494) (-0.650,0.459) (-0.603,0.425) (-0.550,0.389) (-0.500,0.353) (-0.450,0.318) (-0.400,0.283) (-0.354,0.250) (-0.300,0.213) (-0.250,0.177) (-0.220,0.155)}
\EQ{(0.206,0.144) (0.248,0.175) (0.300,0.213) (0.350,0.248) (0.400,0.283) (0.450,0.318) (0.496,0.350) (0.550,0.389) (0.600,0.423) (0.650,0.459) (0.700,0.494) (0.746,0.525) (0.800,0.563) (0.850,0.598) (0.900,0.632) (0.950,0.666) (1.000,0.700) (1.050,0.734) (1.100,0.768) (1.148,0.800) (1.200,0.835) (1.250,0.867) (1.300,0.900) (1.350,0.932) (1.400,0.964) (1.450,0.996) (1.497,1.025) (1.550,1.058) (1.600,1.088) (1.650,1.118) (1.700,1.147) (1.748,1.175) (1.792,1.200) (1.837,1.225) (1.883,1.250) (1.929,1.275) (1.975,1.299) (2.025,1.325) (2.075,1.350) (2.125,1.374) (2.175,1.398) (2.225,1.421) (2.275,1.443) (2.325,1.465) (2.375,1.486) (2.425,1.505) (2.475,1.525) (2.525,1.543) (2.575,1.560) (2.625,1.577) (2.675,1.593) (2.725,1.608) (2.775,1.622) (2.825,1.635) (2.875,1.647) (2.925,1.658) (2.975,1.668) (3.025,1.677) (3.075,1.686) (3.125,1.693) (3.175,1.699) (3.225,1.704) (3.275,1.708) (3.325,1.711) (3.375,1.713) (3.425,1.713) (3.475,1.713) (3.525,1.711) (3.575,1.708) (3.625,1.704) (3.675,1.698) (3.725,1.691) (3.775,1.683) (3.825,1.674) (3.875,1.663) (3.925,1.650) (3.975,1.636) (4.025,1.621) (4.075,1.604) (4.125,1.585) (4.175,1.564) (4.225,1.542) (4.275,1.518) (4.325,1.491) (4.375,1.462) (4.425,1.431) (4.472,1.400) (4.525,1.362) (4.572,1.325) (4.625,1.280) (4.675,1.234) (4.725,1.185) (4.775,1.130) (4.822,1.075) (4.861,1.025) (4.897,0.975) (4.925,0.932) (4.961,0.875) (4.989,0.825) (5.015,0.775) (5.040,0.725) (5.062,0.675) (5.082,0.625) (5.100,0.577) (5.117,0.525) (5.133,0.475) (5.146,0.425) (5.158,0.375) (5.168,0.325) (5.177,0.275) (5.184,0.225) (5.190,0.175) (5.195,0.125) (5.197,0.075) (5.199,0.025) (5.198,-0.050) (5.196,-0.100) (5.193,-0.150) (5.187,-0.200) (5.181,-0.250) (5.173,-0.300) (5.163,-0.350) (5.152,-0.400) (5.140,-0.450) (5.125,-0.500) (5.109,-0.550) (5.092,-0.600) (5.072,-0.650) (5.051,-0.700) (5.028,-0.750) (5.002,-0.800) (4.975,-0.850) (4.946,-0.900) (4.914,-0.950) (4.879,-1.000) (4.842,-1.050) (4.801,-1.100) (4.758,-1.150) (4.710,-1.200) (4.659,-1.250) (4.602,-1.300) (4.550,-1.343) (4.507,-1.375) (4.450,-1.415) (4.400,-1.447) (4.353,-1.475) (4.309,-1.500) (4.260,-1.525) (4.208,-1.550) (4.150,-1.575) (4.100,-1.595) (4.050,-1.613) (4.000,-1.629) (3.950,-1.644) (3.900,-1.657) (3.850,-1.668) (3.800,-1.679) (3.750,-1.687) (3.700,-1.695) (3.650,-1.701) (3.600,-1.706) (3.550,-1.710) (3.500,-1.712) (3.450,-1.713) (3.400,-1.713) (3.350,-1.712) (3.300,-1.710) (3.250,-1.706) (3.200,-1.702) (3.150,-1.696) (3.100,-1.689) (3.050,-1.682) (3.000,-1.673) (2.950,-1.663) (2.900,-1.653) (2.850,-1.641) (2.800,-1.628) (2.750,-1.615) (2.700,-1.600) (2.650,-1.585) (2.600,-1.569) (2.550,-1.552) (2.500,-1.534) (2.450,-1.515) (2.400,-1.496) (2.350,-1.475) (2.300,-1.454) (2.250,-1.432) (2.200,-1.410) (2.150,-1.386) (2.100,-1.362) (2.050,-1.338) (2.000,-1.312) (1.950,-1.286) (1.900,-1.259) (1.850,-1.232) (1.800,-1.204) (1.750,-1.176) (1.705,-1.150) (1.650,-1.118) (1.600,-1.088) (1.550,-1.058) (1.500,-1.027) (1.457,-1.000) (1.400,-0.964) (1.350,-0.932) (1.300,-0.900) (1.250,-0.867) (1.200,-0.835) (1.150,-0.801) (1.100,-0.768) (1.050,-0.734) (1.000,-0.700) (0.950,-0.666) (0.900,-0.632) (0.854,-0.600) (0.800,-0.563) (0.750,-0.528) (0.700,-0.494) (0.650,-0.459) (0.603,-0.425) (0.550,-0.389) (0.500,-0.353) (0.450,-0.318) (0.400,-0.283) (0.354,-0.250) (0.300,-0.213) (0.250,-0.177) (0.206,-0.144)}
\EQ{(-0.220,-0.155) (-0.275,-0.195) (-0.319,-0.225) (-0.375,-0.266) (-0.425,-0.300) (-0.475,-0.336) (-0.525,-0.371) (-0.567,-0.400) (-0.625,-0.441) (-0.674,-0.475) (-0.725,-0.511) (-0.775,-0.546) (-0.817,-0.575) (-0.875,-0.616) (-0.925,-0.650) (-0.975,-0.684) (-1.025,-0.719) (-1.072,-0.750) (-1.125,-0.787) (-1.175,-0.820) (-1.220,-0.850) (-1.275,-0.887) (-1.325,-0.920) (-1.372,-0.950) (-1.425,-0.985) (-1.475,-1.017) (-1.525,-1.048) (-1.568,-1.075) (-1.625,-1.110) (-1.675,-1.141) (-1.725,-1.171) (-1.774,-1.200) (-1.825,-1.230) (-1.875,-1.259) (-1.925,-1.287) (-1.975,-1.314) (-2.025,-1.342) (-2.075,-1.368) (-2.125,-1.394) (-2.175,-1.420) (-2.225,-1.444) (-2.275,-1.468) (-2.325,-1.492) (-2.375,-1.515) (-2.425,-1.537) (-2.475,-1.558) (-2.525,-1.579) (-2.575,-1.599) (-2.625,-1.618) (-2.675,-1.637) (-2.725,-1.654) (-2.775,-1.671) (-2.825,-1.687) (-2.875,-1.702) (-2.925,-1.717) (-2.975,-1.730) (-3.025,-1.743) (-3.075,-1.755) (-3.125,-1.765) (-3.175,-1.775) (-3.225,-1.784) (-3.275,-1.792) (-3.325,-1.799) (-3.375,-1.805) (-3.425,-1.811) (-3.475,-1.815) (-3.525,-1.818) (-3.575,-1.820) (-3.625,-1.820) (-3.675,-1.820) (-3.725,-1.819) (-3.775,-1.816) (-3.825,-1.813) (-3.875,-1.808) (-3.925,-1.802) (-3.975,-1.795) (-4.025,-1.786) (-4.075,-1.776) (-4.125,-1.765) (-4.175,-1.753) (-4.225,-1.739) (-4.275,-1.723) (-4.325,-1.706) (-4.375,-1.688) (-4.425,-1.667) (-4.475,-1.645) (-4.525,-1.622) (-4.575,-1.596) (-4.625,-1.568) (-4.675,-1.538) (-4.725,-1.506) (-4.770,-1.475) (-4.825,-1.435) (-4.868,-1.400) (-4.925,-1.352) (-4.975,-1.305) (-5.025,-1.255) (-5.075,-1.200) (-5.118,-1.150) (-5.150,-1.109) (-5.194,-1.050) (-5.225,-1.004) (-5.259,-0.950) (-5.289,-0.900) (-5.316,-0.850) (-5.341,-0.800) (-5.365,-0.750) (-5.386,-0.700) (-5.406,-0.650) (-5.424,-0.600) (-5.441,-0.550) (-5.456,-0.500) (-5.469,-0.450) (-5.481,-0.400) (-5.491,-0.350) (-5.500,-0.300) (-5.508,-0.250) (-5.514,-0.200) (-5.519,-0.150) (-5.522,-0.100) (-5.524,-0.050) (-5.525,0.000) (-5.524,0.050) (-5.522,0.100) (-5.519,0.150) (-5.514,0.200) (-5.508,0.250) (-5.500,0.300) (-5.491,0.350) (-5.481,0.400) (-5.469,0.450) (-5.456,0.500) (-5.441,0.550) (-5.424,0.600) (-5.406,0.650) (-5.386,0.700) (-5.365,0.750) (-5.341,0.800) (-5.316,0.850) (-5.289,0.900) (-5.259,0.950) (-5.228,1.000) (-5.194,1.050) (-5.157,1.100) (-5.125,1.141) (-5.075,1.200) (-5.030,1.250) (-4.980,1.300) (-4.927,1.350) (-4.875,1.395) (-4.825,1.435) (-4.775,1.472) (-4.725,1.506) (-4.675,1.538) (-4.625,1.568) (-4.575,1.596) (-4.525,1.622) (-4.475,1.645) (-4.425,1.667) (-4.375,1.688) (-4.325,1.706) (-4.275,1.723) (-4.225,1.739) (-4.175,1.753) (-4.125,1.765) (-4.075,1.776) (-4.025,1.786) (-3.975,1.795) (-3.925,1.802) (-3.875,1.808) (-3.825,1.813) (-3.775,1.816) (-3.725,1.819) (-3.675,1.820) (-3.625,1.820) (-3.575,1.820) (-3.525,1.818) (-3.475,1.815) (-3.425,1.811) (-3.375,1.805) (-3.325,1.799) (-3.275,1.792) (-3.225,1.784) (-3.175,1.775) (-3.125,1.765) (-3.075,1.755) (-3.025,1.743) (-2.975,1.730) (-2.925,1.717) (-2.875,1.702) (-2.825,1.687) (-2.775,1.671) (-2.725,1.654) (-2.675,1.637) (-2.625,1.618) (-2.578,1.600) (-2.525,1.579) (-2.475,1.558) (-2.425,1.537) (-2.375,1.515) (-2.325,1.492) (-2.275,1.468) (-2.225,1.444) (-2.175,1.420) (-2.125,1.394) (-2.075,1.368) (-2.025,1.342) (-1.975,1.314) (-1.925,1.287) (-1.875,1.259) (-1.825,1.230) (-1.775,1.201) (-1.725,1.171) (-1.675,1.141) (-1.625,1.110) (-1.575,1.079) (-1.528,1.050) (-1.475,1.017) (-1.425,0.985) (-1.375,0.952) (-1.325,0.920) (-1.275,0.887) (-1.225,0.854) (-1.183,0.825) (-1.125,0.787) (-1.075,0.753) (-1.025,0.719) (-0.975,0.684) (-0.925,0.650) (-0.875,0.616) (-0.825,0.581) (-0.781,0.550) (-0.725,0.511) (-0.675,0.476) (-0.625,0.441) (-0.575,0.406) (-0.532,0.375) (-0.475,0.336) (-0.425,0.300) (-0.375,0.266) (-0.325,0.230) (-0.284,0.200) (-0.225,0.160) (-0.220,0.155)}
\EQ{(0.206,0.144) (0.248,0.175) (0.300,0.213) (0.350,0.248) (0.400,0.283) (0.450,0.318) (0.496,0.350) (0.550,0.389) (0.600,0.424) (0.650,0.459) (0.700,0.494) (0.745,0.525) (0.800,0.564) (0.850,0.598) (0.900,0.633) (0.950,0.667) (0.998,0.700) (1.050,0.736) (1.100,0.770) (1.145,0.800) (1.200,0.837) (1.250,0.870) (1.296,0.900) (1.350,0.936) (1.400,0.968) (1.449,1.000) (1.500,1.032) (1.550,1.064) (1.600,1.095) (1.649,1.125) (1.700,1.156) (1.750,1.186) (1.800,1.215) (1.850,1.244) (1.900,1.273) (1.949,1.300) (1.995,1.325) (2.041,1.350) (2.088,1.375) (2.137,1.400) (2.186,1.425) (2.237,1.450) (2.289,1.475) (2.343,1.500) (2.398,1.525) (2.450,1.548) (2.500,1.569) (2.550,1.589) (2.600,1.609) (2.650,1.628) (2.700,1.646) (2.750,1.663) (2.800,1.679) (2.850,1.695) (2.900,1.710) (2.950,1.724) (3.000,1.737) (3.050,1.749) (3.100,1.760) (3.150,1.771) (3.200,1.780) (3.250,1.788) (3.300,1.796) (3.350,1.803) (3.400,1.808) (3.450,1.813) (3.500,1.816) (3.550,1.819) (3.600,1.820) (3.650,1.820) (3.700,1.820) (3.750,1.818) (3.800,1.815) (3.850,1.811) (3.900,1.805) (3.950,1.799) (4.000,1.791) (4.050,1.782) (4.100,1.771) (4.150,1.759) (4.200,1.746) (4.250,1.731) (4.300,1.715) (4.350,1.697) (4.400,1.678) (4.450,1.657) (4.500,1.634) (4.550,1.609) (4.600,1.582) (4.650,1.553) (4.696,1.525) (4.750,1.489) (4.800,1.453) (4.850,1.415) (4.898,1.375) (4.950,1.329) (5.000,1.280) (5.050,1.228) (5.097,1.175) (5.138,1.125) (5.175,1.076) (5.211,1.025) (5.244,0.975) (5.274,0.925) (5.300,0.880) (5.325,0.832) (5.350,0.782) (5.375,0.726) (5.396,0.675) (5.415,0.625) (5.433,0.575) (5.448,0.525) (5.462,0.475) (5.475,0.425) (5.486,0.375) (5.496,0.325) (5.504,0.275) (5.511,0.225) (5.516,0.175) (5.521,0.125) (5.523,0.075) (5.525,0.025) (5.524,-0.050) (5.522,-0.100) (5.519,-0.150) (5.514,-0.200) (5.508,-0.250) (5.500,-0.300) (5.491,-0.350) (5.481,-0.400) (5.469,-0.450) (5.456,-0.500) (5.441,-0.550) (5.424,-0.600) (5.406,-0.650) (5.386,-0.700) (5.365,-0.750) (5.341,-0.800) (5.316,-0.850) (5.289,-0.900) (5.259,-0.950) (5.228,-1.000) (5.194,-1.050) (5.157,-1.100) (5.125,-1.141) (5.075,-1.200) (5.030,-1.250) (4.980,-1.300) (4.927,-1.350) (4.875,-1.395) (4.825,-1.435) (4.775,-1.472) (4.725,-1.506) (4.675,-1.538) (4.625,-1.568) (4.575,-1.596) (4.525,-1.622) (4.475,-1.645) (4.425,-1.667) (4.375,-1.688) (4.325,-1.706) (4.275,-1.723) (4.225,-1.739) (4.175,-1.753) (4.125,-1.765) (4.075,-1.776) (4.025,-1.786) (3.975,-1.795) (3.925,-1.802) (3.875,-1.808) (3.825,-1.813) (3.775,-1.816) (3.725,-1.819) (3.675,-1.820) (3.625,-1.820) (3.575,-1.820) (3.525,-1.818) (3.475,-1.815) (3.425,-1.811) (3.375,-1.805) (3.325,-1.799) (3.275,-1.792) (3.225,-1.784) (3.175,-1.775) (3.125,-1.765) (3.075,-1.755) (3.025,-1.743) (2.975,-1.730) (2.925,-1.717) (2.875,-1.702) (2.825,-1.687) (2.775,-1.671) (2.725,-1.654) (2.675,-1.637) (2.625,-1.618) (2.578,-1.600) (2.525,-1.579) (2.475,-1.558) (2.425,-1.537) (2.375,-1.515) (2.325,-1.492) (2.275,-1.468) (2.225,-1.444) (2.175,-1.420) (2.125,-1.394) (2.075,-1.368) (2.025,-1.342) (1.975,-1.314) (1.925,-1.287) (1.875,-1.259) (1.825,-1.230) (1.775,-1.201) (1.725,-1.171) (1.675,-1.141) (1.625,-1.110) (1.575,-1.079) (1.528,-1.050) (1.475,-1.017) (1.425,-0.985) (1.375,-0.952) (1.325,-0.920) (1.275,-0.887) (1.225,-0.854) (1.183,-0.825) (1.125,-0.787) (1.075,-0.753) (1.025,-0.719) (0.975,-0.684) (0.925,-0.650) (0.875,-0.616) (0.825,-0.581) (0.781,-0.550) (0.725,-0.511) (0.675,-0.476) (0.625,-0.441) (0.575,-0.406) (0.532,-0.375) (0.475,-0.336) (0.425,-0.300) (0.375,-0.266) (0.325,-0.230) (0.284,-0.200) (0.225,-0.160) (0.206,-0.144)}
\EQ{(-0.220,-0.155) (-0.275,-0.195) (-0.319,-0.225) (-0.375,-0.266) (-0.425,-0.300) (-0.475,-0.336) (-0.525,-0.371) (-0.567,-0.400) (-0.625,-0.442) (-0.673,-0.475) (-0.725,-0.512) (-0.775,-0.547) (-0.817,-0.575) (-0.875,-0.616) (-0.924,-0.650) (-0.975,-0.686) (-1.025,-0.720) (-1.069,-0.750) (-1.125,-0.788) (-1.175,-0.822) (-1.225,-0.856) (-1.275,-0.890) (-1.325,-0.923) (-1.375,-0.956) (-1.425,-0.989) (-1.475,-1.022) (-1.519,-1.050) (-1.575,-1.086) (-1.625,-1.118) (-1.675,-1.149) (-1.725,-1.180) (-1.775,-1.211) (-1.825,-1.242) (-1.875,-1.271) (-1.924,-1.300) (-1.975,-1.330) (-2.025,-1.359) (-2.075,-1.387) (-2.125,-1.415) (-2.175,-1.442) (-2.225,-1.468) (-2.275,-1.495) (-2.325,-1.520) (-2.375,-1.545) (-2.425,-1.570) (-2.475,-1.593) (-2.525,-1.617) (-2.575,-1.639) (-2.625,-1.661) (-2.675,-1.682) (-2.725,-1.703) (-2.775,-1.723) (-2.825,-1.742) (-2.875,-1.761) (-2.925,-1.778) (-2.975,-1.796) (-3.025,-1.812) (-3.075,-1.827) (-3.125,-1.842) (-3.175,-1.856) (-3.225,-1.869) (-3.275,-1.881) (-3.325,-1.893) (-3.375,-1.904) (-3.425,-1.913) (-3.475,-1.922) (-3.525,-1.930) (-3.575,-1.937) (-3.625,-1.944) (-3.675,-1.949) (-3.725,-1.953) (-3.775,-1.957) (-3.825,-1.959) (-3.875,-1.961) (-3.925,-1.961) (-3.975,-1.960) (-4.025,-1.959) (-4.075,-1.956) (-4.125,-1.952) (-4.175,-1.948) (-4.225,-1.942) (-4.275,-1.934) (-4.325,-1.926) (-4.375,-1.917) (-4.425,-1.906) (-4.475,-1.894) (-4.525,-1.880) (-4.575,-1.866) (-4.624,-1.850) (-4.675,-1.832) (-4.725,-1.813) (-4.775,-1.793) (-4.825,-1.770) (-4.875,-1.746) (-4.925,-1.721) (-4.975,-1.693) (-5.025,-1.664) (-5.075,-1.632) (-5.123,-1.600) (-5.175,-1.562) (-5.223,-1.525) (-5.275,-1.482) (-5.325,-1.438) (-5.365,-1.400) (-5.415,-1.350) (-5.450,-1.312) (-5.500,-1.254) (-5.544,-1.200) (-5.575,-1.158) (-5.616,-1.100) (-5.649,-1.050) (-5.675,-1.007) (-5.708,-0.950) (-5.735,-0.900) (-5.759,-0.850) (-5.782,-0.800) (-5.804,-0.750) (-5.823,-0.700) (-5.842,-0.650) (-5.858,-0.600) (-5.874,-0.550) (-5.887,-0.500) (-5.900,-0.450) (-5.911,-0.400) (-5.920,-0.350) (-5.929,-0.300) (-5.936,-0.250) (-5.941,-0.200) (-5.946,-0.150) (-5.949,-0.100) (-5.951,-0.050) (-5.951,0.000) (-5.951,0.050) (-5.949,0.100) (-5.946,0.150) (-5.941,0.200) (-5.936,0.250) (-5.929,0.300) (-5.920,0.350) (-5.911,0.400) (-5.900,0.450) (-5.887,0.500) (-5.874,0.550) (-5.858,0.600) (-5.842,0.650) (-5.823,0.700) (-5.804,0.750) (-5.782,0.800) (-5.759,0.850) (-5.735,0.900) (-5.708,0.950) (-5.679,1.000) (-5.650,1.048) (-5.616,1.100) (-5.581,1.150) (-5.550,1.192) (-5.504,1.250) (-5.461,1.300) (-5.425,1.339) (-5.375,1.390) (-5.325,1.438) (-5.284,1.475) (-5.225,1.524) (-5.175,1.562) (-5.125,1.599) (-5.075,1.632) (-5.025,1.664) (-4.975,1.693) (-4.925,1.721) (-4.875,1.746) (-4.825,1.770) (-4.775,1.793) (-4.725,1.813) (-4.675,1.832) (-4.625,1.850) (-4.575,1.866) (-4.525,1.880) (-4.475,1.894) (-4.425,1.906) (-4.375,1.917) (-4.325,1.926) (-4.275,1.934) (-4.225,1.942) (-4.175,1.948) (-4.125,1.952) (-4.075,1.956) (-4.025,1.959) (-3.975,1.960) (-3.925,1.961) (-3.875,1.961) (-3.825,1.959) (-3.775,1.957) (-3.725,1.953) (-3.675,1.949) (-3.625,1.944) (-3.575,1.937) (-3.525,1.930) (-3.475,1.922) (-3.425,1.913) (-3.375,1.904) (-3.325,1.893) (-3.275,1.881) (-3.225,1.869) (-3.175,1.856) (-3.125,1.842) (-3.075,1.827) (-3.025,1.812) (-2.975,1.796) (-2.925,1.778) (-2.875,1.761) (-2.825,1.742) (-2.775,1.723) (-2.725,1.703) (-2.675,1.682) (-2.625,1.661) (-2.575,1.639) (-2.525,1.617) (-2.475,1.593) (-2.425,1.570) (-2.375,1.545) (-2.325,1.520) (-2.275,1.495) (-2.225,1.468) (-2.175,1.442) (-2.125,1.415) (-2.075,1.387) (-2.025,1.359) (-1.975,1.330) (-1.925,1.301) (-1.881,1.275) (-1.825,1.242) (-1.775,1.211) (-1.725,1.180) (-1.676,1.150) (-1.625,1.118) (-1.575,1.086) (-1.525,1.054) (-1.480,1.025) (-1.425,0.989) (-1.375,0.956) (-1.328,0.925) (-1.275,0.890) (-1.225,0.856) (-1.179,0.825) (-1.125,0.788) (-1.075,0.754) (-1.033,0.725) (-0.975,0.686) (-0.925,0.651) (-0.875,0.616) (-0.825,0.581) (-0.781,0.550) (-0.725,0.512) (-0.675,0.476) (-0.625,0.442) (-0.575,0.406) (-0.531,0.375) (-0.475,0.336) (-0.425,0.300) (-0.375,0.266) (-0.325,0.230) (-0.284,0.200) (-0.225,0.160) (-0.220,0.155)}
\EQ{(0.206,0.144) (0.248,0.175) (0.300,0.213) (0.350,0.248) (0.400,0.283) (0.450,0.318) (0.496,0.350) (0.550,0.389) (0.600,0.424) (0.650,0.459) (0.700,0.494) (0.745,0.525) (0.800,0.564) (0.850,0.599) (0.900,0.634) (0.950,0.668) (0.997,0.700) (1.050,0.737) (1.100,0.771) (1.143,0.800) (1.200,0.839) (1.250,0.873) (1.300,0.906) (1.350,0.940) (1.400,0.973) (1.442,1.000) (1.500,1.038) (1.550,1.070) (1.597,1.100) (1.650,1.134) (1.700,1.165) (1.750,1.196) (1.798,1.225) (1.850,1.257) (1.900,1.286) (1.950,1.316) (2.000,1.344) (2.050,1.373) (2.099,1.400) (2.144,1.425) (2.191,1.450) (2.238,1.475) (2.286,1.500) (2.335,1.525) (2.385,1.550) (2.437,1.575) (2.489,1.600) (2.544,1.625) (2.600,1.650) (2.650,1.672) (2.700,1.693) (2.750,1.713) (2.800,1.733) (2.850,1.751) (2.900,1.770) (2.950,1.787) (3.000,1.804) (3.050,1.820) (3.100,1.835) (3.150,1.849) (3.200,1.863) (3.249,1.875) (3.300,1.887) (3.350,1.898) (3.400,1.909) (3.450,1.918) (3.500,1.926) (3.550,1.934) (3.600,1.941) (3.650,1.946) (3.700,1.951) (3.750,1.955) (3.800,1.958) (3.850,1.960) (3.900,1.961) (3.950,1.961) (4.000,1.960) (4.050,1.958) (4.100,1.954) (4.150,1.950) (4.200,1.945) (4.250,1.938) (4.300,1.930) (4.350,1.922) (4.400,1.911) (4.450,1.900) (4.500,1.887) (4.550,1.873) (4.600,1.858) (4.650,1.841) (4.700,1.823) (4.750,1.803) (4.800,1.782) (4.850,1.759) (4.900,1.734) (4.950,1.707) (5.000,1.679) (5.047,1.650) (5.100,1.616) (5.150,1.581) (5.191,1.550) (5.250,1.503) (5.300,1.461) (5.339,1.425) (5.390,1.375) (5.438,1.325) (5.475,1.284) (5.524,1.225) (5.563,1.175) (5.599,1.125) (5.633,1.075) (5.664,1.025) (5.694,0.975) (5.721,0.925) (5.747,0.875) (5.771,0.825) (5.793,0.775) (5.814,0.725) (5.833,0.675) (5.850,0.625) (5.866,0.575) (5.881,0.525) (5.894,0.475) (5.905,0.425) (5.916,0.375) (5.925,0.325) (5.932,0.275) (5.939,0.225) (5.944,0.175) (5.947,0.125) (5.950,0.075) (5.951,0.025) (5.951,-0.050) (5.949,-0.100) (5.946,-0.150) (5.941,-0.200) (5.936,-0.250) (5.929,-0.300) (5.920,-0.350) (5.911,-0.400) (5.900,-0.450) (5.887,-0.500) (5.874,-0.550) (5.858,-0.600) (5.842,-0.650) (5.823,-0.700) (5.804,-0.750) (5.782,-0.800) (5.759,-0.850) (5.735,-0.900) (5.708,-0.950) (5.679,-1.000) (5.650,-1.048) (5.616,-1.100) (5.581,-1.150) (5.550,-1.192) (5.504,-1.250) (5.461,-1.300) (5.425,-1.339) (5.375,-1.390) (5.325,-1.438) (5.284,-1.475) (5.225,-1.524) (5.175,-1.562) (5.125,-1.599) (5.075,-1.632) (5.025,-1.664) (4.975,-1.693) (4.925,-1.721) (4.875,-1.746) (4.825,-1.770) (4.775,-1.793) (4.725,-1.813) (4.675,-1.832) (4.625,-1.850) (4.575,-1.866) (4.525,-1.880) (4.475,-1.894) (4.425,-1.906) (4.375,-1.917) (4.325,-1.926) (4.275,-1.934) (4.225,-1.942) (4.175,-1.948) (4.125,-1.952) (4.075,-1.956) (4.025,-1.959) (3.975,-1.960) (3.925,-1.961) (3.875,-1.961) (3.825,-1.959) (3.775,-1.957) (3.725,-1.953) (3.675,-1.949) (3.625,-1.944) (3.575,-1.937) (3.525,-1.930) (3.475,-1.922) (3.425,-1.913) (3.375,-1.904) (3.325,-1.893) (3.275,-1.881) (3.225,-1.869) (3.175,-1.856) (3.125,-1.842) (3.075,-1.827) (3.025,-1.812) (2.975,-1.796) (2.925,-1.778) (2.875,-1.761) (2.825,-1.742) (2.775,-1.723) (2.725,-1.703) (2.675,-1.682) (2.625,-1.661) (2.575,-1.639) (2.525,-1.617) (2.475,-1.593) (2.425,-1.570) (2.375,-1.545) (2.325,-1.520) (2.275,-1.495) (2.225,-1.468) (2.175,-1.442) (2.125,-1.415) (2.075,-1.387) (2.025,-1.359) (1.975,-1.330) (1.925,-1.301) (1.881,-1.275) (1.825,-1.242) (1.775,-1.211) (1.725,-1.180) (1.676,-1.150) (1.625,-1.118) (1.575,-1.086) (1.525,-1.054) (1.480,-1.025) (1.425,-0.989) (1.375,-0.956) (1.328,-0.925) (1.275,-0.890) (1.225,-0.856) (1.179,-0.825) (1.125,-0.788) (1.075,-0.754) (1.033,-0.725) (0.975,-0.686) (0.925,-0.651) (0.875,-0.616) (0.825,-0.581) (0.781,-0.550) (0.725,-0.512) (0.675,-0.476) (0.625,-0.442) (0.575,-0.406) (0.531,-0.375) (0.475,-0.336) (0.425,-0.300) (0.375,-0.266) (0.325,-0.230) (0.284,-0.200) (0.225,-0.160) (0.206,-0.144)}
\EQ{(-0.220,-0.155) (-0.275,-0.195) (-0.319,-0.225) (-0.375,-0.266) (-0.425,-0.300) (-0.475,-0.336) (-0.525,-0.371) (-0.567,-0.400) (-0.625,-0.442) (-0.673,-0.475) (-0.725,-0.512) (-0.775,-0.547) (-0.825,-0.582) (-0.875,-0.617) (-0.923,-0.650) (-0.975,-0.687) (-1.025,-0.721) (-1.067,-0.750) (-1.125,-0.790) (-1.175,-0.824) (-1.225,-0.859) (-1.275,-0.893) (-1.323,-0.925) (-1.375,-0.960) (-1.425,-0.994) (-1.472,-1.025) (-1.525,-1.060) (-1.575,-1.093) (-1.624,-1.125) (-1.675,-1.158) (-1.725,-1.190) (-1.775,-1.222) (-1.820,-1.250) (-1.875,-1.285) (-1.925,-1.316) (-1.975,-1.346) (-2.023,-1.375) (-2.075,-1.406) (-2.125,-1.436) (-2.175,-1.465) (-2.225,-1.493) (-2.275,-1.521) (-2.325,-1.549) (-2.372,-1.575) (-2.419,-1.600) (-2.466,-1.625) (-2.514,-1.650) (-2.563,-1.675) (-2.613,-1.700) (-2.664,-1.725) (-2.716,-1.750) (-2.770,-1.775) (-2.825,-1.800) (-2.875,-1.822) (-2.925,-1.844) (-2.975,-1.865) (-3.025,-1.885) (-3.075,-1.904) (-3.125,-1.923) (-3.175,-1.942) (-3.225,-1.959) (-3.275,-1.976) (-3.325,-1.993) (-3.375,-2.008) (-3.425,-2.023) (-3.475,-2.037) (-3.525,-2.051) (-3.575,-2.064) (-3.625,-2.076) (-3.675,-2.087) (-3.725,-2.097) (-3.775,-2.107) (-3.825,-2.116) (-3.875,-2.124) (-3.925,-2.131) (-3.975,-2.138) (-4.025,-2.143) (-4.075,-2.148) (-4.125,-2.152) (-4.175,-2.155) (-4.225,-2.157) (-4.275,-2.158) (-4.325,-2.158) (-4.375,-2.158) (-4.425,-2.156) (-4.475,-2.154) (-4.525,-2.150) (-4.575,-2.146) (-4.625,-2.140) (-4.675,-2.134) (-4.725,-2.126) (-4.775,-2.117) (-4.825,-2.108) (-4.875,-2.097) (-4.925,-2.085) (-4.975,-2.071) (-5.025,-2.057) (-5.075,-2.041) (-5.125,-2.024) (-5.175,-2.006) (-5.225,-1.986) (-5.275,-1.965) (-5.325,-1.942) (-5.375,-1.918) (-5.425,-1.892) (-5.475,-1.864) (-5.525,-1.835) (-5.575,-1.804) (-5.618,-1.775) (-5.675,-1.735) (-5.722,-1.700) (-5.775,-1.658) (-5.814,-1.625) (-5.869,-1.575) (-5.921,-1.525) (-5.970,-1.475) (-6.016,-1.425) (-6.050,-1.385) (-6.098,-1.325) (-6.136,-1.275) (-6.172,-1.225) (-6.200,-1.183) (-6.236,-1.125) (-6.266,-1.075) (-6.294,-1.025) (-6.320,-0.975) (-6.344,-0.925) (-6.367,-0.875) (-6.388,-0.825) (-6.408,-0.775) (-6.426,-0.725) (-6.443,-0.675) (-6.459,-0.625) (-6.473,-0.575) (-6.486,-0.525) (-6.498,-0.475) (-6.509,-0.425) (-6.518,-0.375) (-6.526,-0.325) (-6.533,-0.275) (-6.539,-0.225) (-6.543,-0.175) (-6.547,-0.125) (-6.549,-0.075) (-6.550,-0.025) (-6.550,0.037) (-6.548,0.100) (-6.545,0.150) (-6.541,0.200) (-6.536,0.250) (-6.530,0.300) (-6.522,0.350) (-6.513,0.400) (-6.503,0.450) (-6.492,0.500) (-6.480,0.550) (-6.466,0.600) (-6.451,0.650) (-6.435,0.700) (-6.417,0.750) (-6.398,0.800) (-6.378,0.850) (-6.356,0.900) (-6.332,0.950) (-6.307,1.000) (-6.280,1.050) (-6.251,1.100) (-6.225,1.143) (-6.189,1.200) (-6.154,1.250) (-6.118,1.300) (-6.079,1.350) (-6.037,1.400) (-6.000,1.442) (-5.950,1.496) (-5.900,1.546) (-5.850,1.593) (-5.800,1.637) (-5.754,1.675) (-5.700,1.717) (-5.654,1.750) (-5.600,1.787) (-5.550,1.819) (-5.500,1.850) (-5.456,1.875) (-5.409,1.900) (-5.360,1.925) (-5.307,1.950) (-5.251,1.975) (-5.200,1.996) (-5.150,2.015) (-5.100,2.033) (-5.050,2.049) (-5.000,2.064) (-4.950,2.078) (-4.900,2.091) (-4.850,2.102) (-4.800,2.113) (-4.750,2.122) (-4.700,2.130) (-4.650,2.137) (-4.600,2.143) (-4.550,2.148) (-4.500,2.152) (-4.450,2.155) (-4.400,2.157) (-4.350,2.158) (-4.300,2.158) (-4.250,2.158) (-4.200,2.156) (-4.150,2.153) (-4.100,2.150) (-4.050,2.146) (-4.000,2.140) (-3.950,2.134) (-3.900,2.128) (-3.850,2.120) (-3.800,2.111) (-3.750,2.102) (-3.700,2.092) (-3.650,2.081) (-3.600,2.070) (-3.550,2.057) (-3.500,2.044) (-3.450,2.030) (-3.400,2.016) (-3.350,2.001) (-3.300,1.985) (-3.250,1.968) (-3.200,1.951) (-3.150,1.933) (-3.100,1.914) (-3.050,1.895) (-3.001,1.875) (-2.950,1.854) (-2.900,1.833) (-2.850,1.811) (-2.800,1.789) (-2.750,1.766) (-2.700,1.742) (-2.650,1.718) (-2.600,1.694) (-2.550,1.669) (-2.500,1.643) (-2.450,1.617) (-2.400,1.590) (-2.350,1.563) (-2.300,1.535) (-2.250,1.507) (-2.200,1.479) (-2.150,1.450) (-2.100,1.421) (-2.050,1.391) (-2.000,1.361) (-1.950,1.331) (-1.900,1.300) (-1.850,1.269) (-1.800,1.238) (-1.750,1.206) (-1.702,1.175) (-1.650,1.142) (-1.600,1.109) (-1.550,1.076) (-1.500,1.044) (-1.450,1.010) (-1.400,0.977) (-1.350,0.943) (-1.300,0.910) (-1.250,0.876) (-1.200,0.842) (-1.150,0.807) (-1.103,0.775) (-1.050,0.739) (-1.000,0.704) (-0.950,0.669) (-0.900,0.634) (-0.851,0.600) (-0.800,0.565) (-0.750,0.529) (-0.709,0.500) (-0.650,0.459) (-0.602,0.425) (-0.550,0.389) (-0.500,0.354) (-0.450,0.319) (-0.400,0.283) (-0.354,0.250) (-0.300,0.213) (-0.250,0.177) (-0.220,0.155)}
\EQ{(0.206,0.144) (0.248,0.175) (0.300,0.213) (0.350,0.248) (0.400,0.283) (0.450,0.319) (0.496,0.350) (0.550,0.389) (0.600,0.424) (0.650,0.459) (0.700,0.494) (0.744,0.525) (0.800,0.565) (0.850,0.599) (0.900,0.634) (0.950,0.669) (0.995,0.700) (1.050,0.739) (1.100,0.773) (1.150,0.807) (1.200,0.842) (1.249,0.875) (1.300,0.910) (1.350,0.943) (1.397,0.975) (1.450,1.010) (1.500,1.044) (1.548,1.075) (1.600,1.109) (1.650,1.142) (1.700,1.174) (1.750,1.206) (1.800,1.238) (1.850,1.269) (1.900,1.300) (1.950,1.331) (2.000,1.361) (2.050,1.391) (2.100,1.421) (2.150,1.450) (2.193,1.475) (2.237,1.500) (2.282,1.525) (2.327,1.550) (2.372,1.575) (2.419,1.600) (2.466,1.625) (2.514,1.650) (2.563,1.675) (2.613,1.700) (2.664,1.725) (2.716,1.750) (2.770,1.775) (2.825,1.800) (2.875,1.822) (2.925,1.844) (2.975,1.865) (3.025,1.885) (3.075,1.904) (3.125,1.923) (3.175,1.942) (3.225,1.959) (3.275,1.976) (3.325,1.993) (3.375,2.008) (3.425,2.023) (3.475,2.037) (3.525,2.051) (3.575,2.064) (3.625,2.076) (3.675,2.087) (3.725,2.097) (3.775,2.107) (3.825,2.116) (3.875,2.124) (3.925,2.131) (3.975,2.138) (4.025,2.143) (4.075,2.148) (4.125,2.152) (4.175,2.155) (4.225,2.157) (4.275,2.158) (4.325,2.158) (4.375,2.158) (4.425,2.156) (4.475,2.154) (4.525,2.150) (4.575,2.146) (4.625,2.140) (4.675,2.134) (4.725,2.126) (4.775,2.117) (4.825,2.108) (4.875,2.097) (4.925,2.085) (4.975,2.071) (5.025,2.057) (5.075,2.041) (5.125,2.024) (5.175,2.006) (5.225,1.986) (5.275,1.965) (5.325,1.942) (5.375,1.918) (5.425,1.892) (5.475,1.864) (5.525,1.835) (5.575,1.804) (5.618,1.775) (5.675,1.735) (5.722,1.700) (5.775,1.658) (5.814,1.625) (5.869,1.575) (5.921,1.525) (5.970,1.475) (6.016,1.425) (6.050,1.385) (6.098,1.325) (6.136,1.275) (6.172,1.225) (6.200,1.183) (6.236,1.125) (6.266,1.075) (6.294,1.025) (6.320,0.975) (6.344,0.925) (6.367,0.875) (6.388,0.825) (6.408,0.775) (6.426,0.725) (6.443,0.675) (6.459,0.625) (6.473,0.575) (6.486,0.525) (6.498,0.475) (6.509,0.425) (6.518,0.375) (6.526,0.325) (6.533,0.275) (6.539,0.225) (6.543,0.175) (6.547,0.125) (6.549,0.075) (6.550,0.025) (6.550,-0.037) (6.548,-0.100) (6.545,-0.150) (6.541,-0.200) (6.536,-0.250) (6.530,-0.300) (6.522,-0.350) (6.513,-0.400) (6.503,-0.450) (6.492,-0.500) (6.480,-0.550) (6.466,-0.600) (6.451,-0.650) (6.435,-0.700) (6.417,-0.750) (6.398,-0.800) (6.378,-0.850) (6.356,-0.900) (6.332,-0.950) (6.307,-1.000) (6.280,-1.050) (6.251,-1.100) (6.225,-1.143) (6.189,-1.200) (6.154,-1.250) (6.118,-1.300) (6.079,-1.350) (6.037,-1.400) (6.000,-1.442) (5.950,-1.496) (5.900,-1.546) (5.850,-1.593) (5.800,-1.637) (5.754,-1.675) (5.700,-1.717) (5.654,-1.750) (5.600,-1.787) (5.550,-1.819) (5.500,-1.850) (5.456,-1.875) (5.409,-1.900) (5.360,-1.925) (5.307,-1.950) (5.251,-1.975) (5.200,-1.996) (5.150,-2.015) (5.100,-2.033) (5.050,-2.049) (5.000,-2.064) (4.950,-2.078) (4.900,-2.091) (4.850,-2.102) (4.800,-2.113) (4.750,-2.122) (4.700,-2.130) (4.650,-2.137) (4.600,-2.143) (4.550,-2.148) (4.500,-2.152) (4.450,-2.155) (4.400,-2.157) (4.350,-2.158) (4.300,-2.158) (4.250,-2.158) (4.200,-2.156) (4.150,-2.153) (4.100,-2.150) (4.050,-2.146) (4.000,-2.140) (3.950,-2.134) (3.900,-2.128) (3.850,-2.120) (3.800,-2.111) (3.750,-2.102) (3.700,-2.092) (3.650,-2.081) (3.600,-2.070) (3.550,-2.057) (3.500,-2.044) (3.450,-2.030) (3.400,-2.016) (3.350,-2.001) (3.300,-1.985) (3.250,-1.968) (3.200,-1.951) (3.150,-1.933) (3.100,-1.914) (3.050,-1.895) (3.001,-1.875) (2.950,-1.854) (2.900,-1.833) (2.850,-1.811) (2.800,-1.789) (2.750,-1.766) (2.700,-1.742) (2.650,-1.718) (2.600,-1.694) (2.550,-1.669) (2.500,-1.643) (2.450,-1.617) (2.400,-1.590) (2.350,-1.563) (2.300,-1.535) (2.250,-1.507) (2.200,-1.479) (2.150,-1.450) (2.100,-1.421) (2.050,-1.391) (2.000,-1.361) (1.950,-1.331) (1.900,-1.300) (1.850,-1.269) (1.800,-1.238) (1.750,-1.206) (1.702,-1.175) (1.650,-1.142) (1.600,-1.109) (1.550,-1.076) (1.500,-1.044) (1.450,-1.010) (1.400,-0.977) (1.350,-0.943) (1.300,-0.910) (1.250,-0.876) (1.200,-0.842) (1.150,-0.807) (1.103,-0.775) (1.050,-0.739) (1.000,-0.704) (0.950,-0.669) (0.900,-0.634) (0.851,-0.600) (0.800,-0.565) (0.750,-0.529) (0.709,-0.500) (0.650,-0.459) (0.602,-0.425) (0.550,-0.389) (0.500,-0.354) (0.450,-0.319) (0.400,-0.283) (0.354,-0.250) (0.300,-0.213) (0.250,-0.177) (0.206,-0.144)}
\EQ{(-0.220,-0.155) (-0.275,-0.195) (-0.319,-0.225) (-0.375,-0.266) (-0.424,-0.300) (-0.475,-0.336) (-0.525,-0.371) (-0.567,-0.400) (-0.625,-0.442) (-0.673,-0.475) (-0.725,-0.512) (-0.775,-0.547) (-0.825,-0.583) (-0.875,-0.618) (-0.922,-0.650) (-0.975,-0.688) (-1.025,-0.722) (-1.075,-0.757) (-1.125,-0.792) (-1.173,-0.825) (-1.225,-0.861) (-1.275,-0.896) (-1.318,-0.925) (-1.375,-0.964) (-1.425,-0.998) (-1.475,-1.032) (-1.525,-1.066) (-1.575,-1.100) (-1.625,-1.133) (-1.675,-1.167) (-1.725,-1.200) (-1.775,-1.233) (-1.825,-1.266) (-1.875,-1.298) (-1.925,-1.330) (-1.975,-1.362) (-2.025,-1.394) (-2.074,-1.425) (-2.125,-1.457) (-2.175,-1.488) (-2.225,-1.519) (-2.275,-1.549) (-2.325,-1.580) (-2.375,-1.609) (-2.425,-1.639) (-2.475,-1.668) (-2.525,-1.697) (-2.575,-1.725) (-2.619,-1.750) (-2.665,-1.775) (-2.710,-1.800) (-2.757,-1.825) (-2.803,-1.850) (-2.850,-1.874) (-2.900,-1.900) (-2.949,-1.925) (-2.999,-1.950) (-3.050,-1.975) (-3.100,-1.999) (-3.150,-2.022) (-3.200,-2.045) (-3.250,-2.068) (-3.300,-2.090) (-3.350,-2.111) (-3.400,-2.132) (-3.450,-2.152) (-3.500,-2.172) (-3.550,-2.191) (-3.600,-2.210) (-3.650,-2.228) (-3.700,-2.246) (-3.750,-2.263) (-3.800,-2.279) (-3.850,-2.295) (-3.900,-2.310) (-3.950,-2.324) (-4.000,-2.338) (-4.050,-2.352) (-4.100,-2.364) (-4.150,-2.376) (-4.200,-2.387) (-4.250,-2.398) (-4.300,-2.408) (-4.350,-2.417) (-4.400,-2.425) (-4.450,-2.433) (-4.500,-2.440) (-4.550,-2.447) (-4.600,-2.452) (-4.650,-2.457) (-4.700,-2.462) (-4.750,-2.465) (-4.800,-2.468) (-4.850,-2.469) (-4.900,-2.470) (-4.950,-2.471) (-5.000,-2.470) (-5.050,-2.469) (-5.100,-2.467) (-5.150,-2.464) (-5.200,-2.460) (-5.250,-2.455) (-5.300,-2.449) (-5.350,-2.443) (-5.400,-2.435) (-5.450,-2.427) (-5.500,-2.417) (-5.550,-2.407) (-5.600,-2.395) (-5.650,-2.383) (-5.700,-2.370) (-5.750,-2.355) (-5.800,-2.339) (-5.850,-2.323) (-5.900,-2.305) (-5.950,-2.286) (-6.000,-2.265) (-6.050,-2.244) (-6.100,-2.221) (-6.150,-2.196) (-6.200,-2.171) (-6.250,-2.144) (-6.300,-2.115) (-6.350,-2.085) (-6.400,-2.053) (-6.450,-2.019) (-6.500,-1.983) (-6.545,-1.950) (-6.600,-1.906) (-6.650,-1.864) (-6.695,-1.825) (-6.748,-1.775) (-6.799,-1.725) (-6.846,-1.675) (-6.892,-1.625) (-6.925,-1.586) (-6.974,-1.525) (-7.013,-1.475) (-7.049,-1.425) (-7.083,-1.375) (-7.116,-1.325) (-7.147,-1.275) (-7.175,-1.226) (-7.200,-1.181) (-7.225,-1.134) (-7.250,-1.084) (-7.275,-1.031) (-7.300,-0.975) (-7.320,-0.925) (-7.340,-0.875) (-7.358,-0.825) (-7.375,-0.775) (-7.391,-0.725) (-7.405,-0.675) (-7.419,-0.625) (-7.431,-0.575) (-7.443,-0.525) (-7.453,-0.475) (-7.462,-0.425) (-7.470,-0.375) (-7.477,-0.325) (-7.483,-0.275) (-7.488,-0.225) (-7.492,-0.175) (-7.495,-0.125) (-7.497,-0.075) (-7.498,-0.025) (-7.498,0.050) (-7.496,0.100) (-7.494,0.150) (-7.490,0.200) (-7.486,0.250) (-7.480,0.300) (-7.474,0.350) (-7.466,0.400) (-7.457,0.450) (-7.448,0.500) (-7.437,0.550) (-7.425,0.600) (-7.412,0.650) (-7.398,0.700) (-7.383,0.750) (-7.367,0.800) (-7.349,0.850) (-7.330,0.900) (-7.310,0.950) (-7.289,1.000) (-7.266,1.050) (-7.242,1.100) (-7.217,1.150) (-7.190,1.200) (-7.161,1.250) (-7.131,1.300) (-7.100,1.350) (-7.066,1.400) (-7.031,1.450) (-7.000,1.492) (-6.955,1.550) (-6.913,1.600) (-6.875,1.644) (-6.825,1.698) (-6.775,1.749) (-6.725,1.797) (-6.675,1.843) (-6.625,1.886) (-6.577,1.925) (-6.525,1.965) (-6.477,2.000) (-6.425,2.036) (-6.375,2.069) (-6.325,2.100) (-6.275,2.129) (-6.225,2.157) (-6.175,2.184) (-6.125,2.209) (-6.075,2.232) (-6.025,2.255) (-5.977,2.275) (-5.925,2.295) (-5.875,2.314) (-5.825,2.331) (-5.775,2.347) (-5.725,2.362) (-5.675,2.376) (-5.625,2.389) (-5.575,2.401) (-5.525,2.412) (-5.475,2.422) (-5.425,2.431) (-5.375,2.439) (-5.325,2.446) (-5.275,2.452) (-5.225,2.457) (-5.175,2.462) (-5.125,2.465) (-5.075,2.468) (-5.025,2.470) (-4.975,2.471) (-4.925,2.471) (-4.875,2.470) (-4.825,2.469) (-4.775,2.466) (-4.725,2.463) (-4.675,2.460) (-4.625,2.455) (-4.575,2.450) (-4.525,2.444) (-4.475,2.437) (-4.425,2.430) (-4.375,2.421) (-4.325,2.412) (-4.275,2.403) (-4.225,2.393) (-4.175,2.382) (-4.125,2.370) (-4.075,2.358) (-4.025,2.345) (-3.975,2.331) (-3.925,2.317) (-3.875,2.302) (-3.825,2.287) (-3.775,2.271) (-3.725,2.254) (-3.675,2.237) (-3.625,2.219) (-3.575,2.201) (-3.525,2.182) (-3.475,2.162) (-3.425,2.142) (-3.375,2.122) (-3.325,2.100) (-3.275,2.079) (-3.225,2.057) (-3.175,2.034) (-3.125,2.011) (-3.075,1.987) (-3.025,1.963) (-2.975,1.938) (-2.925,1.913) (-2.875,1.887) (-2.825,1.861) (-2.775,1.835) (-2.725,1.808) (-2.675,1.781) (-2.625,1.753) (-2.575,1.725) (-2.531,1.700) (-2.487,1.675) (-2.425,1.639) (-2.375,1.609) (-2.325,1.580) (-2.276,1.550) (-2.225,1.519) (-2.175,1.488) (-2.125,1.457) (-2.075,1.426) (-2.025,1.394) (-1.975,1.362) (-1.925,1.330) (-1.878,1.300) (-1.825,1.266) (-1.775,1.233) (-1.725,1.200) (-1.675,1.167) (-1.625,1.133) (-1.576,1.100) (-1.525,1.066) (-1.475,1.032) (-1.428,1.000) (-1.375,0.964) (-1.325,0.930) (-1.282,0.900) (-1.225,0.861) (-1.175,0.827) (-1.125,0.792) (-1.075,0.757) (-1.029,0.725) (-0.975,0.688) (-0.925,0.652) (-0.875,0.618) (-0.825,0.583) (-0.779,0.550) (-0.725,0.512) (-0.675,0.477) (-0.625,0.442) (-0.575,0.407) (-0.531,0.375) (-0.475,0.336) (-0.425,0.300) (-0.375,0.266) (-0.325,0.230) (-0.284,0.200) (-0.225,0.160) (-0.220,0.155)}
\EQ{(0.206,0.144) (0.248,0.175) (0.300,0.213) (0.350,0.248) (0.400,0.283) (0.450,0.319) (0.496,0.350) (0.550,0.389) (0.600,0.424) (0.650,0.460) (0.700,0.495) (0.744,0.525) (0.800,0.565) (0.850,0.600) (0.900,0.635) (0.950,0.670) (0.993,0.700) (1.050,0.740) (1.100,0.774) (1.150,0.809) (1.200,0.844) (1.245,0.875) (1.300,0.913) (1.350,0.947) (1.400,0.981) (1.450,1.015) (1.500,1.049) (1.550,1.083) (1.600,1.117) (1.650,1.150) (1.700,1.183) (1.750,1.216) (1.800,1.249) (1.850,1.282) (1.900,1.314) (1.950,1.346) (1.995,1.375) (2.050,1.410) (2.100,1.442) (2.150,1.473) (2.194,1.500) (2.250,1.534) (2.300,1.565) (2.350,1.594) (2.400,1.624) (2.444,1.650) (2.500,1.682) (2.550,1.711) (2.600,1.739) (2.650,1.767) (2.700,1.795) (2.750,1.822) (2.800,1.848) (2.850,1.874) (2.900,1.900) (2.949,1.925) (2.999,1.950) (3.050,1.975) (3.100,1.999) (3.150,2.022) (3.200,2.045) (3.250,2.068) (3.300,2.090) (3.350,2.111) (3.400,2.132) (3.450,2.152) (3.500,2.172) (3.550,2.191) (3.600,2.210) (3.650,2.228) (3.700,2.246) (3.750,2.263) (3.800,2.279) (3.850,2.295) (3.900,2.310) (3.950,2.324) (4.000,2.338) (4.050,2.352) (4.100,2.364) (4.150,2.376) (4.200,2.387) (4.250,2.398) (4.300,2.408) (4.350,2.417) (4.400,2.425) (4.450,2.433) (4.500,2.440) (4.550,2.447) (4.600,2.452) (4.650,2.457) (4.700,2.462) (4.750,2.465) (4.800,2.468) (4.850,2.469) (4.900,2.470) (4.950,2.471) (5.000,2.470) (5.050,2.469) (5.100,2.467) (5.150,2.464) (5.200,2.460) (5.250,2.455) (5.300,2.449) (5.350,2.443) (5.400,2.435) (5.450,2.427) (5.500,2.417) (5.550,2.407) (5.600,2.395) (5.650,2.383) (5.700,2.370) (5.750,2.355) (5.800,2.339) (5.850,2.323) (5.900,2.305) (5.950,2.286) (6.000,2.265) (6.050,2.244) (6.100,2.221) (6.150,2.196) (6.200,2.171) (6.250,2.144) (6.300,2.115) (6.350,2.085) (6.400,2.053) (6.450,2.019) (6.500,1.983) (6.545,1.950) (6.600,1.906) (6.650,1.864) (6.695,1.825) (6.748,1.775) (6.799,1.725) (6.846,1.675) (6.892,1.625) (6.925,1.586) (6.974,1.525) (7.013,1.475) (7.049,1.425) (7.083,1.375) (7.116,1.325) (7.147,1.275) (7.175,1.226) (7.200,1.181) (7.225,1.134) (7.250,1.084) (7.275,1.031) (7.300,0.975) (7.320,0.925) (7.340,0.875) (7.358,0.825) (7.375,0.775) (7.391,0.725) (7.405,0.675) (7.419,0.625) (7.431,0.575) (7.443,0.525) (7.453,0.475) (7.462,0.425) (7.470,0.375) (7.477,0.325) (7.483,0.275) (7.488,0.225) (7.492,0.175) (7.495,0.125) (7.497,0.075) (7.498,0.025) (7.498,-0.050) (7.496,-0.100) (7.494,-0.150) (7.490,-0.200) (7.486,-0.250) (7.480,-0.300) (7.474,-0.350) (7.466,-0.400) (7.457,-0.450) (7.448,-0.500) (7.437,-0.550) (7.425,-0.600) (7.412,-0.650) (7.398,-0.700) (7.383,-0.750) (7.367,-0.800) (7.349,-0.850) (7.330,-0.900) (7.310,-0.950) (7.289,-1.000) (7.266,-1.050) (7.242,-1.100) (7.217,-1.150) (7.190,-1.200) (7.161,-1.250) (7.131,-1.300) (7.100,-1.350) (7.066,-1.400) (7.031,-1.450) (7.000,-1.492) (6.955,-1.550) (6.913,-1.600) (6.875,-1.644) (6.825,-1.698) (6.775,-1.749) (6.725,-1.797) (6.675,-1.843) (6.625,-1.886) (6.577,-1.925) (6.525,-1.965) (6.477,-2.000) (6.425,-2.036) (6.375,-2.069) (6.325,-2.100) (6.275,-2.129) (6.225,-2.157) (6.175,-2.184) (6.125,-2.209) (6.075,-2.232) (6.025,-2.255) (5.977,-2.275) (5.925,-2.295) (5.875,-2.314) (5.825,-2.331) (5.775,-2.347) (5.725,-2.362) (5.675,-2.376) (5.625,-2.389) (5.575,-2.401) (5.525,-2.412) (5.475,-2.422) (5.425,-2.431) (5.375,-2.439) (5.325,-2.446) (5.275,-2.452) (5.225,-2.457) (5.175,-2.462) (5.125,-2.465) (5.075,-2.468) (5.025,-2.470) (4.975,-2.471) (4.925,-2.471) (4.875,-2.470) (4.825,-2.469) (4.775,-2.466) (4.725,-2.463) (4.675,-2.460) (4.625,-2.455) (4.575,-2.450) (4.525,-2.444) (4.475,-2.437) (4.425,-2.430) (4.375,-2.421) (4.325,-2.412) (4.275,-2.403) (4.225,-2.393) (4.175,-2.382) (4.125,-2.370) (4.075,-2.358) (4.025,-2.345) (3.975,-2.331) (3.925,-2.317) (3.875,-2.302) (3.825,-2.287) (3.775,-2.271) (3.725,-2.254) (3.675,-2.237) (3.625,-2.219) (3.575,-2.201) (3.525,-2.182) (3.475,-2.162) (3.425,-2.142) (3.375,-2.122) (3.325,-2.100) (3.275,-2.079) (3.225,-2.057) (3.175,-2.034) (3.125,-2.011) (3.075,-1.987) (3.025,-1.963) (2.975,-1.938) (2.925,-1.913) (2.875,-1.887) (2.825,-1.861) (2.775,-1.835) (2.725,-1.808) (2.675,-1.781) (2.625,-1.753) (2.575,-1.725) (2.531,-1.700) (2.487,-1.675) (2.425,-1.639) (2.375,-1.609) (2.325,-1.580) (2.276,-1.550) (2.225,-1.519) (2.175,-1.488) (2.125,-1.457) (2.075,-1.426) (2.025,-1.394) (1.975,-1.362) (1.925,-1.330) (1.878,-1.300) (1.825,-1.266) (1.775,-1.233) (1.725,-1.200) (1.675,-1.167) (1.625,-1.133) (1.576,-1.100) (1.525,-1.066) (1.475,-1.032) (1.428,-1.000) (1.375,-0.964) (1.325,-0.930) (1.282,-0.900) (1.225,-0.861) (1.175,-0.827) (1.125,-0.792) (1.075,-0.757) (1.029,-0.725) (0.975,-0.688) (0.925,-0.652) (0.875,-0.618) (0.825,-0.583) (0.779,-0.550) (0.725,-0.512) (0.675,-0.477) (0.625,-0.442) (0.575,-0.407) (0.531,-0.375) (0.475,-0.336) (0.425,-0.300) (0.375,-0.266) (0.325,-0.230) (0.284,-0.200) (0.225,-0.160) (0.206,-0.144)}
\EQ{(-0.220,-0.155) (-0.275,-0.195) (-0.319,-0.225) (-0.375,-0.266) (-0.424,-0.300) (-0.475,-0.336) (-0.525,-0.371) (-0.567,-0.400) (-0.625,-0.442) (-0.672,-0.475) (-0.725,-0.513) (-0.775,-0.548) (-0.825,-0.583) (-0.875,-0.618) (-0.921,-0.650) (-0.975,-0.689) (-1.025,-0.724) (-1.075,-0.759) (-1.125,-0.794) (-1.170,-0.825) (-1.225,-0.864) (-1.275,-0.899) (-1.325,-0.934) (-1.375,-0.968) (-1.421,-1.000) (-1.475,-1.038) (-1.525,-1.072) (-1.575,-1.107) (-1.625,-1.141) (-1.674,-1.175) (-1.725,-1.210) (-1.775,-1.244) (-1.821,-1.275) (-1.875,-1.312) (-1.925,-1.346) (-1.969,-1.375) (-2.025,-1.413) (-2.075,-1.446) (-2.118,-1.475) (-2.175,-1.513) (-2.225,-1.546) (-2.270,-1.575) (-2.325,-1.611) (-2.375,-1.644) (-2.424,-1.675) (-2.475,-1.708) (-2.525,-1.740) (-2.575,-1.772) (-2.620,-1.800) (-2.675,-1.834) (-2.725,-1.866) (-2.775,-1.896) (-2.822,-1.925) (-2.875,-1.957) (-2.925,-1.988) (-2.975,-2.017) (-3.025,-2.047) (-3.073,-2.075) (-3.125,-2.105) (-3.175,-2.134) (-3.225,-2.163) (-3.275,-2.191) (-3.325,-2.219) (-3.375,-2.246) (-3.425,-2.274) (-3.474,-2.300) (-3.521,-2.325) (-3.569,-2.350) (-3.617,-2.375) (-3.665,-2.400) (-3.715,-2.425) (-3.765,-2.450) (-3.816,-2.475) (-3.867,-2.500) (-3.920,-2.525) (-3.973,-2.550) (-4.025,-2.574) (-4.075,-2.596) (-4.125,-2.619) (-4.175,-2.640) (-4.225,-2.662) (-4.275,-2.683) (-4.325,-2.703) (-4.375,-2.723) (-4.425,-2.743) (-4.475,-2.762) (-4.525,-2.781) (-4.575,-2.799) (-4.625,-2.817) (-4.675,-2.834) (-4.725,-2.851) (-4.775,-2.867) (-4.825,-2.883) (-4.875,-2.899) (-4.925,-2.914) (-4.975,-2.928) (-5.025,-2.942) (-5.075,-2.956) (-5.125,-2.968) (-5.175,-2.981) (-5.225,-2.993) (-5.275,-3.004) (-5.325,-3.015) (-5.375,-3.025) (-5.425,-3.035) (-5.475,-3.044) (-5.525,-3.053) (-5.575,-3.061) (-5.625,-3.069) (-5.675,-3.075) (-5.725,-3.082) (-5.775,-3.088) (-5.825,-3.093) (-5.875,-3.098) (-5.925,-3.102) (-5.975,-3.105) (-6.025,-3.108) (-6.075,-3.110) (-6.125,-3.112) (-6.175,-3.113) (-6.225,-3.113) (-6.275,-3.113) (-6.325,-3.112) (-6.375,-3.110) (-6.425,-3.108) (-6.475,-3.105) (-6.525,-3.101) (-6.575,-3.097) (-6.625,-3.092) (-6.675,-3.086) (-6.725,-3.080) (-6.775,-3.072) (-6.825,-3.065) (-6.875,-3.056) (-6.925,-3.046) (-6.975,-3.036) (-7.025,-3.025) (-7.075,-3.013) (-7.125,-3.001) (-7.175,-2.987) (-7.225,-2.973) (-7.275,-2.958) (-7.325,-2.942) (-7.375,-2.925) (-7.425,-2.907) (-7.475,-2.888) (-7.525,-2.868) (-7.575,-2.848) (-7.625,-2.826) (-7.675,-2.803) (-7.725,-2.779) (-7.775,-2.754) (-7.825,-2.728) (-7.875,-2.701) (-7.920,-2.675) (-7.975,-2.642) (-8.025,-2.611) (-8.075,-2.579) (-8.117,-2.550) (-8.175,-2.509) (-8.222,-2.475) (-8.275,-2.434) (-8.318,-2.400) (-8.375,-2.352) (-8.425,-2.308) (-8.475,-2.262) (-8.514,-2.225) (-8.550,-2.189) (-8.589,-2.150) (-8.625,-2.111) (-8.675,-2.056) (-8.723,-2.000) (-8.765,-1.950) (-8.800,-1.905) (-8.842,-1.850) (-8.875,-1.804) (-8.912,-1.750) (-8.946,-1.700) (-8.975,-1.654) (-9.008,-1.600) (-9.037,-1.550) (-9.065,-1.500) (-9.092,-1.450) (-9.118,-1.400) (-9.142,-1.350) (-9.165,-1.300) (-9.188,-1.250) (-9.209,-1.200) (-9.229,-1.150) (-9.248,-1.100) (-9.267,-1.050) (-9.284,-1.000) (-9.300,-0.951) (-9.316,-0.900) (-9.330,-0.850) (-9.344,-0.800) (-9.357,-0.750) (-9.368,-0.700) (-9.379,-0.650) (-9.390,-0.600) (-9.399,-0.550) (-9.407,-0.500) (-9.415,-0.450) (-9.422,-0.400) (-9.428,-0.350) (-9.433,-0.300) (-9.437,-0.250) (-9.441,-0.200) (-9.444,-0.150) (-9.446,-0.100) (-9.447,-0.050) (-9.447,0.000) (-9.447,0.050) (-9.446,0.100) (-9.444,0.150) (-9.441,0.200) (-9.437,0.250) (-9.433,0.300) (-9.428,0.350) (-9.422,0.400) (-9.415,0.450) (-9.407,0.500) (-9.399,0.550) (-9.390,0.600) (-9.379,0.650) (-9.368,0.700) (-9.357,0.750) (-9.344,0.800) (-9.330,0.850) (-9.316,0.900) (-9.300,0.950) (-9.284,1.000) (-9.267,1.050) (-9.248,1.100) (-9.229,1.150) (-9.209,1.200) (-9.188,1.250) (-9.165,1.300) (-9.142,1.350) (-9.118,1.400) (-9.092,1.450) (-9.065,1.500) (-9.037,1.550) (-9.008,1.600) (-8.978,1.650) (-8.950,1.693) (-8.912,1.750) (-8.878,1.800) (-8.842,1.850) (-8.804,1.900) (-8.765,1.950) (-8.725,1.998) (-8.680,2.050) (-8.636,2.100) (-8.600,2.138) (-8.564,2.175) (-8.525,2.214) (-8.488,2.250) (-8.450,2.285) (-8.406,2.325) (-8.350,2.373) (-8.300,2.414) (-8.255,2.450) (-8.200,2.491) (-8.153,2.525) (-8.100,2.562) (-8.050,2.595) (-8.003,2.625) (-7.950,2.657) (-7.900,2.686) (-7.850,2.714) (-7.800,2.741) (-7.750,2.767) (-7.700,2.791) (-7.650,2.814) (-7.600,2.837) (-7.550,2.858) (-7.500,2.878) (-7.450,2.898) (-7.400,2.916) (-7.350,2.933) (-7.300,2.950) (-7.250,2.966) (-7.200,2.980) (-7.150,2.994) (-7.100,3.007) (-7.050,3.019) (-7.000,3.031) (-6.950,3.041) (-6.900,3.051) (-6.850,3.060) (-6.800,3.069) (-6.750,3.076) (-6.700,3.083) (-6.650,3.089) (-6.600,3.094) (-6.550,3.099) (-6.500,3.103) (-6.450,3.106) (-6.400,3.109) (-6.350,3.111) (-6.300,3.112) (-6.250,3.113) (-6.200,3.113) (-6.150,3.112) (-6.100,3.111) (-6.050,3.109) (-6.000,3.107) (-5.950,3.103) (-5.900,3.100) (-5.850,3.095) (-5.800,3.090) (-5.750,3.085) (-5.700,3.079) (-5.650,3.072) (-5.600,3.065) (-5.550,3.057) (-5.500,3.049) (-5.450,3.040) (-5.400,3.030) (-5.350,3.020) (-5.300,3.010) (-5.250,2.998) (-5.200,2.987) (-5.151,2.975) (-5.100,2.962) (-5.050,2.949) (-5.000,2.935) (-4.950,2.921) (-4.900,2.906) (-4.850,2.891) (-4.800,2.875) (-4.750,2.859) (-4.700,2.843) (-4.650,2.825) (-4.600,2.808) (-4.550,2.790) (-4.500,2.771) (-4.450,2.752) (-4.400,2.733) (-4.350,2.713) (-4.300,2.693) (-4.250,2.672) (-4.200,2.651) (-4.150,2.629) (-4.100,2.608) (-4.050,2.585) (-4.000,2.562) (-3.950,2.539) (-3.900,2.516) (-3.850,2.492) (-3.800,2.467) (-3.750,2.443) (-3.700,2.418) (-3.650,2.392) (-3.600,2.366) (-3.550,2.340) (-3.500,2.314) (-3.450,2.287) (-3.400,2.260) (-3.350,2.233) (-3.300,2.205) (-3.250,2.177) (-3.203,2.150) (-3.159,2.125) (-3.116,2.100) (-3.050,2.062) (-3.000,2.032) (-2.950,2.002) (-2.904,1.975) (-2.850,1.942) (-2.800,1.912) (-2.750,1.881) (-2.700,1.850) (-2.650,1.819) (-2.600,1.787) (-2.550,1.756) (-2.502,1.725) (-2.450,1.692) (-2.400,1.660) (-2.350,1.627) (-2.300,1.595) (-2.250,1.562) (-2.200,1.529) (-2.156,1.500) (-2.100,1.463) (-2.050,1.430) (-2.006,1.400) (-1.950,1.363) (-1.900,1.329) (-1.858,1.300) (-1.800,1.261) (-1.750,1.227) (-1.700,1.193) (-1.650,1.158) (-1.602,1.125) (-1.550,1.090) (-1.500,1.055) (-1.457,1.025) (-1.400,0.986) (-1.350,0.951) (-1.300,0.916) (-1.250,0.881) (-1.206,0.850) (-1.150,0.811) (-1.100,0.776) (-1.050,0.741) (-1.000,0.706) (-0.956,0.675) (-0.900,0.636) (-0.850,0.600) (-0.800,0.566) (-0.750,0.530) (-0.708,0.500) (-0.650,0.460) (-0.601,0.425) (-0.550,0.389) (-0.500,0.354) (-0.450,0.319) (-0.400,0.283) (-0.354,0.250) (-0.300,0.213) (-0.250,0.177) (-0.220,0.155)}
\EQ{(0.206,0.144) (0.248,0.175) (0.300,0.213) (0.350,0.248) (0.400,0.283) (0.450,0.319) (0.496,0.350) (0.550,0.389) (0.600,0.424) (0.650,0.460) (0.700,0.495) (0.743,0.525) (0.800,0.566) (0.849,0.600) (0.900,0.636) (0.950,0.671) (0.992,0.700) (1.050,0.741) (1.098,0.775) (1.150,0.811) (1.200,0.846) (1.242,0.875) (1.300,0.916) (1.349,0.950) (1.400,0.986) (1.450,1.020) (1.493,1.050) (1.550,1.090) (1.600,1.124) (1.650,1.158) (1.700,1.193) (1.748,1.225) (1.800,1.261) (1.850,1.295) (1.895,1.325) (1.950,1.363) (2.000,1.396) (2.043,1.425) (2.100,1.463) (2.150,1.496) (2.194,1.525) (2.250,1.562) (2.300,1.595) (2.347,1.625) (2.400,1.660) (2.450,1.692) (2.500,1.724) (2.550,1.756) (2.600,1.787) (2.650,1.819) (2.700,1.850) (2.750,1.881) (2.800,1.912) (2.850,1.942) (2.900,1.972) (2.946,2.000) (3.000,2.032) (3.050,2.062) (3.100,2.091) (3.150,2.120) (3.200,2.148) (3.247,2.175) (3.292,2.200) (3.337,2.225) (3.382,2.250) (3.428,2.275) (3.474,2.300) (3.521,2.325) (3.569,2.350) (3.617,2.375) (3.665,2.400) (3.715,2.425) (3.765,2.450) (3.816,2.475) (3.867,2.500) (3.920,2.525) (3.973,2.550) (4.025,2.574) (4.075,2.596) (4.125,2.619) (4.175,2.640) (4.225,2.662) (4.275,2.683) (4.325,2.703) (4.375,2.723) (4.425,2.743) (4.475,2.762) (4.525,2.781) (4.575,2.799) (4.625,2.817) (4.675,2.834) (4.725,2.851) (4.775,2.867) (4.825,2.883) (4.875,2.899) (4.925,2.914) (4.975,2.928) (5.025,2.942) (5.075,2.956) (5.125,2.968) (5.175,2.981) (5.225,2.993) (5.275,3.004) (5.325,3.015) (5.375,3.025) (5.425,3.035) (5.475,3.044) (5.525,3.053) (5.575,3.061) (5.625,3.069) (5.675,3.075) (5.725,3.082) (5.775,3.088) (5.825,3.093) (5.875,3.098) (5.925,3.102) (5.975,3.105) (6.025,3.108) (6.075,3.110) (6.125,3.112) (6.175,3.113) (6.225,3.113) (6.275,3.113) (6.325,3.112) (6.375,3.110) (6.425,3.108) (6.475,3.105) (6.525,3.101) (6.575,3.097) (6.625,3.092) (6.675,3.086) (6.725,3.080) (6.775,3.072) (6.825,3.065) (6.875,3.056) (6.925,3.046) (6.975,3.036) (7.025,3.025) (7.075,3.013) (7.125,3.001) (7.175,2.987) (7.225,2.973) (7.275,2.958) (7.325,2.942) (7.375,2.925) (7.425,2.907) (7.475,2.888) (7.525,2.868) (7.575,2.848) (7.625,2.826) (7.675,2.803) (7.725,2.779) (7.775,2.754) (7.825,2.728) (7.875,2.701) (7.920,2.675) (7.975,2.642) (8.025,2.611) (8.075,2.579) (8.117,2.550) (8.175,2.509) (8.222,2.475) (8.275,2.434) (8.318,2.400) (8.375,2.352) (8.425,2.308) (8.475,2.262) (8.514,2.225) (8.550,2.189) (8.589,2.150) (8.625,2.111) (8.675,2.056) (8.723,2.000) (8.765,1.950) (8.800,1.905) (8.842,1.850) (8.875,1.804) (8.912,1.750) (8.946,1.700) (8.975,1.654) (9.008,1.600) (9.037,1.550) (9.065,1.500) (9.092,1.450) (9.118,1.400) (9.142,1.350) (9.165,1.300) (9.188,1.250) (9.209,1.200) (9.229,1.150) (9.248,1.100) (9.267,1.050) (9.284,1.000) (9.300,0.951) (9.316,0.900) (9.330,0.850) (9.344,0.800) (9.357,0.750) (9.368,0.700) (9.379,0.650) (9.390,0.600) (9.399,0.550) (9.407,0.500) (9.415,0.450) (9.422,0.400) (9.428,0.350) (9.433,0.300) (9.437,0.250) (9.441,0.200) (9.444,0.150) (9.446,0.100) (9.447,0.050) (9.447,0.000) (9.447,-0.050) (9.446,-0.100) (9.444,-0.150) (9.441,-0.200) (9.437,-0.250) (9.433,-0.300) (9.428,-0.350) (9.422,-0.400) (9.415,-0.450) (9.407,-0.500) (9.399,-0.550) (9.390,-0.600) (9.379,-0.650) (9.368,-0.700) (9.357,-0.750) (9.344,-0.800) (9.330,-0.850) (9.316,-0.900) (9.300,-0.950) (9.284,-1.000) (9.267,-1.050) (9.248,-1.100) (9.229,-1.150) (9.209,-1.200) (9.188,-1.250) (9.165,-1.300) (9.142,-1.350) (9.118,-1.400) (9.092,-1.450) (9.065,-1.500) (9.037,-1.550) (9.008,-1.600) (8.978,-1.650) (8.950,-1.693) (8.912,-1.750) (8.878,-1.800) (8.842,-1.850) (8.804,-1.900) (8.765,-1.950) (8.725,-1.998) (8.680,-2.050) (8.636,-2.100) (8.600,-2.138) (8.564,-2.175) (8.525,-2.214) (8.488,-2.250) (8.450,-2.285) (8.406,-2.325) (8.350,-2.373) (8.300,-2.414) (8.255,-2.450) (8.200,-2.491) (8.153,-2.525) (8.100,-2.562) (8.050,-2.595) (8.003,-2.625) (7.950,-2.657) (7.900,-2.686) (7.850,-2.714) (7.800,-2.741) (7.750,-2.767) (7.700,-2.791) (7.650,-2.814) (7.600,-2.837) (7.550,-2.858) (7.500,-2.878) (7.450,-2.898) (7.400,-2.916) (7.350,-2.933) (7.300,-2.950) (7.250,-2.966) (7.200,-2.980) (7.150,-2.994) (7.100,-3.007) (7.050,-3.019) (7.000,-3.031) (6.950,-3.041) (6.900,-3.051) (6.850,-3.060) (6.800,-3.069) (6.750,-3.076) (6.700,-3.083) (6.650,-3.089) (6.600,-3.094) (6.550,-3.099) (6.500,-3.103) (6.450,-3.106) (6.400,-3.109) (6.350,-3.111) (6.300,-3.112) (6.250,-3.113) (6.200,-3.113) (6.150,-3.112) (6.100,-3.111) (6.050,-3.109) (6.000,-3.107) (5.950,-3.103) (5.900,-3.100) (5.850,-3.095) (5.800,-3.090) (5.750,-3.085) (5.700,-3.079) (5.650,-3.072) (5.600,-3.065) (5.550,-3.057) (5.500,-3.049) (5.450,-3.040) (5.400,-3.030) (5.350,-3.020) (5.300,-3.010) (5.250,-2.998) (5.200,-2.987) (5.151,-2.975) (5.100,-2.962) (5.050,-2.949) (5.000,-2.935) (4.950,-2.921) (4.900,-2.906) (4.850,-2.891) (4.800,-2.875) (4.750,-2.859) (4.700,-2.843) (4.650,-2.825) (4.600,-2.808) (4.550,-2.790) (4.500,-2.771) (4.450,-2.752) (4.400,-2.733) (4.350,-2.713) (4.300,-2.693) (4.250,-2.672) (4.200,-2.651) (4.150,-2.629) (4.100,-2.608) (4.050,-2.585) (4.000,-2.562) (3.950,-2.539) (3.900,-2.516) (3.850,-2.492) (3.800,-2.467) (3.750,-2.443) (3.700,-2.418) (3.650,-2.392) (3.600,-2.366) (3.550,-2.340) (3.500,-2.314) (3.450,-2.287) (3.400,-2.260) (3.350,-2.233) (3.300,-2.205) (3.250,-2.177) (3.203,-2.150) (3.159,-2.125) (3.116,-2.100) (3.050,-2.062) (3.000,-2.032) (2.950,-2.002) (2.904,-1.975) (2.850,-1.942) (2.800,-1.912) (2.750,-1.881) (2.700,-1.850) (2.650,-1.819) (2.600,-1.787) (2.550,-1.756) (2.502,-1.725) (2.450,-1.692) (2.400,-1.660) (2.350,-1.627) (2.300,-1.595) (2.250,-1.562) (2.200,-1.529) (2.156,-1.500) (2.100,-1.463) (2.050,-1.430) (2.006,-1.400) (1.950,-1.363) (1.900,-1.329) (1.858,-1.300) (1.800,-1.261) (1.750,-1.227) (1.700,-1.193) (1.650,-1.158) (1.602,-1.125) (1.550,-1.090) (1.500,-1.055) (1.457,-1.025) (1.400,-0.986) (1.350,-0.951) (1.300,-0.916) (1.250,-0.881) (1.206,-0.850) (1.150,-0.811) (1.100,-0.776) (1.050,-0.741) (1.000,-0.706) (0.956,-0.675) (0.900,-0.636) (0.850,-0.600) (0.800,-0.566) (0.750,-0.530) (0.708,-0.500) (0.650,-0.460) (0.601,-0.425) (0.550,-0.389) (0.500,-0.354) (0.450,-0.319) (0.400,-0.283) (0.354,-0.250) (0.300,-0.213) (0.250,-0.177) (0.206,-0.144)}
\EQ{(10.000,-7.071) (9.950,-7.036) (9.900,-7.000) (9.850,-6.965) (9.800,-6.930) (9.758,-6.900) (9.700,-6.859) (9.652,-6.825) (9.600,-6.788) (9.550,-6.753) (9.500,-6.718) (9.450,-6.682) (9.405,-6.650) (9.350,-6.611) (9.300,-6.576) (9.250,-6.541) (9.200,-6.505) (9.157,-6.475) (9.100,-6.435) (9.051,-6.400) (9.000,-6.364) (8.950,-6.329) (8.900,-6.293) (8.850,-6.258) (8.804,-6.225) (8.750,-6.187) (8.700,-6.152) (8.650,-6.117) (8.600,-6.081) (8.556,-6.050) (8.500,-6.010) (8.450,-5.975) (8.400,-5.940) (8.350,-5.904) (8.309,-5.875) (8.250,-5.834) (8.202,-5.800) (8.150,-5.763) (8.100,-5.728) (8.050,-5.692) (8.000,-5.657) (7.955,-5.625) (7.900,-5.586) (7.850,-5.551) (7.800,-5.515) (7.750,-5.480) (7.708,-5.450) (7.650,-5.409) (7.601,-5.375) (7.550,-5.339) (7.500,-5.303) (7.450,-5.268) (7.400,-5.233) (7.354,-5.200) (7.300,-5.162) (7.250,-5.127) (7.200,-5.091) (7.150,-5.056) (7.106,-5.025) (7.050,-4.985) (7.000,-4.950) (6.950,-4.914) (6.900,-4.879) (6.859,-4.850) (6.800,-4.808) (6.753,-4.775) (6.700,-4.738) (6.650,-4.702) (6.600,-4.667) (6.550,-4.632) (6.505,-4.600) (6.450,-4.561) (6.400,-4.525) (6.350,-4.490) (6.300,-4.455) (6.258,-4.425) (6.200,-4.384) (6.152,-4.350) (6.100,-4.313) (6.050,-4.278) (6.000,-4.243) (5.950,-4.207) (5.904,-4.175) (5.850,-4.137) (5.800,-4.101) (5.750,-4.066) (5.700,-4.031) (5.657,-4.000) (5.600,-3.960) (5.551,-3.925) (5.500,-3.889) (5.450,-3.854) (5.400,-3.818) (5.350,-3.783) (5.303,-3.750) (5.250,-3.712) (5.200,-3.677) (5.150,-3.642) (5.100,-3.606) (5.056,-3.575) (5.000,-3.536) (4.950,-3.500) (4.900,-3.465) (4.850,-3.429) (4.808,-3.400) (4.750,-3.359) (4.702,-3.325) (4.650,-3.288) (4.600,-3.253) (4.550,-3.217) (4.500,-3.182) (4.455,-3.150) (4.400,-3.111) (4.350,-3.076) (4.300,-3.041) (4.250,-3.005) (4.207,-2.975) (4.150,-2.935) (4.101,-2.900) (4.050,-2.864) (4.000,-2.828) (3.950,-2.793) (3.900,-2.758) (3.854,-2.725) (3.800,-2.687) (3.750,-2.652) (3.700,-2.616) (3.650,-2.581) (3.606,-2.550) (3.550,-2.510) (3.500,-2.475) (3.450,-2.440) (3.400,-2.404) (3.359,-2.375) (3.300,-2.334) (3.253,-2.300) (3.200,-2.263) (3.150,-2.227) (3.100,-2.192) (3.050,-2.157) (3.005,-2.125) (2.950,-2.086) (2.900,-2.051) (2.850,-2.015) (2.800,-1.980) (2.758,-1.950) (2.700,-1.909) (2.652,-1.875) (2.600,-1.839) (2.550,-1.803) (2.500,-1.768) (2.450,-1.732) (2.404,-1.700) (2.350,-1.662) (2.300,-1.626) (2.250,-1.591) (2.200,-1.556) (2.157,-1.525) (2.100,-1.485) (2.051,-1.450) (2.000,-1.414) (1.950,-1.379) (1.900,-1.344) (1.850,-1.308) (1.803,-1.275) (1.750,-1.238) (1.700,-1.202) (1.650,-1.167) (1.600,-1.131) (1.556,-1.100) (1.500,-1.061) (1.450,-1.025) (1.400,-0.990) (1.350,-0.955) (1.308,-0.925) (1.250,-0.884) (1.202,-0.850) (1.150,-0.813) (1.100,-0.778) (1.050,-0.743) (1.000,-0.707) (0.955,-0.675) (0.900,-0.637) (0.850,-0.601) (0.800,-0.566) (0.750,-0.531) (0.708,-0.500) (0.650,-0.460) (0.601,-0.425) (0.550,-0.389) (0.500,-0.354) (0.450,-0.319) (0.400,-0.283) (0.354,-0.250) (0.300,-0.213) (0.250,-0.177) (0.206,-0.144)}
\EQ{(-0.220,-0.155) (-0.275,-0.195) (-0.319,-0.225) (-0.375,-0.266) (-0.424,-0.300) (-0.475,-0.336) (-0.525,-0.371) (-0.566,-0.400) (-0.625,-0.442) (-0.672,-0.475) (-0.725,-0.513) (-0.775,-0.548) (-0.825,-0.584) (-0.875,-0.619) (-0.920,-0.650) (-0.975,-0.690) (-1.025,-0.725) (-1.075,-0.760) (-1.125,-0.796) (-1.167,-0.825) (-1.225,-0.866) (-1.273,-0.900) (-1.325,-0.937) (-1.375,-0.972) (-1.425,-1.008) (-1.475,-1.043) (-1.520,-1.075) (-1.575,-1.114) (-1.625,-1.149) (-1.675,-1.185) (-1.725,-1.220) (-1.768,-1.250) (-1.825,-1.291) (-1.874,-1.325) (-1.925,-1.361) (-1.975,-1.397) (-2.025,-1.432) (-2.075,-1.467) (-2.121,-1.500) (-2.175,-1.538) (-2.225,-1.573) (-2.275,-1.609) (-2.325,-1.644) (-2.369,-1.675) (-2.425,-1.715) (-2.475,-1.750) (-2.525,-1.786) (-2.575,-1.821) (-2.616,-1.850) (-2.675,-1.892) (-2.722,-1.925) (-2.775,-1.962) (-2.825,-1.998) (-2.875,-2.033) (-2.925,-2.068) (-2.970,-2.100) (-3.025,-2.139) (-3.075,-2.174) (-3.125,-2.210) (-3.175,-2.245) (-3.217,-2.275) (-3.275,-2.316) (-3.323,-2.350) (-3.375,-2.387) (-3.425,-2.422) (-3.475,-2.457) (-3.525,-2.493) (-3.571,-2.525) (-3.625,-2.563) (-3.675,-2.599) (-3.725,-2.634) (-3.775,-2.669) (-3.818,-2.700) (-3.875,-2.740) (-3.924,-2.775) (-3.975,-2.811) (-4.025,-2.846) (-4.066,-2.875) (-4.125,-2.917) (-4.172,-2.950) (-4.225,-2.988) (-4.275,-3.023) (-4.325,-3.058) (-4.375,-3.094) (-4.419,-3.125) (-4.475,-3.164) (-4.525,-3.200) (-4.575,-3.235) (-4.625,-3.270) (-4.667,-3.300) (-4.725,-3.341) (-4.773,-3.375) (-4.825,-3.412) (-4.875,-3.447) (-4.925,-3.483) (-4.975,-3.518) (-5.021,-3.550) (-5.075,-3.589) (-5.125,-3.624) (-5.175,-3.659) (-5.225,-3.695) (-5.268,-3.725) (-5.325,-3.765) (-5.374,-3.800) (-5.425,-3.836) (-5.475,-3.871) (-5.525,-3.907) (-5.575,-3.942) (-5.622,-3.975) (-5.675,-4.013) (-5.725,-4.048) (-5.775,-4.084) (-5.825,-4.119) (-5.869,-4.150) (-5.925,-4.190) (-5.975,-4.225) (-6.025,-4.260) (-6.075,-4.296) (-6.117,-4.325) (-6.175,-4.366) (-6.223,-4.400) (-6.275,-4.437) (-6.325,-4.472) (-6.375,-4.508) (-6.425,-4.543) (-6.470,-4.575) (-6.525,-4.614) (-6.575,-4.649) (-6.625,-4.685) (-6.675,-4.720) (-6.718,-4.750) (-6.775,-4.791) (-6.824,-4.825) (-6.875,-4.861) (-6.925,-4.897) (-6.975,-4.932) (-7.025,-4.967) (-7.071,-5.000) (-7.125,-5.038) (-7.175,-5.073) (-7.225,-5.109) (-7.275,-5.144) (-7.319,-5.175) (-7.375,-5.215) (-7.425,-5.250) (-7.475,-5.286) (-7.525,-5.321) (-7.566,-5.350) (-7.625,-5.392) (-7.672,-5.425) (-7.725,-5.462) (-7.775,-5.498) (-7.825,-5.533) (-7.875,-5.568) (-7.920,-5.600) (-7.975,-5.639) (-8.025,-5.675) (-8.075,-5.710) (-8.125,-5.745) (-8.167,-5.775) (-8.225,-5.816) (-8.273,-5.850) (-8.325,-5.887) (-8.375,-5.922) (-8.425,-5.957) (-8.475,-5.993) (-8.521,-6.025) (-8.575,-6.063) (-8.625,-6.099) (-8.675,-6.134) (-8.725,-6.170) (-8.768,-6.200) (-8.825,-6.240) (-8.874,-6.275) (-8.925,-6.311) (-8.975,-6.346) (-9.025,-6.382) (-9.075,-6.417) (-9.122,-6.450) (-9.175,-6.488) (-9.225,-6.523) (-9.275,-6.558) (-9.325,-6.594) (-9.369,-6.625) (-9.425,-6.665) (-9.475,-6.700) (-9.525,-6.735) (-9.575,-6.771) (-9.617,-6.800) (-9.675,-6.841) (-9.723,-6.875) (-9.775,-6.912) (-9.825,-6.947) (-9.875,-6.983) (-9.925,-7.018) (-9.970,-7.050) (-10.000,-7.071)}
\EQ{(0.206,0.144) (0.248,0.175) (0.300,0.213) (0.350,0.248) (0.400,0.283) (0.450,0.319) (0.496,0.350) (0.550,0.389) (0.600,0.424) (0.650,0.460) (0.700,0.495) (0.743,0.525) (0.800,0.566) (0.849,0.600) (0.900,0.637) (0.950,0.672) (1.000,0.707) (1.050,0.743) (1.096,0.775) (1.150,0.813) (1.200,0.849) (1.250,0.884) (1.300,0.919) (1.344,0.950) (1.400,0.990) (1.450,1.025) (1.500,1.061) (1.550,1.096) (1.591,1.125) (1.650,1.167) (1.697,1.200) (1.750,1.238) (1.800,1.273) (1.850,1.308) (1.900,1.344) (1.945,1.375) (2.000,1.414) (2.050,1.450) (2.100,1.485) (2.150,1.520) (2.192,1.550) (2.250,1.591) (2.298,1.625) (2.350,1.662) (2.400,1.697) (2.450,1.732) (2.500,1.768) (2.546,1.800) (2.600,1.839) (2.650,1.874) (2.700,1.909) (2.750,1.945) (2.793,1.975) (2.850,2.015) (2.899,2.050) (2.950,2.086) (3.000,2.121) (3.050,2.157) (3.100,2.192) (3.147,2.225) (3.200,2.263) (3.250,2.298) (3.300,2.334) (3.350,2.369) (3.394,2.400) (3.450,2.440) (3.500,2.475) (3.550,2.510) (3.600,2.546) (3.642,2.575) (3.700,2.616) (3.748,2.650) (3.800,2.687) (3.850,2.722) (3.900,2.758) (3.950,2.793) (3.995,2.825) (4.050,2.864) (4.100,2.899) (4.150,2.935) (4.200,2.970) (4.243,3.000) (4.300,3.041) (4.349,3.075) (4.400,3.111) (4.450,3.147) (4.500,3.182) (4.550,3.217) (4.596,3.250) (4.650,3.288) (4.700,3.323) (4.750,3.359) (4.800,3.394) (4.844,3.425) (4.900,3.465) (4.950,3.500) (5.000,3.536) (5.050,3.571) (5.091,3.600) (5.150,3.642) (5.197,3.675) (5.250,3.712) (5.300,3.748) (5.350,3.783) (5.400,3.818) (5.445,3.850) (5.500,3.889) (5.550,3.924) (5.600,3.960) (5.650,3.995) (5.692,4.025) (5.750,4.066) (5.798,4.100) (5.850,4.137) (5.900,4.172) (5.950,4.207) (6.000,4.243) (6.046,4.275) (6.100,4.313) (6.150,4.349) (6.200,4.384) (6.250,4.419) (6.293,4.450) (6.350,4.490) (6.399,4.525) (6.450,4.561) (6.500,4.596) (6.550,4.632) (6.600,4.667) (6.647,4.700) (6.700,4.738) (6.750,4.773) (6.800,4.808) (6.850,4.844) (6.894,4.875) (6.950,4.914) (7.000,4.950) (7.050,4.985) (7.100,5.020) (7.142,5.050) (7.200,5.091) (7.248,5.125) (7.300,5.162) (7.350,5.197) (7.400,5.233) (7.450,5.268) (7.495,5.300) (7.550,5.339) (7.600,5.374) (7.650,5.409) (7.700,5.445) (7.743,5.475) (7.800,5.515) (7.849,5.550) (7.900,5.586) (7.950,5.622) (8.000,5.657) (8.050,5.692) (8.096,5.725) (8.150,5.763) (8.200,5.798) (8.250,5.834) (8.300,5.869) (8.344,5.900) (8.400,5.940) (8.450,5.975) (8.500,6.010) (8.550,6.046) (8.591,6.075) (8.650,6.117) (8.697,6.150) (8.750,6.187) (8.800,6.223) (8.850,6.258) (8.900,6.293) (8.945,6.325) (9.000,6.364) (9.050,6.399) (9.100,6.435) (9.150,6.470) (9.192,6.500) (9.250,6.541) (9.298,6.575) (9.350,6.611) (9.400,6.647) (9.450,6.682) (9.500,6.718) (9.546,6.750) (9.600,6.788) (9.650,6.824) (9.700,6.859) (9.750,6.894) (9.793,6.925) (9.850,6.965) (9.899,7.000) (9.950,7.036) (10.000,7.071)}
\EQ{(-10.000,7.071) (-9.950,7.036) (-9.900,7.000) (-9.850,6.965) (-9.800,6.930) (-9.758,6.900) (-9.700,6.859) (-9.652,6.825) (-9.600,6.788) (-9.550,6.753) (-9.500,6.718) (-9.450,6.682) (-9.405,6.650) (-9.350,6.611) (-9.300,6.576) (-9.250,6.541) (-9.200,6.505) (-9.157,6.475) (-9.100,6.435) (-9.051,6.400) (-9.000,6.364) (-8.950,6.329) (-8.900,6.293) (-8.850,6.258) (-8.804,6.225) (-8.750,6.187) (-8.700,6.152) (-8.650,6.117) (-8.600,6.081) (-8.556,6.050) (-8.500,6.010) (-8.450,5.975) (-8.400,5.940) (-8.350,5.904) (-8.309,5.875) (-8.250,5.834) (-8.202,5.800) (-8.150,5.763) (-8.100,5.728) (-8.050,5.692) (-8.000,5.657) (-7.955,5.625) (-7.900,5.586) (-7.850,5.551) (-7.800,5.515) (-7.750,5.480) (-7.708,5.450) (-7.650,5.409) (-7.601,5.375) (-7.550,5.339) (-7.500,5.303) (-7.450,5.268) (-7.400,5.233) (-7.354,5.200) (-7.300,5.162) (-7.250,5.127) (-7.200,5.091) (-7.150,5.056) (-7.106,5.025) (-7.050,4.985) (-7.000,4.950) (-6.950,4.914) (-6.900,4.879) (-6.859,4.850) (-6.800,4.808) (-6.753,4.775) (-6.700,4.738) (-6.650,4.702) (-6.600,4.667) (-6.550,4.632) (-6.505,4.600) (-6.450,4.561) (-6.400,4.525) (-6.350,4.490) (-6.300,4.455) (-6.258,4.425) (-6.200,4.384) (-6.152,4.350) (-6.100,4.313) (-6.050,4.278) (-6.000,4.243) (-5.950,4.207) (-5.904,4.175) (-5.850,4.137) (-5.800,4.101) (-5.750,4.066) (-5.700,4.031) (-5.657,4.000) (-5.600,3.960) (-5.551,3.925) (-5.500,3.889) (-5.450,3.854) (-5.400,3.818) (-5.350,3.783) (-5.303,3.750) (-5.250,3.712) (-5.200,3.677) (-5.150,3.642) (-5.100,3.606) (-5.056,3.575) (-5.000,3.536) (-4.950,3.500) (-4.900,3.465) (-4.850,3.429) (-4.808,3.400) (-4.750,3.359) (-4.702,3.325) (-4.650,3.288) (-4.600,3.253) (-4.550,3.217) (-4.500,3.182) (-4.455,3.150) (-4.400,3.111) (-4.350,3.076) (-4.300,3.041) (-4.250,3.005) (-4.207,2.975) (-4.150,2.935) (-4.101,2.900) (-4.050,2.864) (-4.000,2.828) (-3.950,2.793) (-3.900,2.758) (-3.854,2.725) (-3.800,2.687) (-3.750,2.652) (-3.700,2.616) (-3.650,2.581) (-3.606,2.550) (-3.550,2.510) (-3.500,2.475) (-3.450,2.440) (-3.400,2.404) (-3.359,2.375) (-3.300,2.334) (-3.253,2.300) (-3.200,2.263) (-3.150,2.227) (-3.100,2.192) (-3.050,2.157) (-3.005,2.125) (-2.950,2.086) (-2.900,2.051) (-2.850,2.015) (-2.800,1.980) (-2.758,1.950) (-2.700,1.909) (-2.652,1.875) (-2.600,1.839) (-2.550,1.803) (-2.500,1.768) (-2.450,1.732) (-2.404,1.700) (-2.350,1.662) (-2.300,1.626) (-2.250,1.591) (-2.200,1.556) (-2.157,1.525) (-2.100,1.485) (-2.051,1.450) (-2.000,1.414) (-1.950,1.379) (-1.900,1.344) (-1.850,1.308) (-1.803,1.275) (-1.750,1.238) (-1.700,1.202) (-1.650,1.167) (-1.600,1.131) (-1.556,1.100) (-1.500,1.061) (-1.450,1.025) (-1.400,0.990) (-1.350,0.955) (-1.308,0.925) (-1.250,0.884) (-1.202,0.850) (-1.150,0.813) (-1.100,0.778) (-1.050,0.743) (-1.000,0.707) (-0.955,0.675) (-0.900,0.637) (-0.850,0.601) (-0.800,0.566) (-0.750,0.531) (-0.708,0.500) (-0.650,0.460) (-0.601,0.425) (-0.550,0.389) (-0.500,0.354) (-0.450,0.319) (-0.400,0.283) (-0.354,0.250) (-0.300,0.213) (-0.250,0.177) (-0.220,0.155)}
\EQ{(4.306,-10.000) (4.350,-9.950) (4.392,-9.900) (4.425,-9.860) (4.473,-9.800) (4.513,-9.750) (4.550,-9.701) (4.588,-9.650) (4.624,-9.600) (4.660,-9.550) (4.695,-9.500) (4.725,-9.455) (4.761,-9.400) (4.793,-9.350) (4.824,-9.300) (4.855,-9.250) (4.884,-9.200) (4.913,-9.150) (4.941,-9.100) (4.969,-9.050) (4.995,-9.000) (5.021,-8.950) (5.046,-8.900) (5.071,-8.850) (5.095,-8.800) (5.118,-8.750) (5.140,-8.700) (5.162,-8.650) (5.183,-8.600) (5.203,-8.550) (5.223,-8.500) (5.242,-8.450) (5.261,-8.400) (5.279,-8.350) (5.296,-8.300) (5.313,-8.250) (5.329,-8.200) (5.344,-8.150) (5.359,-8.100) (5.373,-8.050) (5.387,-8.000) (5.400,-7.951) (5.413,-7.900) (5.425,-7.850) (5.436,-7.800) (5.447,-7.750) (5.457,-7.700) (5.467,-7.650) (5.476,-7.600) (5.484,-7.550) (5.492,-7.500) (5.500,-7.450) (5.507,-7.400) (5.513,-7.350) (5.519,-7.300) (5.524,-7.250) (5.529,-7.200) (5.533,-7.150) (5.537,-7.100) (5.540,-7.050) (5.542,-7.000) (5.544,-6.950) (5.545,-6.900) (5.546,-6.850) (5.547,-6.800) (5.546,-6.750) (5.546,-6.700) (5.544,-6.650) (5.543,-6.600) (5.540,-6.550) (5.537,-6.500) (5.534,-6.450) (5.530,-6.400) (5.525,-6.350) (5.520,-6.300) (5.515,-6.250) (5.509,-6.200) (5.502,-6.150) (5.495,-6.100) (5.487,-6.050) (5.478,-6.000) (5.469,-5.950) (5.460,-5.900) (5.450,-5.850) (5.439,-5.800) (5.428,-5.750) (5.416,-5.700) (5.404,-5.650) (5.391,-5.600) (5.378,-5.550) (5.363,-5.500) (5.349,-5.450) (5.334,-5.400) (5.318,-5.350) (5.301,-5.300) (5.284,-5.250) (5.267,-5.200) (5.249,-5.150) (5.230,-5.100) (5.210,-5.050) (5.190,-5.000) (5.170,-4.950) (5.150,-4.903) (5.127,-4.850) (5.104,-4.800) (5.081,-4.750) (5.057,-4.700) (5.033,-4.650) (5.008,-4.600) (4.982,-4.550) (4.956,-4.500) (4.929,-4.450) (4.901,-4.400) (4.875,-4.354) (4.844,-4.300) (4.814,-4.250) (4.784,-4.200) (4.753,-4.150) (4.725,-4.107) (4.688,-4.050) (4.655,-4.000) (4.625,-3.955) (4.587,-3.900) (4.551,-3.850) (4.516,-3.800) (4.479,-3.750) (4.441,-3.700) (4.403,-3.650) (4.364,-3.600) (4.325,-3.551) (4.284,-3.500) (4.250,-3.458) (4.201,-3.400) (4.159,-3.350) (4.125,-3.311) (4.075,-3.254) (4.026,-3.200) (3.981,-3.150) (3.935,-3.100) (3.900,-3.063) (3.864,-3.025) (3.825,-2.985) (3.775,-2.934) (3.725,-2.883) (3.675,-2.834) (3.625,-2.785) (3.575,-2.736) (3.537,-2.700) (3.500,-2.665) (3.457,-2.625) (3.403,-2.575) (3.350,-2.526) (3.300,-2.481) (3.250,-2.437) (3.209,-2.400) (3.151,-2.350) (3.100,-2.306) (3.050,-2.263) (3.006,-2.225) (2.950,-2.178) (2.900,-2.137) (2.856,-2.100) (2.800,-2.054) (2.750,-2.014) (2.702,-1.975) (2.650,-1.933) (2.609,-1.900) (2.550,-1.854) (2.500,-1.814) (2.450,-1.775) (2.400,-1.737) (2.353,-1.700) (2.300,-1.660) (2.255,-1.625) (2.200,-1.583) (2.156,-1.550) (2.100,-1.508) (2.056,-1.475) (2.000,-1.433) (1.956,-1.400) (1.900,-1.359) (1.854,-1.325) (1.800,-1.285) (1.752,-1.250) (1.700,-1.212) (1.650,-1.175) (1.600,-1.139) (1.550,-1.103) (1.500,-1.067) (1.450,-1.031) (1.408,-1.000) (1.350,-0.959) (1.304,-0.925) (1.250,-0.887) (1.200,-0.851) (1.150,-0.815) (1.100,-0.780) (1.059,-0.750) (1.000,-0.708) (0.953,-0.675) (0.900,-0.637) (0.850,-0.602) (0.800,-0.566) (0.750,-0.531) (0.707,-0.500) (0.650,-0.460) (0.601,-0.425) (0.550,-0.389) (0.500,-0.354) (0.450,-0.319) (0.400,-0.283) (0.354,-0.250) (0.300,-0.213) (0.250,-0.177) (0.206,-0.144)}
\EQ{(-0.220,-0.155) (-0.275,-0.195) (-0.319,-0.225) (-0.375,-0.266) (-0.424,-0.300) (-0.475,-0.336) (-0.525,-0.372) (-0.566,-0.400) (-0.625,-0.442) (-0.672,-0.475) (-0.725,-0.513) (-0.775,-0.549) (-0.825,-0.584) (-0.875,-0.620) (-0.918,-0.650) (-0.975,-0.691) (-1.024,-0.725) (-1.075,-0.762) (-1.125,-0.797) (-1.175,-0.833) (-1.225,-0.869) (-1.269,-0.900) (-1.325,-0.941) (-1.373,-0.975) (-1.425,-1.013) (-1.475,-1.049) (-1.525,-1.085) (-1.575,-1.121) (-1.625,-1.157) (-1.675,-1.194) (-1.718,-1.225) (-1.775,-1.267) (-1.820,-1.300) (-1.875,-1.340) (-1.922,-1.375) (-1.975,-1.415) (-2.023,-1.450) (-2.075,-1.489) (-2.123,-1.525) (-2.175,-1.564) (-2.222,-1.600) (-2.275,-1.640) (-2.320,-1.675) (-2.375,-1.717) (-2.418,-1.750) (-2.475,-1.795) (-2.525,-1.834) (-2.575,-1.873) (-2.625,-1.913) (-2.671,-1.950) (-2.725,-1.993) (-2.764,-2.025) (-2.825,-2.075) (-2.875,-2.116) (-2.916,-2.150) (-2.975,-2.199) (-3.025,-2.241) (-3.065,-2.275) (-3.123,-2.325) (-3.175,-2.371) (-3.225,-2.415) (-3.265,-2.450) (-3.321,-2.500) (-3.375,-2.549) (-3.425,-2.595) (-3.475,-2.642) (-3.525,-2.689) (-3.563,-2.725) (-3.600,-2.760) (-3.641,-2.800) (-3.692,-2.850) (-3.742,-2.900) (-3.791,-2.950) (-3.840,-3.000) (-3.875,-3.037) (-3.911,-3.075) (-3.950,-3.117) (-4.000,-3.171) (-4.049,-3.225) (-4.093,-3.275) (-4.137,-3.325) (-4.175,-3.369) (-4.222,-3.425) (-4.264,-3.475) (-4.300,-3.519) (-4.345,-3.575) (-4.384,-3.625) (-4.422,-3.675) (-4.460,-3.725) (-4.497,-3.775) (-4.534,-3.825) (-4.569,-3.875) (-4.600,-3.919) (-4.638,-3.975) (-4.672,-4.025) (-4.700,-4.068) (-4.737,-4.125) (-4.768,-4.175) (-4.799,-4.225) (-4.825,-4.269) (-4.858,-4.325) (-4.887,-4.375) (-4.915,-4.425) (-4.942,-4.475) (-4.969,-4.525) (-4.995,-4.575) (-5.020,-4.625) (-5.045,-4.675) (-5.069,-4.725) (-5.093,-4.775) (-5.116,-4.825) (-5.138,-4.875) (-5.159,-4.925) (-5.180,-4.975) (-5.200,-5.024) (-5.220,-5.075) (-5.239,-5.125) (-5.258,-5.175) (-5.275,-5.223) (-5.293,-5.275) (-5.310,-5.325) (-5.325,-5.373) (-5.341,-5.425) (-5.356,-5.475) (-5.371,-5.525) (-5.384,-5.575) (-5.398,-5.625) (-5.410,-5.675) (-5.422,-5.725) (-5.434,-5.775) (-5.445,-5.825) (-5.455,-5.875) (-5.465,-5.925) (-5.474,-5.975) (-5.483,-6.025) (-5.491,-6.075) (-5.498,-6.125) (-5.505,-6.175) (-5.512,-6.225) (-5.518,-6.275) (-5.523,-6.325) (-5.528,-6.375) (-5.532,-6.425) (-5.536,-6.475) (-5.539,-6.525) (-5.542,-6.575) (-5.544,-6.625) (-5.545,-6.675) (-5.546,-6.725) (-5.547,-6.775) (-5.547,-6.825) (-5.546,-6.875) (-5.545,-6.925) (-5.543,-6.975) (-5.541,-7.025) (-5.538,-7.075) (-5.535,-7.125) (-5.531,-7.175) (-5.526,-7.225) (-5.522,-7.275) (-5.516,-7.325) (-5.510,-7.375) (-5.503,-7.425) (-5.496,-7.475) (-5.489,-7.525) (-5.480,-7.575) (-5.471,-7.625) (-5.462,-7.675) (-5.452,-7.725) (-5.442,-7.775) (-5.430,-7.825) (-5.419,-7.875) (-5.407,-7.925) (-5.394,-7.975) (-5.380,-8.025) (-5.366,-8.075) (-5.352,-8.125) (-5.337,-8.175) (-5.321,-8.225) (-5.305,-8.275) (-5.288,-8.325) (-5.270,-8.375) (-5.252,-8.425) (-5.233,-8.475) (-5.213,-8.525) (-5.193,-8.575) (-5.173,-8.625) (-5.151,-8.675) (-5.129,-8.725) (-5.106,-8.775) (-5.083,-8.825) (-5.059,-8.875) (-5.034,-8.925) (-5.008,-8.975) (-4.982,-9.025) (-4.955,-9.075) (-4.927,-9.125) (-4.900,-9.173) (-4.870,-9.225) (-4.840,-9.275) (-4.809,-9.325) (-4.777,-9.375) (-4.750,-9.417) (-4.711,-9.475) (-4.677,-9.525) (-4.642,-9.575) (-4.606,-9.625) (-4.575,-9.668) (-4.532,-9.725) (-4.500,-9.766) (-4.453,-9.825) (-4.413,-9.875) (-4.375,-9.920) (-4.328,-9.975) (-4.306,-10.000)}
\EQ{(0.206,0.144) (0.248,0.175) (0.300,0.213) (0.350,0.248) (0.400,0.283) (0.450,0.319) (0.495,0.350) (0.550,0.389) (0.600,0.424) (0.650,0.460) (0.700,0.495) (0.742,0.525) (0.800,0.566) (0.848,0.600) (0.900,0.637) (0.950,0.673) (1.000,0.708) (1.050,0.744) (1.094,0.775) (1.150,0.815) (1.199,0.850) (1.250,0.887) (1.300,0.923) (1.350,0.959) (1.400,0.995) (1.443,1.025) (1.500,1.067) (1.546,1.100) (1.600,1.139) (1.650,1.175) (1.700,1.212) (1.750,1.248) (1.800,1.285) (1.850,1.322) (1.900,1.359) (1.950,1.396) (2.000,1.433) (2.050,1.470) (2.100,1.508) (2.150,1.546) (2.200,1.583) (2.250,1.621) (2.300,1.660) (2.350,1.698) (2.400,1.737) (2.450,1.775) (2.500,1.814) (2.546,1.850) (2.600,1.893) (2.640,1.925) (2.700,1.973) (2.750,2.014) (2.795,2.050) (2.850,2.095) (2.900,2.137) (2.946,2.175) (3.000,2.220) (3.050,2.263) (3.094,2.300) (3.150,2.349) (3.200,2.393) (3.250,2.437) (3.293,2.475) (3.349,2.525) (3.400,2.572) (3.450,2.618) (3.500,2.665) (3.537,2.700) (3.575,2.736) (3.615,2.775) (3.666,2.825) (3.717,2.875) (3.767,2.925) (3.815,2.975) (3.864,3.025) (3.900,3.063) (3.935,3.100) (3.975,3.144) (4.025,3.198) (4.071,3.250) (4.115,3.300) (4.150,3.340) (4.200,3.399) (4.243,3.450) (4.275,3.489) (4.325,3.550) (4.364,3.600) (4.400,3.646) (4.441,3.700) (4.475,3.745) (4.516,3.800) (4.550,3.848) (4.587,3.900) (4.621,3.950) (4.650,3.992) (4.688,4.050) (4.721,4.100) (4.750,4.146) (4.784,4.200) (4.814,4.250) (4.844,4.300) (4.873,4.350) (4.900,4.398) (4.925,4.443) (4.950,4.489) (4.975,4.536) (5.000,4.585) (5.025,4.634) (5.050,4.685) (5.075,4.737) (5.100,4.791) (5.125,4.846) (5.149,4.900) (5.170,4.950) (5.190,5.000) (5.210,5.050) (5.230,5.100) (5.249,5.150) (5.267,5.200) (5.284,5.250) (5.301,5.300) (5.318,5.350) (5.334,5.400) (5.349,5.450) (5.363,5.500) (5.378,5.550) (5.391,5.600) (5.404,5.650) (5.416,5.700) (5.428,5.750) (5.439,5.800) (5.450,5.850) (5.460,5.900) (5.469,5.950) (5.478,6.000) (5.487,6.050) (5.495,6.100) (5.502,6.150) (5.509,6.200) (5.515,6.250) (5.520,6.300) (5.525,6.350) (5.530,6.400) (5.534,6.450) (5.537,6.500) (5.540,6.550) (5.543,6.600) (5.544,6.650) (5.546,6.700) (5.546,6.750) (5.547,6.800) (5.546,6.850) (5.545,6.900) (5.544,6.950) (5.542,7.000) (5.540,7.050) (5.537,7.100) (5.533,7.150) (5.529,7.200) (5.524,7.250) (5.519,7.300) (5.513,7.350) (5.507,7.400) (5.500,7.450) (5.492,7.500) (5.484,7.550) (5.476,7.600) (5.467,7.650) (5.457,7.700) (5.447,7.750) (5.436,7.800) (5.425,7.850) (5.413,7.900) (5.400,7.950) (5.387,8.000) (5.373,8.050) (5.359,8.100) (5.344,8.150) (5.329,8.200) (5.313,8.250) (5.296,8.300) (5.279,8.350) (5.261,8.400) (5.242,8.450) (5.223,8.500) (5.203,8.550) (5.183,8.600) (5.162,8.650) (5.140,8.700) (5.118,8.750) (5.095,8.800) (5.071,8.850) (5.046,8.900) (5.021,8.950) (4.995,9.000) (4.969,9.050) (4.941,9.100) (4.913,9.150) (4.884,9.200) (4.855,9.250) (4.825,9.299) (4.793,9.350) (4.761,9.400) (4.728,9.450) (4.700,9.492) (4.660,9.550) (4.625,9.599) (4.588,9.650) (4.551,9.700) (4.513,9.750) (4.475,9.798) (4.433,9.850) (4.400,9.890) (4.350,9.949) (4.306,10.000)}
\EQ{(-4.306,10.000) (-4.350,9.950) (-4.392,9.900) (-4.425,9.860) (-4.473,9.800) (-4.513,9.750) (-4.550,9.701) (-4.588,9.650) (-4.624,9.600) (-4.660,9.550) (-4.695,9.500) (-4.725,9.455) (-4.761,9.400) (-4.793,9.350) (-4.824,9.300) (-4.855,9.250) (-4.884,9.200) (-4.913,9.150) (-4.941,9.100) (-4.969,9.050) (-4.995,9.000) (-5.021,8.950) (-5.046,8.900) (-5.071,8.850) (-5.095,8.800) (-5.118,8.750) (-5.140,8.700) (-5.162,8.650) (-5.183,8.600) (-5.203,8.550) (-5.223,8.500) (-5.242,8.450) (-5.261,8.400) (-5.279,8.350) (-5.296,8.300) (-5.313,8.250) (-5.329,8.200) (-5.344,8.150) (-5.359,8.100) (-5.373,8.050) (-5.387,8.000) (-5.400,7.951) (-5.413,7.900) (-5.425,7.850) (-5.436,7.800) (-5.447,7.750) (-5.457,7.700) (-5.467,7.650) (-5.476,7.600) (-5.484,7.550) (-5.492,7.500) (-5.500,7.450) (-5.507,7.400) (-5.513,7.350) (-5.519,7.300) (-5.524,7.250) (-5.529,7.200) (-5.533,7.150) (-5.537,7.100) (-5.540,7.050) (-5.542,7.000) (-5.544,6.950) (-5.545,6.900) (-5.546,6.850) (-5.547,6.800) (-5.546,6.750) (-5.546,6.700) (-5.544,6.650) (-5.543,6.600) (-5.540,6.550) (-5.537,6.500) (-5.534,6.450) (-5.530,6.400) (-5.525,6.350) (-5.520,6.300) (-5.515,6.250) (-5.509,6.200) (-5.502,6.150) (-5.495,6.100) (-5.487,6.050) (-5.478,6.000) (-5.469,5.950) (-5.460,5.900) (-5.450,5.850) (-5.439,5.800) (-5.428,5.750) (-5.416,5.700) (-5.404,5.650) (-5.391,5.600) (-5.378,5.550) (-5.363,5.500) (-5.349,5.450) (-5.334,5.400) (-5.318,5.350) (-5.301,5.300) (-5.284,5.250) (-5.267,5.200) (-5.249,5.150) (-5.230,5.100) (-5.210,5.050) (-5.190,5.000) (-5.170,4.950) (-5.150,4.903) (-5.127,4.850) (-5.104,4.800) (-5.081,4.750) (-5.057,4.700) (-5.033,4.650) (-5.008,4.600) (-4.982,4.550) (-4.956,4.500) (-4.929,4.450) (-4.901,4.400) (-4.875,4.354) (-4.844,4.300) (-4.814,4.250) (-4.784,4.200) (-4.753,4.150) (-4.725,4.107) (-4.688,4.050) (-4.655,4.000) (-4.625,3.955) (-4.587,3.900) (-4.551,3.850) (-4.516,3.800) (-4.479,3.750) (-4.441,3.700) (-4.403,3.650) (-4.364,3.600) (-4.325,3.551) (-4.284,3.500) (-4.250,3.458) (-4.201,3.400) (-4.159,3.350) (-4.125,3.311) (-4.075,3.254) (-4.026,3.200) (-3.981,3.150) (-3.935,3.100) (-3.900,3.063) (-3.864,3.025) (-3.825,2.985) (-3.775,2.934) (-3.725,2.883) (-3.675,2.834) (-3.625,2.785) (-3.575,2.736) (-3.537,2.700) (-3.500,2.665) (-3.457,2.625) (-3.403,2.575) (-3.350,2.526) (-3.300,2.481) (-3.250,2.437) (-3.209,2.400) (-3.151,2.350) (-3.100,2.306) (-3.050,2.263) (-3.006,2.225) (-2.950,2.178) (-2.900,2.137) (-2.856,2.100) (-2.800,2.054) (-2.750,2.014) (-2.702,1.975) (-2.650,1.933) (-2.609,1.900) (-2.550,1.854) (-2.500,1.814) (-2.450,1.775) (-2.400,1.737) (-2.353,1.700) (-2.300,1.660) (-2.255,1.625) (-2.200,1.583) (-2.156,1.550) (-2.100,1.508) (-2.056,1.475) (-2.000,1.433) (-1.956,1.400) (-1.900,1.359) (-1.854,1.325) (-1.800,1.285) (-1.752,1.250) (-1.700,1.212) (-1.650,1.175) (-1.600,1.139) (-1.550,1.103) (-1.500,1.067) (-1.450,1.031) (-1.408,1.000) (-1.350,0.959) (-1.304,0.925) (-1.250,0.887) (-1.200,0.851) (-1.150,0.815) (-1.100,0.780) (-1.059,0.750) (-1.000,0.708) (-0.953,0.675) (-0.900,0.637) (-0.850,0.602) (-0.800,0.566) (-0.750,0.531) (-0.707,0.500) (-0.650,0.460) (-0.601,0.425) (-0.550,0.389) (-0.500,0.354) (-0.450,0.319) (-0.400,0.283) (-0.354,0.250) (-0.300,0.213) (-0.250,0.177) (-0.220,0.155)}
\EQ{(-0.220,-0.155) (-0.275,-0.195) (-0.319,-0.225) (-0.375,-0.266) (-0.424,-0.300) (-0.475,-0.337) (-0.525,-0.372) (-0.575,-0.407) (-0.625,-0.443) (-0.671,-0.475) (-0.725,-0.514) (-0.775,-0.549) (-0.825,-0.585) (-0.875,-0.620) (-0.917,-0.650) (-0.975,-0.692) (-1.022,-0.725) (-1.075,-0.763) (-1.125,-0.799) (-1.175,-0.835) (-1.225,-0.872) (-1.275,-0.908) (-1.325,-0.944) (-1.368,-0.975) (-1.425,-1.017) (-1.470,-1.050) (-1.525,-1.091) (-1.571,-1.125) (-1.625,-1.166) (-1.671,-1.200) (-1.725,-1.241) (-1.770,-1.275) (-1.825,-1.317) (-1.869,-1.350) (-1.925,-1.394) (-1.965,-1.425) (-2.025,-1.472) (-2.075,-1.512) (-2.123,-1.550) (-2.175,-1.592) (-2.216,-1.625) (-2.275,-1.674) (-2.325,-1.715) (-2.367,-1.750) (-2.425,-1.799) (-2.475,-1.843) (-2.525,-1.886) (-2.570,-1.925) (-2.625,-1.975) (-2.675,-2.020) (-2.725,-2.066) (-2.762,-2.100) (-2.800,-2.136) (-2.841,-2.175) (-2.893,-2.225) (-2.943,-2.275) (-2.993,-2.325) (-3.042,-2.375) (-3.090,-2.425) (-3.125,-2.463) (-3.160,-2.500) (-3.200,-2.545) (-3.249,-2.600) (-3.293,-2.650) (-3.336,-2.700) (-3.375,-2.747) (-3.418,-2.800) (-3.450,-2.840) (-3.497,-2.900) (-3.535,-2.950) (-3.572,-3.000) (-3.608,-3.050) (-3.644,-3.100) (-3.675,-3.146) (-3.712,-3.200) (-3.744,-3.250) (-3.775,-3.299) (-3.807,-3.350) (-3.837,-3.400) (-3.866,-3.450) (-3.894,-3.500) (-3.922,-3.550) (-3.948,-3.600) (-3.974,-3.650) (-3.999,-3.700) (-4.023,-3.750) (-4.047,-3.800) (-4.069,-3.850) (-4.091,-3.900) (-4.112,-3.950) (-4.132,-4.000) (-4.152,-4.050) (-4.170,-4.100) (-4.188,-4.150) (-4.206,-4.200) (-4.222,-4.250) (-4.238,-4.300) (-4.253,-4.350) (-4.267,-4.400) (-4.280,-4.450) (-4.293,-4.500) (-4.305,-4.550) (-4.316,-4.600) (-4.327,-4.650) (-4.337,-4.700) (-4.346,-4.750) (-4.354,-4.800) (-4.362,-4.850) (-4.369,-4.900) (-4.376,-4.950) (-4.381,-5.000) (-4.386,-5.050) (-4.391,-5.100) (-4.394,-5.150) (-4.397,-5.200) (-4.400,-5.250) (-4.401,-5.300) (-4.402,-5.350) (-4.402,-5.400) (-4.402,-5.450) (-4.401,-5.500) (-4.399,-5.550) (-4.396,-5.600) (-4.393,-5.650) (-4.389,-5.700) (-4.385,-5.750) (-4.380,-5.800) (-4.374,-5.850) (-4.367,-5.900) (-4.360,-5.950) (-4.352,-6.000) (-4.343,-6.050) (-4.334,-6.100) (-4.323,-6.150) (-4.313,-6.200) (-4.301,-6.250) (-4.289,-6.300) (-4.276,-6.350) (-4.262,-6.400) (-4.247,-6.450) (-4.232,-6.500) (-4.216,-6.550) (-4.200,-6.598) (-4.182,-6.650) (-4.163,-6.700) (-4.144,-6.750) (-4.125,-6.798) (-4.104,-6.850) (-4.082,-6.900) (-4.060,-6.950) (-4.036,-7.000) (-4.012,-7.050) (-3.987,-7.100) (-3.961,-7.150) (-3.934,-7.200) (-3.907,-7.250) (-3.878,-7.300) (-3.850,-7.347) (-3.818,-7.400) (-3.786,-7.450) (-3.753,-7.500) (-3.725,-7.542) (-3.684,-7.600) (-3.650,-7.648) (-3.611,-7.700) (-3.575,-7.747) (-3.533,-7.800) (-3.500,-7.841) (-3.450,-7.900) (-3.406,-7.950) (-3.361,-8.000) (-3.325,-8.039) (-3.275,-8.092) (-3.225,-8.142) (-3.175,-8.191) (-3.125,-8.238) (-3.085,-8.275) (-3.029,-8.325) (-2.975,-8.372) (-2.925,-8.414) (-2.880,-8.450) (-2.825,-8.494) (-2.784,-8.525) (-2.725,-8.569) (-2.682,-8.600) (-2.625,-8.640) (-2.575,-8.674) (-2.525,-8.707) (-2.475,-8.739) (-2.425,-8.770) (-2.375,-8.800) (-2.325,-8.829) (-2.275,-8.858) (-2.225,-8.886) (-2.175,-8.912) (-2.125,-8.938) (-2.075,-8.964) (-2.025,-8.988) (-1.975,-9.012) (-1.925,-9.035) (-1.875,-9.057) (-1.825,-9.078) (-1.775,-9.099) (-1.725,-9.119) (-1.675,-9.139) (-1.625,-9.158) (-1.577,-9.175) (-1.525,-9.193) (-1.475,-9.210) (-1.425,-9.226) (-1.375,-9.242) (-1.325,-9.257) (-1.275,-9.272) (-1.225,-9.285) (-1.175,-9.298) (-1.125,-9.311) (-1.075,-9.323) (-1.025,-9.335) (-0.975,-9.345) (-0.925,-9.356) (-0.875,-9.365) (-0.825,-9.375) (-0.775,-9.383) (-0.725,-9.391) (-0.675,-9.399) (-0.625,-9.406) (-0.575,-9.412) (-0.525,-9.418) (-0.475,-9.423) (-0.425,-9.428) (-0.375,-9.432) (-0.325,-9.436) (-0.275,-9.439) (-0.225,-9.442) (-0.175,-9.444) (-0.125,-9.446) (-0.075,-9.447) (-0.025,-9.447) (0.050,-9.447) (0.100,-9.446) (0.150,-9.445) (0.200,-9.443) (0.250,-9.441) (0.300,-9.438) (0.350,-9.434) (0.400,-9.430) (0.450,-9.426) (0.500,-9.421) (0.550,-9.415) (0.600,-9.409) (0.650,-9.402) (0.700,-9.395) (0.750,-9.387) (0.800,-9.379) (0.850,-9.370) (0.900,-9.361) (0.950,-9.351) (1.000,-9.340) (1.050,-9.329) (1.100,-9.317) (1.150,-9.305) (1.200,-9.292) (1.250,-9.278) (1.300,-9.264) (1.349,-9.250) (1.400,-9.234) (1.450,-9.218) (1.500,-9.202) (1.550,-9.185) (1.600,-9.167) (1.650,-9.148) (1.700,-9.129) (1.750,-9.109) (1.800,-9.089) (1.850,-9.068) (1.900,-9.046) (1.946,-9.025) (2.000,-9.000) (2.050,-8.976) (2.100,-8.951) (2.150,-8.925) (2.198,-8.900) (2.244,-8.875) (2.289,-8.850) (2.333,-8.825) (2.400,-8.785) (2.450,-8.755) (2.497,-8.725) (2.550,-8.691) (2.600,-8.657) (2.647,-8.625) (2.700,-8.587) (2.750,-8.551) (2.800,-8.513) (2.849,-8.475) (2.900,-8.434) (2.942,-8.400) (3.000,-8.351) (3.050,-8.307) (3.100,-8.262) (3.139,-8.225) (3.192,-8.175) (3.242,-8.125) (3.291,-8.075) (3.338,-8.025) (3.375,-7.985) (3.425,-7.929) (3.471,-7.875) (3.513,-7.825) (3.550,-7.779) (3.592,-7.725) (3.625,-7.682) (3.667,-7.625) (3.700,-7.578) (3.736,-7.525) (3.770,-7.475) (3.800,-7.428) (3.833,-7.375) (3.863,-7.325) (3.892,-7.275) (3.921,-7.225) (3.948,-7.175) (3.974,-7.125) (4.000,-7.075) (4.024,-7.025) (4.048,-6.975) (4.071,-6.925) (4.093,-6.875) (4.114,-6.825) (4.134,-6.775) (4.154,-6.725) (4.173,-6.675) (4.191,-6.625) (4.208,-6.575) (4.224,-6.525) (4.240,-6.475) (4.255,-6.425) (4.269,-6.375) (4.282,-6.325) (4.295,-6.275) (4.307,-6.225) (4.318,-6.175) (4.329,-6.125) (4.338,-6.075) (4.347,-6.025) (4.356,-5.975) (4.364,-5.925) (4.370,-5.875) (4.377,-5.825) (4.382,-5.775) (4.387,-5.725) (4.391,-5.675) (4.395,-5.625) (4.398,-5.575) (4.400,-5.525) (4.401,-5.475) (4.402,-5.425) (4.402,-5.375) (4.402,-5.325) (4.401,-5.275) (4.399,-5.225) (4.396,-5.175) (4.393,-5.125) (4.389,-5.075) (4.384,-5.025) (4.379,-4.975) (4.373,-4.925) (4.366,-4.875) (4.358,-4.825) (4.350,-4.775) (4.342,-4.725) (4.332,-4.675) (4.322,-4.625) (4.311,-4.575) (4.299,-4.525) (4.287,-4.475) (4.274,-4.425) (4.260,-4.375) (4.245,-4.325) (4.230,-4.275) (4.214,-4.225) (4.197,-4.175) (4.180,-4.125) (4.161,-4.075) (4.142,-4.025) (4.122,-3.975) (4.102,-3.925) (4.080,-3.875) (4.058,-3.825) (4.035,-3.775) (4.011,-3.725) (3.987,-3.675) (3.961,-3.625) (3.935,-3.575) (3.908,-3.525) (3.880,-3.475) (3.852,-3.425) (3.825,-3.380) (3.792,-3.325) (3.760,-3.275) (3.728,-3.225) (3.700,-3.183) (3.661,-3.125) (3.626,-3.075) (3.590,-3.025) (3.554,-2.975) (3.516,-2.925) (3.478,-2.875) (3.438,-2.825) (3.400,-2.778) (3.357,-2.725) (3.315,-2.675) (3.275,-2.629) (3.227,-2.575) (3.182,-2.525) (3.136,-2.475) (3.100,-2.436) (3.050,-2.384) (3.000,-2.332) (2.950,-2.282) (2.900,-2.232) (2.850,-2.184) (2.800,-2.136) (2.762,-2.100) (2.725,-2.066) (2.681,-2.025) (2.626,-1.975) (2.575,-1.930) (2.525,-1.886) (2.484,-1.850) (2.426,-1.800) (2.375,-1.757) (2.325,-1.715) (2.277,-1.675) (2.225,-1.633) (2.185,-1.600) (2.125,-1.552) (2.075,-1.512) (2.029,-1.475) (1.975,-1.433) (1.933,-1.400) (1.875,-1.355) (1.825,-1.317) (1.775,-1.279) (1.725,-1.241) (1.675,-1.203) (1.625,-1.166) (1.575,-1.128) (1.525,-1.091) (1.475,-1.054) (1.425,-1.017) (1.375,-0.981) (1.333,-0.950) (1.275,-0.908) (1.230,-0.875) (1.175,-0.835) (1.126,-0.800) (1.075,-0.763) (1.025,-0.727) (0.975,-0.692) (0.925,-0.656) (0.882,-0.625) (0.825,-0.585) (0.777,-0.550) (0.725,-0.514) (0.675,-0.478) (0.625,-0.443) (0.575,-0.407) (0.531,-0.375) (0.475,-0.337) (0.425,-0.301) (0.375,-0.266) (0.325,-0.230) (0.284,-0.200) (0.225,-0.160) (0.206,-0.144)}
\EQ{(0.206,0.144) (0.248,0.175) (0.300,0.213) (0.350,0.248) (0.400,0.283) (0.450,0.319) (0.495,0.350) (0.550,0.390) (0.600,0.425) (0.650,0.460) (0.700,0.496) (0.742,0.525) (0.800,0.567) (0.847,0.600) (0.900,0.638) (0.950,0.674) (1.000,0.710) (1.050,0.745) (1.092,0.775) (1.150,0.817) (1.196,0.850) (1.250,0.890) (1.299,0.925) (1.350,0.963) (1.400,0.999) (1.450,1.036) (1.500,1.073) (1.550,1.110) (1.600,1.147) (1.650,1.184) (1.700,1.222) (1.750,1.260) (1.800,1.298) (1.850,1.336) (1.900,1.374) (1.950,1.413) (1.997,1.450) (2.050,1.492) (2.092,1.525) (2.150,1.572) (2.200,1.612) (2.246,1.650) (2.300,1.694) (2.350,1.736) (2.396,1.775) (2.450,1.821) (2.500,1.864) (2.541,1.900) (2.598,1.950) (2.650,1.997) (2.700,2.043) (2.750,2.089) (2.789,2.125) (2.825,2.160) (2.867,2.200) (2.918,2.250) (2.968,2.300) (3.017,2.350) (3.066,2.400) (3.113,2.450) (3.150,2.490) (3.200,2.545) (3.249,2.600) (3.293,2.650) (3.336,2.700) (3.375,2.747) (3.418,2.800) (3.450,2.840) (3.497,2.900) (3.535,2.950) (3.572,3.000) (3.608,3.050) (3.644,3.100) (3.675,3.146) (3.712,3.200) (3.744,3.250) (3.775,3.299) (3.807,3.350) (3.837,3.400) (3.866,3.450) (3.894,3.500) (3.922,3.550) (3.948,3.600) (3.974,3.650) (3.999,3.700) (4.023,3.750) (4.047,3.800) (4.069,3.850) (4.091,3.900) (4.112,3.950) (4.132,4.000) (4.152,4.050) (4.170,4.100) (4.188,4.150) (4.206,4.200) (4.222,4.250) (4.238,4.300) (4.253,4.350) (4.267,4.400) (4.280,4.450) (4.293,4.500) (4.305,4.550) (4.316,4.600) (4.327,4.650) (4.337,4.700) (4.346,4.750) (4.354,4.800) (4.362,4.850) (4.369,4.900) (4.376,4.950) (4.381,5.000) (4.386,5.050) (4.391,5.100) (4.394,5.150) (4.397,5.200) (4.400,5.250) (4.401,5.300) (4.402,5.350) (4.402,5.400) (4.402,5.450) (4.401,5.500) (4.399,5.550) (4.396,5.600) (4.393,5.650) (4.389,5.700) (4.385,5.750) (4.380,5.800) (4.374,5.850) (4.367,5.900) (4.360,5.950) (4.352,6.000) (4.343,6.050) (4.334,6.100) (4.323,6.150) (4.313,6.200) (4.301,6.250) (4.289,6.300) (4.276,6.350) (4.262,6.400) (4.247,6.450) (4.232,6.500) (4.216,6.550) (4.200,6.598) (4.182,6.650) (4.163,6.700) (4.144,6.750) (4.125,6.798) (4.104,6.850) (4.082,6.900) (4.060,6.950) (4.036,7.000) (4.012,7.050) (3.987,7.100) (3.961,7.150) (3.934,7.200) (3.907,7.250) (3.878,7.300) (3.850,7.347) (3.818,7.400) (3.786,7.450) (3.753,7.500) (3.725,7.542) (3.684,7.600) (3.650,7.648) (3.611,7.700) (3.575,7.747) (3.533,7.800) (3.500,7.841) (3.450,7.900) (3.406,7.950) (3.361,8.000) (3.325,8.039) (3.275,8.092) (3.225,8.142) (3.175,8.191) (3.125,8.238) (3.085,8.275) (3.029,8.325) (2.975,8.372) (2.925,8.414) (2.880,8.450) (2.825,8.494) (2.784,8.525) (2.725,8.569) (2.682,8.600) (2.625,8.640) (2.575,8.674) (2.525,8.707) (2.475,8.739) (2.425,8.770) (2.375,8.800) (2.325,8.829) (2.275,8.858) (2.225,8.886) (2.175,8.912) (2.125,8.938) (2.075,8.964) (2.025,8.988) (1.975,9.012) (1.925,9.035) (1.875,9.057) (1.825,9.078) (1.775,9.099) (1.725,9.119) (1.675,9.139) (1.625,9.158) (1.577,9.175) (1.525,9.193) (1.475,9.210) (1.425,9.226) (1.375,9.242) (1.325,9.257) (1.275,9.272) (1.225,9.285) (1.175,9.298) (1.125,9.311) (1.075,9.323) (1.025,9.335) (0.975,9.345) (0.925,9.356) (0.875,9.365) (0.825,9.375) (0.775,9.383) (0.725,9.391) (0.675,9.399) (0.625,9.406) (0.575,9.412) (0.525,9.418) (0.475,9.423) (0.425,9.428) (0.375,9.432) (0.325,9.436) (0.275,9.439) (0.225,9.442) (0.175,9.444) (0.125,9.446) (0.075,9.447) (0.025,9.447) (-0.050,9.447) (-0.100,9.446) (-0.150,9.445) (-0.200,9.443) (-0.250,9.441) (-0.300,9.438) (-0.350,9.434) (-0.400,9.430) (-0.450,9.426) (-0.500,9.421) (-0.550,9.415) (-0.600,9.409) (-0.650,9.402) (-0.700,9.395) (-0.750,9.387) (-0.800,9.379) (-0.850,9.370) (-0.900,9.361) (-0.950,9.351) (-1.000,9.340) (-1.050,9.329) (-1.100,9.317) (-1.150,9.305) (-1.200,9.292) (-1.250,9.278) (-1.300,9.264) (-1.349,9.250) (-1.400,9.234) (-1.450,9.218) (-1.500,9.202) (-1.550,9.185) (-1.600,9.167) (-1.650,9.148) (-1.700,9.129) (-1.750,9.109) (-1.800,9.089) (-1.850,9.068) (-1.900,9.046) (-1.946,9.025) (-2.000,9.000) (-2.050,8.976) (-2.100,8.951) (-2.150,8.925) (-2.198,8.900) (-2.244,8.875) (-2.289,8.850) (-2.333,8.825) (-2.400,8.785) (-2.450,8.755) (-2.497,8.725) (-2.550,8.691) (-2.600,8.657) (-2.647,8.625) (-2.700,8.587) (-2.750,8.551) (-2.800,8.513) (-2.849,8.475) (-2.900,8.434) (-2.942,8.400) (-3.000,8.351) (-3.050,8.307) (-3.100,8.262) (-3.139,8.225) (-3.192,8.175) (-3.242,8.125) (-3.291,8.075) (-3.338,8.025) (-3.375,7.985) (-3.425,7.929) (-3.471,7.875) (-3.513,7.825) (-3.550,7.779) (-3.592,7.725) (-3.625,7.682) (-3.667,7.625) (-3.700,7.578) (-3.736,7.525) (-3.770,7.475) (-3.800,7.428) (-3.833,7.375) (-3.863,7.325) (-3.892,7.275) (-3.921,7.225) (-3.948,7.175) (-3.974,7.125) (-4.000,7.075) (-4.024,7.025) (-4.048,6.975) (-4.071,6.925) (-4.093,6.875) (-4.114,6.825) (-4.134,6.775) (-4.154,6.725) (-4.173,6.675) (-4.191,6.625) (-4.208,6.575) (-4.224,6.525) (-4.240,6.475) (-4.255,6.425) (-4.269,6.375) (-4.282,6.325) (-4.295,6.275) (-4.307,6.225) (-4.318,6.175) (-4.329,6.125) (-4.338,6.075) (-4.347,6.025) (-4.356,5.975) (-4.364,5.925) (-4.370,5.875) (-4.377,5.825) (-4.382,5.775) (-4.387,5.725) (-4.391,5.675) (-4.395,5.625) (-4.398,5.575) (-4.400,5.525) (-4.401,5.475) (-4.402,5.425) (-4.402,5.375) (-4.402,5.325) (-4.401,5.275) (-4.399,5.225) (-4.396,5.175) (-4.393,5.125) (-4.389,5.075) (-4.384,5.025) (-4.379,4.975) (-4.373,4.925) (-4.366,4.875) (-4.358,4.825) (-4.350,4.775) (-4.342,4.725) (-4.332,4.675) (-4.322,4.625) (-4.311,4.575) (-4.299,4.525) (-4.287,4.475) (-4.274,4.425) (-4.260,4.375) (-4.245,4.325) (-4.230,4.275) (-4.214,4.225) (-4.197,4.175) (-4.180,4.125) (-4.161,4.075) (-4.142,4.025) (-4.122,3.975) (-4.102,3.925) (-4.080,3.875) (-4.058,3.825) (-4.035,3.775) (-4.011,3.725) (-3.987,3.675) (-3.961,3.625) (-3.935,3.575) (-3.908,3.525) (-3.880,3.475) (-3.852,3.425) (-3.825,3.380) (-3.792,3.325) (-3.760,3.275) (-3.728,3.225) (-3.700,3.183) (-3.661,3.125) (-3.626,3.075) (-3.590,3.025) (-3.554,2.975) (-3.516,2.925) (-3.478,2.875) (-3.438,2.825) (-3.400,2.778) (-3.357,2.725) (-3.315,2.675) (-3.275,2.629) (-3.227,2.575) (-3.182,2.525) (-3.136,2.475) (-3.100,2.436) (-3.050,2.384) (-3.000,2.332) (-2.950,2.282) (-2.900,2.232) (-2.850,2.184) (-2.800,2.136) (-2.762,2.100) (-2.725,2.066) (-2.681,2.025) (-2.626,1.975) (-2.575,1.930) (-2.525,1.886) (-2.484,1.850) (-2.426,1.800) (-2.375,1.757) (-2.325,1.715) (-2.277,1.675) (-2.225,1.633) (-2.185,1.600) (-2.125,1.552) (-2.075,1.512) (-2.029,1.475) (-1.975,1.433) (-1.933,1.400) (-1.875,1.355) (-1.825,1.317) (-1.775,1.279) (-1.725,1.241) (-1.675,1.203) (-1.625,1.166) (-1.575,1.128) (-1.525,1.091) (-1.475,1.054) (-1.425,1.017) (-1.375,0.981) (-1.333,0.950) (-1.275,0.908) (-1.230,0.875) (-1.175,0.835) (-1.126,0.800) (-1.075,0.763) (-1.025,0.727) (-0.975,0.692) (-0.925,0.656) (-0.882,0.625) (-0.825,0.585) (-0.777,0.550) (-0.725,0.514) (-0.675,0.478) (-0.625,0.443) (-0.575,0.407) (-0.531,0.375) (-0.475,0.337) (-0.425,0.301) (-0.375,0.266) (-0.325,0.230) (-0.284,0.200) (-0.225,0.160) (-0.220,0.155)}
\EQ{(-0.220,-0.155) (-0.275,-0.195) (-0.319,-0.225) (-0.375,-0.266) (-0.424,-0.300) (-0.475,-0.337) (-0.525,-0.372) (-0.575,-0.407) (-0.625,-0.443) (-0.671,-0.475) (-0.725,-0.514) (-0.775,-0.549) (-0.825,-0.585) (-0.875,-0.621) (-0.916,-0.650) (-0.975,-0.693) (-1.020,-0.725) (-1.075,-0.765) (-1.124,-0.800) (-1.175,-0.838) (-1.225,-0.874) (-1.275,-0.911) (-1.325,-0.948) (-1.375,-0.985) (-1.425,-1.022) (-1.475,-1.060) (-1.525,-1.098) (-1.575,-1.136) (-1.625,-1.174) (-1.675,-1.213) (-1.723,-1.250) (-1.775,-1.291) (-1.818,-1.325) (-1.875,-1.370) (-1.925,-1.411) (-1.973,-1.450) (-2.025,-1.493) (-2.075,-1.535) (-2.122,-1.575) (-2.175,-1.621) (-2.225,-1.664) (-2.265,-1.700) (-2.321,-1.750) (-2.375,-1.800) (-2.425,-1.846) (-2.475,-1.894) (-2.525,-1.942) (-2.575,-1.991) (-2.625,-2.042) (-2.675,-2.093) (-2.725,-2.146) (-2.775,-2.200) (-2.820,-2.250) (-2.864,-2.300) (-2.900,-2.342) (-2.949,-2.400) (-2.989,-2.450) (-3.025,-2.495) (-3.068,-2.550) (-3.100,-2.593) (-3.142,-2.650) (-3.175,-2.697) (-3.212,-2.750) (-3.245,-2.800) (-3.275,-2.846) (-3.309,-2.900) (-3.340,-2.950) (-3.369,-3.000) (-3.398,-3.050) (-3.425,-3.099) (-3.450,-3.146) (-3.475,-3.195) (-3.500,-3.245) (-3.525,-3.298) (-3.549,-3.350) (-3.571,-3.400) (-3.592,-3.450) (-3.612,-3.500) (-3.631,-3.550) (-3.650,-3.600) (-3.667,-3.650) (-3.684,-3.700) (-3.700,-3.750) (-3.715,-3.800) (-3.729,-3.850) (-3.742,-3.900) (-3.755,-3.950) (-3.767,-4.000) (-3.777,-4.050) (-3.787,-4.100) (-3.797,-4.150) (-3.805,-4.200) (-3.813,-4.250) (-3.820,-4.300) (-3.826,-4.350) (-3.831,-4.400) (-3.835,-4.450) (-3.839,-4.500) (-3.842,-4.550) (-3.844,-4.600) (-3.845,-4.650) (-3.846,-4.700) (-3.846,-4.750) (-3.845,-4.800) (-3.843,-4.850) (-3.840,-4.900) (-3.837,-4.950) (-3.833,-5.000) (-3.828,-5.050) (-3.822,-5.100) (-3.816,-5.150) (-3.808,-5.200) (-3.800,-5.250) (-3.791,-5.300) (-3.782,-5.350) (-3.771,-5.400) (-3.760,-5.450) (-3.747,-5.500) (-3.734,-5.550) (-3.721,-5.600) (-3.706,-5.650) (-3.690,-5.700) (-3.674,-5.750) (-3.656,-5.800) (-3.638,-5.850) (-3.619,-5.900) (-3.600,-5.948) (-3.578,-6.000) (-3.556,-6.050) (-3.533,-6.100) (-3.509,-6.150) (-3.484,-6.200) (-3.458,-6.250) (-3.431,-6.300) (-3.403,-6.350) (-3.375,-6.399) (-3.344,-6.450) (-3.313,-6.500) (-3.280,-6.550) (-3.250,-6.594) (-3.211,-6.650) (-3.175,-6.699) (-3.137,-6.750) (-3.100,-6.797) (-3.057,-6.850) (-3.015,-6.900) (-2.975,-6.946) (-2.926,-7.000) (-2.879,-7.050) (-2.830,-7.100) (-2.779,-7.150) (-2.727,-7.200) (-2.675,-7.247) (-2.625,-7.291) (-2.585,-7.325) (-2.525,-7.375) (-2.475,-7.414) (-2.429,-7.450) (-2.375,-7.490) (-2.326,-7.525) (-2.275,-7.560) (-2.225,-7.594) (-2.176,-7.625) (-2.125,-7.657) (-2.075,-7.688) (-2.025,-7.717) (-1.975,-7.745) (-1.925,-7.772) (-1.875,-7.799) (-1.825,-7.824) (-1.775,-7.849) (-1.725,-7.873) (-1.675,-7.896) (-1.625,-7.918) (-1.575,-7.939) (-1.525,-7.959) (-1.475,-7.979) (-1.425,-7.998) (-1.375,-8.016) (-1.325,-8.034) (-1.276,-8.050) (-1.225,-8.066) (-1.175,-8.082) (-1.125,-8.096) (-1.075,-8.110) (-1.025,-8.123) (-0.975,-8.136) (-0.925,-8.148) (-0.875,-8.159) (-0.825,-8.170) (-0.775,-8.179) (-0.725,-8.189) (-0.675,-8.197) (-0.625,-8.205) (-0.575,-8.213) (-0.525,-8.219) (-0.475,-8.225) (-0.425,-8.231) (-0.375,-8.236) (-0.325,-8.240) (-0.275,-8.244) (-0.225,-8.247) (-0.175,-8.249) (-0.125,-8.251) (-0.075,-8.252) (-0.025,-8.253) (0.050,-8.253) (0.100,-8.252) (0.150,-8.250) (0.200,-8.248) (0.250,-8.245) (0.300,-8.242) (0.350,-8.238) (0.400,-8.234) (0.450,-8.228) (0.500,-8.223) (0.550,-8.216) (0.600,-8.209) (0.650,-8.201) (0.700,-8.193) (0.750,-8.184) (0.800,-8.175) (0.850,-8.164) (0.900,-8.153) (0.950,-8.142) (1.000,-8.130) (1.050,-8.117) (1.100,-8.103) (1.150,-8.089) (1.200,-8.074) (1.250,-8.058) (1.300,-8.042) (1.350,-8.025) (1.400,-8.007) (1.450,-7.989) (1.500,-7.969) (1.548,-7.950) (1.600,-7.928) (1.650,-7.907) (1.700,-7.884) (1.750,-7.861) (1.800,-7.837) (1.850,-7.812) (1.900,-7.786) (1.950,-7.759) (2.000,-7.731) (2.050,-7.702) (2.096,-7.675) (2.150,-7.642) (2.200,-7.610) (2.250,-7.577) (2.300,-7.543) (2.350,-7.508) (2.395,-7.475) (2.450,-7.434) (2.493,-7.400) (2.550,-7.354) (2.600,-7.313) (2.644,-7.275) (2.700,-7.225) (2.750,-7.178) (2.800,-7.130) (2.850,-7.080) (2.900,-7.028) (2.949,-6.975) (2.993,-6.925) (3.036,-6.875) (3.075,-6.828) (3.117,-6.775) (3.150,-6.733) (3.193,-6.675) (3.225,-6.630) (3.263,-6.575) (3.296,-6.525) (3.325,-6.480) (3.359,-6.425) (3.389,-6.375) (3.418,-6.325) (3.445,-6.275) (3.472,-6.225) (3.497,-6.175) (3.521,-6.125) (3.545,-6.075) (3.567,-6.025) (3.589,-5.975) (3.609,-5.925) (3.629,-5.875) (3.647,-5.825) (3.665,-5.775) (3.682,-5.725) (3.698,-5.675) (3.713,-5.625) (3.728,-5.575) (3.741,-5.525) (3.754,-5.475) (3.765,-5.425) (3.776,-5.375) (3.787,-5.325) (3.796,-5.275) (3.804,-5.225) (3.812,-5.175) (3.819,-5.125) (3.825,-5.075) (3.830,-5.025) (3.835,-4.975) (3.839,-4.925) (3.842,-4.875) (3.844,-4.825) (3.845,-4.775) (3.846,-4.725) (3.846,-4.675) (3.845,-4.625) (3.843,-4.575) (3.841,-4.525) (3.837,-4.475) (3.833,-4.425) (3.828,-4.375) (3.823,-4.325) (3.816,-4.275) (3.809,-4.225) (3.801,-4.175) (3.792,-4.125) (3.783,-4.075) (3.772,-4.025) (3.761,-3.975) (3.749,-3.925) (3.736,-3.875) (3.722,-3.825) (3.708,-3.775) (3.692,-3.725) (3.676,-3.675) (3.659,-3.625) (3.641,-3.575) (3.622,-3.525) (3.602,-3.475) (3.582,-3.425) (3.560,-3.375) (3.538,-3.325) (3.514,-3.275) (3.490,-3.225) (3.465,-3.175) (3.439,-3.125) (3.412,-3.075) (3.384,-3.025) (3.355,-2.975) (3.325,-2.926) (3.294,-2.875) (3.262,-2.825) (3.229,-2.775) (3.200,-2.733) (3.160,-2.675) (3.125,-2.627) (3.087,-2.575) (3.050,-2.527) (3.009,-2.475) (2.975,-2.432) (2.928,-2.375) (2.886,-2.325) (2.850,-2.284) (2.800,-2.228) (2.752,-2.175) (2.705,-2.125) (2.657,-2.075) (2.609,-2.025) (2.559,-1.975) (2.508,-1.925) (2.456,-1.875) (2.403,-1.825) (2.350,-1.777) (2.300,-1.731) (2.250,-1.687) (2.209,-1.650) (2.151,-1.600) (2.100,-1.556) (2.050,-1.514) (2.003,-1.475) (1.950,-1.431) (1.900,-1.391) (1.850,-1.350) (1.800,-1.311) (1.755,-1.275) (1.700,-1.232) (1.659,-1.200) (1.600,-1.155) (1.550,-1.117) (1.500,-1.079) (1.450,-1.041) (1.400,-1.004) (1.350,-0.966) (1.300,-0.929) (1.250,-0.893) (1.200,-0.856) (1.158,-0.825) (1.100,-0.783) (1.055,-0.750) (1.000,-0.711) (0.951,-0.675) (0.900,-0.639) (0.850,-0.603) (0.800,-0.567) (0.750,-0.532) (0.706,-0.500) (0.650,-0.461) (0.600,-0.425) (0.550,-0.390) (0.500,-0.354) (0.450,-0.319) (0.400,-0.284) (0.354,-0.250) (0.300,-0.213) (0.250,-0.177) (0.206,-0.144)}
\EQ{(0.206,0.144) (0.248,0.175) (0.300,0.213) (0.350,0.248) (0.400,0.284) (0.450,0.319) (0.495,0.350) (0.550,0.390) (0.600,0.425) (0.650,0.461) (0.700,0.496) (0.742,0.525) (0.800,0.567) (0.846,0.600) (0.900,0.639) (0.950,0.675) (1.000,0.711) (1.050,0.747) (1.100,0.783) (1.150,0.819) (1.192,0.850) (1.250,0.893) (1.294,0.925) (1.350,0.966) (1.396,1.000) (1.450,1.041) (1.495,1.075) (1.550,1.117) (1.594,1.150) (1.650,1.193) (1.691,1.225) (1.750,1.271) (1.800,1.311) (1.850,1.350) (1.900,1.391) (1.942,1.425) (2.000,1.472) (2.050,1.514) (2.093,1.550) (2.150,1.599) (2.200,1.643) (2.250,1.687) (2.293,1.725) (2.348,1.775) (2.400,1.823) (2.450,1.870) (2.500,1.918) (2.550,1.967) (2.600,2.017) (2.650,2.068) (2.700,2.120) (2.750,2.173) (2.797,2.225) (2.842,2.275) (2.875,2.313) (2.925,2.372) (2.969,2.425) (3.009,2.475) (3.049,2.525) (3.087,2.575) (3.124,2.625) (3.160,2.675) (3.195,2.725) (3.225,2.770) (3.262,2.825) (3.294,2.875) (3.325,2.925) (3.355,2.975) (3.384,3.025) (3.412,3.075) (3.439,3.125) (3.465,3.175) (3.490,3.225) (3.514,3.275) (3.538,3.325) (3.560,3.375) (3.582,3.425) (3.602,3.475) (3.622,3.525) (3.641,3.575) (3.659,3.625) (3.675,3.673) (3.692,3.725) (3.708,3.775) (3.722,3.825) (3.736,3.875) (3.749,3.925) (3.761,3.975) (3.772,4.025) (3.783,4.075) (3.792,4.125) (3.801,4.175) (3.809,4.225) (3.816,4.275) (3.823,4.325) (3.828,4.375) (3.833,4.425) (3.837,4.475) (3.841,4.525) (3.843,4.575) (3.845,4.625) (3.846,4.675) (3.846,4.725) (3.845,4.775) (3.844,4.825) (3.842,4.875) (3.839,4.925) (3.835,4.975) (3.830,5.025) (3.825,5.075) (3.819,5.125) (3.812,5.175) (3.804,5.225) (3.796,5.275) (3.787,5.325) (3.776,5.375) (3.765,5.425) (3.754,5.475) (3.741,5.525) (3.728,5.575) (3.713,5.625) (3.698,5.675) (3.682,5.725) (3.665,5.775) (3.647,5.825) (3.629,5.875) (3.609,5.925) (3.589,5.975) (3.567,6.025) (3.545,6.075) (3.521,6.125) (3.497,6.175) (3.472,6.225) (3.445,6.275) (3.418,6.325) (3.389,6.375) (3.359,6.425) (3.328,6.475) (3.300,6.519) (3.263,6.575) (3.229,6.625) (3.193,6.675) (3.156,6.725) (3.125,6.765) (3.077,6.825) (3.036,6.875) (3.000,6.917) (2.950,6.974) (2.903,7.025) (2.855,7.075) (2.805,7.125) (2.753,7.175) (2.700,7.225) (2.650,7.269) (2.600,7.313) (2.555,7.350) (2.500,7.395) (2.450,7.434) (2.400,7.471) (2.350,7.508) (2.300,7.543) (2.253,7.575) (2.200,7.610) (2.150,7.642) (2.100,7.673) (2.054,7.700) (2.000,7.731) (1.950,7.759) (1.900,7.786) (1.850,7.812) (1.800,7.837) (1.750,7.861) (1.700,7.884) (1.650,7.907) (1.600,7.928) (1.550,7.949) (1.500,7.969) (1.450,7.989) (1.400,8.007) (1.350,8.025) (1.300,8.042) (1.250,8.058) (1.200,8.074) (1.150,8.089) (1.100,8.103) (1.050,8.117) (1.000,8.130) (0.950,8.142) (0.900,8.153) (0.850,8.164) (0.800,8.175) (0.750,8.184) (0.700,8.193) (0.650,8.201) (0.600,8.209) (0.550,8.216) (0.500,8.223) (0.450,8.228) (0.400,8.234) (0.350,8.238) (0.300,8.242) (0.250,8.245) (0.200,8.248) (0.150,8.250) (0.100,8.252) (0.050,8.253) (0.000,8.253) (-0.050,8.253) (-0.100,8.252) (-0.150,8.250) (-0.200,8.248) (-0.250,8.245) (-0.300,8.242) (-0.350,8.238) (-0.400,8.234) (-0.450,8.228) (-0.500,8.223) (-0.550,8.216) (-0.600,8.209) (-0.650,8.201) (-0.700,8.193) (-0.750,8.184) (-0.800,8.175) (-0.850,8.164) (-0.900,8.153) (-0.950,8.142) (-1.000,8.130) (-1.050,8.117) (-1.100,8.103) (-1.150,8.089) (-1.200,8.074) (-1.250,8.058) (-1.300,8.042) (-1.350,8.025) (-1.400,8.007) (-1.450,7.989) (-1.500,7.969) (-1.548,7.950) (-1.600,7.928) (-1.650,7.907) (-1.700,7.884) (-1.750,7.861) (-1.800,7.837) (-1.850,7.812) (-1.900,7.786) (-1.950,7.759) (-2.000,7.731) (-2.050,7.702) (-2.096,7.675) (-2.150,7.642) (-2.200,7.610) (-2.250,7.577) (-2.300,7.543) (-2.350,7.508) (-2.395,7.475) (-2.450,7.434) (-2.493,7.400) (-2.550,7.354) (-2.600,7.313) (-2.644,7.275) (-2.700,7.225) (-2.750,7.178) (-2.800,7.130) (-2.850,7.080) (-2.900,7.028) (-2.949,6.975) (-2.993,6.925) (-3.036,6.875) (-3.075,6.828) (-3.117,6.775) (-3.150,6.733) (-3.193,6.675) (-3.225,6.630) (-3.263,6.575) (-3.296,6.525) (-3.325,6.480) (-3.359,6.425) (-3.389,6.375) (-3.418,6.325) (-3.445,6.275) (-3.472,6.225) (-3.497,6.175) (-3.521,6.125) (-3.545,6.075) (-3.567,6.025) (-3.589,5.975) (-3.609,5.925) (-3.629,5.875) (-3.647,5.825) (-3.665,5.775) (-3.682,5.725) (-3.698,5.675) (-3.713,5.625) (-3.728,5.575) (-3.741,5.525) (-3.754,5.475) (-3.765,5.425) (-3.776,5.375) (-3.787,5.325) (-3.796,5.275) (-3.804,5.225) (-3.812,5.175) (-3.819,5.125) (-3.825,5.075) (-3.830,5.025) (-3.835,4.975) (-3.839,4.925) (-3.842,4.875) (-3.844,4.825) (-3.845,4.775) (-3.846,4.725) (-3.846,4.675) (-3.845,4.625) (-3.843,4.575) (-3.841,4.525) (-3.837,4.475) (-3.833,4.425) (-3.828,4.375) (-3.823,4.325) (-3.816,4.275) (-3.809,4.225) (-3.801,4.175) (-3.792,4.125) (-3.783,4.075) (-3.772,4.025) (-3.761,3.975) (-3.749,3.925) (-3.736,3.875) (-3.722,3.825) (-3.708,3.775) (-3.692,3.725) (-3.676,3.675) (-3.659,3.625) (-3.641,3.575) (-3.622,3.525) (-3.602,3.475) (-3.582,3.425) (-3.560,3.375) (-3.538,3.325) (-3.514,3.275) (-3.490,3.225) (-3.465,3.175) (-3.439,3.125) (-3.412,3.075) (-3.384,3.025) (-3.355,2.975) (-3.325,2.926) (-3.294,2.875) (-3.262,2.825) (-3.229,2.775) (-3.200,2.733) (-3.160,2.675) (-3.125,2.627) (-3.087,2.575) (-3.050,2.527) (-3.009,2.475) (-2.975,2.432) (-2.928,2.375) (-2.886,2.325) (-2.850,2.284) (-2.800,2.228) (-2.752,2.175) (-2.705,2.125) (-2.657,2.075) (-2.609,2.025) (-2.559,1.975) (-2.508,1.925) (-2.456,1.875) (-2.403,1.825) (-2.350,1.777) (-2.300,1.731) (-2.250,1.687) (-2.209,1.650) (-2.151,1.600) (-2.100,1.556) (-2.050,1.514) (-2.003,1.475) (-1.950,1.431) (-1.900,1.391) (-1.850,1.350) (-1.800,1.311) (-1.755,1.275) (-1.700,1.232) (-1.659,1.200) (-1.600,1.155) (-1.550,1.117) (-1.500,1.079) (-1.450,1.041) (-1.400,1.004) (-1.350,0.966) (-1.300,0.929) (-1.250,0.893) (-1.200,0.856) (-1.158,0.825) (-1.100,0.783) (-1.055,0.750) (-1.000,0.711) (-0.951,0.675) (-0.900,0.639) (-0.850,0.603) (-0.800,0.567) (-0.750,0.532) (-0.706,0.500) (-0.650,0.461) (-0.600,0.425) (-0.550,0.390) (-0.500,0.354) (-0.450,0.319) (-0.400,0.284) (-0.354,0.250) (-0.300,0.213) (-0.250,0.177) (-0.220,0.155)}
\EQ{(-0.220,-0.155) (-0.275,-0.195) (-0.319,-0.225) (-0.375,-0.266) (-0.424,-0.300) (-0.475,-0.337) (-0.525,-0.372) (-0.575,-0.407) (-0.625,-0.443) (-0.671,-0.475) (-0.725,-0.514) (-0.775,-0.550) (-0.825,-0.586) (-0.875,-0.622) (-0.925,-0.658) (-0.975,-0.694) (-1.019,-0.725) (-1.075,-0.767) (-1.121,-0.800) (-1.175,-0.840) (-1.223,-0.875) (-1.275,-0.914) (-1.323,-0.950) (-1.375,-0.989) (-1.422,-1.025) (-1.475,-1.066) (-1.520,-1.100) (-1.575,-1.143) (-1.616,-1.175) (-1.675,-1.222) (-1.725,-1.263) (-1.771,-1.300) (-1.825,-1.344) (-1.875,-1.386) (-1.921,-1.425) (-1.975,-1.472) (-2.025,-1.515) (-2.064,-1.550) (-2.120,-1.600) (-2.174,-1.650) (-2.225,-1.698) (-2.275,-1.746) (-2.325,-1.795) (-2.375,-1.845) (-2.425,-1.896) (-2.475,-1.949) (-2.522,-2.000) (-2.567,-2.050) (-2.611,-2.100) (-2.650,-2.146) (-2.695,-2.200) (-2.735,-2.250) (-2.774,-2.300) (-2.812,-2.350) (-2.849,-2.400) (-2.885,-2.450) (-2.919,-2.500) (-2.950,-2.546) (-2.985,-2.600) (-3.016,-2.650) (-3.046,-2.700) (-3.075,-2.750) (-3.100,-2.794) (-3.125,-2.841) (-3.150,-2.889) (-3.175,-2.939) (-3.200,-2.991) (-3.225,-3.045) (-3.249,-3.100) (-3.270,-3.150) (-3.290,-3.200) (-3.309,-3.250) (-3.326,-3.300) (-3.343,-3.350) (-3.359,-3.400) (-3.374,-3.450) (-3.389,-3.500) (-3.402,-3.550) (-3.414,-3.600) (-3.426,-3.650) (-3.436,-3.700) (-3.446,-3.750) (-3.455,-3.800) (-3.462,-3.850) (-3.469,-3.900) (-3.476,-3.950) (-3.481,-4.000) (-3.485,-4.050) (-3.489,-4.100) (-3.491,-4.150) (-3.493,-4.200) (-3.494,-4.250) (-3.494,-4.300) (-3.493,-4.350) (-3.492,-4.400) (-3.489,-4.450) (-3.486,-4.500) (-3.482,-4.550) (-3.477,-4.600) (-3.471,-4.650) (-3.464,-4.700) (-3.456,-4.750) (-3.447,-4.800) (-3.438,-4.850) (-3.428,-4.900) (-3.416,-4.950) (-3.404,-5.000) (-3.391,-5.050) (-3.377,-5.100) (-3.362,-5.150) (-3.346,-5.200) (-3.329,-5.250) (-3.311,-5.300) (-3.292,-5.350) (-3.272,-5.400) (-3.251,-5.450) (-3.229,-5.500) (-3.206,-5.550) (-3.182,-5.600) (-3.157,-5.650) (-3.131,-5.700) (-3.103,-5.750) (-3.075,-5.799) (-3.044,-5.850) (-3.013,-5.900) (-2.981,-5.950) (-2.950,-5.995) (-2.912,-6.050) (-2.875,-6.100) (-2.837,-6.150) (-2.800,-6.196) (-2.755,-6.250) (-2.712,-6.300) (-2.675,-6.342) (-2.625,-6.396) (-2.575,-6.447) (-2.525,-6.496) (-2.475,-6.544) (-2.425,-6.589) (-2.384,-6.625) (-2.325,-6.675) (-2.275,-6.716) (-2.231,-6.750) (-2.175,-6.792) (-2.130,-6.825) (-2.075,-6.863) (-2.025,-6.897) (-1.982,-6.925) (-1.925,-6.961) (-1.875,-6.991) (-1.825,-7.020) (-1.775,-7.048) (-1.725,-7.074) (-1.676,-7.100) (-1.625,-7.125) (-1.575,-7.149) (-1.525,-7.172) (-1.475,-7.194) (-1.425,-7.215) (-1.375,-7.236) (-1.325,-7.255) (-1.275,-7.274) (-1.225,-7.292) (-1.175,-7.309) (-1.125,-7.325) (-1.075,-7.340) (-1.025,-7.355) (-0.975,-7.369) (-0.925,-7.382) (-0.875,-7.395) (-0.825,-7.406) (-0.775,-7.417) (-0.725,-7.427) (-0.675,-7.437) (-0.625,-7.446) (-0.575,-7.454) (-0.525,-7.461) (-0.475,-7.468) (-0.425,-7.474) (-0.375,-7.479) (-0.325,-7.484) (-0.275,-7.488) (-0.225,-7.492) (-0.175,-7.494) (-0.125,-7.496) (-0.075,-7.498) (-0.025,-7.498) (0.050,-7.498) (0.100,-7.497) (0.150,-7.495) (0.200,-7.493) (0.250,-7.490) (0.300,-7.486) (0.350,-7.482) (0.400,-7.477) (0.450,-7.471) (0.500,-7.465) (0.550,-7.458) (0.599,-7.450) (0.650,-7.441) (0.700,-7.432) (0.750,-7.422) (0.800,-7.412) (0.850,-7.400) (0.900,-7.388) (0.950,-7.376) (1.000,-7.362) (1.050,-7.348) (1.100,-7.333) (1.150,-7.317) (1.200,-7.300) (1.250,-7.283) (1.300,-7.265) (1.350,-7.246) (1.400,-7.226) (1.450,-7.205) (1.500,-7.183) (1.550,-7.161) (1.600,-7.137) (1.650,-7.113) (1.700,-7.088) (1.750,-7.061) (1.800,-7.034) (1.850,-7.005) (1.900,-6.976) (1.950,-6.945) (2.000,-6.913) (2.050,-6.880) (2.094,-6.850) (2.150,-6.811) (2.198,-6.775) (2.250,-6.735) (2.295,-6.700) (2.350,-6.654) (2.400,-6.611) (2.441,-6.575) (2.495,-6.525) (2.547,-6.475) (2.597,-6.425) (2.644,-6.375) (2.690,-6.325) (2.725,-6.286) (2.775,-6.227) (2.817,-6.175) (2.850,-6.133) (2.893,-6.075) (2.925,-6.031) (2.964,-5.975) (2.997,-5.925) (3.025,-5.881) (3.060,-5.825) (3.089,-5.775) (3.117,-5.725) (3.144,-5.675) (3.170,-5.625) (3.195,-5.575) (3.218,-5.525) (3.241,-5.475) (3.262,-5.425) (3.282,-5.375) (3.302,-5.325) (3.320,-5.275) (3.338,-5.225) (3.354,-5.175) (3.370,-5.125) (3.384,-5.075) (3.398,-5.025) (3.410,-4.975) (3.422,-4.925) (3.433,-4.875) (3.443,-4.825) (3.452,-4.775) (3.460,-4.725) (3.467,-4.675) (3.474,-4.625) (3.479,-4.575) (3.484,-4.525) (3.488,-4.475) (3.491,-4.425) (3.493,-4.375) (3.494,-4.325) (3.494,-4.275) (3.494,-4.225) (3.492,-4.175) (3.490,-4.125) (3.487,-4.075) (3.483,-4.025) (3.478,-3.975) (3.473,-3.925) (3.466,-3.875) (3.459,-3.825) (3.450,-3.775) (3.441,-3.725) (3.431,-3.675) (3.420,-3.625) (3.408,-3.575) (3.395,-3.525) (3.382,-3.475) (3.367,-3.425) (3.352,-3.375) (3.335,-3.325) (3.318,-3.275) (3.300,-3.227) (3.280,-3.175) (3.260,-3.125) (3.238,-3.075) (3.216,-3.025) (3.193,-2.975) (3.168,-2.925) (3.143,-2.875) (3.117,-2.825) (3.089,-2.775) (3.061,-2.725) (3.031,-2.675) (3.001,-2.625) (2.969,-2.575) (2.936,-2.525) (2.902,-2.475) (2.867,-2.425) (2.831,-2.375) (2.800,-2.334) (2.755,-2.275) (2.715,-2.225) (2.675,-2.176) (2.632,-2.125) (2.589,-2.075) (2.550,-2.031) (2.500,-1.976) (2.452,-1.925) (2.404,-1.875) (2.355,-1.825) (2.305,-1.775) (2.254,-1.725) (2.201,-1.675) (2.150,-1.628) (2.100,-1.582) (2.050,-1.537) (2.008,-1.500) (1.950,-1.450) (1.900,-1.407) (1.850,-1.365) (1.802,-1.325) (1.750,-1.283) (1.710,-1.250) (1.650,-1.202) (1.600,-1.163) (1.552,-1.125) (1.500,-1.085) (1.455,-1.050) (1.400,-1.008) (1.357,-0.975) (1.300,-0.933) (1.257,-0.900) (1.200,-0.858) (1.155,-0.825) (1.100,-0.785) (1.053,-0.750) (1.000,-0.712) (0.950,-0.676) (0.900,-0.640) (0.850,-0.604) (0.800,-0.568) (0.750,-0.532) (0.706,-0.500) (0.650,-0.461) (0.600,-0.425) (0.550,-0.390) (0.500,-0.354) (0.450,-0.319) (0.400,-0.284) (0.354,-0.250) (0.300,-0.213) (0.250,-0.177) (0.206,-0.144)}
\EQ{(0.206,0.144) (0.248,0.175) (0.300,0.213) (0.350,0.248) (0.400,0.284) (0.450,0.319) (0.495,0.350) (0.550,0.390) (0.600,0.425) (0.650,0.461) (0.700,0.496) (0.741,0.525) (0.800,0.568) (0.845,0.600) (0.900,0.640) (0.949,0.675) (1.000,0.712) (1.050,0.748) (1.100,0.785) (1.150,0.821) (1.200,0.858) (1.250,0.895) (1.300,0.933) (1.350,0.970) (1.400,1.008) (1.450,1.046) (1.500,1.085) (1.550,1.124) (1.600,1.163) (1.647,1.200) (1.700,1.242) (1.741,1.275) (1.800,1.324) (1.850,1.365) (1.892,1.400) (1.950,1.450) (2.000,1.493) (2.050,1.537) (2.092,1.575) (2.147,1.625) (2.200,1.674) (2.250,1.722) (2.300,1.770) (2.350,1.820) (2.400,1.871) (2.450,1.923) (2.499,1.975) (2.545,2.025) (2.589,2.075) (2.625,2.117) (2.674,2.175) (2.715,2.225) (2.750,2.269) (2.793,2.325) (2.825,2.367) (2.867,2.425) (2.900,2.472) (2.936,2.525) (2.969,2.575) (3.000,2.624) (3.031,2.675) (3.061,2.725) (3.089,2.775) (3.117,2.825) (3.143,2.875) (3.168,2.925) (3.193,2.975) (3.216,3.025) (3.238,3.075) (3.260,3.125) (3.280,3.175) (3.299,3.225) (3.318,3.275) (3.335,3.325) (3.352,3.375) (3.367,3.425) (3.382,3.475) (3.395,3.525) (3.408,3.575) (3.420,3.625) (3.431,3.675) (3.441,3.725) (3.450,3.775) (3.459,3.825) (3.466,3.875) (3.473,3.925) (3.478,3.975) (3.483,4.025) (3.487,4.075) (3.490,4.125) (3.492,4.175) (3.494,4.225) (3.494,4.275) (3.494,4.325) (3.493,4.375) (3.491,4.425) (3.488,4.475) (3.484,4.525) (3.479,4.575) (3.474,4.625) (3.467,4.675) (3.460,4.725) (3.452,4.775) (3.443,4.825) (3.433,4.875) (3.422,4.925) (3.410,4.975) (3.398,5.025) (3.384,5.075) (3.370,5.125) (3.354,5.175) (3.338,5.225) (3.320,5.275) (3.302,5.325) (3.282,5.375) (3.262,5.425) (3.241,5.475) (3.218,5.525) (3.195,5.575) (3.170,5.625) (3.144,5.675) (3.117,5.725) (3.089,5.775) (3.060,5.825) (3.029,5.875) (3.000,5.921) (2.964,5.975) (2.929,6.025) (2.900,6.066) (2.856,6.125) (2.825,6.165) (2.776,6.225) (2.734,6.275) (2.700,6.314) (2.650,6.369) (2.600,6.422) (2.550,6.472) (2.500,6.520) (2.450,6.567) (2.400,6.611) (2.355,6.650) (2.300,6.696) (2.250,6.735) (2.200,6.774) (2.150,6.811) (2.100,6.846) (2.058,6.875) (2.000,6.913) (1.950,6.945) (1.901,6.975) (1.850,7.005) (1.800,7.034) (1.750,7.061) (1.700,7.088) (1.650,7.113) (1.600,7.137) (1.550,7.161) (1.500,7.183) (1.450,7.205) (1.402,7.225) (1.350,7.246) (1.300,7.265) (1.250,7.283) (1.201,7.300) (1.150,7.317) (1.100,7.333) (1.050,7.348) (1.000,7.362) (0.950,7.376) (0.900,7.388) (0.850,7.400) (0.800,7.412) (0.750,7.422) (0.700,7.432) (0.650,7.441) (0.600,7.450) (0.550,7.458) (0.500,7.465) (0.450,7.471) (0.400,7.477) (0.350,7.482) (0.300,7.486) (0.250,7.490) (0.200,7.493) (0.150,7.495) (0.100,7.497) (0.050,7.498) (0.000,7.498) (-0.050,7.498) (-0.100,7.497) (-0.150,7.495) (-0.200,7.493) (-0.250,7.490) (-0.300,7.486) (-0.350,7.482) (-0.400,7.477) (-0.450,7.471) (-0.500,7.465) (-0.550,7.458) (-0.599,7.450) (-0.650,7.441) (-0.700,7.432) (-0.750,7.422) (-0.800,7.412) (-0.850,7.400) (-0.900,7.388) (-0.950,7.376) (-1.000,7.362) (-1.050,7.348) (-1.100,7.333) (-1.150,7.317) (-1.200,7.300) (-1.250,7.283) (-1.300,7.265) (-1.350,7.246) (-1.400,7.226) (-1.450,7.205) (-1.500,7.183) (-1.550,7.161) (-1.600,7.137) (-1.650,7.113) (-1.700,7.088) (-1.750,7.061) (-1.800,7.034) (-1.850,7.005) (-1.900,6.976) (-1.950,6.945) (-2.000,6.913) (-2.050,6.880) (-2.094,6.850) (-2.150,6.811) (-2.198,6.775) (-2.250,6.735) (-2.295,6.700) (-2.350,6.654) (-2.400,6.611) (-2.441,6.575) (-2.495,6.525) (-2.547,6.475) (-2.597,6.425) (-2.644,6.375) (-2.690,6.325) (-2.725,6.286) (-2.775,6.227) (-2.817,6.175) (-2.850,6.133) (-2.893,6.075) (-2.925,6.031) (-2.964,5.975) (-2.997,5.925) (-3.025,5.881) (-3.060,5.825) (-3.089,5.775) (-3.117,5.725) (-3.144,5.675) (-3.170,5.625) (-3.195,5.575) (-3.218,5.525) (-3.241,5.475) (-3.262,5.425) (-3.282,5.375) (-3.302,5.325) (-3.320,5.275) (-3.338,5.225) (-3.354,5.175) (-3.370,5.125) (-3.384,5.075) (-3.398,5.025) (-3.410,4.975) (-3.422,4.925) (-3.433,4.875) (-3.443,4.825) (-3.452,4.775) (-3.460,4.725) (-3.467,4.675) (-3.474,4.625) (-3.479,4.575) (-3.484,4.525) (-3.488,4.475) (-3.491,4.425) (-3.493,4.375) (-3.494,4.325) (-3.494,4.275) (-3.494,4.225) (-3.492,4.175) (-3.490,4.125) (-3.487,4.075) (-3.483,4.025) (-3.478,3.975) (-3.473,3.925) (-3.466,3.875) (-3.459,3.825) (-3.450,3.775) (-3.441,3.725) (-3.431,3.675) (-3.420,3.625) (-3.408,3.575) (-3.395,3.525) (-3.382,3.475) (-3.367,3.425) (-3.352,3.375) (-3.335,3.325) (-3.318,3.275) (-3.300,3.227) (-3.280,3.175) (-3.260,3.125) (-3.238,3.075) (-3.216,3.025) (-3.193,2.975) (-3.168,2.925) (-3.143,2.875) (-3.117,2.825) (-3.089,2.775) (-3.061,2.725) (-3.031,2.675) (-3.001,2.625) (-2.969,2.575) (-2.936,2.525) (-2.902,2.475) (-2.867,2.425) (-2.831,2.375) (-2.800,2.334) (-2.755,2.275) (-2.715,2.225) (-2.675,2.176) (-2.632,2.125) (-2.589,2.075) (-2.550,2.031) (-2.500,1.976) (-2.452,1.925) (-2.404,1.875) (-2.355,1.825) (-2.305,1.775) (-2.254,1.725) (-2.201,1.675) (-2.150,1.628) (-2.100,1.582) (-2.050,1.537) (-2.008,1.500) (-1.950,1.450) (-1.900,1.407) (-1.850,1.365) (-1.802,1.325) (-1.750,1.283) (-1.710,1.250) (-1.650,1.202) (-1.600,1.163) (-1.552,1.125) (-1.500,1.085) (-1.455,1.050) (-1.400,1.008) (-1.357,0.975) (-1.300,0.933) (-1.257,0.900) (-1.200,0.858) (-1.155,0.825) (-1.100,0.785) (-1.053,0.750) (-1.000,0.712) (-0.950,0.676) (-0.900,0.640) (-0.850,0.604) (-0.800,0.568) (-0.750,0.532) (-0.706,0.500) (-0.650,0.461) (-0.600,0.425) (-0.550,0.390) (-0.500,0.354) (-0.450,0.319) (-0.400,0.284) (-0.354,0.250) (-0.300,0.213) (-0.250,0.177) (-0.220,0.155)}
\EQ{(-0.220,-0.155) (-0.275,-0.195) (-0.319,-0.225) (-0.375,-0.266) (-0.424,-0.300) (-0.475,-0.337) (-0.525,-0.372) (-0.575,-0.408) (-0.625,-0.443) (-0.671,-0.475) (-0.725,-0.515) (-0.775,-0.550) (-0.825,-0.586) (-0.875,-0.622) (-0.925,-0.659) (-0.975,-0.695) (-1.017,-0.725) (-1.075,-0.768) (-1.119,-0.800) (-1.175,-0.842) (-1.219,-0.875) (-1.275,-0.917) (-1.319,-0.950) (-1.375,-0.994) (-1.416,-1.025) (-1.475,-1.071) (-1.525,-1.111) (-1.574,-1.150) (-1.625,-1.191) (-1.667,-1.225) (-1.725,-1.274) (-1.775,-1.316) (-1.815,-1.350) (-1.872,-1.400) (-1.925,-1.447) (-1.975,-1.492) (-2.025,-1.538) (-2.064,-1.575) (-2.117,-1.625) (-2.167,-1.675) (-2.217,-1.725) (-2.265,-1.775) (-2.300,-1.812) (-2.335,-1.850) (-2.375,-1.894) (-2.424,-1.950) (-2.467,-2.000) (-2.500,-2.040) (-2.548,-2.100) (-2.587,-2.150) (-2.624,-2.200) (-2.661,-2.250) (-2.696,-2.300) (-2.725,-2.343) (-2.762,-2.400) (-2.794,-2.450) (-2.824,-2.500) (-2.850,-2.545) (-2.875,-2.589) (-2.900,-2.635) (-2.925,-2.683) (-2.950,-2.733) (-2.975,-2.786) (-3.000,-2.841) (-3.025,-2.899) (-3.046,-2.950) (-3.065,-3.000) (-3.083,-3.050) (-3.100,-3.099) (-3.117,-3.150) (-3.132,-3.200) (-3.146,-3.250) (-3.159,-3.300) (-3.171,-3.350) (-3.183,-3.400) (-3.193,-3.450) (-3.202,-3.500) (-3.211,-3.550) (-3.218,-3.600) (-3.224,-3.650) (-3.230,-3.700) (-3.235,-3.750) (-3.238,-3.800) (-3.241,-3.850) (-3.243,-3.900) (-3.244,-3.950) (-3.244,-4.000) (-3.243,-4.050) (-3.241,-4.100) (-3.238,-4.150) (-3.234,-4.200) (-3.229,-4.250) (-3.224,-4.300) (-3.217,-4.350) (-3.210,-4.400) (-3.201,-4.450) (-3.192,-4.500) (-3.182,-4.550) (-3.170,-4.600) (-3.158,-4.650) (-3.145,-4.700) (-3.130,-4.750) (-3.115,-4.800) (-3.099,-4.850) (-3.081,-4.900) (-3.063,-4.950) (-3.043,-5.000) (-3.022,-5.050) (-3.001,-5.100) (-2.978,-5.150) (-2.954,-5.200) (-2.928,-5.250) (-2.902,-5.300) (-2.875,-5.348) (-2.845,-5.400) (-2.814,-5.450) (-2.782,-5.500) (-2.750,-5.549) (-2.715,-5.600) (-2.678,-5.650) (-2.640,-5.700) (-2.601,-5.750) (-2.560,-5.800) (-2.525,-5.840) (-2.475,-5.896) (-2.425,-5.949) (-2.375,-6.000) (-2.325,-6.049) (-2.275,-6.096) (-2.225,-6.140) (-2.184,-6.175) (-2.125,-6.224) (-2.075,-6.263) (-2.026,-6.300) (-1.975,-6.337) (-1.925,-6.372) (-1.875,-6.406) (-1.825,-6.438) (-1.775,-6.469) (-1.725,-6.499) (-1.679,-6.525) (-1.634,-6.550) (-1.587,-6.575) (-1.538,-6.600) (-1.487,-6.625) (-1.433,-6.650) (-1.377,-6.675) (-1.325,-6.697) (-1.275,-6.718) (-1.225,-6.737) (-1.175,-6.756) (-1.125,-6.773) (-1.075,-6.790) (-1.025,-6.806) (-0.975,-6.821) (-0.925,-6.835) (-0.875,-6.849) (-0.825,-6.861) (-0.775,-6.873) (-0.725,-6.884) (-0.675,-6.895) (-0.625,-6.904) (-0.575,-6.913) (-0.525,-6.921) (-0.475,-6.928) (-0.425,-6.935) (-0.375,-6.941) (-0.325,-6.946) (-0.275,-6.950) (-0.225,-6.954) (-0.175,-6.956) (-0.125,-6.959) (-0.075,-6.960) (-0.025,-6.961) (0.050,-6.960) (0.100,-6.959) (0.150,-6.958) (0.200,-6.955) (0.250,-6.952) (0.300,-6.948) (0.350,-6.943) (0.400,-6.938) (0.450,-6.932) (0.500,-6.925) (0.550,-6.917) (0.600,-6.909) (0.650,-6.899) (0.700,-6.890) (0.750,-6.879) (0.800,-6.867) (0.850,-6.855) (0.900,-6.842) (0.950,-6.828) (1.000,-6.814) (1.050,-6.798) (1.100,-6.782) (1.150,-6.765) (1.200,-6.747) (1.250,-6.728) (1.300,-6.708) (1.350,-6.687) (1.400,-6.665) (1.450,-6.642) (1.500,-6.619) (1.550,-6.594) (1.600,-6.568) (1.650,-6.541) (1.700,-6.513) (1.750,-6.484) (1.800,-6.454) (1.845,-6.425) (1.900,-6.389) (1.950,-6.355) (1.992,-6.325) (2.050,-6.282) (2.092,-6.250) (2.150,-6.204) (2.200,-6.162) (2.242,-6.125) (2.297,-6.075) (2.350,-6.025) (2.400,-5.975) (2.448,-5.925) (2.494,-5.875) (2.538,-5.825) (2.575,-5.782) (2.621,-5.725) (2.660,-5.675) (2.697,-5.625) (2.732,-5.575) (2.766,-5.525) (2.799,-5.475) (2.830,-5.425) (2.859,-5.375) (2.888,-5.325) (2.915,-5.275) (2.941,-5.225) (2.966,-5.175) (2.989,-5.125) (3.012,-5.075) (3.033,-5.025) (3.053,-4.975) (3.072,-4.925) (3.090,-4.875) (3.107,-4.825) (3.123,-4.775) (3.138,-4.725) (3.151,-4.675) (3.164,-4.625) (3.176,-4.575) (3.187,-4.525) (3.197,-4.475) (3.206,-4.425) (3.214,-4.375) (3.221,-4.325) (3.227,-4.275) (3.232,-4.225) (3.236,-4.175) (3.239,-4.125) (3.242,-4.075) (3.243,-4.025) (3.244,-3.975) (3.243,-3.925) (3.242,-3.875) (3.240,-3.825) (3.237,-3.775) (3.232,-3.725) (3.227,-3.675) (3.221,-3.625) (3.214,-3.575) (3.207,-3.525) (3.198,-3.475) (3.188,-3.425) (3.177,-3.375) (3.165,-3.325) (3.153,-3.275) (3.139,-3.225) (3.124,-3.175) (3.109,-3.125) (3.092,-3.075) (3.075,-3.027) (3.056,-2.975) (3.036,-2.925) (3.015,-2.875) (2.993,-2.825) (2.970,-2.775) (2.946,-2.725) (2.921,-2.675) (2.895,-2.625) (2.867,-2.575) (2.839,-2.525) (2.809,-2.475) (2.778,-2.425) (2.750,-2.381) (2.713,-2.325) (2.678,-2.275) (2.643,-2.225) (2.606,-2.175) (2.575,-2.135) (2.528,-2.075) (2.488,-2.025) (2.450,-1.980) (2.402,-1.925) (2.358,-1.875) (2.312,-1.825) (2.275,-1.785) (2.225,-1.733) (2.175,-1.683) (2.125,-1.633) (2.075,-1.585) (2.038,-1.550) (2.000,-1.515) (1.957,-1.475) (1.901,-1.425) (1.850,-1.381) (1.800,-1.337) (1.756,-1.300) (1.700,-1.253) (1.650,-1.212) (1.605,-1.175) (1.550,-1.131) (1.500,-1.091) (1.450,-1.052) (1.400,-1.013) (1.351,-0.975) (1.300,-0.936) (1.252,-0.900) (1.200,-0.861) (1.152,-0.825) (1.100,-0.787) (1.051,-0.750) (1.000,-0.713) (0.950,-0.677) (0.900,-0.640) (0.850,-0.604) (0.800,-0.568) (0.750,-0.532) (0.705,-0.500) (0.650,-0.461) (0.600,-0.425) (0.550,-0.390) (0.500,-0.354) (0.450,-0.319) (0.400,-0.284) (0.354,-0.250) (0.300,-0.213) (0.250,-0.177) (0.206,-0.144)}
\EQ{(0.206,0.144) (0.248,0.175) (0.300,0.213) (0.350,0.248) (0.400,0.284) (0.450,0.319) (0.495,0.350) (0.550,0.390) (0.600,0.425) (0.650,0.461) (0.700,0.497) (0.750,0.532) (0.800,0.568) (0.845,0.600) (0.900,0.640) (0.948,0.675) (1.000,0.713) (1.050,0.750) (1.100,0.787) (1.150,0.823) (1.200,0.861) (1.250,0.898) (1.300,0.936) (1.350,0.974) (1.400,1.013) (1.448,1.050) (1.500,1.091) (1.543,1.125) (1.600,1.171) (1.650,1.212) (1.697,1.250) (1.750,1.295) (1.800,1.337) (1.844,1.375) (1.900,1.424) (1.950,1.469) (2.000,1.515) (2.038,1.550) (2.075,1.585) (2.117,1.625) (2.167,1.675) (2.217,1.725) (2.265,1.775) (2.300,1.812) (2.335,1.850) (2.375,1.894) (2.424,1.950) (2.467,2.000) (2.500,2.040) (2.548,2.100) (2.587,2.150) (2.624,2.200) (2.661,2.250) (2.696,2.300) (2.725,2.343) (2.762,2.400) (2.794,2.450) (2.824,2.500) (2.850,2.545) (2.875,2.589) (2.900,2.635) (2.925,2.683) (2.950,2.733) (2.975,2.786) (3.000,2.841) (3.025,2.899) (3.046,2.950) (3.065,3.000) (3.083,3.050) (3.100,3.099) (3.117,3.150) (3.132,3.200) (3.146,3.250) (3.159,3.300) (3.171,3.350) (3.183,3.400) (3.193,3.450) (3.202,3.500) (3.211,3.550) (3.218,3.600) (3.224,3.650) (3.230,3.700) (3.235,3.750) (3.238,3.800) (3.241,3.850) (3.243,3.900) (3.244,3.950) (3.244,4.000) (3.243,4.050) (3.241,4.100) (3.238,4.150) (3.234,4.200) (3.229,4.250) (3.224,4.300) (3.217,4.350) (3.210,4.400) (3.201,4.450) (3.192,4.500) (3.182,4.550) (3.170,4.600) (3.158,4.650) (3.145,4.700) (3.130,4.750) (3.115,4.800) (3.099,4.850) (3.081,4.900) (3.063,4.950) (3.043,5.000) (3.022,5.050) (3.001,5.100) (2.978,5.150) (2.954,5.200) (2.928,5.250) (2.902,5.300) (2.875,5.348) (2.845,5.400) (2.814,5.450) (2.782,5.500) (2.750,5.549) (2.715,5.600) (2.678,5.650) (2.640,5.700) (2.601,5.750) (2.560,5.800) (2.525,5.840) (2.475,5.896) (2.425,5.949) (2.375,6.000) (2.325,6.049) (2.275,6.096) (2.225,6.140) (2.184,6.175) (2.125,6.224) (2.075,6.263) (2.026,6.300) (1.975,6.337) (1.925,6.372) (1.875,6.406) (1.825,6.438) (1.775,6.469) (1.725,6.499) (1.679,6.525) (1.634,6.550) (1.587,6.575) (1.538,6.600) (1.487,6.625) (1.433,6.650) (1.377,6.675) (1.325,6.697) (1.275,6.718) (1.225,6.737) (1.175,6.756) (1.125,6.773) (1.075,6.790) (1.025,6.806) (0.975,6.821) (0.925,6.835) (0.875,6.849) (0.825,6.861) (0.775,6.873) (0.725,6.884) (0.675,6.895) (0.625,6.904) (0.575,6.913) (0.525,6.921) (0.475,6.928) (0.425,6.935) (0.375,6.941) (0.325,6.946) (0.275,6.950) (0.225,6.954) (0.175,6.956) (0.125,6.959) (0.075,6.960) (0.025,6.961) (-0.050,6.960) (-0.100,6.959) (-0.150,6.958) (-0.200,6.955) (-0.250,6.952) (-0.300,6.948) (-0.350,6.943) (-0.400,6.938) (-0.450,6.932) (-0.500,6.925) (-0.550,6.917) (-0.600,6.909) (-0.650,6.899) (-0.700,6.890) (-0.750,6.879) (-0.800,6.867) (-0.850,6.855) (-0.900,6.842) (-0.950,6.828) (-1.000,6.814) (-1.050,6.798) (-1.100,6.782) (-1.150,6.765) (-1.200,6.747) (-1.250,6.728) (-1.300,6.708) (-1.350,6.687) (-1.400,6.665) (-1.450,6.642) (-1.500,6.619) (-1.550,6.594) (-1.600,6.568) (-1.650,6.541) (-1.700,6.513) (-1.750,6.484) (-1.800,6.454) (-1.845,6.425) (-1.900,6.389) (-1.950,6.355) (-1.992,6.325) (-2.050,6.282) (-2.092,6.250) (-2.150,6.204) (-2.200,6.162) (-2.242,6.125) (-2.297,6.075) (-2.350,6.025) (-2.400,5.975) (-2.448,5.925) (-2.494,5.875) (-2.538,5.825) (-2.575,5.782) (-2.621,5.725) (-2.660,5.675) (-2.697,5.625) (-2.732,5.575) (-2.766,5.525) (-2.799,5.475) (-2.830,5.425) (-2.859,5.375) (-2.888,5.325) (-2.915,5.275) (-2.941,5.225) (-2.966,5.175) (-2.989,5.125) (-3.012,5.075) (-3.033,5.025) (-3.053,4.975) (-3.072,4.925) (-3.090,4.875) (-3.107,4.825) (-3.123,4.775) (-3.138,4.725) (-3.151,4.675) (-3.164,4.625) (-3.176,4.575) (-3.187,4.525) (-3.197,4.475) (-3.206,4.425) (-3.214,4.375) (-3.221,4.325) (-3.227,4.275) (-3.232,4.225) (-3.236,4.175) (-3.239,4.125) (-3.242,4.075) (-3.243,4.025) (-3.244,3.975) (-3.243,3.925) (-3.242,3.875) (-3.240,3.825) (-3.237,3.775) (-3.232,3.725) (-3.227,3.675) (-3.221,3.625) (-3.214,3.575) (-3.207,3.525) (-3.198,3.475) (-3.188,3.425) (-3.177,3.375) (-3.165,3.325) (-3.153,3.275) (-3.139,3.225) (-3.124,3.175) (-3.109,3.125) (-3.092,3.075) (-3.075,3.027) (-3.056,2.975) (-3.036,2.925) (-3.015,2.875) (-2.993,2.825) (-2.970,2.775) (-2.946,2.725) (-2.921,2.675) (-2.895,2.625) (-2.867,2.575) (-2.839,2.525) (-2.809,2.475) (-2.778,2.425) (-2.750,2.381) (-2.713,2.325) (-2.678,2.275) (-2.643,2.225) (-2.606,2.175) (-2.575,2.135) (-2.528,2.075) (-2.488,2.025) (-2.450,1.980) (-2.402,1.925) (-2.358,1.875) (-2.312,1.825) (-2.275,1.785) (-2.225,1.733) (-2.175,1.683) (-2.125,1.633) (-2.075,1.585) (-2.038,1.550) (-2.000,1.515) (-1.957,1.475) (-1.901,1.425) (-1.850,1.381) (-1.800,1.337) (-1.756,1.300) (-1.700,1.253) (-1.650,1.212) (-1.605,1.175) (-1.550,1.131) (-1.500,1.091) (-1.450,1.052) (-1.400,1.013) (-1.351,0.975) (-1.300,0.936) (-1.252,0.900) (-1.200,0.861) (-1.152,0.825) (-1.100,0.787) (-1.051,0.750) (-1.000,0.713) (-0.950,0.677) (-0.900,0.640) (-0.850,0.604) (-0.800,0.568) (-0.750,0.532) (-0.705,0.500) (-0.650,0.461) (-0.600,0.425) (-0.550,0.390) (-0.500,0.354) (-0.450,0.319) (-0.400,0.284) (-0.354,0.250) (-0.300,0.213) (-0.250,0.177) (-0.220,0.155)}
\EQ{(-0.220,-0.155) (-0.275,-0.195) (-0.319,-0.225) (-0.375,-0.266) (-0.424,-0.300) (-0.475,-0.337) (-0.525,-0.372) (-0.575,-0.408) (-0.625,-0.443) (-0.670,-0.475) (-0.725,-0.515) (-0.774,-0.550) (-0.825,-0.587) (-0.875,-0.623) (-0.925,-0.659) (-0.975,-0.696) (-1.025,-0.733) (-1.075,-0.770) (-1.116,-0.800) (-1.175,-0.844) (-1.216,-0.875) (-1.275,-0.920) (-1.325,-0.959) (-1.375,-0.998) (-1.425,-1.037) (-1.472,-1.075) (-1.525,-1.118) (-1.565,-1.150) (-1.625,-1.200) (-1.675,-1.242) (-1.713,-1.275) (-1.771,-1.325) (-1.825,-1.374) (-1.875,-1.419) (-1.925,-1.466) (-1.962,-1.500) (-2.000,-1.537) (-2.038,-1.575) (-2.075,-1.612) (-2.112,-1.650) (-2.150,-1.690) (-2.200,-1.744) (-2.250,-1.799) (-2.294,-1.850) (-2.336,-1.900) (-2.375,-1.948) (-2.417,-2.000) (-2.450,-2.044) (-2.491,-2.100) (-2.525,-2.147) (-2.561,-2.200) (-2.594,-2.250) (-2.625,-2.299) (-2.656,-2.350) (-2.685,-2.400) (-2.713,-2.450) (-2.740,-2.500) (-2.765,-2.550) (-2.790,-2.600) (-2.813,-2.650) (-2.835,-2.700) (-2.856,-2.750) (-2.875,-2.799) (-2.894,-2.850) (-2.912,-2.900) (-2.928,-2.950) (-2.944,-3.000) (-2.958,-3.050) (-2.972,-3.100) (-2.984,-3.150) (-2.995,-3.200) (-3.005,-3.250) (-3.015,-3.300) (-3.023,-3.350) (-3.030,-3.400) (-3.036,-3.450) (-3.041,-3.500) (-3.045,-3.550) (-3.049,-3.600) (-3.051,-3.650) (-3.052,-3.700) (-3.052,-3.750) (-3.052,-3.800) (-3.050,-3.850) (-3.047,-3.900) (-3.044,-3.950) (-3.039,-4.000) (-3.033,-4.050) (-3.027,-4.100) (-3.019,-4.150) (-3.010,-4.200) (-3.001,-4.250) (-2.990,-4.300) (-2.978,-4.350) (-2.965,-4.400) (-2.951,-4.450) (-2.937,-4.500) (-2.921,-4.550) (-2.903,-4.600) (-2.885,-4.650) (-2.866,-4.700) (-2.845,-4.750) (-2.825,-4.796) (-2.800,-4.850) (-2.776,-4.900) (-2.751,-4.950) (-2.725,-4.998) (-2.700,-5.042) (-2.666,-5.100) (-2.635,-5.150) (-2.602,-5.200) (-2.568,-5.250) (-2.533,-5.300) (-2.500,-5.344) (-2.456,-5.400) (-2.415,-5.450) (-2.375,-5.497) (-2.328,-5.550) (-2.281,-5.600) (-2.231,-5.650) (-2.180,-5.700) (-2.125,-5.750) (-2.075,-5.794) (-2.025,-5.836) (-1.977,-5.875) (-1.925,-5.915) (-1.878,-5.950) (-1.825,-5.987) (-1.775,-6.021) (-1.731,-6.050) (-1.675,-6.085) (-1.625,-6.115) (-1.575,-6.143) (-1.525,-6.170) (-1.475,-6.197) (-1.425,-6.222) (-1.375,-6.246) (-1.325,-6.269) (-1.275,-6.291) (-1.225,-6.311) (-1.175,-6.331) (-1.126,-6.350) (-1.075,-6.368) (-1.025,-6.385) (-0.975,-6.401) (-0.925,-6.417) (-0.875,-6.431) (-0.825,-6.445) (-0.775,-6.457) (-0.725,-6.469) (-0.675,-6.480) (-0.625,-6.490) (-0.575,-6.499) (-0.525,-6.508) (-0.475,-6.516) (-0.425,-6.523) (-0.375,-6.529) (-0.325,-6.534) (-0.275,-6.539) (-0.225,-6.543) (-0.175,-6.546) (-0.125,-6.548) (-0.075,-6.549) (-0.025,-6.550) (0.047,-6.550) (0.100,-6.549) (0.150,-6.547) (0.200,-6.544) (0.250,-6.541) (0.300,-6.537) (0.350,-6.532) (0.400,-6.526) (0.450,-6.519) (0.500,-6.512) (0.550,-6.504) (0.600,-6.495) (0.650,-6.485) (0.700,-6.475) (0.750,-6.463) (0.800,-6.451) (0.850,-6.438) (0.900,-6.424) (0.950,-6.409) (1.000,-6.393) (1.050,-6.377) (1.100,-6.359) (1.150,-6.341) (1.200,-6.322) (1.250,-6.301) (1.300,-6.280) (1.350,-6.257) (1.400,-6.234) (1.450,-6.209) (1.500,-6.184) (1.550,-6.157) (1.600,-6.129) (1.650,-6.100) (1.700,-6.069) (1.750,-6.038) (1.800,-6.004) (1.842,-5.975) (1.900,-5.934) (1.945,-5.900) (2.000,-5.857) (2.050,-5.816) (2.097,-5.775) (2.150,-5.728) (2.200,-5.681) (2.250,-5.632) (2.300,-5.580) (2.350,-5.525) (2.394,-5.475) (2.436,-5.425) (2.475,-5.376) (2.514,-5.325) (2.550,-5.276) (2.586,-5.225) (2.619,-5.175) (2.650,-5.126) (2.681,-5.075) (2.710,-5.025) (2.737,-4.975) (2.763,-4.925) (2.788,-4.875) (2.812,-4.825) (2.834,-4.775) (2.856,-4.725) (2.875,-4.676) (2.894,-4.625) (2.912,-4.575) (2.929,-4.525) (2.944,-4.475) (2.959,-4.425) (2.972,-4.375) (2.984,-4.325) (2.995,-4.275) (3.006,-4.225) (3.015,-4.175) (3.023,-4.125) (3.030,-4.075) (3.036,-4.025) (3.041,-3.975) (3.046,-3.925) (3.049,-3.875) (3.051,-3.825) (3.052,-3.775) (3.052,-3.725) (3.052,-3.675) (3.050,-3.625) (3.047,-3.575) (3.044,-3.525) (3.039,-3.475) (3.033,-3.425) (3.026,-3.375) (3.019,-3.325) (3.010,-3.275) (3.000,-3.225) (2.990,-3.175) (2.978,-3.125) (2.965,-3.075) (2.951,-3.025) (2.936,-2.975) (2.920,-2.925) (2.903,-2.875) (2.885,-2.825) (2.866,-2.775) (2.845,-2.725) (2.825,-2.677) (2.801,-2.625) (2.778,-2.575) (2.753,-2.525) (2.726,-2.475) (2.700,-2.427) (2.675,-2.383) (2.641,-2.325) (2.610,-2.275) (2.578,-2.225) (2.544,-2.175) (2.509,-2.125) (2.475,-2.078) (2.436,-2.025) (2.400,-1.979) (2.357,-1.925) (2.315,-1.875) (2.275,-1.828) (2.228,-1.775) (2.183,-1.725) (2.136,-1.675) (2.100,-1.638) (2.063,-1.600) (2.025,-1.562) (1.987,-1.525) (1.950,-1.489) (1.909,-1.450) (1.854,-1.400) (1.800,-1.351) (1.750,-1.307) (1.700,-1.264) (1.655,-1.225) (1.600,-1.179) (1.550,-1.138) (1.504,-1.100) (1.450,-1.057) (1.410,-1.025) (1.350,-0.978) (1.300,-0.940) (1.250,-0.901) (1.200,-0.863) (1.150,-0.826) (1.100,-0.788) (1.050,-0.751) (1.000,-0.714) (0.950,-0.678) (0.900,-0.641) (0.850,-0.605) (0.809,-0.575) (0.750,-0.533) (0.705,-0.500) (0.650,-0.461) (0.600,-0.425) (0.550,-0.390) (0.500,-0.354) (0.450,-0.319) (0.400,-0.284) (0.354,-0.250) (0.300,-0.213) (0.250,-0.177) (0.206,-0.144)}
\EQ{(0.206,0.144) (0.248,0.175) (0.300,0.213) (0.350,0.248) (0.400,0.284) (0.450,0.319) (0.495,0.350) (0.550,0.390) (0.600,0.425) (0.650,0.461) (0.700,0.497) (0.750,0.533) (0.800,0.569) (0.844,0.600) (0.900,0.641) (0.947,0.675) (1.000,0.714) (1.049,0.750) (1.100,0.788) (1.149,0.825) (1.200,0.863) (1.249,0.900) (1.300,0.940) (1.346,0.975) (1.400,1.018) (1.441,1.050) (1.500,1.097) (1.550,1.138) (1.595,1.175) (1.650,1.221) (1.700,1.264) (1.742,1.300) (1.799,1.350) (1.850,1.396) (1.900,1.442) (1.950,1.489) (1.987,1.525) (2.025,1.562) (2.063,1.600) (2.100,1.638) (2.136,1.675) (2.175,1.717) (2.225,1.771) (2.273,1.825) (2.315,1.875) (2.350,1.917) (2.397,1.975) (2.436,2.025) (2.473,2.075) (2.509,2.125) (2.544,2.175) (2.575,2.221) (2.610,2.275) (2.641,2.325) (2.671,2.375) (2.699,2.425) (2.725,2.472) (2.750,2.520) (2.775,2.570) (2.800,2.622) (2.824,2.675) (2.845,2.725) (2.866,2.775) (2.885,2.825) (2.903,2.875) (2.920,2.925) (2.936,2.975) (2.951,3.025) (2.965,3.075) (2.978,3.125) (2.990,3.175) (3.000,3.225) (3.010,3.275) (3.019,3.325) (3.026,3.375) (3.033,3.425) (3.039,3.475) (3.044,3.525) (3.047,3.575) (3.050,3.625) (3.052,3.675) (3.052,3.725) (3.052,3.775) (3.051,3.825) (3.049,3.875) (3.046,3.925) (3.041,3.975) (3.036,4.025) (3.030,4.075) (3.023,4.125) (3.015,4.175) (3.006,4.225) (2.995,4.275) (2.984,4.325) (2.972,4.375) (2.959,4.425) (2.944,4.475) (2.929,4.525) (2.912,4.575) (2.894,4.625) (2.876,4.675) (2.856,4.725) (2.834,4.775) (2.812,4.825) (2.788,4.875) (2.763,4.925) (2.737,4.975) (2.710,5.025) (2.681,5.075) (2.651,5.125) (2.619,5.175) (2.586,5.225) (2.551,5.275) (2.514,5.325) (2.476,5.375) (2.436,5.425) (2.400,5.468) (2.350,5.525) (2.305,5.575) (2.256,5.625) (2.206,5.675) (2.153,5.725) (2.100,5.773) (2.050,5.816) (2.008,5.850) (1.950,5.896) (1.900,5.934) (1.850,5.970) (1.806,6.000) (1.750,6.038) (1.700,6.069) (1.650,6.100) (1.600,6.129) (1.550,6.157) (1.500,6.184) (1.450,6.209) (1.400,6.234) (1.350,6.257) (1.300,6.280) (1.253,6.300) (1.200,6.322) (1.150,6.341) (1.100,6.359) (1.050,6.377) (1.000,6.393) (0.950,6.409) (0.900,6.424) (0.850,6.438) (0.800,6.451) (0.750,6.463) (0.700,6.475) (0.650,6.485) (0.600,6.495) (0.550,6.504) (0.500,6.512) (0.450,6.519) (0.400,6.526) (0.350,6.532) (0.300,6.537) (0.250,6.541) (0.200,6.544) (0.150,6.547) (0.100,6.549) (0.050,6.550) (0.000,6.550) (-0.050,6.550) (-0.100,6.549) (-0.150,6.547) (-0.200,6.544) (-0.250,6.541) (-0.300,6.537) (-0.350,6.532) (-0.400,6.526) (-0.450,6.519) (-0.500,6.512) (-0.550,6.504) (-0.600,6.495) (-0.650,6.485) (-0.700,6.475) (-0.750,6.463) (-0.800,6.451) (-0.850,6.438) (-0.900,6.424) (-0.950,6.409) (-1.000,6.393) (-1.050,6.377) (-1.100,6.359) (-1.150,6.341) (-1.200,6.322) (-1.250,6.301) (-1.300,6.280) (-1.350,6.257) (-1.400,6.234) (-1.450,6.209) (-1.500,6.184) (-1.550,6.157) (-1.600,6.129) (-1.650,6.100) (-1.700,6.069) (-1.750,6.038) (-1.800,6.004) (-1.842,5.975) (-1.900,5.934) (-1.945,5.900) (-2.000,5.857) (-2.050,5.816) (-2.097,5.775) (-2.150,5.728) (-2.200,5.681) (-2.250,5.632) (-2.300,5.580) (-2.350,5.525) (-2.394,5.475) (-2.436,5.425) (-2.475,5.376) (-2.514,5.325) (-2.550,5.276) (-2.586,5.225) (-2.619,5.175) (-2.650,5.126) (-2.681,5.075) (-2.710,5.025) (-2.737,4.975) (-2.763,4.925) (-2.788,4.875) (-2.812,4.825) (-2.834,4.775) (-2.856,4.725) (-2.875,4.676) (-2.894,4.625) (-2.912,4.575) (-2.929,4.525) (-2.944,4.475) (-2.959,4.425) (-2.972,4.375) (-2.984,4.325) (-2.995,4.275) (-3.006,4.225) (-3.015,4.175) (-3.023,4.125) (-3.030,4.075) (-3.036,4.025) (-3.041,3.975) (-3.046,3.925) (-3.049,3.875) (-3.051,3.825) (-3.052,3.775) (-3.052,3.725) (-3.052,3.675) (-3.050,3.625) (-3.047,3.575) (-3.044,3.525) (-3.039,3.475) (-3.033,3.425) (-3.026,3.375) (-3.019,3.325) (-3.010,3.275) (-3.000,3.225) (-2.990,3.175) (-2.978,3.125) (-2.965,3.075) (-2.951,3.025) (-2.936,2.975) (-2.920,2.925) (-2.903,2.875) (-2.885,2.825) (-2.866,2.775) (-2.845,2.725) (-2.825,2.677) (-2.801,2.625) (-2.778,2.575) (-2.753,2.525) (-2.726,2.475) (-2.700,2.427) (-2.675,2.383) (-2.641,2.325) (-2.610,2.275) (-2.578,2.225) (-2.544,2.175) (-2.509,2.125) (-2.475,2.078) (-2.436,2.025) (-2.400,1.979) (-2.357,1.925) (-2.315,1.875) (-2.275,1.828) (-2.228,1.775) (-2.183,1.725) (-2.136,1.675) (-2.100,1.638) (-2.063,1.600) (-2.025,1.562) (-1.987,1.525) (-1.950,1.489) (-1.909,1.450) (-1.854,1.400) (-1.800,1.351) (-1.750,1.307) (-1.700,1.264) (-1.655,1.225) (-1.600,1.179) (-1.550,1.138) (-1.504,1.100) (-1.450,1.057) (-1.410,1.025) (-1.350,0.978) (-1.300,0.940) (-1.250,0.901) (-1.200,0.863) (-1.150,0.826) (-1.100,0.788) (-1.050,0.751) (-1.000,0.714) (-0.950,0.678) (-0.900,0.641) (-0.850,0.605) (-0.809,0.575) (-0.750,0.533) (-0.705,0.500) (-0.650,0.461) (-0.600,0.425) (-0.550,0.390) (-0.500,0.354) (-0.450,0.319) (-0.400,0.284) (-0.354,0.250) (-0.300,0.213) (-0.250,0.177) (-0.220,0.155)}
\EQ{(-0.220,-0.155) (-0.275,-0.195) (-0.319,-0.225) (-0.375,-0.266) (-0.424,-0.300) (-0.475,-0.337) (-0.525,-0.372) (-0.575,-0.408) (-0.625,-0.444) (-0.670,-0.475) (-0.725,-0.515) (-0.774,-0.550) (-0.825,-0.587) (-0.875,-0.624) (-0.925,-0.660) (-0.975,-0.697) (-1.025,-0.734) (-1.075,-0.771) (-1.125,-0.809) (-1.175,-0.847) (-1.225,-0.885) (-1.275,-0.924) (-1.325,-0.963) (-1.372,-1.000) (-1.425,-1.043) (-1.465,-1.075) (-1.525,-1.124) (-1.575,-1.166) (-1.615,-1.200) (-1.672,-1.250) (-1.725,-1.297) (-1.775,-1.342) (-1.825,-1.389) (-1.863,-1.425) (-1.900,-1.461) (-1.940,-1.500) (-1.990,-1.550) (-2.025,-1.587) (-2.061,-1.625) (-2.100,-1.668) (-2.150,-1.724) (-2.194,-1.775) (-2.236,-1.825) (-2.275,-1.874) (-2.315,-1.925) (-2.350,-1.972) (-2.388,-2.025) (-2.423,-2.075) (-2.456,-2.125) (-2.488,-2.175) (-2.519,-2.225) (-2.548,-2.275) (-2.575,-2.324) (-2.600,-2.370) (-2.625,-2.419) (-2.650,-2.470) (-2.675,-2.524) (-2.697,-2.575) (-2.718,-2.625) (-2.737,-2.675) (-2.756,-2.725) (-2.773,-2.775) (-2.789,-2.825) (-2.804,-2.875) (-2.818,-2.925) (-2.830,-2.975) (-2.842,-3.025) (-2.852,-3.075) (-2.862,-3.125) (-2.870,-3.175) (-2.877,-3.225) (-2.884,-3.275) (-2.889,-3.325) (-2.893,-3.375) (-2.896,-3.425) (-2.898,-3.475) (-2.899,-3.525) (-2.899,-3.575) (-2.898,-3.625) (-2.896,-3.675) (-2.893,-3.725) (-2.889,-3.775) (-2.884,-3.825) (-2.878,-3.875) (-2.871,-3.925) (-2.862,-3.975) (-2.853,-4.025) (-2.842,-4.075) (-2.831,-4.125) (-2.818,-4.175) (-2.804,-4.225) (-2.789,-4.275) (-2.773,-4.325) (-2.756,-4.375) (-2.738,-4.425) (-2.718,-4.475) (-2.697,-4.525) (-2.675,-4.575) (-2.652,-4.625) (-2.627,-4.675) (-2.601,-4.725) (-2.575,-4.772) (-2.550,-4.815) (-2.514,-4.875) (-2.482,-4.925) (-2.450,-4.972) (-2.413,-5.025) (-2.376,-5.075) (-2.337,-5.125) (-2.300,-5.170) (-2.253,-5.225) (-2.208,-5.275) (-2.161,-5.325) (-2.125,-5.362) (-2.086,-5.400) (-2.032,-5.450) (-1.976,-5.500) (-1.925,-5.543) (-1.885,-5.575) (-1.825,-5.621) (-1.775,-5.658) (-1.725,-5.693) (-1.678,-5.725) (-1.625,-5.759) (-1.575,-5.790) (-1.525,-5.819) (-1.475,-5.847) (-1.425,-5.874) (-1.375,-5.900) (-1.325,-5.924) (-1.275,-5.947) (-1.225,-5.970) (-1.175,-5.991) (-1.125,-6.011) (-1.075,-6.030) (-1.025,-6.048) (-0.975,-6.065) (-0.925,-6.081) (-0.875,-6.096) (-0.825,-6.111) (-0.775,-6.124) (-0.725,-6.137) (-0.675,-6.148) (-0.625,-6.159) (-0.575,-6.169) (-0.525,-6.178) (-0.475,-6.186) (-0.425,-6.193) (-0.375,-6.200) (-0.325,-6.205) (-0.275,-6.210) (-0.225,-6.214) (-0.175,-6.217) (-0.125,-6.220) (-0.075,-6.221) (-0.025,-6.222) (0.050,-6.222) (0.100,-6.221) (0.150,-6.219) (0.200,-6.216) (0.250,-6.212) (0.300,-6.208) (0.350,-6.203) (0.400,-6.196) (0.450,-6.190) (0.500,-6.182) (0.550,-6.173) (0.600,-6.164) (0.650,-6.154) (0.700,-6.142) (0.750,-6.130) (0.800,-6.117) (0.850,-6.104) (0.900,-6.089) (0.950,-6.073) (1.000,-6.057) (1.050,-6.039) (1.100,-6.020) (1.150,-6.001) (1.200,-5.980) (1.250,-5.959) (1.300,-5.936) (1.350,-5.912) (1.400,-5.887) (1.450,-5.861) (1.500,-5.833) (1.550,-5.805) (1.599,-5.775) (1.650,-5.743) (1.700,-5.710) (1.750,-5.676) (1.800,-5.640) (1.850,-5.602) (1.900,-5.563) (1.946,-5.525) (2.000,-5.479) (2.050,-5.434) (2.100,-5.386) (2.137,-5.350) (2.175,-5.311) (2.225,-5.257) (2.275,-5.200) (2.317,-5.150) (2.350,-5.108) (2.394,-5.050) (2.425,-5.008) (2.465,-4.950) (2.498,-4.900) (2.525,-4.857) (2.559,-4.800) (2.587,-4.750) (2.614,-4.700) (2.639,-4.650) (2.663,-4.600) (2.686,-4.550) (2.708,-4.500) (2.728,-4.450) (2.747,-4.400) (2.765,-4.350) (2.781,-4.300) (2.797,-4.250) (2.811,-4.200) (2.825,-4.150) (2.837,-4.100) (2.848,-4.050) (2.858,-4.000) (2.867,-3.950) (2.874,-3.900) (2.881,-3.850) (2.887,-3.800) (2.891,-3.750) (2.895,-3.700) (2.898,-3.650) (2.899,-3.600) (2.900,-3.550) (2.899,-3.500) (2.897,-3.450) (2.895,-3.400) (2.891,-3.350) (2.887,-3.300) (2.881,-3.250) (2.874,-3.200) (2.866,-3.150) (2.857,-3.100) (2.847,-3.050) (2.836,-3.000) (2.824,-2.950) (2.811,-2.900) (2.796,-2.850) (2.781,-2.800) (2.764,-2.750) (2.747,-2.700) (2.728,-2.650) (2.708,-2.600) (2.687,-2.550) (2.664,-2.500) (2.640,-2.450) (2.616,-2.400) (2.589,-2.350) (2.562,-2.300) (2.533,-2.250) (2.503,-2.200) (2.475,-2.155) (2.440,-2.100) (2.406,-2.050) (2.375,-2.007) (2.334,-1.950) (2.300,-1.906) (2.256,-1.850) (2.215,-1.800) (2.175,-1.752) (2.129,-1.700) (2.084,-1.650) (2.050,-1.613) (2.014,-1.575) (1.975,-1.535) (1.925,-1.485) (1.889,-1.450) (1.850,-1.413) (1.811,-1.375) (1.756,-1.325) (1.700,-1.275) (1.650,-1.231) (1.600,-1.188) (1.556,-1.150) (1.500,-1.104) (1.450,-1.063) (1.404,-1.025) (1.350,-0.983) (1.309,-0.950) (1.250,-0.904) (1.200,-0.866) (1.150,-0.828) (1.100,-0.790) (1.050,-0.753) (1.000,-0.716) (0.950,-0.679) (0.900,-0.642) (0.850,-0.606) (0.809,-0.575) (0.750,-0.533) (0.705,-0.500) (0.650,-0.461) (0.600,-0.425) (0.550,-0.390) (0.500,-0.354) (0.450,-0.319) (0.400,-0.284) (0.354,-0.250) (0.300,-0.213) (0.250,-0.177) (0.206,-0.144)}
\EQ{(0.206,0.144) (0.248,0.175) (0.300,0.213) (0.350,0.248) (0.400,0.284) (0.450,0.319) (0.495,0.350) (0.550,0.390) (0.600,0.425) (0.650,0.461) (0.700,0.497) (0.750,0.533) (0.800,0.569) (0.843,0.600) (0.900,0.642) (0.946,0.675) (1.000,0.716) (1.047,0.750) (1.100,0.790) (1.147,0.825) (1.200,0.866) (1.245,0.900) (1.300,0.943) (1.341,0.975) (1.400,1.022) (1.450,1.063) (1.496,1.100) (1.550,1.145) (1.600,1.188) (1.644,1.225) (1.700,1.275) (1.750,1.320) (1.800,1.366) (1.837,1.400) (1.875,1.436) (1.915,1.475) (1.965,1.525) (2.000,1.561) (2.038,1.600) (2.075,1.640) (2.125,1.695) (2.173,1.750) (2.215,1.800) (2.250,1.843) (2.296,1.900) (2.334,1.950) (2.370,2.000) (2.400,2.042) (2.440,2.100) (2.472,2.150) (2.500,2.195) (2.533,2.250) (2.562,2.300) (2.589,2.350) (2.616,2.400) (2.640,2.450) (2.664,2.500) (2.687,2.550) (2.708,2.600) (2.728,2.650) (2.747,2.700) (2.764,2.750) (2.781,2.800) (2.796,2.850) (2.811,2.900) (2.824,2.950) (2.836,3.000) (2.847,3.050) (2.857,3.100) (2.866,3.150) (2.874,3.200) (2.881,3.250) (2.887,3.300) (2.891,3.350) (2.895,3.400) (2.897,3.450) (2.899,3.500) (2.900,3.550) (2.899,3.600) (2.898,3.650) (2.895,3.700) (2.891,3.750) (2.887,3.800) (2.881,3.850) (2.874,3.900) (2.867,3.950) (2.858,4.000) (2.848,4.050) (2.837,4.100) (2.825,4.150) (2.811,4.200) (2.797,4.250) (2.781,4.300) (2.765,4.350) (2.747,4.400) (2.728,4.450) (2.708,4.500) (2.686,4.550) (2.663,4.600) (2.639,4.650) (2.614,4.700) (2.587,4.750) (2.559,4.800) (2.529,4.850) (2.500,4.897) (2.465,4.950) (2.431,5.000) (2.400,5.043) (2.357,5.100) (2.325,5.140) (2.275,5.200) (2.231,5.250) (2.185,5.300) (2.150,5.337) (2.112,5.375) (2.075,5.410) (2.032,5.450) (1.976,5.500) (1.925,5.543) (1.885,5.575) (1.825,5.621) (1.775,5.658) (1.725,5.693) (1.678,5.725) (1.625,5.759) (1.575,5.790) (1.525,5.819) (1.475,5.847) (1.425,5.874) (1.375,5.900) (1.325,5.924) (1.275,5.947) (1.225,5.970) (1.175,5.991) (1.125,6.011) (1.075,6.030) (1.025,6.048) (0.975,6.065) (0.925,6.081) (0.875,6.096) (0.825,6.111) (0.775,6.124) (0.725,6.137) (0.675,6.148) (0.625,6.159) (0.575,6.169) (0.525,6.178) (0.475,6.186) (0.425,6.193) (0.375,6.200) (0.325,6.205) (0.275,6.210) (0.225,6.214) (0.175,6.217) (0.125,6.220) (0.075,6.221) (0.025,6.222) (-0.050,6.222) (-0.100,6.221) (-0.150,6.219) (-0.200,6.216) (-0.250,6.212) (-0.300,6.208) (-0.350,6.203) (-0.400,6.196) (-0.450,6.190) (-0.500,6.182) (-0.550,6.173) (-0.600,6.164) (-0.650,6.154) (-0.700,6.142) (-0.750,6.130) (-0.800,6.117) (-0.850,6.104) (-0.900,6.089) (-0.950,6.073) (-1.000,6.057) (-1.050,6.039) (-1.100,6.020) (-1.150,6.001) (-1.200,5.980) (-1.250,5.959) (-1.300,5.936) (-1.350,5.912) (-1.400,5.887) (-1.450,5.861) (-1.500,5.833) (-1.550,5.805) (-1.599,5.775) (-1.650,5.743) (-1.700,5.710) (-1.750,5.676) (-1.800,5.640) (-1.850,5.602) (-1.900,5.563) (-1.946,5.525) (-2.000,5.479) (-2.050,5.434) (-2.100,5.386) (-2.137,5.350) (-2.175,5.311) (-2.225,5.257) (-2.275,5.200) (-2.317,5.150) (-2.350,5.108) (-2.394,5.050) (-2.425,5.008) (-2.465,4.950) (-2.498,4.900) (-2.525,4.857) (-2.559,4.800) (-2.587,4.750) (-2.614,4.700) (-2.639,4.650) (-2.663,4.600) (-2.686,4.550) (-2.708,4.500) (-2.728,4.450) (-2.747,4.400) (-2.765,4.350) (-2.781,4.300) (-2.797,4.250) (-2.811,4.200) (-2.825,4.150) (-2.837,4.100) (-2.848,4.050) (-2.858,4.000) (-2.867,3.950) (-2.874,3.900) (-2.881,3.850) (-2.887,3.800) (-2.891,3.750) (-2.895,3.700) (-2.898,3.650) (-2.899,3.600) (-2.900,3.550) (-2.899,3.500) (-2.897,3.450) (-2.895,3.400) (-2.891,3.350) (-2.887,3.300) (-2.881,3.250) (-2.874,3.200) (-2.866,3.150) (-2.857,3.100) (-2.847,3.050) (-2.836,3.000) (-2.824,2.950) (-2.811,2.900) (-2.796,2.850) (-2.781,2.800) (-2.764,2.750) (-2.747,2.700) (-2.728,2.650) (-2.708,2.600) (-2.687,2.550) (-2.664,2.500) (-2.640,2.450) (-2.616,2.400) (-2.589,2.350) (-2.562,2.300) (-2.533,2.250) (-2.503,2.200) (-2.475,2.155) (-2.440,2.100) (-2.406,2.050) (-2.375,2.007) (-2.334,1.950) (-2.300,1.906) (-2.256,1.850) (-2.215,1.800) (-2.175,1.752) (-2.129,1.700) (-2.084,1.650) (-2.050,1.613) (-2.014,1.575) (-1.975,1.535) (-1.925,1.485) (-1.889,1.450) (-1.850,1.413) (-1.811,1.375) (-1.756,1.325) (-1.700,1.275) (-1.650,1.231) (-1.600,1.188) (-1.556,1.150) (-1.500,1.104) (-1.450,1.063) (-1.404,1.025) (-1.350,0.983) (-1.309,0.950) (-1.250,0.904) (-1.200,0.866) (-1.150,0.828) (-1.100,0.790) (-1.050,0.753) (-1.000,0.716) (-0.950,0.679) (-0.900,0.642) (-0.850,0.606) (-0.809,0.575) (-0.750,0.533) (-0.705,0.500) (-0.650,0.461) (-0.600,0.425) (-0.550,0.390) (-0.500,0.354) (-0.450,0.319) (-0.400,0.284) (-0.354,0.250) (-0.300,0.213) (-0.250,0.177) (-0.220,0.155)}
\EQ{(-0.220,-0.155) (-0.275,-0.195) (-0.319,-0.225) (-0.375,-0.266) (-0.424,-0.300) (-0.475,-0.337) (-0.525,-0.372) (-0.575,-0.408) (-0.625,-0.444) (-0.670,-0.475) (-0.725,-0.516) (-0.773,-0.550) (-0.825,-0.588) (-0.875,-0.624) (-0.925,-0.661) (-0.975,-0.698) (-1.025,-0.735) (-1.075,-0.773) (-1.125,-0.811) (-1.175,-0.849) (-1.225,-0.888) (-1.273,-0.925) (-1.325,-0.967) (-1.367,-1.000) (-1.425,-1.048) (-1.475,-1.089) (-1.518,-1.125) (-1.575,-1.174) (-1.625,-1.218) (-1.675,-1.263) (-1.715,-1.300) (-1.769,-1.350) (-1.821,-1.400) (-1.871,-1.450) (-1.920,-1.500) (-1.967,-1.550) (-2.013,-1.600) (-2.050,-1.641) (-2.100,-1.699) (-2.142,-1.750) (-2.175,-1.791) (-2.221,-1.850) (-2.258,-1.900) (-2.293,-1.950) (-2.325,-1.996) (-2.361,-2.050) (-2.392,-2.100) (-2.422,-2.150) (-2.450,-2.199) (-2.475,-2.245) (-2.500,-2.292) (-2.525,-2.343) (-2.550,-2.396) (-2.574,-2.450) (-2.595,-2.500) (-2.615,-2.550) (-2.633,-2.600) (-2.650,-2.648) (-2.667,-2.700) (-2.682,-2.750) (-2.696,-2.800) (-2.708,-2.850) (-2.720,-2.900) (-2.730,-2.950) (-2.739,-3.000) (-2.747,-3.050) (-2.754,-3.100) (-2.760,-3.150) (-2.765,-3.200) (-2.769,-3.250) (-2.771,-3.300) (-2.773,-3.350) (-2.773,-3.400) (-2.773,-3.450) (-2.771,-3.500) (-2.768,-3.550) (-2.764,-3.600) (-2.759,-3.650) (-2.753,-3.700) (-2.746,-3.750) (-2.738,-3.800) (-2.729,-3.850) (-2.718,-3.900) (-2.706,-3.950) (-2.694,-4.000) (-2.680,-4.050) (-2.664,-4.100) (-2.648,-4.150) (-2.630,-4.200) (-2.612,-4.250) (-2.592,-4.300) (-2.570,-4.350) (-2.547,-4.400) (-2.525,-4.446) (-2.500,-4.496) (-2.475,-4.542) (-2.450,-4.587) (-2.412,-4.650) (-2.381,-4.700) (-2.350,-4.746) (-2.312,-4.800) (-2.275,-4.850) (-2.237,-4.900) (-2.200,-4.945) (-2.153,-5.000) (-2.108,-5.050) (-2.061,-5.100) (-2.025,-5.136) (-1.984,-5.175) (-1.930,-5.225) (-1.875,-5.273) (-1.825,-5.315) (-1.781,-5.350) (-1.725,-5.392) (-1.679,-5.425) (-1.625,-5.462) (-1.575,-5.495) (-1.527,-5.525) (-1.475,-5.556) (-1.425,-5.585) (-1.375,-5.612) (-1.325,-5.638) (-1.275,-5.662) (-1.225,-5.686) (-1.175,-5.708) (-1.125,-5.729) (-1.075,-5.749) (-1.025,-5.769) (-0.975,-5.787) (-0.925,-5.804) (-0.875,-5.819) (-0.825,-5.834) (-0.775,-5.848) (-0.725,-5.862) (-0.675,-5.874) (-0.625,-5.885) (-0.575,-5.895) (-0.525,-5.905) (-0.475,-5.913) (-0.425,-5.921) (-0.375,-5.928) (-0.325,-5.934) (-0.275,-5.939) (-0.225,-5.943) (-0.175,-5.946) (-0.125,-5.949) (-0.075,-5.950) (-0.025,-5.951) (0.050,-5.951) (0.100,-5.950) (0.150,-5.948) (0.200,-5.945) (0.250,-5.941) (0.300,-5.936) (0.350,-5.931) (0.400,-5.924) (0.450,-5.917) (0.500,-5.909) (0.550,-5.900) (0.600,-5.890) (0.650,-5.879) (0.700,-5.868) (0.750,-5.855) (0.800,-5.842) (0.850,-5.827) (0.900,-5.812) (0.950,-5.795) (1.000,-5.778) (1.050,-5.759) (1.100,-5.740) (1.150,-5.719) (1.200,-5.697) (1.249,-5.675) (1.300,-5.650) (1.350,-5.625) (1.397,-5.600) (1.442,-5.575) (1.486,-5.550) (1.550,-5.511) (1.600,-5.479) (1.643,-5.450) (1.700,-5.410) (1.748,-5.375) (1.800,-5.335) (1.843,-5.300) (1.900,-5.252) (1.950,-5.207) (2.000,-5.160) (2.036,-5.125) (2.075,-5.085) (2.125,-5.031) (2.175,-4.975) (2.217,-4.925) (2.250,-4.883) (2.294,-4.825) (2.325,-4.782) (2.364,-4.725) (2.397,-4.675) (2.425,-4.629) (2.457,-4.575) (2.484,-4.525) (2.511,-4.475) (2.535,-4.425) (2.559,-4.375) (2.581,-4.325) (2.602,-4.275) (2.621,-4.225) (2.639,-4.175) (2.656,-4.125) (2.672,-4.075) (2.687,-4.025) (2.700,-3.975) (2.712,-3.925) (2.723,-3.875) (2.733,-3.825) (2.742,-3.775) (2.750,-3.725) (2.757,-3.675) (2.762,-3.625) (2.766,-3.575) (2.770,-3.525) (2.772,-3.475) (2.773,-3.425) (2.773,-3.375) (2.772,-3.325) (2.770,-3.275) (2.767,-3.225) (2.763,-3.175) (2.757,-3.125) (2.751,-3.075) (2.743,-3.025) (2.735,-2.975) (2.725,-2.926) (2.714,-2.875) (2.702,-2.825) (2.689,-2.775) (2.674,-2.725) (2.659,-2.675) (2.642,-2.625) (2.625,-2.577) (2.605,-2.525) (2.585,-2.475) (2.563,-2.425) (2.541,-2.375) (2.517,-2.325) (2.491,-2.275) (2.464,-2.225) (2.436,-2.175) (2.407,-2.125) (2.376,-2.075) (2.344,-2.025) (2.311,-1.975) (2.276,-1.925) (2.240,-1.875) (2.202,-1.825) (2.162,-1.775) (2.125,-1.729) (2.079,-1.675) (2.036,-1.625) (2.000,-1.586) (1.950,-1.532) (1.900,-1.480) (1.850,-1.429) (1.800,-1.380) (1.750,-1.333) (1.700,-1.286) (1.661,-1.250) (1.604,-1.200) (1.550,-1.153) (1.500,-1.110) (1.458,-1.075) (1.400,-1.027) (1.350,-0.987) (1.304,-0.950) (1.250,-0.907) (1.209,-0.875) (1.150,-0.830) (1.100,-0.792) (1.050,-0.754) (1.000,-0.717) (0.950,-0.680) (0.900,-0.643) (0.850,-0.606) (0.808,-0.575) (0.750,-0.534) (0.704,-0.500) (0.650,-0.462) (0.600,-0.425) (0.550,-0.390) (0.500,-0.354) (0.450,-0.319) (0.400,-0.284) (0.354,-0.250) (0.300,-0.213) (0.250,-0.177) (0.206,-0.144)}
\EQ{(0.206,0.144) (0.248,0.175) (0.300,0.213) (0.350,0.248) (0.400,0.284) (0.450,0.319) (0.495,0.350) (0.550,0.390) (0.599,0.425) (0.650,0.462) (0.700,0.497) (0.750,0.534) (0.800,0.570) (0.842,0.600) (0.900,0.643) (0.944,0.675) (1.000,0.717) (1.045,0.750) (1.100,0.792) (1.144,0.825) (1.200,0.868) (1.241,0.900) (1.300,0.947) (1.350,0.987) (1.398,1.025) (1.450,1.068) (1.500,1.110) (1.547,1.150) (1.600,1.196) (1.650,1.241) (1.688,1.275) (1.725,1.309) (1.769,1.350) (1.821,1.400) (1.871,1.450) (1.920,1.500) (1.967,1.550) (2.013,1.600) (2.050,1.641) (2.100,1.699) (2.142,1.750) (2.175,1.791) (2.221,1.850) (2.258,1.900) (2.293,1.950) (2.325,1.996) (2.361,2.050) (2.392,2.100) (2.422,2.150) (2.450,2.199) (2.475,2.245) (2.500,2.292) (2.525,2.343) (2.550,2.396) (2.574,2.450) (2.595,2.500) (2.615,2.550) (2.633,2.600) (2.650,2.648) (2.667,2.700) (2.682,2.750) (2.696,2.800) (2.708,2.850) (2.720,2.900) (2.730,2.950) (2.739,3.000) (2.747,3.050) (2.754,3.100) (2.760,3.150) (2.765,3.200) (2.769,3.250) (2.771,3.300) (2.773,3.350) (2.773,3.400) (2.773,3.450) (2.771,3.500) (2.768,3.550) (2.764,3.600) (2.759,3.650) (2.753,3.700) (2.746,3.750) (2.738,3.800) (2.729,3.850) (2.718,3.900) (2.706,3.950) (2.694,4.000) (2.680,4.050) (2.664,4.100) (2.648,4.150) (2.630,4.200) (2.612,4.250) (2.592,4.300) (2.570,4.350) (2.547,4.400) (2.525,4.446) (2.500,4.496) (2.475,4.542) (2.450,4.587) (2.412,4.650) (2.381,4.700) (2.350,4.746) (2.312,4.800) (2.275,4.850) (2.237,4.900) (2.200,4.945) (2.153,5.000) (2.108,5.050) (2.061,5.100) (2.025,5.136) (1.984,5.175) (1.930,5.225) (1.875,5.273) (1.825,5.315) (1.781,5.350) (1.725,5.392) (1.679,5.425) (1.625,5.462) (1.575,5.495) (1.527,5.525) (1.475,5.556) (1.425,5.585) (1.375,5.612) (1.325,5.638) (1.275,5.662) (1.225,5.686) (1.175,5.708) (1.125,5.729) (1.075,5.749) (1.025,5.769) (0.975,5.787) (0.925,5.804) (0.875,5.819) (0.825,5.834) (0.775,5.848) (0.725,5.862) (0.675,5.874) (0.625,5.885) (0.575,5.895) (0.525,5.905) (0.475,5.913) (0.425,5.921) (0.375,5.928) (0.325,5.934) (0.275,5.939) (0.225,5.943) (0.175,5.946) (0.125,5.949) (0.075,5.950) (0.025,5.951) (-0.050,5.951) (-0.100,5.950) (-0.150,5.948) (-0.200,5.945) (-0.250,5.941) (-0.300,5.936) (-0.350,5.931) (-0.400,5.924) (-0.450,5.917) (-0.500,5.909) (-0.550,5.900) (-0.600,5.890) (-0.650,5.879) (-0.700,5.868) (-0.750,5.855) (-0.800,5.842) (-0.850,5.827) (-0.900,5.812) (-0.950,5.795) (-1.000,5.778) (-1.050,5.759) (-1.100,5.740) (-1.150,5.719) (-1.200,5.697) (-1.249,5.675) (-1.300,5.650) (-1.350,5.625) (-1.397,5.600) (-1.442,5.575) (-1.486,5.550) (-1.550,5.511) (-1.600,5.479) (-1.643,5.450) (-1.700,5.410) (-1.748,5.375) (-1.800,5.335) (-1.843,5.300) (-1.900,5.252) (-1.950,5.207) (-2.000,5.160) (-2.036,5.125) (-2.075,5.085) (-2.125,5.031) (-2.175,4.975) (-2.217,4.925) (-2.250,4.883) (-2.294,4.825) (-2.325,4.782) (-2.364,4.725) (-2.397,4.675) (-2.425,4.629) (-2.457,4.575) (-2.484,4.525) (-2.511,4.475) (-2.535,4.425) (-2.559,4.375) (-2.581,4.325) (-2.602,4.275) (-2.621,4.225) (-2.639,4.175) (-2.656,4.125) (-2.672,4.075) (-2.687,4.025) (-2.700,3.975) (-2.712,3.925) (-2.723,3.875) (-2.733,3.825) (-2.742,3.775) (-2.750,3.725) (-2.757,3.675) (-2.762,3.625) (-2.766,3.575) (-2.770,3.525) (-2.772,3.475) (-2.773,3.425) (-2.773,3.375) (-2.772,3.325) (-2.770,3.275) (-2.767,3.225) (-2.763,3.175) (-2.757,3.125) (-2.751,3.075) (-2.743,3.025) (-2.735,2.975) (-2.725,2.926) (-2.714,2.875) (-2.702,2.825) (-2.689,2.775) (-2.674,2.725) (-2.659,2.675) (-2.642,2.625) (-2.625,2.577) (-2.605,2.525) (-2.585,2.475) (-2.563,2.425) (-2.541,2.375) (-2.517,2.325) (-2.491,2.275) (-2.464,2.225) (-2.436,2.175) (-2.407,2.125) (-2.376,2.075) (-2.344,2.025) (-2.311,1.975) (-2.276,1.925) (-2.240,1.875) (-2.202,1.825) (-2.162,1.775) (-2.125,1.729) (-2.079,1.675) (-2.036,1.625) (-2.000,1.586) (-1.950,1.532) (-1.900,1.480) (-1.850,1.429) (-1.800,1.380) (-1.750,1.333) (-1.700,1.286) (-1.661,1.250) (-1.604,1.200) (-1.550,1.153) (-1.500,1.110) (-1.458,1.075) (-1.400,1.027) (-1.350,0.987) (-1.304,0.950) (-1.250,0.907) (-1.209,0.875) (-1.150,0.830) (-1.100,0.792) (-1.050,0.754) (-1.000,0.717) (-0.950,0.680) (-0.900,0.643) (-0.850,0.606) (-0.808,0.575) (-0.750,0.534) (-0.704,0.500) (-0.650,0.462) (-0.600,0.425) (-0.550,0.390) (-0.500,0.354) (-0.450,0.319) (-0.400,0.284) (-0.354,0.250) (-0.300,0.213) (-0.250,0.177) (-0.220,0.155)}
\EQ{(-0.220,-0.155) (-0.275,-0.195) (-0.319,-0.225) (-0.375,-0.266) (-0.424,-0.300) (-0.475,-0.337) (-0.525,-0.372) (-0.575,-0.408) (-0.625,-0.444) (-0.669,-0.475) (-0.725,-0.516) (-0.773,-0.550) (-0.825,-0.589) (-0.875,-0.625) (-0.925,-0.662) (-0.975,-0.699) (-1.025,-0.737) (-1.075,-0.774) (-1.125,-0.813) (-1.173,-0.850) (-1.225,-0.891) (-1.269,-0.925) (-1.325,-0.970) (-1.375,-1.011) (-1.422,-1.050) (-1.475,-1.095) (-1.525,-1.139) (-1.567,-1.175) (-1.622,-1.225) (-1.675,-1.274) (-1.725,-1.322) (-1.775,-1.370) (-1.825,-1.421) (-1.875,-1.473) (-1.924,-1.525) (-1.969,-1.575) (-2.012,-1.625) (-2.050,-1.671) (-2.094,-1.725) (-2.125,-1.766) (-2.170,-1.825) (-2.200,-1.868) (-2.239,-1.925) (-2.272,-1.975) (-2.300,-2.020) (-2.333,-2.075) (-2.362,-2.125) (-2.389,-2.175) (-2.414,-2.225) (-2.439,-2.275) (-2.462,-2.325) (-2.483,-2.375) (-2.504,-2.425) (-2.523,-2.475) (-2.541,-2.525) (-2.558,-2.575) (-2.573,-2.625) (-2.587,-2.675) (-2.600,-2.725) (-2.612,-2.775) (-2.623,-2.825) (-2.632,-2.875) (-2.640,-2.925) (-2.647,-2.975) (-2.654,-3.025) (-2.658,-3.075) (-2.662,-3.125) (-2.665,-3.175) (-2.666,-3.225) (-2.667,-3.275) (-2.666,-3.325) (-2.664,-3.375) (-2.661,-3.425) (-2.657,-3.475) (-2.651,-3.525) (-2.645,-3.575) (-2.637,-3.625) (-2.629,-3.675) (-2.619,-3.725) (-2.608,-3.775) (-2.595,-3.825) (-2.582,-3.875) (-2.567,-3.925) (-2.551,-3.975) (-2.534,-4.025) (-2.515,-4.075) (-2.496,-4.125) (-2.475,-4.174) (-2.452,-4.225) (-2.428,-4.275) (-2.402,-4.325) (-2.375,-4.375) (-2.350,-4.420) (-2.317,-4.475) (-2.285,-4.525) (-2.251,-4.575) (-2.216,-4.625) (-2.179,-4.675) (-2.140,-4.725) (-2.100,-4.773) (-2.055,-4.825) (-2.009,-4.875) (-1.960,-4.925) (-1.909,-4.975) (-1.855,-5.025) (-1.800,-5.073) (-1.750,-5.114) (-1.704,-5.150) (-1.650,-5.191) (-1.602,-5.225) (-1.550,-5.260) (-1.500,-5.293) (-1.450,-5.323) (-1.405,-5.350) (-1.361,-5.375) (-1.314,-5.400) (-1.265,-5.425) (-1.214,-5.450) (-1.160,-5.475) (-1.103,-5.500) (-1.050,-5.522) (-1.000,-5.541) (-0.950,-5.559) (-0.900,-5.577) (-0.850,-5.593) (-0.800,-5.608) (-0.750,-5.622) (-0.700,-5.635) (-0.650,-5.647) (-0.600,-5.659) (-0.550,-5.669) (-0.500,-5.678) (-0.450,-5.687) (-0.400,-5.694) (-0.350,-5.701) (-0.300,-5.707) (-0.250,-5.711) (-0.200,-5.715) (-0.150,-5.718) (-0.100,-5.721) (-0.050,-5.722) (0.000,-5.722) (0.050,-5.722) (0.100,-5.721) (0.150,-5.718) (0.200,-5.715) (0.250,-5.711) (0.300,-5.707) (0.350,-5.701) (0.400,-5.694) (0.450,-5.687) (0.500,-5.678) (0.550,-5.669) (0.600,-5.659) (0.650,-5.647) (0.700,-5.635) (0.750,-5.622) (0.800,-5.608) (0.850,-5.593) (0.900,-5.577) (0.950,-5.559) (1.000,-5.541) (1.050,-5.522) (1.100,-5.501) (1.150,-5.480) (1.200,-5.457) (1.250,-5.433) (1.300,-5.407) (1.350,-5.381) (1.400,-5.353) (1.447,-5.325) (1.500,-5.293) (1.550,-5.260) (1.600,-5.226) (1.650,-5.191) (1.700,-5.153) (1.750,-5.114) (1.797,-5.075) (1.850,-5.029) (1.900,-4.983) (1.950,-4.935) (2.000,-4.884) (2.050,-4.830) (2.098,-4.775) (2.140,-4.725) (2.175,-4.680) (2.216,-4.625) (2.250,-4.577) (2.285,-4.525) (2.317,-4.475) (2.347,-4.425) (2.375,-4.376) (2.400,-4.330) (2.425,-4.281) (2.450,-4.229) (2.474,-4.175) (2.496,-4.125) (2.515,-4.075) (2.534,-4.025) (2.551,-3.975) (2.567,-3.925) (2.582,-3.875) (2.595,-3.825) (2.608,-3.775) (2.619,-3.725) (2.629,-3.675) (2.637,-3.625) (2.645,-3.575) (2.651,-3.525) (2.657,-3.475) (2.661,-3.425) (2.664,-3.375) (2.666,-3.325) (2.667,-3.275) (2.666,-3.225) (2.665,-3.175) (2.662,-3.125) (2.658,-3.075) (2.654,-3.025) (2.647,-2.975) (2.640,-2.925) (2.632,-2.875) (2.623,-2.825) (2.612,-2.775) (2.600,-2.725) (2.587,-2.675) (2.573,-2.625) (2.558,-2.575) (2.541,-2.525) (2.523,-2.475) (2.504,-2.425) (2.483,-2.375) (2.462,-2.325) (2.439,-2.275) (2.414,-2.225) (2.389,-2.175) (2.362,-2.125) (2.333,-2.075) (2.303,-2.025) (2.275,-1.980) (2.239,-1.925) (2.205,-1.875) (2.175,-1.833) (2.132,-1.775) (2.100,-1.733) (2.053,-1.675) (2.012,-1.625) (1.975,-1.583) (1.925,-1.527) (1.877,-1.475) (1.829,-1.425) (1.780,-1.375) (1.729,-1.325) (1.676,-1.275) (1.625,-1.228) (1.575,-1.183) (1.525,-1.139) (1.481,-1.100) (1.425,-1.053) (1.375,-1.011) (1.331,-0.975) (1.275,-0.930) (1.225,-0.891) (1.175,-0.851) (1.125,-0.813) (1.076,-0.775) (1.025,-0.737) (0.976,-0.700) (0.925,-0.662) (0.875,-0.625) (0.825,-0.589) (0.775,-0.552) (0.725,-0.516) (0.675,-0.480) (0.635,-0.450) (0.575,-0.408) (0.530,-0.375) (0.475,-0.337) (0.425,-0.301) (0.375,-0.266) (0.325,-0.230) (0.284,-0.200) (0.225,-0.160) (0.206,-0.144)}
\EQ{(0.206,0.144) (0.248,0.175) (0.300,0.213) (0.350,0.248) (0.400,0.284) (0.450,0.319) (0.495,0.350) (0.550,0.390) (0.599,0.425) (0.650,0.462) (0.700,0.498) (0.750,0.534) (0.800,0.570) (0.842,0.600) (0.900,0.644) (0.943,0.675) (1.000,0.718) (1.043,0.750) (1.100,0.794) (1.141,0.825) (1.200,0.871) (1.250,0.910) (1.300,0.950) (1.350,0.991) (1.392,1.025) (1.450,1.074) (1.500,1.117) (1.538,1.150) (1.595,1.200) (1.649,1.250) (1.700,1.298) (1.750,1.346) (1.800,1.395) (1.850,1.447) (1.900,1.499) (1.946,1.550) (1.990,1.600) (2.025,1.641) (2.074,1.700) (2.113,1.750) (2.150,1.799) (2.188,1.850) (2.222,1.900) (2.256,1.950) (2.288,2.000) (2.318,2.050) (2.348,2.100) (2.375,2.149) (2.400,2.197) (2.425,2.247) (2.450,2.299) (2.473,2.350) (2.494,2.400) (2.514,2.450) (2.532,2.500) (2.549,2.550) (2.565,2.600) (2.580,2.650) (2.594,2.700) (2.606,2.750) (2.617,2.800) (2.627,2.850) (2.636,2.900) (2.644,2.950) (2.651,3.000) (2.656,3.050) (2.660,3.100) (2.664,3.150) (2.666,3.200) (2.667,3.250) (2.666,3.300) (2.665,3.350) (2.663,3.400) (2.659,3.450) (2.654,3.500) (2.648,3.550) (2.641,3.600) (2.633,3.650) (2.624,3.700) (2.613,3.750) (2.602,3.800) (2.589,3.850) (2.575,3.899) (2.559,3.950) (2.543,4.000) (2.525,4.050) (2.506,4.100) (2.485,4.150) (2.463,4.200) (2.440,4.250) (2.415,4.300) (2.389,4.350) (2.361,4.400) (2.332,4.450) (2.301,4.500) (2.268,4.550) (2.234,4.600) (2.200,4.647) (2.160,4.700) (2.125,4.743) (2.077,4.800) (2.032,4.850) (1.985,4.900) (1.935,4.950) (1.882,5.000) (1.826,5.050) (1.775,5.094) (1.725,5.134) (1.675,5.172) (1.625,5.209) (1.575,5.244) (1.528,5.275) (1.475,5.308) (1.425,5.338) (1.375,5.367) (1.325,5.394) (1.275,5.420) (1.225,5.445) (1.175,5.468) (1.125,5.491) (1.075,5.512) (1.025,5.532) (0.976,5.550) (0.925,5.568) (0.875,5.585) (0.826,5.600) (0.775,5.615) (0.725,5.629) (0.675,5.641) (0.625,5.653) (0.575,5.664) (0.525,5.674) (0.475,5.683) (0.425,5.691) (0.375,5.698) (0.325,5.704) (0.275,5.709) (0.225,5.713) (0.175,5.717) (0.125,5.720) (0.075,5.721) (0.025,5.722) (-0.050,5.722) (-0.100,5.721) (-0.150,5.718) (-0.200,5.715) (-0.250,5.711) (-0.300,5.707) (-0.350,5.701) (-0.400,5.694) (-0.450,5.687) (-0.500,5.678) (-0.550,5.669) (-0.600,5.659) (-0.650,5.647) (-0.700,5.635) (-0.750,5.622) (-0.800,5.608) (-0.850,5.593) (-0.900,5.577) (-0.950,5.559) (-1.000,5.541) (-1.050,5.522) (-1.100,5.501) (-1.150,5.480) (-1.200,5.457) (-1.250,5.433) (-1.300,5.407) (-1.350,5.381) (-1.400,5.353) (-1.447,5.325) (-1.500,5.293) (-1.550,5.260) (-1.600,5.226) (-1.650,5.191) (-1.700,5.153) (-1.750,5.114) (-1.797,5.075) (-1.850,5.029) (-1.900,4.983) (-1.950,4.935) (-2.000,4.884) (-2.050,4.830) (-2.098,4.775) (-2.140,4.725) (-2.175,4.680) (-2.216,4.625) (-2.250,4.577) (-2.285,4.525) (-2.317,4.475) (-2.347,4.425) (-2.375,4.376) (-2.400,4.330) (-2.425,4.281) (-2.450,4.229) (-2.474,4.175) (-2.496,4.125) (-2.515,4.075) (-2.534,4.025) (-2.551,3.975) (-2.567,3.925) (-2.582,3.875) (-2.595,3.825) (-2.608,3.775) (-2.619,3.725) (-2.629,3.675) (-2.637,3.625) (-2.645,3.575) (-2.651,3.525) (-2.657,3.475) (-2.661,3.425) (-2.664,3.375) (-2.666,3.325) (-2.667,3.275) (-2.666,3.225) (-2.665,3.175) (-2.662,3.125) (-2.658,3.075) (-2.654,3.025) (-2.647,2.975) (-2.640,2.925) (-2.632,2.875) (-2.623,2.825) (-2.612,2.775) (-2.600,2.725) (-2.587,2.675) (-2.573,2.625) (-2.558,2.575) (-2.541,2.525) (-2.523,2.475) (-2.504,2.425) (-2.483,2.375) (-2.462,2.325) (-2.439,2.275) (-2.414,2.225) (-2.389,2.175) (-2.362,2.125) (-2.333,2.075) (-2.303,2.025) (-2.275,1.980) (-2.239,1.925) (-2.205,1.875) (-2.175,1.833) (-2.132,1.775) (-2.100,1.733) (-2.053,1.675) (-2.012,1.625) (-1.975,1.583) (-1.925,1.527) (-1.877,1.475) (-1.829,1.425) (-1.780,1.375) (-1.729,1.325) (-1.676,1.275) (-1.625,1.228) (-1.575,1.183) (-1.525,1.139) (-1.481,1.100) (-1.425,1.053) (-1.375,1.011) (-1.331,0.975) (-1.275,0.930) (-1.225,0.891) (-1.175,0.851) (-1.125,0.813) (-1.076,0.775) (-1.025,0.737) (-0.976,0.700) (-0.925,0.662) (-0.875,0.625) (-0.825,0.589) (-0.775,0.552) (-0.725,0.516) (-0.675,0.480) (-0.635,0.450) (-0.575,0.408) (-0.530,0.375) (-0.475,0.337) (-0.425,0.301) (-0.375,0.266) (-0.325,0.230) (-0.284,0.200) (-0.225,0.160) (-0.220,0.155)}
\EQ{(-0.220,-0.155) (-0.275,-0.195) (-0.319,-0.225) (-0.375,-0.266) (-0.424,-0.300) (-0.475,-0.337) (-0.525,-0.372) (-0.575,-0.408) (-0.625,-0.444) (-0.669,-0.475) (-0.725,-0.516) (-0.772,-0.550) (-0.825,-0.589) (-0.874,-0.625) (-0.925,-0.663) (-0.975,-0.700) (-1.025,-0.738) (-1.074,-0.775) (-1.125,-0.815) (-1.171,-0.850) (-1.225,-0.893) (-1.265,-0.925) (-1.325,-0.974) (-1.375,-1.016) (-1.416,-1.050) (-1.473,-1.100) (-1.525,-1.146) (-1.575,-1.191) (-1.612,-1.225) (-1.650,-1.261) (-1.691,-1.300) (-1.741,-1.350) (-1.790,-1.400) (-1.825,-1.438) (-1.860,-1.475) (-1.900,-1.520) (-1.947,-1.575) (-1.989,-1.625) (-2.025,-1.670) (-2.067,-1.725) (-2.100,-1.770) (-2.139,-1.825) (-2.173,-1.875) (-2.205,-1.925) (-2.236,-1.975) (-2.265,-2.025) (-2.293,-2.075) (-2.319,-2.125) (-2.344,-2.175) (-2.368,-2.225) (-2.390,-2.275) (-2.411,-2.325) (-2.431,-2.375) (-2.449,-2.425) (-2.466,-2.475) (-2.481,-2.525) (-2.496,-2.575) (-2.509,-2.625) (-2.521,-2.675) (-2.532,-2.725) (-2.541,-2.775) (-2.549,-2.825) (-2.557,-2.875) (-2.562,-2.925) (-2.567,-2.975) (-2.571,-3.025) (-2.573,-3.075) (-2.574,-3.125) (-2.574,-3.175) (-2.573,-3.225) (-2.571,-3.275) (-2.568,-3.325) (-2.563,-3.375) (-2.557,-3.425) (-2.550,-3.475) (-2.542,-3.525) (-2.533,-3.575) (-2.522,-3.625) (-2.510,-3.675) (-2.497,-3.725) (-2.483,-3.775) (-2.467,-3.825) (-2.450,-3.875) (-2.432,-3.925) (-2.413,-3.975) (-2.392,-4.025) (-2.369,-4.075) (-2.345,-4.125) (-2.320,-4.175) (-2.293,-4.225) (-2.264,-4.275) (-2.234,-4.325) (-2.202,-4.375) (-2.168,-4.425) (-2.133,-4.475) (-2.100,-4.518) (-2.055,-4.575) (-2.013,-4.625) (-1.975,-4.668) (-1.925,-4.722) (-1.875,-4.772) (-1.825,-4.820) (-1.775,-4.866) (-1.735,-4.900) (-1.675,-4.949) (-1.625,-4.988) (-1.575,-5.025) (-1.525,-5.060) (-1.475,-5.093) (-1.425,-5.124) (-1.375,-5.155) (-1.325,-5.183) (-1.275,-5.210) (-1.225,-5.236) (-1.175,-5.261) (-1.125,-5.284) (-1.075,-5.306) (-1.025,-5.327) (-0.975,-5.346) (-0.925,-5.365) (-0.875,-5.382) (-0.825,-5.398) (-0.775,-5.414) (-0.725,-5.428) (-0.675,-5.441) (-0.625,-5.453) (-0.575,-5.464) (-0.525,-5.474) (-0.475,-5.484) (-0.425,-5.492) (-0.375,-5.499) (-0.325,-5.506) (-0.275,-5.511) (-0.225,-5.516) (-0.175,-5.519) (-0.125,-5.522) (-0.075,-5.524) (-0.025,-5.525) (0.050,-5.524) (0.100,-5.523) (0.150,-5.521) (0.200,-5.518) (0.250,-5.513) (0.300,-5.508) (0.350,-5.503) (0.400,-5.496) (0.450,-5.488) (0.500,-5.479) (0.550,-5.469) (0.600,-5.459) (0.650,-5.447) (0.700,-5.434) (0.750,-5.421) (0.800,-5.406) (0.850,-5.390) (0.900,-5.374) (0.950,-5.356) (1.000,-5.337) (1.050,-5.316) (1.100,-5.295) (1.150,-5.272) (1.197,-5.250) (1.247,-5.225) (1.294,-5.200) (1.340,-5.175) (1.400,-5.140) (1.450,-5.109) (1.500,-5.076) (1.550,-5.042) (1.600,-5.006) (1.642,-4.975) (1.700,-4.929) (1.750,-4.887) (1.792,-4.850) (1.846,-4.800) (1.898,-4.750) (1.946,-4.700) (1.991,-4.650) (2.025,-4.611) (2.075,-4.550) (2.114,-4.500) (2.150,-4.451) (2.185,-4.400) (2.218,-4.350) (2.249,-4.300) (2.275,-4.257) (2.300,-4.212) (2.325,-4.165) (2.350,-4.115) (2.375,-4.062) (2.400,-4.005) (2.423,-3.950) (2.442,-3.900) (2.459,-3.850) (2.475,-3.801) (2.490,-3.750) (2.504,-3.700) (2.516,-3.650) (2.528,-3.600) (2.538,-3.550) (2.546,-3.500) (2.554,-3.450) (2.560,-3.400) (2.566,-3.350) (2.570,-3.300) (2.572,-3.250) (2.574,-3.200) (2.575,-3.150) (2.574,-3.100) (2.572,-3.050) (2.569,-3.000) (2.565,-2.950) (2.560,-2.900) (2.553,-2.850) (2.545,-2.800) (2.537,-2.750) (2.526,-2.700) (2.515,-2.650) (2.502,-2.600) (2.489,-2.550) (2.474,-2.500) (2.457,-2.450) (2.440,-2.400) (2.421,-2.350) (2.401,-2.300) (2.379,-2.250) (2.356,-2.200) (2.332,-2.150) (2.306,-2.100) (2.279,-2.050) (2.251,-2.000) (2.221,-1.950) (2.189,-1.900) (2.156,-1.850) (2.125,-1.805) (2.086,-1.750) (2.050,-1.702) (2.009,-1.650) (1.975,-1.608) (1.926,-1.550) (1.882,-1.500) (1.837,-1.450) (1.800,-1.411) (1.765,-1.375) (1.725,-1.334) (1.675,-1.285) (1.639,-1.250) (1.600,-1.214) (1.558,-1.175) (1.502,-1.125) (1.450,-1.080) (1.400,-1.037) (1.356,-1.000) (1.300,-0.954) (1.250,-0.913) (1.202,-0.875) (1.150,-0.834) (1.107,-0.800) (1.050,-0.757) (1.008,-0.725) (0.950,-0.682) (0.908,-0.650) (0.850,-0.607) (0.807,-0.575) (0.750,-0.534) (0.703,-0.500) (0.650,-0.462) (0.600,-0.426) (0.550,-0.390) (0.500,-0.355) (0.450,-0.319) (0.400,-0.284) (0.354,-0.250) (0.300,-0.213) (0.250,-0.177) (0.206,-0.144)}
\EQ{(0.206,0.144) (0.248,0.175) (0.300,0.213) (0.350,0.248) (0.400,0.284) (0.450,0.319) (0.495,0.350) (0.550,0.390) (0.599,0.425) (0.650,0.462) (0.700,0.498) (0.750,0.534) (0.800,0.571) (0.841,0.600) (0.900,0.644) (0.942,0.675) (1.000,0.719) (1.041,0.750) (1.100,0.795) (1.150,0.834) (1.200,0.873) (1.250,0.913) (1.296,0.950) (1.350,0.995) (1.400,1.037) (1.445,1.075) (1.500,1.123) (1.550,1.168) (1.600,1.214) (1.639,1.250) (1.675,1.285) (1.716,1.325) (1.765,1.375) (1.800,1.411) (1.837,1.450) (1.875,1.492) (1.925,1.549) (1.968,1.600) (2.000,1.639) (2.048,1.700) (2.086,1.750) (2.122,1.800) (2.156,1.850) (2.189,1.900) (2.221,1.950) (2.250,1.999) (2.275,2.043) (2.300,2.088) (2.325,2.136) (2.350,2.187) (2.375,2.241) (2.400,2.299) (2.421,2.350) (2.440,2.400) (2.457,2.450) (2.474,2.500) (2.489,2.550) (2.502,2.600) (2.515,2.650) (2.526,2.700) (2.537,2.750) (2.545,2.800) (2.553,2.850) (2.560,2.900) (2.565,2.950) (2.569,3.000) (2.572,3.050) (2.574,3.100) (2.575,3.150) (2.574,3.200) (2.572,3.250) (2.570,3.300) (2.566,3.350) (2.560,3.400) (2.554,3.450) (2.546,3.500) (2.538,3.550) (2.528,3.600) (2.516,3.650) (2.504,3.700) (2.490,3.750) (2.475,3.800) (2.459,3.850) (2.442,3.900) (2.423,3.950) (2.402,4.000) (2.381,4.050) (2.357,4.100) (2.333,4.150) (2.307,4.200) (2.279,4.250) (2.250,4.299) (2.218,4.350) (2.185,4.400) (2.151,4.450) (2.114,4.500) (2.075,4.550) (2.034,4.600) (2.000,4.640) (1.950,4.695) (1.900,4.747) (1.850,4.797) (1.800,4.843) (1.750,4.887) (1.705,4.925) (1.650,4.969) (1.609,5.000) (1.550,5.042) (1.502,5.075) (1.450,5.109) (1.400,5.140) (1.350,5.169) (1.300,5.197) (1.250,5.223) (1.200,5.249) (1.150,5.272) (1.100,5.295) (1.050,5.316) (1.000,5.337) (0.950,5.356) (0.900,5.374) (0.850,5.390) (0.800,5.406) (0.750,5.421) (0.700,5.434) (0.650,5.447) (0.600,5.459) (0.550,5.469) (0.500,5.479) (0.450,5.488) (0.400,5.496) (0.350,5.503) (0.300,5.508) (0.250,5.513) (0.200,5.518) (0.150,5.521) (0.100,5.523) (0.050,5.524) (0.000,5.525) (-0.050,5.524) (-0.100,5.523) (-0.150,5.521) (-0.200,5.518) (-0.250,5.513) (-0.300,5.508) (-0.350,5.503) (-0.400,5.496) (-0.450,5.488) (-0.500,5.479) (-0.550,5.469) (-0.600,5.459) (-0.650,5.447) (-0.700,5.434) (-0.750,5.421) (-0.800,5.406) (-0.850,5.390) (-0.900,5.374) (-0.950,5.356) (-1.000,5.337) (-1.050,5.316) (-1.100,5.295) (-1.150,5.272) (-1.197,5.250) (-1.247,5.225) (-1.294,5.200) (-1.340,5.175) (-1.400,5.140) (-1.450,5.109) (-1.500,5.076) (-1.550,5.042) (-1.600,5.006) (-1.642,4.975) (-1.700,4.929) (-1.750,4.887) (-1.792,4.850) (-1.846,4.800) (-1.898,4.750) (-1.946,4.700) (-1.991,4.650) (-2.025,4.611) (-2.075,4.550) (-2.114,4.500) (-2.150,4.451) (-2.185,4.400) (-2.218,4.350) (-2.249,4.300) (-2.275,4.257) (-2.300,4.212) (-2.325,4.165) (-2.350,4.115) (-2.375,4.062) (-2.400,4.005) (-2.423,3.950) (-2.442,3.900) (-2.459,3.850) (-2.475,3.801) (-2.490,3.750) (-2.504,3.700) (-2.516,3.650) (-2.528,3.600) (-2.538,3.550) (-2.546,3.500) (-2.554,3.450) (-2.560,3.400) (-2.566,3.350) (-2.570,3.300) (-2.572,3.250) (-2.574,3.200) (-2.575,3.150) (-2.574,3.100) (-2.572,3.050) (-2.569,3.000) (-2.565,2.950) (-2.560,2.900) (-2.553,2.850) (-2.545,2.800) (-2.537,2.750) (-2.526,2.700) (-2.515,2.650) (-2.502,2.600) (-2.489,2.550) (-2.474,2.500) (-2.457,2.450) (-2.440,2.400) (-2.421,2.350) (-2.401,2.300) (-2.379,2.250) (-2.356,2.200) (-2.332,2.150) (-2.306,2.100) (-2.279,2.050) (-2.251,2.000) (-2.221,1.950) (-2.189,1.900) (-2.156,1.850) (-2.125,1.805) (-2.086,1.750) (-2.050,1.702) (-2.009,1.650) (-1.975,1.608) (-1.926,1.550) (-1.882,1.500) (-1.837,1.450) (-1.800,1.411) (-1.765,1.375) (-1.725,1.334) (-1.675,1.285) (-1.639,1.250) (-1.600,1.214) (-1.558,1.175) (-1.502,1.125) (-1.450,1.080) (-1.400,1.037) (-1.356,1.000) (-1.300,0.954) (-1.250,0.913) (-1.202,0.875) (-1.150,0.834) (-1.107,0.800) (-1.050,0.757) (-1.008,0.725) (-0.950,0.682) (-0.908,0.650) (-0.850,0.607) (-0.807,0.575) (-0.750,0.534) (-0.703,0.500) (-0.650,0.462) (-0.600,0.426) (-0.550,0.390) (-0.500,0.355) (-0.450,0.319) (-0.400,0.284) (-0.354,0.250) (-0.300,0.213) (-0.250,0.177) (-0.220,0.155)}
}
  \marcoMMA{10}{5}{$x$ (u.\,arb.)}{$z$ (u.\,arb.)}
\end{tikzpicture}
\end{minipage}
\caption{Equipotenciales, a intervalos iguales de potencial, de los tres primeros términos del desarrollo multipolar con simetría de revolución alrededor de $z$:
monopolar (esferas $r=\text{cte}$, que se apiñan hacia el centro porque $\phi\propto1/r$), dipolar ($r^{2}\propto\cos\theta$, dos lóbulos de signo opuesto separados por
el plano $z=0$) y cuadrupolar ($r^{3}\propto3\cos^{2}\theta-1$, cuatro lóbulos separados por el cono nodal $\cos^{2}\theta=1/3$, las rectas oblicuas, hacia el que convergen las curvas; la región en blanco, con forma de trébol, es la de potencial más intenso que el último nivel trazado). Recreación
de~\cite{FIGURAS-T2}.}
\label{fig:equipotenciales}
\end{figure}
